\documentclass[11pt,a4paper]{article}
\usepackage[T1]{fontenc}
\usepackage[utf8]{inputenc}
\usepackage{lmodern}
\usepackage[margin=27mm,headheight=14pt]{geometry}
\usepackage{amsmath,amssymb,amsthm,mathtools,bm}
\usepackage{mathrsfs}
\usepackage{microtype,booktabs,longtable,array,enumitem,aliascnt}
\usepackage{fancyhdr}
\usepackage[breaklinks=true,colorlinks=true,linkcolor=black,citecolor=black,urlcolor=black]{hyperref}
\usepackage[nameinlink,noabbrev]{cleveref}
\hypersetup{pdftitle={Renewal Yang--Mills: Exact Shell Mixing, V16},pdfauthor={Research continuation},pdfsubject={Rigorous finite Gaussian response identities and conditional cofinal criteria}}
\setlength{\emergencystretch}{3em}
\setlength{\parskip}{3pt}
\allowdisplaybreaks[2]
\raggedbottom
\pagestyle{fancy}
\fancyhf{}
\fancyhead[L]{\small Yang--Mills: exact shell-response continuation}
\fancyhead[R]{\small V16}
\fancyfoot[C]{\thepage}
\newtheorem{theorem}{Theorem}[section]
\newaliascnt{proposition}{theorem}
\newtheorem{proposition}[proposition]{Proposition}
\aliascntresetthe{proposition}
\crefname{proposition}{proposition}{propositions}
\Crefname{proposition}{Proposition}{Propositions}
\newaliascnt{lemma}{theorem}
\newtheorem{lemma}[lemma]{Lemma}
\aliascntresetthe{lemma}
\crefname{lemma}{lemma}{lemmas}
\Crefname{lemma}{Lemma}{Lemmas}
\newaliascnt{corollary}{theorem}
\newtheorem{corollary}[corollary]{Corollary}
\aliascntresetthe{corollary}
\crefname{corollary}{corollary}{corollaries}
\Crefname{corollary}{Corollary}{Corollaries}
\theoremstyle{definition}
\newaliascnt{definition}{theorem}
\newtheorem{definition}[definition]{Definition}
\aliascntresetthe{definition}
\crefname{definition}{definition}{definitions}
\Crefname{definition}{Definition}{Definitions}
\newaliascnt{example}{theorem}
\newtheorem{example}[example]{Example}
\aliascntresetthe{example}
\crefname{example}{example}{examples}
\Crefname{example}{Example}{Examples}
\newaliascnt{criterion}{theorem}
\newtheorem{criterion}[criterion]{Criterion}
\aliascntresetthe{criterion}
\crefname{criterion}{criterion}{criteria}
\Crefname{criterion}{Criterion}{Criteria}
\theoremstyle{remark}
\newaliascnt{remark}{theorem}
\newtheorem{remark}[remark]{Remark}
\aliascntresetthe{remark}
\crefname{remark}{remark}{remarks}
\Crefname{remark}{Remark}{Remarks}
\newaliascnt{scope}{theorem}
\newtheorem{scope}[scope]{Scope}
\aliascntresetthe{scope}
\crefname{scope}{scope}{scope statements}
\Crefname{scope}{Scope}{Scope statements}
\DeclareMathOperator{\Tr}{Tr}
\DeclareMathOperator{\Cov}{Cov}
\DeclareMathOperator{\Var}{Var}
\DeclareMathOperator{\Ran}{Ran}
\DeclareMathOperator{\diag}{diag}
\DeclareMathOperator{\Rea}{Re}
\newcommand{\R}{\mathbb R}
\newcommand{\C}{\mathbb C}
\newcommand{\E}{\mathbb E}
\newcommand{\T}{\mathbb T}
\newcommand{\dd}{\mathrm d}
\newcommand{\ip}[2]{\langle #1,#2\rangle}
\newcommand{\norm}[1]{\lVert #1\rVert}
\newcommand{\abs}[1]{\lvert #1\rvert}
\newcommand{\HS}{\mathrm{HS}}
\newcommand{\Id}{I}
\newcommand{\bext}{\mathfrak b}
\newcommand{\src}[2]{\cite[#2]{#1}}
\title{Renewal Yang--Mills Programme\\[4pt]
Exact shell mixing, inherited vertices,\\and the complete Gaussian response\\[8pt]
\large Research notes --- Version 16}
\author{Research continuation}
\date{14 September 2026}
\usepackage{etoolbox}
\makeatletter
\patchcmd{\l@section}{1.5em}{2.5em}{}{}
\renewcommand*\l@subsection{\@dottedtocline{2}{2.5em}{3.0em}}
\renewcommand{\@pnumwidth}{2.8em}
\renewcommand{\@tocrmarg}{3.8em}
\makeatother
\begin{document}
\maketitle
\begin{abstract}
The supplied programme already separates compact construction,
renormalization, and physical spectral closure. Its current local-repair
lineage also rules out a temperature-uniform hidden diffusion gap in one
of the actual conditional fibres. We therefore address a different,
upstream quantitative object: the composition of the Gaussian principal
response along the explicit nested covariance tower. A positive shell
resolution does not make the quadratic background determinant additive
in the same-shell bubble terms. We derive the missing mixed-shell terms,
identify their exact Schur and Gaussian-martingale ownership, and express
their total as a commutator energy for genuine Gaussian innovations. A
uniformly analytic example shows that their omission can change the
coefficient accumulated per scale. The actual response normalization in
the closure monograph is retained in a separate two-momentum formula.
We then prove a complete shell-pair fixed-profile criterion, an explicit
finite-band error bound, and a polarized background-comparison identity.
The finite identities are unconditional under their stated matrix
hypotheses; the cofinal Yang--Mills estimates and the physical mass gap
are not claimed. An accompanying exact-arithmetic checker verifies the
new algebra without importing unprovided numerical certificates.
\end{abstract}
\tableofcontents

\section{Context, source audit, and the selected continuation}
\label{sec:v1-context}

\subsection{The inherited programme and its actual frontier}
The basis for this continuation is the supplied Grand Tensor V075,
Yang--Mills closure monograph V3, the completed f12 V100 archive, and the
active V16 local-repair lineage \cite{GT75,YM3,F12,F13}. The purpose is not
to restate their conditional closure theorem or to restart an already
excluded spectral route. The following distinctions determine the new
work.

The Grand Tensor's half-slab construction identifies the physical
boundary channel, its Doob law, and its relation to the transfer edge.
Its action-flow results also distinguish source-relative control from a
uniform scalar gap. These are retained interfaces, not new results here
\src{GT75}{\texttt{subsec:YM-half-slab-hierarchy},
\texttt{prop:YM-divergence-source-flow}}.

The closure monograph provides a positive nested Gaussian block-average
tower, harmonic extensions, and the finite tree-gauge representative
map. It explicitly leaves the higher-shell complete response, its
uniform momentum estimates, and fixed/harmonic background matching
unpopulated. Its final physical contraction hypothesis is independent
of these Gaussian issues \src{YM3}{Chapters on the explicit Gaussian
tower, the physical spectral endpoint, and integrated conditional
closure}.

The f12 archive already contains compact completion, conditional energy
selection, direct heat-bath escape criteria, and a detailed analysis of
noncentral conditional wells. Reproving the coordinate-improvement
criterion would add no acquisition \src{F12}{V18, V74, V89--V100}.
The active V16 lineage has gone further: its V8 proves a secondary-well
obstruction to a temperature-uniform hidden diffusion gap; V13--V16
construct genuine visible-cylinder repair floors without imposing a
hidden global gap. The final V16 theorem still requires its explicit
source-clearance and visible-curvature tests; it does not cover their
complement \src{F13}{V8, V13--V16}.

Consequently, two tempting continuations are not selected: further
attempts to prove that same false hidden diffusion estimate, and an
unjustified promotion of the V16 fixed-event floor to an unrestricted
interface or physical gap. The selected task is instead to make the
Gaussian response being accumulated in the renormalization calculation
compositionally complete.

\subsection{The precise new question}
At a fixed finite regulator, write a covariance as a positive sum
\[
                    C=\sum_{j=0}^{N-1}\Gamma_j.
\]
The tadpole is linear in $C$, but the two-vertex bubble is quadratic.
Thus the expression
\[
 \sum_j\left\{\frac12\Tr(\Gamma_j Q)
                  -\frac12\Tr(\Gamma_jR\Gamma_jR)\right\}
\]
is not generally the second background derivative of the determinant
for the summed covariance. The mixed contractions with two different
shells must either occur explicitly or be carried by the inherited
vertices of the effective action. The word ``complete'' in the
monograph's hypotheses already permits the latter. What is not supplied
by the displayed same-shell functional alone is its exact composition.

This distinction preserves the determinant telescope in
\src{YM3}{\texttt{eq:det-telescope}}. That identity applies to the
\emph{actual conditional determinants}. The work below identifies
concretely what their second jets must contain. Nor is a new discovery
of Gaussian elimination claimed: the acquisition is an explicit,
nonduplicating shell-response ledger for this programme, together with
criteria that can be applied to its already constructed tower.

\subsection{Proof and verification convention}
Every theorem proved below is either a finite-dimensional identity or an
implication with its analytic hypotheses written explicitly. The large
certificate scripts cited inside the supplied manuscripts were not
included as executable attachments; their reported numerical outputs
are not independently recertified here. The new checker concerns only
the new identities and examples. The source audit used structural and
dependency inspection of all four files and detailed reading of the
relevant constructions and terminal statements; it is not a line-by-line
independent verification of every proof in those manuscripts.

The physical target remains a nontrivial four-dimensional Yang--Mills
theory with the required reconstruction properties and a positive
physical Hamiltonian gap. The official problem formulation is retained
as the external target \cite{Clay}. Gauge-invariant renormalization
calculations such as \cite{Morris,AGM} provide external context for the
need to compare complete regularized responses, not an imported proof
of the present regulator's coefficient or continuum limit.

\section{One finite Gaussian law and two distinct parameters}
\label{sec:v1-finite}

Let $\mathsf X$ and $\mathsf Z_j$ be finite-dimensional real Hilbert
spaces, and let
\[
 U_j:\mathsf Z_j\longrightarrow\mathsf X,
 \qquad \Gamma_j=U_jU_j^*,\qquad
 U=(U_0,\ldots,U_{N-1}),\qquad C=UU^*.
\]
No $\Gamma_j$ is required to be invertible, and their ranges need not be
orthogonal in $\mathsf X$. If the $z_j$ are independent standard Gaussian
vectors, then $X=\sum_jU_jz_j$ has covariance $C$. This supplies a common
law rather than a collection of unrelated shell measures.

Let $V(t)$ be a $C^2$ self-adjoint matrix family on $\mathsf X$, with
\begin{equation}\label{eq:v1-Vjet}
             V(0)=0,\qquad V'(0)=R,\qquad V''(0)=Q.
\end{equation}
Here $t$ is a \emph{background amplitude}. It is not the external momentum
$q$ used later, not a renormalization index, and not physical time. Set
\begin{equation}\label{eq:v1-K}
 K(t)=I+U^*V(t)U,
 \qquad
 \mathcal W(t)=-\log\E\exp\left[-\tfrac12\ip{X}{V(t)X}\right].
\end{equation}
At this fixed regulator $K(t)$ is positive for sufficiently small real
$t$, since $K(0)=I$. Define the innovation-space banks
\begin{equation}\label{eq:v1-banks}
 A_{ij}=U_i^*RU_j,\qquad B_{ij}=U_i^*QU_j,
 \qquad A=U^*RU,\quad B=U^*QU.
\end{equation}
The adjoint notation allows the same trace identities to be used after
complexifying a real Fourier representation. The stochastic statements
are proved for the underlying real Gaussian law; no complex-Gaussian
normalization is silently substituted.

\begin{lemma}[Normalizer on a possibly degenerate covariance]
\label{lem:v1-normalizer}
Under the preceding hypotheses,
\begin{equation}\label{eq:v1-normalizer}
 \mathcal W(t)=\frac12\log\det K(t),\qquad
 \mathcal W''(0)=\frac12\Tr B-\frac12\Tr A^2
               =\frac12\Tr(CQ)-\frac12\Tr(CRCR).
\end{equation}
\end{lemma}
\begin{proof}
The joint vector $z=(z_0,\ldots,z_{N-1})$ is standard Gaussian. Combining
its reference density with the exponential in \cref{eq:v1-K} gives the
precision matrix $K(t)$. Orthogonal diagonalization and the elementary
one-dimensional Gaussian integral give the determinant formula.
For any differentiable invertible matrix $K$,
\[
 \partial_t\log\det K=\Tr(K^{-1}K'),\qquad
 \partial_t K^{-1}=-K^{-1}K'K^{-1}.
\]
These follow by differentiating the determinant and $K^{-1}K=I$.
Differentiating once more at $K(0)=I$ proves the first response formula.
Cyclicity of the finite trace and $UU^*=C$ prove the second. No inverse
of $C$ or of an individual shell was used.
\end{proof}

\begin{scope}[The principal determinant is not the full interacting law]
The matrices in \cref{eq:v1-Vjet} must be the derivatives of the declared
quadratic precision. A linear fluctuation source, a background-dependent
scalar action, a Haar amplitude, quotient factors, or nonquadratic
interactions are additional terms unless they have already been included
by a proved reduction. For example, a source $t\ip{h}{X}$ in the
exponent contributes $-\ip{h}{Ch}$ to the second derivative of the
negative log normalizer. The formulas below do not declare such terms
zero. They settle the composition of the specified Gaussian quadratic
component.
\end{scope}

\section{The missing mixed-shell contribution and a sharp obstruction}
\label{sec:v1-mixed}

\begin{theorem}[Exact shell-pair decomposition]
\label{thm:v1-mixed}
Set
\begin{equation}\label{eq:v1-diagonal}
 d_j=\frac12\Tr B_{jj}-\frac12\Tr A_{jj}^2,
 \qquad c_{ij}=\norm{A_{ij}}_\HS^2\quad(i\ne j).
\end{equation}
Then
\begin{equation}\label{eq:v1-mixed}
 \boxed{\displaystyle
 \mathcal W''(0)=\sum_{j=0}^{N-1}d_j
                      -\sum_{0\le i<j<N}c_{ij}.}
\end{equation}
Moreover,
\begin{equation}\label{eq:v1-paircov}
 c_{ij}=\Tr(\Gamma_iR\Gamma_jR)
       =\norm{U_i^*RU_j}_\HS^2\ge0.
\end{equation}
The same-shell formula is exact for this amplitude response if and only
if $U_i^*RU_j=0$ for every $i\ne j$.
\end{theorem}
\begin{proof}
Self-adjointness gives $A_{ji}=A_{ij}^*$. Block multiplication yields
\[
 \Tr A^2=\sum_j\Tr A_{jj}^2
               +2\sum_{i<j}\Tr(A_{ij}A_{ij}^*).
\]
Insert this identity into \cref{eq:v1-normalizer}. Cyclicity gives
\cref{eq:v1-paircov}. Every omitted summand is nonnegative, so its total
vanishes exactly when each block vanishes.
\end{proof}

The independence of the $z_j$ does not imply the vanishing of the
$A_{ij}$. Independence removes a two-point covariance between $z_i$ and
$z_j$; it does not remove the interaction
$\ip{z_i}{A_{ij}z_j}$ or the expectation of its square. This is the
specific distinction needed when using independent Gaussian innovations
inside a nonlinear background response.

\begin{example}[Uniform analytic control does not make the omission sublinear]
\label{ex:v1-chain}
For $n\ge1$, take $\mathsf X=\R^{n+1}$,
$U_j=e_j$ for $0\le j\le n$, $Q=0$, and
\[
 R_n=\sum_{j=0}^{n-1}(e_je_{j+1}^*+e_{j+1}e_j^*).
\]
The shells are independent orthogonal rank-one components of $C=I$.
Nevertheless,
\begin{equation}\label{eq:v1-chain}
 \sum_{j=0}^{n}d_j=0,\qquad
             \mathcal W_n''(0)=-n.
\end{equation}
The family $I+tR_n$ is positive for $|t|<1/2$, uniformly in $n$.
\end{example}
\begin{proof}
Every diagonal block of $R_n$ vanishes. Its only nonzero off-diagonal
blocks are the $n$ nearest-neighbour pairs, each of squared norm one.
Thus \cref{thm:v1-mixed} gives \cref{eq:v1-chain}. For any real $x$,
\[
 |\ip{x}{R_nx}|\le2\sum_{j=0}^{n-1}|x_jx_{j+1}|
       \le2\norm{x}^2.
\]
Hence $I+tR_n\succeq(1-2|t|)I$. The positive lower bound supplies the
stated common analytic neighbourhood.
\end{proof}

This is a counterexample to a Gaussian \emph{composition implication},
not a counterexample to Yang--Mills or to the monograph's conditional
closure theorem. It shows that the unproved assertion ``mixed shells
only contribute a bounded endpoint'' cannot be justified by positivity,
independence, or a common analytic radius alone.

\begin{corollary}[Diagonal cards do not determine the complete response]
\label{cor:v1-cards}
Even complete knowledge of all individual shell covariances and all
same-shell response cards does not determine the summed determinant
response without an off-diagonal vertex or equivalent inherited-action
input.
\end{corollary}
\begin{proof}
Take $C=I_2$, $U_0=e_0$, $U_1=e_1$, $Q=0$, and
$R_\alpha=\left(\begin{smallmatrix}0&\alpha\\\alpha&0\end{smallmatrix}\right)$.
Every same-shell card is zero for every $\alpha$, while the complete
response is $-\alpha^2$.
\end{proof}

\section{Exact elimination and nonduplicating ownership}
\label{sec:v1-owner}

Fix the elimination order $0,1,\ldots,N-1$. For a given $j$, write
$E_j=\{0,\ldots,j-1\}$. Matrices indexed by an empty set contribute zero.
The effective precision on block $j$ after eliminating $E_j$ is
\begin{equation}\label{eq:v1-Schurj}
 S_j(t)=K_{jj}(t)-K_{jE_j}(t)K_{E_jE_j}(t)^{-1}K_{E_jj}(t).
\end{equation}
The inversion is performed near $t=0$, where the eliminated block is
positive. It is an inversion on independent innovation coordinates,
not an inversion of a singular physical covariance.

\begin{theorem}[Inherited second vertex and exact shell ownership]
\label{thm:v1-owner}
For every $j$,
\begin{align}
 S_j(0)&=I,\qquad S_j'(0)=A_{jj},\label{eq:v1-Sjet1}\\
 S_j''(0)&=B_{jj}-2\sum_{i<j}A_{ji}A_{ij}.\label{eq:v1-Sjet2}
\end{align}
Consequently, if $\mathcal W_j(t)=\tfrac12\log\det S_j(t)$, then
\begin{equation}\label{eq:v1-owner}
 \boxed{\displaystyle
 \mathcal W_j''(0)=d_j-\sum_{i<j}\norm{A_{ji}}_\HS^2,
 \qquad \mathcal W(t)=\sum_j\mathcal W_j(t).}
\end{equation}
Each unordered mixed-shell pair is charged exactly once, to its
later-eliminated endpoint. The total is unchanged under any permutation
of the elimination order, although individual owned terms may change.
\end{theorem}
\begin{proof}
At zero background, $K(0)=I$, so the off-diagonal blocks of $K(0)$ vanish.
In the second derivative of the product subtracted in
\cref{eq:v1-Schurj}, the only surviving terms differentiate each outer
off-diagonal block once. There are two such product-rule terms. Derivatives
of the middle inverse are multiplied by an undifferentiated zero outer
block and vanish. This gives
\[
 \left.\partial_t^2\bigl(K_{jE}K_{EE}^{-1}K_{Ej}\bigr)\right|_0
      =2A_{jE}A_{Ej}=2\sum_{i<j}A_{ji}A_{ij},
\]
proving \cref{eq:v1-Sjet1,eq:v1-Sjet2}. Applying the determinant derivative
formula to $S_j$ gives its owned response.

For a positive block matrix,
\[
 \det K_{\le j,\le j}=\det K_{E_jE_j}\det S_j.
\]
Indeed, multiplication on the left and right by block triangular matrices
of determinant one block-diagonalizes it with diagonal blocks
$K_{E_jE_j}$ and $S_j$. Iteration proves the normalizer identity. Every
pair has exactly one later endpoint. Reordering gives the same original
determinant, which proves the final assertion.
\end{proof}

\begin{corollary}[The retained action carries the unspent mixed response]
\label{cor:v1-retained}
For a division into eliminated innovations $E$ and unintegrated
innovations $F$, the exact retained precision obeys
\begin{equation}\label{eq:v1-retained}
 \left(K_{FF}-K_{FE}K_{EE}^{-1}K_{EF}\right)''\Big|_0
                      =B_{FF}-2A_{FE}A_{EF}.
\end{equation}
The eliminated normalizer contains mixed pairs internal to $E$.
Pairs between $E$ and $F$ enter the retained second vertex and must not
also be counted as already integrated scalar bubbles. If the retained
variables are later integrated, those pairs are charged at that later
step.
\end{corollary}
\begin{proof}
The same product-rule calculation proves \cref{eq:v1-retained}. The
block determinant factorization separates the eliminated normalizer
from the retained Gaussian integral. Applying \cref{thm:v1-owner} to the
latter gives the allocation described.
\end{proof}

In particular, an unintegrated retained action is not an additional
Gaussian determinant. If one elects to integrate it, its determinant is
not generally background independent. Any endpoint declared invisible
to $\bext$ must be justified in the same normalization. These distinctions
prevent both omission and double charging.

\subsection{The same ownership from the actual Gaussian score}

\begin{theorem}[Orthogonal martingale ownership of the bubble]
\label{thm:v1-score}
For independent real standard Gaussian $z_j$, let
\[
 \mathcal Q=\tfrac12\bigl(\ip{z}{Az}-\Tr A\bigr),\qquad
 \mathcal F_j=\sigma(z_0,\ldots,z_j).
\]
Then $\mathcal Q=\sum_jM_j$, where
\begin{equation}\label{eq:v1-Mj}
 M_j=\tfrac12\bigl(\ip{z_j}{A_{jj}z_j}-\Tr A_{jj}\bigr)
                       +\sum_{i<j}\ip{z_j}{A_{ji}z_i}.
\end{equation}
The $M_j$ are orthogonal martingale differences and
\begin{equation}\label{eq:v1-scorevar}
 \E M_j^2=\tfrac12\Tr A_{jj}^2+\sum_{i<j}\norm{A_{ji}}_\HS^2,
 \qquad
 \mathcal W_j''(0)=\tfrac12\Tr B_{jj}-\E M_j^2.
\end{equation}
\end{theorem}
\begin{proof}
Expanding the self-adjoint quadratic form groups its diagonal term and
each mixed pair by its larger index, giving \cref{eq:v1-Mj}.
The conditional expectation of $M_j$ given $\mathcal F_{j-1}$ is zero:
$z_j$ has mean zero, covariance $I$, and is independent of the preceding
variables. The martingale-difference property proves orthogonality.

For a real standard Gaussian vector $z$,
\[
 \E(z_az_bz_cz_d)=\delta_{ab}\delta_{cd}
                  +\delta_{ac}\delta_{bd}+\delta_{ad}\delta_{bc}.
\]
This follows by independence, vanishing odd moments, and
$\E z_a^2=1$, $\E z_a^4=3$, the latter obtained by integrating the
one-dimensional Gaussian density by parts. Applying it to the centered
diagonal quadratic form gives variance $\tfrac12\Tr A_{jj}^2$.
For $i<j$, independence gives
$\E\ip{z_j}{A_{ji}z_i}^2=\Tr(A_{ji}A_{ji}^*)$.
Different $i$ give zero cross expectation. The diagonal quadratic term
has zero covariance with every bilinear term, by vanishing odd moments.
This proves the variance formula; \cref{thm:v1-owner} proves the last
identity.
\end{proof}

Thus the response ledger has both a Schur realization and a same-law
conditional-score realization. It does not obtain a stochastic physical
clock from Gaussian elimination, and it does not use a spectral gap of
an auxiliary hidden diffusion.

\section{Genuine innovations and the commutator energy}
\label{sec:v1-commutator}

The preceding results allow arbitrary positive covariance sums. The
specific nested Gaussian tower has additional orthogonality in its
\emph{Gaussian Hilbert space}, rather than in the unweighted fine-edge
coordinates.

\begin{lemma}[Whitening a nested explained-covariance tower]
\label{lem:v1-projectors}
Let $C_0\succ0$ on a finite Gaussian quotient. For nested retained linear
observations with explained covariances $\mathcal C_j$, assume
$\mathcal C_0=C_0$ and that the retained Gaussian linear subspaces decrease
with $j$. There are decreasing orthogonal projections $\mathcal P_j$
with $\mathcal P_0=I$ such that
\[
 \mathcal C_j=C_0^{1/2}\mathcal P_jC_0^{1/2}.
\]
Hence, for $j<N$, the innovations and the retained tail have the form
\begin{equation}\label{eq:v1-Pj}
 \Gamma_j=C_0^{1/2}P_jC_0^{1/2},\quad
 P_j=\mathcal P_j-\mathcal P_{j+1},\qquad
 \mathcal C_N=C_0^{1/2}P_NC_0^{1/2},\quad P_N=\mathcal P_N,
\end{equation}
where the $P_0,\ldots,P_N$ are mutually orthogonal and sum to $I$.
\end{lemma}
\begin{proof}
For an observation $L_jX$, its explained covariance is
\[
 C_0L_j^*(L_jC_0L_j^*)^+L_jC_0.
\]
Whitening gives the orthogonal projection onto
$\Ran(C_0^{1/2}L_j^*)$; the formula remains valid with redundant
observations because the Moore--Penrose inverse removes their null
space. Nested observation spaces give nested projections. For two
nested orthogonal projections $P\succeq Q$, $PQ=QP=Q$, so $P-Q$ is an
orthogonal projection. Successive differences are orthogonal and their
sum telescopes to $I$.
\end{proof}

If $C_0$ is singular because of gauge directions, this statement is made
on its support. It asserts no inversion of a toron or gauge-null mode
outside that support. The monograph's shell identities are the relevant
input for applying this lemma \src{YM3}{\texttt{thm:shells}}.

\begin{theorem}[Exact innovation-commutator defect]
\label{thm:v1-commutator}
Let $P_0,\ldots,P_N$ be the resolution in \cref{lem:v1-projectors}, and
put $\mathsf A=C_0^{1/2}RC_0^{1/2}$. For the complete Gaussian integral,
including the retained tail when that tail is actually integrated,
\begin{equation}\label{eq:v1-commutator}
 \boxed{\displaystyle
 \mathcal W''(0)-\mathcal W''_{\rm same}(0)
 =-\sum_{i<j}\norm{P_i\mathsf AP_j}_\HS^2
 =-\frac14\sum_j\norm{[\mathsf A,P_j]}_\HS^2.}
\end{equation}
Here $\mathcal W''_{\rm same}$ keeps the tadpole and the diagonal
innovation bubbles only. The defect vanishes if and only if
$\mathsf A$ preserves every innovation subspace.
\end{theorem}
\begin{proof}
Use $U_j=C_0^{1/2}|_{\Ran P_j}$ in \cref{thm:v1-mixed}. This identifies
the mixed bank with $P_i\mathsf AP_j$ on the appropriate ranges.
For each $j$, the commutator has off-diagonal blocks
$P_i\mathsf AP_j$ and $-P_j\mathsf AP_i$, $i\ne j$, and zero diagonal
blocks. Orthogonality of the block Hilbert--Schmidt norm gives
\[
 \norm{[\mathsf A,P_j]}_\HS^2
             =2\sum_{i\ne j}\norm{P_i\mathsf AP_j}_\HS^2.
\]
Summing counts each unordered pair four times. Vanishing is equivalent
to all off-diagonal blocks being zero, which is equivalent to commuting
with every $P_j$.
\end{proof}

This gives a concrete alternative to assuming that shell mixing is
small: estimate the commutator banks of the background precision with
the actual Gaussian conditional projections. It is an algebraic
reduction, not a bound on those banks for the full Yang--Mills tower.

\subsection{A two-block geometric interpretation}

\begin{proposition}[Background rotation of a retained Gaussian subspace]
\label{prop:v1-projector-motion}
Let $K(t)=I+t\mathsf A+O(t^2)$ be a $C^2$ positive self-adjoint
precision family, and let $L$ have full row rank. The explained-subspace projection
in standard Gaussian coordinates is
\begin{equation}\label{eq:v1-Pt}
 P(t)=K(t)^{-1/2}L^*
       \bigl(LK(t)^{-1}L^*\bigr)^{-1}LK(t)^{-1/2}.
\end{equation}
Writing $P=P(0)$, one has
\begin{align}
 P'(0)&=-\tfrac12\{(I-P)\mathsf AP+P\mathsf A(I-P)\},
                         \label{eq:v1-Pprime}\\
 \norm{P\mathsf A(I-P)}_\HS^2&=2\norm{P'(0)}_\HS^2.
                         \label{eq:v1-Pmotion}
\end{align}
Thus the omitted mixed bubble at this two-block cut is exactly twice
the squared speed of this projection.
\end{proposition}
\begin{proof}
An invertible, $t$-independent transformation of the rows of $L$ does not
change \cref{eq:v1-Pt}. We may therefore arrange $LL^*=I$, so $P=L^*L$.
Near $t=0$, functional calculus or the convergent binomial series gives
$(K^{-1/2})'(0)=-\mathsf A/2$, while
$(LK^{-1}L^*)'(0)=-L\mathsf AL^*$.
Differentiate \cref{eq:v1-Pt} to obtain
\[
 P'(0)=-\tfrac12\mathsf AP-\tfrac12P\mathsf A+P\mathsf AP,
\]
which is \cref{eq:v1-Pprime}. The two off-diagonal blocks are orthogonal
in Hilbert--Schmidt norm and are adjoints. Their squared norms give
\cref{eq:v1-Pmotion}.
\end{proof}

The projection in this proposition is a Gaussian fluctuation subspace.
It is not the physical background field $b_j^{\rm harm}$. The two
constructions must not be conflated. Also, if the observation map itself
depends on $t$, its derivatives contribute additional terms to
\cref{eq:v1-Pprime}; the proposition deliberately assumes fixed $L$.

\section{Common representatives, regrouping, and data sufficiency}
\label{sec:v1-frames}

\begin{proposition}[Common-frame invariance of all vertex banks]
\label{prop:v1-frame}
Let $S$ be an invertible change of the common physical carrier, and set
\[
 U'_j=SU_j,\qquad R'=S^{-*}RS^{-1},\qquad Q'=S^{-*}QS^{-1}.
\]
Then $A'_{ij}=A_{ij}$ and $B'_{ij}=B_{ij}$ for all $i,j$. Independent
orthogonal changes of the innovation bases conjugate the diagonal banks
and transform mixed banks by left and right orthogonal factors; their
traces and Hilbert--Schmidt norms are unchanged.
\end{proposition}
\begin{proof}
Substitution cancels $S^*S^{-*}$ and $S^{-1}S$. Orthogonal innovation
basis changes give $A_{ij}\mapsto O_i^*A_{ij}O_j$ and similarly for $B$.
Trace and Hilbert--Schmidt invariance prove the rest.
\end{proof}

This statement is not a license to move only a covariance and leave its
vertices in another frame. In particular, the tree-gauge maps of the
monograph have to be assembled into compatible maps from the shell
innovations into one declared fine-edge carrier. An isolated Wilson
vertex need not annihilate a gauge displacement at a nonzero background;
the monograph explicitly warns about representative dependence of
partially assembled cards \src{YM3}{\texttt{prop:representative-dependence}}.
A nonlinear gauge change can additionally transport the measure and
higher vertices. \Cref{prop:v1-frame} applies after that common object has
been fixed; it does not replace that construction.

\begin{proposition}[Regrouping retains every pair exactly once]
\label{prop:v1-regroup}
Partition the shell indices into disjoint groups. Computing the complete
response first within each group and then between the grouped
covariances gives the same result as \cref{thm:v1-mixed}. Replacing each
group by only the sum of its original same-shell cards does not in
general have this property.
\end{proposition}
\begin{proof}
Expand $\Tr(CRCR)$ using the partition. An unordered shell pair either
has both endpoints in one group or has endpoints in two different
groups, never both. The first class enters a grouped self-response; the
second enters a grouped mixed response. The tadpole is linear. Omitting
the first class inside any group gives the defect of
\cref{thm:v1-mixed} restricted to that group.
\end{proof}

For implementation, it is therefore sufficient to retain $U_i^*RU_j$
for the relevant pairs and $U_j^*QU_j$ for the tadpoles. Same-shell
corner cards alone cannot reconstruct the mixed bank by
\cref{cor:v1-cards}. Full-edge vertex data and compatible innovation maps,
when actually supplied, do determine it by matrix multiplication.

\section{Insertion into the monograph's two-momentum normalization}
\label{sec:v1-Wilson}

The coefficient $-1/2$ in \cref{eq:v1-normalizer} is the normalization of
a real Gaussian determinant with the derivatives in
\cref{eq:v1-Vjet}. It must not be silently substituted for the monograph's
reported colour and vertex normalization. We now preserve the latter
exactly.

For a covariance $\Xi(k)$ in one common representative, write
\begin{equation}\label{eq:v1-Omega}
 \Omega(\Xi;k,q)=\frac12\Tr\bigl[\Xi(k)Q_0(k,q)\bigr]
 -\frac13\Tr\bigl[\Xi(k+qe_1)R_0(k,q)\Xi(k)R_0(k,q)^*\bigr].
\end{equation}
This is the functional in
\src{YM3}{\texttt{eq:shell-functional}}, with its momentum dependence
made explicit. The same background prescription is used in every
occurrence of $Q_0,R_0$. For a finite shell family in a common lifted
carrier, put
\[
 \Xi_{<j}(k)=\sum_{i<j}\Xi_i(k),\qquad
 \Xi_{<N}(k)=\sum_{i<N}\Xi_i(k),
\]
and define the ordered mixed bubble
\begin{equation}\label{eq:v1-Bij}
 \mathcal B_{ij}(k,q)=\Tr\bigl[
 \Xi_i(k+qe_1)R_0(k,q)\Xi_j(k)R_0(k,q)^*\bigr].
\end{equation}

\begin{theorem}[Complete prefix response in the stated Wilson normalization]
\label{thm:v1-Wilson}
At every momentum where the displayed matrices are defined,
\begin{equation}\label{eq:v1-Omega-total}
 \boxed{\displaystyle
 \Omega(\Xi_{<N};k,q)
 =\sum_{j<N}\Omega(\Xi_j;k,q)
     -\frac13\sum_{i<j<N}\{\mathcal B_{ij}(k,q)+\mathcal B_{ji}(k,q)\}.}
\end{equation}
An exact ordered increment is
\begin{align}
 \Omega_j^{\rm own}(k,q)
 &=\Omega(\Xi_{<j+1};k,q)-\Omega(\Xi_{<j};k,q)\notag\\
 &=\Omega(\Xi_j;k,q)
       -\frac13\sum_{i<j}(\mathcal B_{ij}+\mathcal B_{ji})(k,q).
                                  \label{eq:v1-Omega-owner}
\end{align}
If the second external-momentum derivative can be passed through the
integral, the downward coefficient satisfies
\begin{align}
 b_{<N}^{\downarrow}
 &:=\frac32\int_{\T^4}
        \partial_q^2\Omega(\Xi_{<N};k,q)\big|_{q=0}\,\dd k,
                                         \label{eq:v1-bcomplete}\\
 b_{<N}^{\downarrow}-\sum_{j<N}b_{j,\rm same}^{\downarrow}
 &=-\frac12\sum_{i<j<N}\int_{\T^4}
       \partial_q^2(\mathcal B_{ij}+\mathcal B_{ji})(k,0)\,\dd k.
                                         \label{eq:v1-bdefect}
\end{align}
\end{theorem}
\begin{proof}
The first term of \cref{eq:v1-Omega} is linear in the covariance. Expanding
both covariance factors of the second term gives the ordered pairs
$(i,j)$. Separate the diagonal pairs and group the two orderings of
each off-diagonal pair. Subtract consecutive prefix formulas to obtain
\cref{eq:v1-Omega-owner}. Finite sums commute with differentiation.
The asserted integral passage then gives \cref{eq:v1-bdefect}, with
$(3/2)(-1/3)=-1/2$.
\end{proof}

\begin{scope}[What this corrects, and what it does not identify]
\Cref{thm:v1-Wilson} is a finite algebraic completion of the displayed
functional, not a proof that this functional alone is the complete
Yang--Mills principal action. It requires one common background and
compatible covariance carriers. The scale-by-scale maps $T_{{\rm g},j}$
are not presumed to make a sum of matrices in unrelated coordinate
spaces meaningful. The full Haar, quotient, scalar, and harmonic
background contributions remain those required by the monograph.
After they are supplied, the identity tells exactly how its indicated
two-vertex component composes.
\end{scope}

\subsection{A derivative criterion without inverse shell covariances}
Suppose $\Xi_i(k)=W_i(k)W_i(k)^*$ on a chosen chart, and set
\[
 Z_{ij}(k,q)=W_i(k+qe_1)^*R_0(k,q)W_j(k).
\]
Then $\mathcal B_{ij}=\norm{Z_{ij}}_\HS^2\ge0$. These factors may be
rectangular; no inverse shell covariance is used. Smooth factors are
an explicit input, not a consequence of mere positivity at a changing
rank.

\begin{lemma}[Complete momentum jet and an integrable majorant]
\label{lem:v1-jet}
If $Z(k,q)$ is $C^2$ in real $q$, then
\begin{equation}\label{eq:v1-jet}
 \partial_q^2\norm{Z}_\HS^2
       =2\norm{Z'}_\HS^2+2\Rea\Tr(Z''Z^*).
\end{equation}
Suppose on $|q|<q_0$ that
$\norm{\partial_q^rZ(k,q)}_\HS\le a_r(k)$, $r=0,1,2$, and
\begin{equation}\label{eq:v1-majorant}
 a_0^2+2a_0a_1+2a_1^2+2a_0a_2\in L^1(\T^4).
\end{equation}
Then the integral of $\norm{Z(k,q)}_\HS^2$ is twice differentiable and its
second derivative is the integral of \cref{eq:v1-jet}. The absolute
second derivative is bounded by $2a_1^2+2a_0a_2$.
\end{lemma}
\begin{proof}
Differentiate $\Tr(ZZ^*)$ twice and combine the two conjugate terms.
The trace Cauchy--Schwarz inequality gives the stated bounds for the
function and its first two derivatives. The fundamental theorem of
calculus in $q$ and dominated convergence justify both differentiations
under the integral.
\end{proof}

Nonnegativity of $\mathcal B_{ij}(k,q)$ does not determine the sign of
its second derivative at zero. In particular, the nonpositive
amplitude-response defect in \cref{thm:v1-mixed} does not by itself
establish the sign of the extracted $F^2$ coefficient. Near singular
Fourier or toron charts, \cref{eq:v1-majorant} or an alternative regular
representation must actually be proved. The fixed covariance convergence
rate alone does not provide this derivative stock.

\section{Closed-word ownership to every order of the quadratic determinant}
\label{sec:v1-words}

\begin{theorem}[Maximum-shell ownership of Gaussian determinant words]
\label{thm:v1-words}
Let $K(t)=I+M(t)$ be a finite positive matrix family, analytic near zero,
with $M(0)=0$ and $\norm{M(t)}<1$. Let $M_{\le j}$ denote its principal
block restriction to shells $0,\ldots,j$. The exact owned normalizer is
\begin{equation}\label{eq:v1-word-owner}
 \mathcal W_j(t)=\frac12\sum_{r\ge1}\frac{(-1)^{r+1}}r
 \left\{\Tr M_{\le j}(t)^r-\Tr M_{<j}(t)^r\right\}.
\end{equation}
In its block-word expansion, precisely the closed words whose maximum
shell index is $j$ occur. Summing over $j$ gives
$\tfrac12\log\det K(t)$ with no omitted or additionally duplicated word.
\end{theorem}
\begin{proof}
For $\norm{M}<1$, the matrix logarithm series converges in operator norm.
At fixed finite dimension it may be traced term by term. The same holds
for every principal restriction because its norm is at most
$\norm{M}$. The trace of the $r$th power is the sum
\[
 \sum_{i_1,\ldots,i_r}\Tr
       (M_{i_1i_2}M_{i_2i_3}\cdots M_{i_ri_1}).
\]
Subtracting the restriction with all indices below $j$ removes exactly
those words whose maximum is below $j$. All remaining indices are at
most $j$, so the surviving maximum is $j$. Summing the differences
telescopes. By block determinant factorization, the difference of the
two traced logarithms is also the logarithm of the Schur precision in
\cref{eq:v1-Schurj}. The factor $1/r$ is the one already present in the
ordinary trace-log series; no extra division by the number of occurrences
of $j$ is introduced.
\end{proof}

This theorem concerns all background orders of a \emph{quadratic
Gaussian determinant}. It is not an all-loop construction of the
interacting Yang--Mills measure. Likewise, norm convergence at each
finite regulator is not a regulator-uniform trace bound. Its utility
is a deterministic ownership rule that can be retained when higher
background jets of the Gaussian principal object are acquired.

\section{A complete shell-pair fixed profile and a finite-band certificate}
\label{sec:v1-profile}

The scalar matching problem must use the complete owned response.
Let $\bext$ be the fixed linear extraction of the declared normalized
background coefficient, after the required momentum integration. Absorb
all prescribed colour and orientation constants into the following
scalars. Write the finite integrated response as
\begin{equation}\label{eq:v1-profile-ledger}
 \mathfrak D_N=\sum_{j=0}^{N-1}d_j
                    -\sum_{j=1}^{N-1}\sum_{\ell=1}^j c_{j,j-\ell}.
\end{equation}
Here $d_j$ is the extracted same-shell contribution and $c_{j,i}$ is the
\emph{complete} extracted unordered mixed-pair debit. Thus $c_{j,i}$
need not be nonnegative after extraction. Equation
\cref{eq:v1-profile-ledger} follows from \cref{thm:v1-owner} or, in the
monograph normalization, from \cref{thm:v1-Wilson}. In this section it is
the exact scalar ledger, not an approximation.

\begin{theorem}[Cofinal coefficient from a complete pair profile]
\label{thm:v1-profile}
Assume $0<\rho,\theta<1$ and constants $0\le A_0,B_0,D_0<\infty$ such that
\begin{align}
 |d_j-d_*|&\le A_0\rho^j,\label{eq:v1-profile1}\\
 |c_{j,j-\ell}-c_\ell|&\le B_0\rho^{j-\ell}\theta^\ell
                         &&(1\le\ell\le j),\label{eq:v1-profile2}\\
 |c_\ell|&\le D_0\theta^\ell &&(\ell\ge1).\label{eq:v1-profile3}
\end{align}
Then the series below is absolutely convergent, and
\begin{equation}\label{eq:v1-bstar}
 \boxed{\displaystyle
 b_*^{\rm complete}=d_*-\sum_{\ell\ge1}c_\ell,\qquad
 \mathfrak D_N=Nb_*^{\rm complete}+O(1).}
\end{equation}
More precisely,
\begin{equation}\label{eq:v1-profile-bound}
 \left|\mathfrak D_N-Nb_*^{\rm complete}\right|
 \le \frac{A_0}{1-\rho}
    +\frac{B_0\theta}{(1-\rho)(1-\theta)}
    +\frac{D_0\theta}{(1-\theta)^2}.
\end{equation}
If the actual complete principal response differs from
$\mathfrak D_N$ by $e_N=o(N)$, the same Ces\`aro coefficient survives.
\end{theorem}
\begin{proof}
Absolute convergence follows from \cref{eq:v1-profile3}. Subtract
$N(d_*-\sum_{\ell\ge1}c_\ell)$ from
\cref{eq:v1-profile-ledger}. The resulting terms are
\[
 \sum_{j<N}(d_j-d_*)
 -\sum_{j<N}\sum_{\ell=1}^j(c_{j,j-\ell}-c_\ell)
 +\sum_{j<N}\sum_{\ell>j}c_\ell.
\]
The first has absolute value at most $A_0/(1-\rho)$. In the second use
$i=j-\ell\ge0$ and sum the two geometric series; its bound is
$B_0\theta/((1-\rho)(1-\theta))$. The last sum is bounded by
\[
 D_0\sum_{\ell\ge1}\min(\ell,N)\theta^\ell
                \le \frac{D_0\theta}{(1-\theta)^2}.
\]
This proves \cref{eq:v1-profile-bound}. Dividing an added $e_N=o(N)$ by
$N$ proves the last assertion.
\end{proof}

The distinction from a purely diagonal fixed-point criterion is the
sum of the entire limiting shell-pair profile. In
\cref{ex:v1-chain}, $d_*=0$, $c_1=1$, and the other $c_\ell$ vanish.
The missing coefficient is therefore order one per shell, consistently
with the exact example.

\begin{corollary}[Explicit finite-band and finite-level errors]
\label{cor:v1-band}
Under \cref{thm:v1-profile}, define
\[
 b_*^{(L)}=d_*-\sum_{\ell=1}^{L}c_\ell,\qquad
 b_{j,L}=d_j-\sum_{\ell=1}^{L}c_{j,j-\ell}\quad(j\ge L).
\]
Then
\begin{align}
 |b_*^{\rm complete}-b_*^{(L)}|
       &\le \frac{D_0\theta^{L+1}}{1-\theta},\label{eq:v1-band-tail}\\
 |b_{j,L}-b_*^{\rm complete}|
       &\le A_0\rho^j
          +B_0\sum_{\ell=1}^{L}\rho^{j-\ell}\theta^\ell
          +\frac{D_0\theta^{L+1}}{1-\theta}.
                                             \label{eq:v1-band-level}
\end{align}
In particular, certified intervals
$d_*\in[d_-,d_+]$ and $c_\ell\in[c_{\ell,-},c_{\ell,+}]$ for
$1\le\ell\le L$ give the interval
\begin{equation}\label{eq:v1-certified-interval}
 \left[d_- -\sum_{\ell=1}^{L}c_{\ell,+}-\varepsilon_L,\quad
       d_+ -\sum_{\ell=1}^{L}c_{\ell,-}+\varepsilon_L\right],
 \qquad \varepsilon_L=\frac{D_0\theta^{L+1}}{1-\theta},
\end{equation}
for the complete limiting coefficient.
\end{corollary}
\begin{proof}
Sum the geometric tail for \cref{eq:v1-band-tail}. Subtract the two
finite-band sums and apply \cref{eq:v1-profile1,eq:v1-profile2} to obtain
\cref{eq:v1-band-level}. Interval subtraction and the tail enclosure
prove \cref{eq:v1-certified-interval}.
\end{proof}

\begin{criterion}[What would populate the Gaussian matching interface]
\label{crit:v1-YM}
To apply this result to the monograph's Gaussian matching problem, the
$d_j,c_{j,i}$ must come from its actual common representative and one
complete background prescription. One must prove the relevant momentum
and shell-separation bounds, account for all other principal terms and
endpoint contributions, and identify
\[
       b_*^{\rm complete}
          =b_{\rm YM}=\frac{22C_A}{3(4\pi)^2}\log2
\]
in the monograph's upward convention. Under these inputs,
\cref{thm:v1-profile} supplies the Gaussian coefficient part of its
averaged determinant matching card. None of these quantitative
Yang--Mills inputs is inferred from the abstract geometric-rate
hypotheses above.
\end{criterion}

\section{A complete fixed/harmonic background comparison}
\label{sec:v1-background}

The finite background discrepancy in the monograph is already a reason
not to identify the fixed and harmonic prescriptions by notation. We
now give a directly computable polarized defect for their complete
Gaussian banks, including the mixed shell entries.

Let $(A_h(q),B_h(q))$ and $(A_f(q),B_f(q))$ be the first and second
background-amplitude precision jets on a \emph{common identified
innovation space}, with all four banks $C^2$ in real $q$. All dependence on the external momentum, background
aliases, representative maps, and covariance factors is included in
these banks. Set
\begin{equation}\label{eq:v1-background-vars}
 D=A_h-A_f,\qquad S=A_h+A_f,\qquad E_B=B_h-B_f.
\end{equation}
The same construction works at a fixed finite torus before replacing
sums by integrals. Define the amplitude responses
\[
 F_h(q)=\tfrac12\Tr B_h(q)-\tfrac12\Tr A_h(q)^2,
 \qquad F_f(q)=\tfrac12\Tr B_f(q)-\tfrac12\Tr A_f(q)^2.
\]

\begin{theorem}[Polarized defect with all external-momentum derivatives]
\label{thm:v1-background}
For the complete banks,
\begin{align}
 F_h-F_f&=\tfrac12\Tr E_B-\tfrac12\Tr(DS),
                                          \label{eq:v1-bgpolar}\\
 \partial_q^2(F_h-F_f)
   &=\tfrac12\Tr E_B''
       -\tfrac12\Tr(D''S+2D'S'+DS'').
                                          \label{eq:v1-bgjet}
\end{align}
Consequently,
\begin{equation}\label{eq:v1-bg-bound}
 \abs{\partial_q^2(F_h-F_f)}
 \le\tfrac12\norm{E_B''}_1
    +\tfrac12\sum_{r=0}^2\binom2r
                    \norm{D^{(r)}}_\HS\norm{S^{(2-r)}}_\HS.
\end{equation}
Here $\norm{\cdot}_1$ is the trace norm. The formulas require no
commutativity of $A_h$ and $A_f$.
\end{theorem}
\begin{proof}
Expand
$(A_h-A_f)(A_h+A_f)=A_h^2-A_f^2+[A_h,A_f]$.
The commutator has zero trace, giving \cref{eq:v1-bgpolar}. Differentiate
it twice, retaining both first-derivative terms, to obtain
\cref{eq:v1-bgjet}. The inequalities
$|\Tr X|\le\norm X_1$ and
$|\Tr XY|\le\norm X_\HS\norm Y_\HS$ prove
\cref{eq:v1-bg-bound}.
\end{proof}

This identity also states exactly why a background alias that is merely
$O(q)$ need not be invisible in an $F^2$ calculation: it can contribute
through $D'S'$ in \cref{eq:v1-bgjet}. This interpretation agrees with,
and does not replace, the monograph's already documented background
warning. The new point is an explicit whole-bank defect to be evaluated
or bounded after the mixed-shell completion.

\begin{corollary}[A sufficient sublinear comparison certificate]
\label{cor:v1-bg-Cesaro}
Suppose the complete banks for the first $N$ integrated shells admit the
same prescribed normalization and the needed differentiated-integral
passage. If the integral of the right-hand side of
\cref{eq:v1-bg-bound}, including the fixed extraction multiplier, is
$o(N)$, then the extracted fixed/harmonic difference is $o(N)$.
If its decomposition into actual conditional shell increments is exact,
its averaged shell discrepancy tends to zero.
\end{corollary}
\begin{proof}
Integrate \cref{eq:v1-bg-bound} and divide by $N$. Exact conditional
normalization telescopes the whole-bank difference into the sum of its
owned increments. No diagonal-only replacement is made.
\end{proof}

The hypothesis is sufficient, not automatically necessary: signed
cancellation may yield a smaller defect than its norm bound. In
particular, a failure of this sufficient estimate alone would not
refute background compatibility. Conversely, the existence of a
formal coboundary notation does not establish the estimate. This
comparison also leaves any separate Haar, scalar, or nonlinear
background terms in their existing ledger.

\section{Executed exact checks and their scope}
\label{sec:v1-checks}

The new file \path{verify_ym_shell_mixing_V1.py} provides both a reusable
finite response ledger and independent algebraic checks. It uses exact
rational and symbolic arithmetic. In particular, the Schur checks use a
generic inverse-series recursion for ordinary Taylor coefficients,
rather than substituting the already derived owner formula as their
only calculation. A direct polynomial determinant is an additional
cross-check.

For the nonorthogonal finite example in the script, the complete
response, same-shell response, and omitted mixed debit are respectively
\begin{equation}\label{eq:v1-check-values}
 -\frac{4507}{12},\qquad -\frac{2005}{9},\qquad \frac{5501}{36},
 \qquad
 -\frac{4507}{12}=-\frac{2005}{9}-\frac{5501}{36}.
\end{equation}
The executed report contains ten passing test families: the degenerate
shell covariance identity and direct determinant; all six orders of a
three-block elimination; the retained second-vertex correction; the
orthogonal linear-drift family at dimensions two through twelve; the
complex-Hermitian commutator and projection-motion factors; full second
external-momentum derivatives; both ordered mixed terms in the
monograph's $1/2,-1/3$ normalization; the polarized background jet;
closed-word ownership through sixth order; and the cofinal profile error
bound at thirty finite lengths.

The command
\begin{quote}\small\ttfamily
python verify\_ym\_shell\_mixing\_V1.py --output verification.json
\end{quote}
reruns these checks. The function \texttt{response\_ledger} accepts
Hermitian precision jets and declared innovation-block sizes, with an
optional elimination order. Supplying the actual Yang--Mills banks is a
separate acquisition; the script does not disguise the rational test
matrices as those banks.

\begin{scope}[What PASS establishes]
The executed report supports the new finite algebra and the stated
rational examples. It is not a Lean formalization, a certification of
the older manuscript scripts, a quadrature of the actual full-torus
Yang--Mills response, a numerical value for the limiting anomaly, or a
proof of the physical mass gap. The general statements rest on the
proofs above, not on an inference from a finite list of checks.
\end{scope}

\section{What has advanced and the next actual acquisition}
\label{sec:v1-next}

The present work closes a specific finite composition issue. The
second-order Gaussian normalizer can now be assembled from one common
innovation bank with exact ownership. The mixed term is not an
unspecified remainder: it is a known off-diagonal contraction, an
inherited Schur vertex, a Gaussian score variance, and, for genuine
innovations, a commutator energy. These are equivalent descriptions of
the same contribution, not four independently chargeable errors.

The selected continuation is therefore more specific than another
request for a full-torus scalar. Its next finite target is the actual
\emph{adjacent-shell} bank, including its true external-momentum jets,
in a common representation and preferably the actual harmonic
background. For a declared shell factorization the object is
\begin{equation}\label{eq:v1-nextbank}
 U_{j-1}(k+qe_1)^*R_0(b_j^{\rm harm};k,q)U_j(k),
 \qquad \partial_q^r\big|_{q=0},\quad r=0,1,2,
\end{equation}
with all alias and colour identifications explicit. The notation in
\cref{eq:v1-nextbank} is a target schema: constructing compatible
$U_j$ and the full harmonic vertex germ is part of its acquisition.
The first certificate should compare the sum of owned increments to
the derivative of a direct two-step Gaussian normalizer. A match of
same-shell corner cards alone is not that test.

After this finite identity is populated, the analytic target is a
shell-separation bound for the extracted mixed rows and the momentum
majorants, sufficient for \cref{thm:v1-profile}. The correction to the
limiting scalar is then the convergent pair profile rather than a
presumed endpoint. The scalar still has to match $b_{\rm YM}$ together
with the other complete principal contributions. This supplies a
concrete route to the renormalization interface rather than assuming
its answer.

On the spectral side, the V16 visible-cylinder theorem remains a valid
local acquisition in its stated regime. This note does not extend it
through the near-active-source or inconclusive-curvature complement.
Even a complete repair inequality would still require the correct
physical-law incidence and the calibrated physical clock. On the
construction side, interacting compact continuation, local continuum
existence, reconstruction, and noncollapse remain the independent
obligations already stated in the closure monograph.

\begin{scope}[Position after this version]
The new finite identities and counterexample are proved. The shell-pair
Ces\`aro theorem and background-comparison theorem are proved with their
explicit hypotheses. No numerical estimate for the actual mixed
Yang--Mills bank, no verification of the shell-profile hypotheses for
that bank, no complete value of $b_*^{\rm complete}$, and no positive
physical Hamiltonian mass gap are claimed. The reason for the remaining
boundary is the identified missing estimates and construction inputs,
not the historical status of the problem.
\end{scope}

\subsection{Append-only continuation convention}
This is a new lineage beginning at V1. The supplied manuscripts are
unchanged. A subsequent requested version is to retain this mathematical
body and append its new sections after the V1 end marker, before
\verb|\end{document}|, with a filename ending in \texttt{\_V2.tex}, then
\texttt{\_V3.tex}, and so on. Existing labels beginning
\texttt{v1-} remain fixed. Corrections, if needed, are to be stated
explicitly in an appended correction section rather than silently
rewriting a prior proof. The accompanying provenance record identifies
the exact input and output files by hashes.

\begin{thebibliography}{99}
\bibitem{GT75}
\emph{Renewal Grand Tensor Monograph}, V075, supplied manuscript,
\path{Renewal_Grand_Tensor_Monograph_V075(6).tex}.
Relevant interface: the section ``Yang--Mills: transfer gaps and
zero-free half-planes'', especially the half-slab hierarchy and
source-relative action-flow results.

\bibitem{YM3}
\emph{Renewal Yang--Mills Closure Monograph}, V3, supplied manuscript,
\path{Renewal_Yang_Mills_Closure_Monograph_V3 (1).tex}.
Relevant labels include \texttt{thm:shells},
\texttt{eq:shell-functional}, \texttt{eq:nested-tree-cov},
\texttt{eq:background-Cesaro}, \texttt{eq:det-telescope}, and
\texttt{hyp:contraction}.

\bibitem{F12}
\emph{Renewal Yang--Mills Mass Gap Research Notes}, f12 V100 archive,
supplied manuscript,
\path{Renewal_Yang_Mills_Mass_Gap_Research_Notes_V100_f12(1).tex}.
The finite compact completion and late conditional spectral programme
are inherited, not recertified by the present checker.

\bibitem{F13}
\emph{Renewal Yang--Mills Mass Gap Research Notes}, active V16 lineage
(internal labels \texttt{f13v1}--\texttt{f13v16}), supplied manuscript,
\path{Renewal_Yang_Mills_Mass_Gap_Research_Notes_V16.tex}.
Relevant results: V8's secondary-well diffusion obstruction and
V13--V16's visible-cylinder repair theorems.

\bibitem{Clay}
Clay Mathematics Institute, \emph{Yang--Mills \& the Mass Gap}, official
problem overview and link to the problem description by A.~Jaffe and
E.~Witten. Consulted 13 September 2026.
\url{https://www.claymath.org/millennium/yang-mills-the-maths-gap/}.

\bibitem{Morris}
T.~R. Morris, \emph{A gauge invariant exact renormalization group II},
JHEP \textbf{0012} (2000), 012, arXiv:hep-th/0006064.
\url{https://arxiv.org/abs/hep-th/0006064}.
Cited for external context only; no theorem from this work is used as
an unverified identification of the present compact response.

\bibitem{AGM}
S.~Arnone, A.~Gatti and T.~R. Morris,
\emph{Towards a manifestly gauge invariant and universal calculus for
Yang--Mills theory}, Acta Phys. Slov. \textbf{52} (2002), 621--634,
arXiv:hep-th/0209130.
\url{https://arxiv.org/abs/hep-th/0209130}.
\end{thebibliography}

% END OF VERSION 1 MATHEMATICAL BODY.
% Append subsequent requested versions below this marker; retain V1 above.
\clearpage
\section{V2: all-level harmonic jets and a populated adjacent-shell Wilson test}
\label{sec:v2-context}

This continuation retains the V1 mathematical body. It addresses the next
acquisition in \cref{eq:v1-nextbank}, rather than proving another abstract
shell-ownership identity. There are two new parts. First, the axial
harmonic lift is solved at every level of the supplied block-average
Gaussian tower, including a common tree representative and its second
external-momentum jet. Second, two successive innovations of that tower
are inserted into an explicitly reconstructed fine-link Wilson vertex.
An exact nonzero mixed contraction, its first two external-momentum
derivatives, and an independent two-step determinant check are obtained.

\paragraph{Nonduplication and scope of the populated calculation.}
The one-step axial minimizing lift and its density symbol are already
proved in f12 V2, and are not new results of this version
\src{F12}{\texttt{thm:f12v2-lift}, \texttt{thm:f12v2-measure}}.
The new harmonic results concern all dyadic levels, a uniform analytic
bound, and an exact replacement of the complete second jet. The finite
Wilson calculation below is the canonical fine-link Wilson component
in a stated exponential chart and a common ordered tree. It uses actual
innovations of the supplied Gaussian block-average rule, not arbitrary
positive test matrices. It does \emph{not} claim that the separately
required nonlinear compact-block tangent, Haar, quotient, or measure
vertices have been included. Nor is the common ordered tree silently
identified with every nested representative convention in the source.

The use of automated finite lattice vertices is not itself a novelty
claim; background-field vertex algorithms are established external
context \cite{HammantV2,LWV2}. The proofs and finite certificates here
are derived from the action and block rule displayed below. No external
renormalization theorem is used to identify their coefficient with the
programme's full principal response.

\section{The axial block-average tower in closed form}
\label{sec:v2-axis}

Let $L=2^n$, $n\ge0$, and consider coarse momentum $qe_1$ and transverse
polarization $e_2$. The axes may subsequently be relabelled. The scalar
transverse covariance is denoted $g_L(q)$, with the inherited convention
$g_1(q)=(4\sin^2(q/2))^{-1}$. This is the covariance convention of
\src{YM3}{\texttt{eq:tower-rec}}, before any separate colour conversion.
Put
\begin{equation}\label{eq:v2-denom}
 z=e^{\mathrm i q},\quad
 D_L(z)=6L^2z+(L^2-1)(z-1)^2,\quad
 S_L(q)=g_L(q)^{-1}.
\end{equation}
The inverse in this section is initially at $q\ne0$ in the transverse
axial mode. Its analytic small-$q$ expression is not an inverse of a
physical zero mode at $q=0$.

\begin{theorem}[Every-level axial covariance]
\label{thm:v2-axis-cov}
For every dyadic $L$ and real $q\notin2\pi\mathbb Z$,
\begin{align}
 g_L(q)&=\frac1{4\sin^2(q/2)}-\frac{1-L^{-2}}6,
                                                   \label{eq:v2-g}\\
 S_L(q)&=-\frac{6L^2(z-1)^2}{D_L(z)}.                \label{eq:v2-S}
\end{align}
In particular, the supplied limiting axial covariance is approached with
exact difference $g_L-g_\infty=1/(6L^2)$.
\end{theorem}
\begin{proof}
On the axial polarization, the only nonzero retained aliases in one
blocking step have momenta $q/2$ and $q/2+\pi$. Their block coefficients
are $1+e^{\mathrm i q/2}$ and $1-e^{\mathrm i q/2}$. Thus
\[
 g_{2L}(q)=\frac1{16}\bigl(
 |1+e^{\mathrm i q/2}|^2g_L(q/2)
 +|1-e^{\mathrm i q/2}|^2g_L(q/2+\pi)\bigr).
\]
Writing $x=q/4$ and $g_L(p)=(4\sin^2(p/2))^{-1}-a_L$, this becomes
\[
 \frac1{16}(\cot^2x+\tan^2x)-\frac{a_L}{4}
 =\frac1{4\sin^2(q/2)}-\frac18-\frac{a_L}{4}.
\]
The recursion $a_{2L}=1/8+a_L/4$, $a_1=0$, has solution
$a_L=(1-L^{-2})/6$. Converting the sine expression to $z$ gives
\cref{eq:v2-S}. Its limit agrees with the inherited fixed-point formula;
that limiting formula is not claimed anew.
\end{proof}

\subsection{An exact all-level minimizing profile}
Use the ordered macrocell $\{0,\ldots,L-1\}^4$. In the transverse
representative, the harmonic background has only its second link
component nonzero and is constant in the other coordinates:
\[
 a_2(s;q)=h_{L,s_1}(q),\qquad a_\mu=0\quad(\mu\ne2).
\]
Extend $h$ by the Bloch rule $h_{s+L}=zh_s$. The retained block constraint
is $\sum_{s=0}^{L-1}h_s=1$. In these units its energy is
$L^3\sum_s|h_{s+1}-h_s|^2$.

\begin{theorem}[Rational harmonic profile at every level]
\label{thm:v2-harmonic}
Set $\kappa_L=S_L/L^3$ and
\begin{align}
 A_L(z)&=\frac{3\{(L-1)+(L+1)z\}}{D_L(z)},\label{eq:v2-AB}\\
 B_L(z)&=\frac{6(z-1)(L-1+z)}{L D_L(z)}.\notag
\end{align}
Then the unique transverse minimizing profile at $z\ne1$ is
\begin{equation}\label{eq:v2-h}
 \boxed{\displaystyle
 h_{L,s}(q)=A_L(z)+sB_L(z)-\frac{\kappa_L(z)}2s(s-1),
 \qquad 0\le s<L.}
\end{equation}
It extends analytically to $q=0$, where $h_{L,s}(0)=1/L$. Moreover,
\begin{equation}\label{eq:v2-euler}
 2h_s-h_{s+1}-h_{s-1}=\frac{S_L}{L^3},\qquad
 \sum_{s=0}^{L-1}h_s=1,\qquad
 L^3\sum_s|h_{s+1}-h_s|^2=S_L
\end{equation}
for real $q\ne0$, with Bloch boundary terms retained.
\end{theorem}
\begin{proof}
The second difference of \cref{eq:v2-h} is $\kappa_L$ in the interior.
The two boundary identities needed to extend that equation are
\[
 A_L+LB_L-\frac{\kappa_L}2L(L-1)=zA_L,\qquad
 B_L+\kappa_L=z^{-1}\{B_L-\kappa_L(L-1)\}.
\]
Substitution of \cref{eq:v2-S,eq:v2-AB} proves both. Summing the
quadratic polynomial using
$\sum s=L(L-1)/2$ and $\sum s(s-1)=L(L-1)(L-2)/3$ gives the block
constraint exactly.

For $z\ne1$ on the unit circle, the Bloch difference has no constant
kernel, so its quadratic energy is strictly positive. The Euler equation
with the affine constraint therefore identifies the unique minimizer.
Multiplying it by $\overline h_s$, summing with the unitary Bloch boundary
condition, and using $\sum\overline h_s=1$, gives the energy identity.
The full supplied Fourier minimizer is this same profile: its nonzero
retained aliases have momentum only in direction one, and $e_2$ is
transverse for each of them. All other retained alias weights vanish.
Thus the scalar constrained problem is the axial restriction of that
minimizer, not a replacement block rule. Analytic continuation of the
rational expression gives the stated value at zero.
\end{proof}

For $L=2$, \cref{eq:v2-h} gives
\[
 h_{2,0}=\frac{1+3z}{z^2+6z+1},\qquad
 h_{2,1}=\frac{z(3+z)}{z^2+6z+1}.
\]
These are exactly f12 V2's $\alpha$ and $\beta$, respectively. This is a
regression to the inherited one-step result, not an additional result
counted twice.

\section{A common tree representative and an exact second-jet substitution}
\label{sec:v2-tree}

Choose the ordered rooted tree whose path from the macrocell root first
moves in direction one, then two, then three, then four. For the preceding
field the tree potential is $\phi(s)=s_2h_{L,s_1}$. Its scalar block
average is
\[
 \mu_L=\frac{L-1}{2L},
\]
because $\sum h_s=1$. The coarse longitudinal part created by subtracting
$d\phi$ is removed by the residual cell-root gauge. The resulting
common representative is
\begin{align}
 b_{L,2}(s;q)&=L h_{L,s_1}(q)\,1_{\{s_2=L-1\}},\label{eq:v2-treefield}\\
 b_{L,1}(s;q)&=-s_2\{h_{L,s_1+1}(q)-h_{L,s_1}(q)\}
       +\mu_L(z-1)1_{\{s_1=L-1\}},\notag\\
 b_{L,3}&=b_{L,4}=0.\notag
\end{align}
The expression at $s_1=L-1$ uses $h_L=zh_0$. Translation to other
macrocells includes their Bloch phase.

\begin{proposition}[The tree constraints and the external zero limit]
\label{prop:v2-tree}
For real $q\ne0$, \cref{eq:v2-treefield} is the unique representative
of this fine gauge class with the declared tree links zero and its
coarse longitudinal link zero. Its retained transverse link is $e_2$.
Formula \cref{eq:v2-treefield} is regular at $q=0$ and contains no
uncancelled $1/(e^{\mathrm i q}-1)$ factor.
\end{proposition}
\begin{proof}
Tree integration gives the displayed $\phi$. The block Ward identity
sends $d\phi$ to the coarse gradient of its scalar block average, namely
$(z-1)\mu_Le_1$. A cellwise constant potential has gradient
$(z-1)e_1$ on the direction-one macrocell boundary, and its block is that
same coarse gradient. Adding the final term of \cref{eq:v2-treefield}
therefore removes precisely the unwanted longitudinal link. It changes
neither the curl nor the tree constraint. Two tree-zero representatives
in the same gauge class differ only by such a cell-root gauge, whose
coarse longitudinal component is nonzero at $q\ne0$ unless its
coefficient vanishes. This proves uniqueness. The explicit formula,
rather than an inverse of a zero-momentum gauge symbol, provides the
zero limit.
\end{proof}

\subsection{The plane-wave replacement is not the old fixed background}
Define the normalized fine plane wave
\begin{equation}\label{eq:v2-plane}
 h^{\rm pw}_{L,s}(q)=
 \frac{e^{\mathrm i qs/L}}{\sum_{r=0}^{L-1}e^{\mathrm i qr/L}}
 =e^{\mathrm i q(s-(L-1)/2)/L}
       \frac{\sin(q/(2L))}{\sin(q/2)}.
\end{equation}
Its zero limit is $1/L$ and its block constraint is exact. Let
$b_L^{\rm pw}$ be obtained by applying precisely
\cref{eq:v2-treefield} to this profile, with the same $\mu_L$.

\begin{theorem}[Exact all-level equality of harmonic and plane-wave second jets]
\label{thm:v2-jet-replacement}
For every dyadic $L$, every link in the macrocell, and $r=0,1,2$,
\begin{equation}\label{eq:v2-jet-equality}
 \boxed{\displaystyle
 \left.\partial_q^r b_L^{\rm harm}(s;q)\right|_0
 =\left.\partial_q^r b_L^{\rm pw}(s;q)\right|_0.}
\end{equation}
Write $x_s=(s_1-(L-1)/2)/L$ and $a_L=(1-L^{-2})/12$.
The nonzero component jets, as true derivatives rather than Taylor
coefficients, are
\begin{align}
 b_{L,2}^{(0)}&=1_{\{s_2=L-1\}},\notag\\
 b_{L,2}^{(1)}&=\mathrm i x_s1_{\{s_2=L-1\}},\notag\\
 b_{L,2}^{(2)}&=(a_L-x_s^2)1_{\{s_2=L-1\}},\label{eq:v2-treejets}\\
 b_{L,1}^{(0)}&=0,\notag\\
 b_{L,1}^{(1)}&=-\frac{\mathrm i s_2}{L^2}
                  +\mathrm i\mu_L1_{\{s_1=L-1\}},\notag\\
 b_{L,1}^{(2)}&=\frac{s_2(2s_1-L+2)}{L^3}
                  -\mu_L1_{\{s_1=L-1\}}.\notag
\end{align}
\end{theorem}
\begin{proof}
For both profiles, direct expansion gives
\begin{equation}\label{eq:v2-hjets}
 h_{L,s}(q)=\frac1L\left\{1+\mathrm i x_s q+
             \left(\frac{1-L^{-2}}{24}-\frac{x_s^2}{2}\right)q^2
                    \right\}+O_L(q^3).
\end{equation}
For the harmonic profile, this follows by substituting
$z=1+\mathrm i q-q^2/2+O(q^3)$ into \cref{eq:v2-h}. For the plane wave,
it follows from the sine ratio and phase in \cref{eq:v2-plane}.
The boundary continuations also have equal second jets, since both use
multiplication by $z$. Substituting into the linear expression
\cref{eq:v2-treefield} proves \cref{eq:v2-jet-equality}. For the first
component, the plane wave satisfies
$h^{\rm pw}_{s+1}-h^{\rm pw}_s=(e^{\mathrm i q/L}-1)h^{\rm pw}_s$,
including the last edge. Expanding that product and the residual
$\mu_L(z-1)$ gives \cref{eq:v2-treejets}.
\end{proof}

\begin{remark}[Which aliases have disappeared, and which have not]
Before the tree transformation, the source's harmonic alias formula is
\[
 M_m(q)=L^{-4}\frac{\overline{D^{\rm seg}_L(p_m)}}
                       {4\sin^2(p_m/2)}S_L(q),\qquad
 p_m=\frac{q+2\pi m}{L},\quad
 D^{\rm seg}_L(p)=\sum_{r=0}^{L-1}e^{\mathrm i rp}.
\]
For $m\not\equiv0\pmod L$, the numerator is $O_L(q)$, the denominator
has a nonzero limit, and $S_L(q)=O(q^2)$. Hence $M_m=O_L(q^3)$.
The constraint implies
$M_0-(D^{\rm seg}_L(q/L))^{-1}=O_L(q^4)$.
The tree operation nevertheless creates nonconstant fine aliases of
order $q$ and $q^2$. Those are retained in
\cref{eq:v2-treefield,eq:v2-treejets}; the theorem does not delete them.
It replaces the harmonic transverse input \emph{before the same tree
operation}. It does not identify this new plane-wave prescription with
$b^{\rm fix}$ in the monograph, whose reported corner discrepancy is
unchanged.
\end{remark}

\begin{theorem}[Uniform external analytic stock]
\label{thm:v2-uniform}
For every $0<a<\operatorname{arcosh}2$, there is a constant $C_a$,
independent of dyadic $L$, such that the harmonic profile is holomorphic
on $|\operatorname{Im}q|<a$ and
\begin{equation}\label{eq:v2-uniform}
 \sup_s L|h_{L,s}(q)|
 +\sup_s L^2|h_{L,s+1}(q)-h_{L,s}(q)|
 +\sup_{s,\mu}|b_{L,\mu}(s;q)|\le C_a
\end{equation}
on every closed narrower strip. All fixed-order external derivatives
have corresponding uniform bounds on still narrower strips. There are
also constants $C,r>0$, independent of $L$, such that
\begin{equation}\label{eq:v2-uniform-defect}
 \sup_{s,\mu}|b^{\rm harm}_{L,\mu}(s;q)
                   -b^{\rm pw}_{L,\mu}(s;q)|\le C|q|^3
 \qquad(|q|\le r).
\end{equation}
\end{theorem}
\begin{proof}
For $q=x+\mathrm i y$,
\[
 \operatorname{Re}(D_L(e^{\mathrm i q})/e^{\mathrm i q})
 =2(2L^2+1)+2(L^2-1)\operatorname{Re}\cos q
 \ge2(2-\cosh a)L^2+2(1+\cosh a)
\]
when $|y|\le a$. Thus the denominator does not vanish and has a uniform
lower bound of order $L^2$ after the harmless factor $|z|$ is accounted
for. Formula \cref{eq:v2-h} has numerator of order $L$, uniformly for
$0\le s<L$. Its exact forward difference is $B_L-\kappa_Ls$, including
the Bloch boundary, and has size $O_a(L^{-2})$. This proves the first
two bounds; \cref{eq:v2-treefield} proves the third. Cauchy's formula on
circles contained in a larger strip gives the derivative bounds.

On a fixed complex disk of sufficiently small radius, both $Lh^{\rm harm}$
and $Lh^{\rm pw}$ are uniformly bounded. For the latter this follows
from the sine-ratio formula, with its removable zero at the origin.
Their first three Taylor coefficients agree by \cref{eq:v2-hjets}.
Cauchy's coefficient estimate and the geometric tail of the Taylor
series give $L\sup_s|h^{\rm harm}-h^{\rm pw}|\le C|q|^3$ on a smaller
disk. Their block averages agree exactly, so the residual root-gauge
terms cancel. The second component of the tree difference is $L$ times
the profile difference. The first is at most $(L-1)$ times a difference
of two such values, with a bounded boundary Bloch factor. This proves
\cref{eq:v2-uniform-defect}.
\end{proof}

These are pointwise link and external-momentum bounds, not uniform
Hilbert--Schmidt trace bounds over a growing macrocell. In particular,
they do not control the internal loop-momentum toron charts.

\begin{corollary}[Second-jet substitution in complete response banks]
\label{cor:v2-bank-substitution}
At a fixed regulator and a regular internal momentum, suppose the
response banks and their common fluctuation representatives depend
$C^2$ on the external momentum and smoothly on the displayed background.
Then substituting $b_L^{\rm pw}$ for $b_L^{\rm harm}$ leaves every
second external-momentum jet unchanged. This applies to all mixed-shell
entries and the complete finite response, not only to diagonal cards.
If the differentiated-integral passage is separately justified, the
extracted second-order coefficient is also unchanged.
\end{corollary}
\begin{proof}
The chain rule through order two uses only the first three jets of the
background, which agree by \cref{thm:v2-jet-replacement}. Apply it to
all entries before assembling the complete traces. Integrating equal
jets gives the last statement under the stated passage hypothesis.
\end{proof}

\section{Two actual Gaussian innovations in an explicit Wilson chart}
\label{sec:v2-finite}

We now populate a finite component of the adjacent-shell target.
Take a side-four macrocell and the continuous Bloch phase
\begin{equation}\label{eq:v2-phases}
 w=\frac{3+4\mathrm i}{5},\qquad
 z=w^4=-\frac{527+336\mathrm i}{625},\qquad p_r=w\mathrm i^r,
 \quad 0\le r<4.
\end{equation}
The internal momentum is in direction two; the fluctuation polarization
is direction three. All other fine momentum coordinates are zero.
This is an exact continuous Bloch-symbol point. It is not claimed to be
a root-of-unity momentum of an ordinary periodic finite torus.

The scalar four-alias covariance, first block, and two-step block are
\begin{align}
 C&=\frac1{256}\diag_r|p_r-1|^{-2},\label{eq:v2-C}\\
 B_{sr}&=1_{\{r\equiv s\ (2)\}}(1+p_r),\quad s=0,1,\notag\\
 W_r&=1+p_r+p_r^2+p_r^3.\notag
\end{align}
In particular $H=C^{-1}=(256/5)\diag(4,18,16,2)$. Define
\begin{align}
 C^{[1]}&=CB^*(BCB^*)^{-1}BC,\quad
 C^{[2]}=CW^*(WCW^*)^{-1}WC,\label{eq:v2-shells}\\
 \Gamma_0&=C-C^{[1]},\qquad \Gamma_1=C^{[1]}-C^{[2]}.\notag
\end{align}
The factors $1/256=16^{-2}$ and $B_{sr}=1+p_r$ are inherited directly
from the two-step tower and its transverse axial restriction. No
independent Gaussian reference has been introduced.

\begin{proposition}[Exact ranges and energy-normalized witnesses]
\label{prop:v2-innovations}
The shell ranks are two and one. Bases of their ranges are
\begin{equation}\label{eq:v2-V}
 V_0=\begin{pmatrix}1&0\\0&1\\-2\mathrm i&0\\0&-\mathrm i/3\end{pmatrix},
 \qquad
 V_1=\begin{pmatrix}
 1\\-51(1+\mathrm i)/2624\\\mathrm i/8\\1377(1-\mathrm i)/2624
 \end{pmatrix}.
\end{equation}
They obey $V_0^*HV_1=0$ and
$\Gamma_j=V_j(V_j^*HV_j)^{-1}V_j^*$. In particular,
\begin{align}
 u_0:=\Gamma_0e_0&=
 (5/17408,\,0,\,-5\mathrm i/8704,\,0)^{\mathsf T},\label{eq:v2-witness}\\
 u_1:=\Gamma_1e_0&=
 (82/22525,\,-3(1+\mathrm i)/42400,\,
          41\mathrm i/90100,\,81(1-\mathrm i)/42400)^{\mathsf T},\notag\\
 u_0^*Hu_0&=5/17408,\qquad u_1^*Hu_1=82/22525.\notag
\end{align}
\end{proposition}
\begin{proof}
Substitute \cref{eq:v2-phases} into \cref{eq:v2-C}. The two rows of $B$
are supported on disjoint alias pairs, so $BCB^*$ is diagonal and
invertible. Also $W$ is the subsequent block of $B$, giving nested
observations. Direct multiplication verifies $BV_0=0$, $WV_1=0$,
$V_1\in\Ran C^{[1]}$, and the displayed energy orthogonality.
Since the first kernel has dimension two and the second innovation
dimension one, these columns span exactly the two ranges. Conditioning
in the $H$ metric gives the stated covariance forms; substitution gives
\cref{eq:v2-witness}. Equivalently these are checked by the rational
identities $\Gamma_jH\Gamma_j=\Gamma_j$ and
$\Gamma_0H\Gamma_1=0$.
\end{proof}

\subsection{Fine Wilson action and the common representative}
Let $E_1,E_2,E_3$ be the imaginary unit quaternions, so that
$E_aE_b=-\delta_{ab}I+\epsilon_{abc}E_c$ and
$\operatorname{Tr}I=2$. The action used in this finite component is
\begin{equation}\label{eq:v2-Wilson-action}
 S_{\rm W}(A)=\frac12\sum_{s,\mu<\nu}
  \left[2-\operatorname{Tr}\left(
 e^{A_\mu(s)}e^{A_\nu(s+e_\mu)}
 e^{-A_\mu(s+e_\nu)}e^{-A_\nu(s)}\right)\right].
\end{equation}
Its one-colour quadratic Hessian is the fine curl Gram, exactly the
normalization underlying \cref{eq:v2-C}. Multiplying this entire action
by $\alpha>0$ and its Gaussian covariance by $\alpha^{-1}$ leaves the
whitened amplitude-vertex banks unchanged. A separate background or
colour convention must still be matched explicitly.

At zero external momentum, the background is
$b_2=E_3\,1_{\{s_2=3\}}$, with other components zero. This is the
$L=4$, $q=0$ value of \cref{eq:v2-treefield}. For an alias vector $u$ put
$h_u(x)=\sum_rp_r^xu_r$. Its common ordered-tree representative is
\begin{align}
 (Tu)_3(s)&=4h_u(s_2)1_{\{s_3=3\}},\label{eq:v2-T}\\
 (Tu)_2(s)&=-s_3\{h_u(s_2+1)-h_u(s_2)\}
       +\frac38\left(\sum_{x=0}^3h_u(x)\right)(z-1)1_{\{s_2=3\}},\notag
\end{align}
with other components zero. For both innovations the final sum is zero,
so the residual term vanishes. The full expression is retained when
constructing a matrix on all four aliases. It is a fine gauge change and
preserves the zero-background curl Gram $H$.

This declares a common representative for the Gaussian innovations and
reconstructs the Wilson vertices in that same representative. It is not
an operation that changes only the covariance and leaves an incompatible
vertex bank fixed. Nonlinear changes of representative or nonlinear
compact-block completion would have additional vertices, which are
outside the component presently evaluated.

\section{An exact nonzero mixed bubble and a two-step determinant certificate}
\label{sec:v2-response}

Use two charged fluctuation colours: $Tu\,E_1$ and $Tv\,E_2$. The
Hermitian first and second background-amplitude Hessian derivatives
on their eight-dimensional complexified alias carrier are
\begin{equation}\label{eq:v2-RQ}
 R=\begin{pmatrix}0&K\\K^*&0\end{pmatrix},\qquad
 Q=\begin{pmatrix}J&0\\0&J\end{pmatrix},
\end{equation}
where
\begin{align}
 K&=\frac{128}{5}
 \begin{pmatrix}
 8\mathrm i&7+7\mathrm i&6&-1+\mathrm i\\
 -7+7\mathrm i&6\mathrm i&-1-\mathrm i&-8\\
 -6&1-\mathrm i&-8\mathrm i&-7-7\mathrm i\\
 1+\mathrm i&8&7-7\mathrm i&-6\mathrm i
 \end{pmatrix},\label{eq:v2-K}\\
 J&=256\begin{pmatrix}
 6&1-\mathrm i&4\mathrm i&5+5\mathrm i\\
 1+\mathrm i&-1&-2+2\mathrm i&3\mathrm i\\
 -4\mathrm i&-2-2\mathrm i&0&2-2\mathrm i\\
 5-5\mathrm i&-3\mathrm i&2+2\mathrm i&7
 \end{pmatrix}.\label{eq:v2-J}
\end{align}
Here $K^*=-K$ and $J^*=J$.

\begin{lemma}[Reconstruction of the displayed Wilson vertices]
\label{lem:v2-vertices}
Equations \cref{eq:v2-RQ,eq:v2-K,eq:v2-J} are the derivatives
$D^3S_{\rm W}[b,\cdot,\cdot]$ and
$D^4S_{\rm W}[b,b,\cdot,\cdot]$ in the declared chart and representative.
\end{lemma}
\begin{proof}
For a plaquette, write its four oriented Lie inputs in their order as
$X_0,X_1,X_2,X_3$. The cubic part of the action is
\[
 -\frac14\operatorname{Tr}\left[
   \left(\sum_iX_i\right)\sum_{i<j}[X_i,X_j]\right].
\]
This follows by multiplying the four exponentials through degree three,
or by using the first two terms of their product logarithm; for
$\mathfrak{su}(2)$ the trace of the cubic power of a single Lie element
is zero. If the three direction fields have oriented scalar entries
$b_i,a_i,c_i$ and colours $E_3,E_1,E_2$, respectively, their trilinear
coefficient is
\begin{equation}\label{eq:v2-cubic-local}
 b_F\sum_{i<j}(a_ic_j-c_ia_j)
 -a_F\sum_{i<j}(b_ic_j-c_ib_j)
 +c_F\sum_{i<j}(b_ia_j-a_ib_j),
 \quad b_F=\sum_i b_i,
\end{equation}
with analogous definitions of $a_F,c_F$.
At zero external momentum the only contributing plaquettes have
orientations $2,3$ and $s_2=3$, where $b=(1,0,-1,0)$.
For their two fluctuation legs, the local third and fourth derivative
matrices are
\begin{equation}\label{eq:v2-local-T}
 T_3=\begin{pmatrix}0&-1&0&1\\1&0&1&2\\0&-1&0&1\\-1&-2&-1&0\end{pmatrix},
 \quad
 T_4=\begin{pmatrix}
 -2/3&-4/3&-2/3&-4/3\\
 -4/3&0&-4/3&-4\\
 -2/3&-4/3&-2/3&-4/3\\
 -4/3&-4&-4/3&0
 \end{pmatrix}.
\end{equation}
The fourth derivative can be obtained without a logarithmic truncation:
expand the ordered product with total degree four. For four labelled
directions, sum over their $24$ orders and all nondecreasing assignments
to the four links, dividing by the factorial of each link's occupancy.
The scalar quaternion product supplies $-\operatorname{Tr}/2$.
For colours $(3,3,1,1)$ and $b=(1,0,-1,0)$ this finite sum is precisely
$T_4$. The same prescription with three directions gives $T_3$ and
\cref{eq:v2-cubic-local}.

Let $P_{x,y}$ map a four-alias vector to the oriented entries of its
$2,3$ plaquette at $(s_2,s_3)=(x,y)$ using \cref{eq:v2-T}. Then
\[
 K=16\sum_{y=0}^3P_{3,y}^*T_3P_{3,y},\qquad
 J=16\sum_{y=0}^3P_{3,y}^*T_4P_{3,y}.
\]
The factor $16$ counts the independent first and fourth coordinates.
Substitution of the four phases in \cref{eq:v2-phases} gives
\cref{eq:v2-K,eq:v2-J}. As a normalization check, the corresponding
quadratic assembly is
$16\sum_{x,y}P_{x,y}^*\boldsymbol 1\boldsymbol 1^{\mathsf T}P_{x,y}=H$.
This derivation retains the order of all noncommuting group factors.
\end{proof}

\begin{theorem}[A populated adjacent-shell mixed response]
\label{thm:v2-populated}
For the two charged copies of each actual innovation,
\begin{align}
 c_{01}&=2\Tr(\Gamma_0K\Gamma_1K^*)
                =\frac{74253125}{51495754}>0,\label{eq:v2-c01}\\
 d_0&=\Tr(\Gamma_0J)-\Tr(\Gamma_0K\Gamma_0K^*)
                =\frac{3648709}{1943236},\notag\\
 d_1&=\Tr(\Gamma_1J)-\Tr(\Gamma_1K\Gamma_1K^*)
                =\frac{22719254375}{1364637481},\notag\\
 \mathcal W''(0)&=d_0+d_1-c_{01}=\frac{191959}{11236}.\label{eq:v2-total}
\end{align}
Moreover a single energy-normalized mixed matrix entry already obeys
\begin{equation}\label{eq:v2-oneentry}
 u_0^*Ku_1=-\frac{25\mathrm i}{61268},\qquad
 \frac{|u_0^*Ku_1|^2}{(u_0^*Hu_0)(u_1^*Hu_1)}
                    =\frac{100000}{627997}>0.
\end{equation}
\end{theorem}
\begin{proof}
Insert \cref{eq:v2-shells} and the explicit matrices
\cref{eq:v2-K,eq:v2-J}; all operations are in $\mathbb Q(\mathrm i)$.
For the covariance $\diag(\Gamma_j,\Gamma_j)$, the two off-diagonal
colour blocks of $R$ contribute the displayed factor two to the mixed
trace. Applying \cref{thm:v1-mixed} gives the last equality in
\cref{eq:v2-total}. Direct multiplication by the two columns
\cref{eq:v2-witness} proves \cref{eq:v2-oneentry}. Thus the nonzero
mixed contraction is established by fully specified finite matrices,
not inferred from generic noncommutativity or from a floating-point fit.
\end{proof}

\subsection{Independent determinant assembly}
Let $T$ be the $8$-by-$6$ inclusion whose first four columns are $V_0$
in each of the two colour copies, and whose last two columns are $V_1$
in those copies. Put
\[
 D=T^*\diag(H,H)T,\quad R_T=T^*RT,\quad Q_T=T^*QT,
 \quad K_T(t)=I+tD^{-1}R_T+\tfrac12t^2D^{-1}Q_T.
\]
The omitted $O(t^3)$ terms of the actual Wilson Hessian cannot change
this determinant's second jet. Although the left-normalized matrix
$K_T$ need not be Hermitian, its determinant is the ratio of the positive
Gaussian precision determinant to $\det D$ near zero.

Direct expansion by permutations, rather than by the trace response
formula, gives
\begin{align}
 \det K_T(t)&=1+\frac{191959}{11236}t^2+O(t^3),\label{eq:v2-detcheck}\\
 \det(K_T(t)_{<4,<4})&=1+\frac{3648709}{1943236}t^2+O(t^3).\notag
\end{align}
The inherited second jet of the last two coordinates gives the remaining
owned response
\[
 \mathcal W_1''(0)=d_1-c_{01}=\frac{41503093125}{2729274962}.
\]
The sum agrees with the direct six-dimensional normalizer. The tail
$C^{[2]}$ has not been integrated and its determinant is not counted.

These numbers use $\frac12\log\det$ on one oriented complexified Bloch
fibre. A conjugate fibre is restored with its own multiplicity in a real
Fourier law. Neither that multiplicity nor the monograph's separate
$1/2,-1/3$ response convention is inserted silently. The active axial
polarization is a subspace of the full fine-field problem; its compressed
Gaussian normalizer is not the whole Wilson normalizer. Its nonzero
mixed matrix element, however, is a genuine entry of the corresponding
fine-Wilson innovation bank in the declared common representative.

\section{True external-momentum jets of one normalized mixed entry}
\label{sec:v2-momentum}

The preceding test can be continued off the external zero rather than
freezing the zero-momentum shell. The background momentum is now $qe_1$,
the left internal momentum is shifted by $qe_1$, and the right internal
momentum remains that of \cref{eq:v2-phases}. The left fluctuation is in
the first innovation and the right fluctuation in the second. The
background has coarse polarization $E_3e_2$; the fluctuations have
polarization $e_3$ and colours $E_1,E_2$.

For $(r,s)\in\{0,2\}^2$, set
\begin{align}
 u_r(q)&=\mathrm i^re^{\mathrm i q/4},\qquad v_s=w\mathrm i^s,\notag\\
 c_{rs}(q)&=\frac1{256\{4-u_r-u_r^{-1}-v_s-v_s^{-1}\}},\notag\\
 b_{rs}(q)&=\tfrac12(1+u_r)(1+v_s),\qquad
 m(q)=\sum_{r,s}|b_{rs}(q)|^2c_{rs}(q),\label{eq:v2-movingcolumn}\\
 a_{rs}(q)&=1_{\{(r,s)=(0,0)\}}c_{00}(q)
     -\frac{c_{rs}(q)\overline{b_{rs}(q)}b_{00}(q)c_{00}(q)}{m(q)}.
 \notag
\end{align}
Conjugation here and in the jets is at real $q$. All denominators are
nonzero in a neighbourhood of zero at the selected internal phase.
This is exactly the $e_{00}$ column of the first innovation covariance,
not a freely chosen continuation of its value at zero. Its support has
four aliases because the first block preserves parity in each of the
two momentum coordinates. The other transverse momentum directions
remain zero; their alias weights decouple in this restriction.

The right column is the four-component $u_1$ in \cref{eq:v2-witness}.
Let
\[
 h_a(x,y;q)=\sum_{r,s\in\{0,2\}}a_{rs}(q)u_r(q)^xv_s^y,
 \qquad h_c(x,y)=\sum_{s=0}^3(u_1)_s p_s^y.
\]
Both columns have zero final block, so their scalar tree potentials
$\phi(s)=s_3h(s_1,s_2)$ have zero block average. Their common tree fields
are therefore
\begin{equation}\label{eq:v2-movingtree}
 A_3=4h(s_1,s_2)1_{\{s_3=3\}},\qquad
 A_1=-s_3\Delta_1h,\qquad A_2=-s_3\Delta_2h,
\end{equation}
with the appropriate Bloch continuations. The background is the
$L=4$ harmonic field in \cref{eq:v2-treefield}; it may be replaced
through order two by \cref{eq:v2-plane} and the \emph{same} tree operation.

Define the Hermitian polarized cubic entry and its normalized square by
\begin{equation}\label{eq:v2-rbeta}
 r(q)=D^3S_{\rm W}\bigl[b_4(q),\overline{A_a(q)}E_1,A_cE_2\bigr],
 \qquad
 \beta(q)=\frac{|r(q)|^2}{\nu_0(q)\nu_1},\quad
 \nu_0(q)=a_{00}(q),\quad \nu_1=\frac{82}{22525}.
\end{equation}
The bar on the left leg is the usual Hermitian Bloch polarization. The
fields in the trilinear form carry the stated matching momenta.
Since $\Gamma H\Gamma=\Gamma$, $\nu_0(q)$ and $\nu_1$ are precisely the
squared energy norms of the unnormalized innovation columns. They are
positive near zero. Thus $\beta$ is a squared matrix entry between unit
vectors in the Gaussian innovation spaces, not a raw unnormalized
coordinate contraction.

\begin{theorem}[Exact adjacent-shell entry through its second external jet]
\label{thm:v2-momentum}
The fine-Wilson component just defined has Taylor expansions
\begin{align}
 r(q)={}&-\frac{25\mathrm i}{61268}-\frac{75}{490144}q
 +\left(\frac{205}{5881728}
       +\frac{208693\mathrm i}{3999575040}\right)q^2+O(q^3),
                                                   \label{eq:v2-rjet}\\
 \nu_0(q)={}&\frac5{17408}+\frac{1115}{170459136}q^2+O(q^3),\notag\\
 \beta(q)={}&\frac{100000}{627997}
                   -\frac{8439475}{384334164}q^2+O(q^3).
                                                   \label{eq:v2-betajet}
\end{align}
In particular, the true derivatives are
\begin{equation}\label{eq:v2-beta-true}
 \boxed{\displaystyle
 \beta(0)=\frac{100000}{627997},\qquad
 \beta'(0)=0,\qquad
 \beta''(0)=-\frac{8439475}{192167082}.}
\end{equation}
\end{theorem}
\begin{proof}
All required algebra can be performed in
$\mathbb Q(\mathrm i)[q]/(q^3)$, where complex conjugation fixes real $q$.
Expansion of \cref{eq:v2-movingcolumn} gives, in the order
$(0,0),(0,2),(2,0),(2,2)$,
\begin{equation}\label{eq:v2-ajet}
 \begin{pmatrix}a_{00}\\a_{02}\\a_{20}\\a_{22}\end{pmatrix}
 =\begin{pmatrix}5/17408\\-5\mathrm i/8704\\0\\0\end{pmatrix}
 +q\begin{pmatrix}0\\0\\-5\mathrm i/52224\\5/156672\end{pmatrix}
 +q^2\begin{pmatrix}1115/170459136\\1265\mathrm i/85229568\\0\\0\end{pmatrix}
 \pmod{q^3}.
\end{equation}
These coefficients obey the exact differentiated first-block constraint
$\sum b_{rs}a_{rs}=0$ in that ring. In particular the newly populated
$r=2$ aliases may not be frozen at their zero values.

Substitute \cref{eq:v2-ajet}, the four right coefficients in
\cref{eq:v2-witness}, and the background jets
\cref{eq:v2-treejets} into the finite fields
\cref{eq:v2-movingtree}. For each $s_1,s_2,s_3\in\{0,1,2,3\}$ and
$\mu<\nu$ in $\{1,2,3\}$ form the ordered plaquette entries
\[
 (F_\mu(s),\ F_\nu(s+e_\mu),\ -F_\mu(s+e_\nu),\ -F_\nu(s)).
\]
Insert the background, conjugate left field, and right field into
\cref{eq:v2-cubic-local}, then sum. The fourth coordinate gives a factor
four. Planes involving direction four vanish identically, since the
fields are independent of that coordinate and their fourth components
are zero. This explicit $64\cdot3$ plaquette sum gives the first line
of \cref{eq:v2-rjet}; its constant agrees independently with
$u_0^*Ku_1$ in \cref{eq:v2-oneentry}. The first entry of
\cref{eq:v2-ajet} gives $\nu_0(q)$.

Finally multiply $r(q)\overline{r(q)}$ and divide by $\nu_0(q)\nu_1$ in the
same truncated ring, whose constant denominator is positive. This gives
\cref{eq:v2-betajet}. The coefficient of $q^2$ is half the true second
derivative, yielding \cref{eq:v2-beta-true}. Each input, finite summation
range, orientation, conjugation, and rational coefficient has been
specified; the accompanying checker executes these identities exactly.
No limit inferred from a finite-difference fit enters the proof.
\end{proof}

\begin{scope}[A genuine populated entry, not the full coefficient]
The result supplies a nonzero adjacent-shell Wilson component and its
actual external second jet, with the harmonic background jet, moving
left covariance, normalization, and common tree all differentiated.
It is not a full mixed-shell Hilbert--Schmidt trace over every
polarization and alias, not a full-torus integral, and not a coefficient
of the complete compact pushed action. A rank-one entry can lower-bound
a squared norm at a fixed momentum, but its second derivative does not
lower-bound the second derivative of that norm. Here $\beta\ge0$ near
zero while $\beta''(0)<0$, explicitly illustrating the sign distinction
in V1. Neither sign is a mass-gap test.
\end{scope}

\section{Verification, remaining acquisition, and version record}
\label{sec:v2-verification}

The bundle contains a checker for the symbolic all-level identities,
the two finite Wilson certificates, and the direct determinant
comparison. Its rational vertex generator uses ordered exponential
products independently of the simplified cubic formula. Its direct
six-dimensional determinant is expanded by permutations. Its momentum
calculation works with ordinary Taylor coefficients in a truncated
Gaussian-rational polynomial ring and converts the final $q^2$
coefficient to the true second derivative explicitly.

The tests include the every-level dyadic covariance recursion; the
symbolic harmonic Euler equations, constraint and boundary conditions;
the $L=2$ regression; the harmonic and plane-wave jet equality; the
local quaternion vertices; the actual shell ranges and energy norms;
the full active-sector mixed response and determinant; and the moving
normalized mixed entry through order two. The inherited V1 checker is
also retained and rerun. These are exact finite algebraic checks,
not a Lean formalization and not a certification of unprovided scripts
mentioned in the older manuscripts. The new general analytic bounds
rest on the proofs above, not on sampling finitely many levels.

\paragraph{The next concrete acquisition.}
The harmonic background on the axial external germ is now available
without a new minimizing solve at every level. The next calculation
should use \cref{eq:v2-treejets} to assemble \emph{all} the adjacent-shell
Wilson entries in one common representative, with the complete internal
momentum dependence. The populated entry is a regression target for
that expansion. The nonlinear compact-block tangent and all principal
Haar and quotient contributions then have to be included or matched by
proved transport identities before comparing with the monograph's
normalization. The internal momentum majorants and shell-separation
estimates needed by \cref{thm:v1-profile} remain to be proved for that
complete bank. These specific tasks are not replaced by the all-level
external analytic bound.

\paragraph{Physical endpoint.}
The finite acquisitions do not supply interacting continuum existence,
noncollapsed local probes, or a uniform contraction at fixed physical
time. Those remain the separate requirements in the supplied closure
monograph. The hidden-diffusion obstruction and the precisely scoped
V16 visible-repair theorem remain unchanged. No physical Hamiltonian
mass gap is claimed by this version.

\paragraph{Append-only record.}
All V1 mathematical text, labels, and its bibliography are retained
above without modification. Only the preamble's version display is
updated. The new body begins at \cref{sec:v2-context}. The provenance
file gives the SHA--256 hashes of the original V1 source, its preserved
body, V2, and the executable checks. A later requested V3 is to append
after the V2 end marker, preserving both previous mathematical bodies.

\begin{thebibliography}{99}
\bibitem[H10]{HammantV2}
T.~C. Hammant, R.~R. Horgan, C.~J. Monahan, A.~G. Hart, E.~H. M\"uller,
A.~Gray, K.~Sivalingham and G.~M. von Hippel,
\emph{Improved automated lattice perturbation theory in background
field gauge}, arXiv:1011.2696 (2010).
\url{https://arxiv.org/abs/1011.2696}.
External methodological context only; no unverified vertex data are
imported from it.
\bibitem[LW95]{LWV2}
M.~L\"uscher and P.~Weisz,
\emph{Background field technique and renormalization in lattice gauge
theory}, Nuclear Physics B \textbf{452} (1995), 213--233,
arXiv:hep-lat/9504006.
\url{https://arxiv.org/abs/hep-lat/9504006}.
Cited for background-field context, not as a proof of the present
nonperturbative construction or mass gap.
\end{thebibliography}

% END OF VERSION 2 MATHEMATICAL BODY.
% Append subsequent requested versions below this marker; retain V1 and V2 above.

\clearpage
\section{V3: complete finite alias traces and removal of their internal soft pole}
\label{sec:v3-context}

The next acquisition in \cref{sec:v2-verification} is a complete bank,
not another normalized matrix entry. This version constructs all entries
on a declared sixteen-dimensional alias carrier, differentiates both
ordered adjacent-shell bubbles, and treats the internal zero-momentum
chart without inverting its soft fine eigenvalue. The resulting reduced
symbol is real analytic over its entire two-dimensional internal momentum
torus. This supplies the differentiated-integral passage for that symbol
and an exact reciprocity relation between its two orderings.

The distinction between the carrier and the full Yang--Mills problem is
essential. We retain all four aliases in each of directions one and two,
set the fine momenta in directions three and four to zero, and use the
transverse polarization $e_3$. This is an invariant covariance subspace
of the supplied Gaussian tower. The Wilson vertex can connect it to
other polarization and alias subspaces, which are not included in the
present trace. Thus ``complete'' below means complete on this specified
sixteen-dimensional carrier, not the unrestricted four-dimensional bank.
The two-dimensional integral introduced below is an integral of a reduced
symbol, not the result of integrating out the missing momentum components.

\paragraph{What is inherited and what is acquired.}
The action, common tree, all-level axial harmonic second jets, and one
normalized moving entry are inherited from V2. The mixed-shell ownership
rule is inherited from V1. Neither is reproved as a new acquisition.
The new finite data comprise the full $16$-by-$16$ cubic bank through
second external order, two rank-$12$ and rank-$3$ innovation covariances,
their complete trace derivatives at a regular internal point, and their
regularized trace derivatives at the internal soft point. The analytic
argument identifies why this last computation is legitimate for the
conditional innovations, rather than an arbitrary prescription for an
infinite free covariance. The scalar trace derivative at the regular
point is positive although V2's normalized entry has negative curvature.
This concretely resolves the insufficiency of extrapolating that entry.

The construction remains the fine-link Wilson component of
\cref{eq:v2-Wilson-action}, with its stated exponential chart. Nonlinear
compact-block tangent terms, Haar and quotient terms, other principal
vertices, and the calibration to the monograph's extracted coefficient
are separate. Automated background-field vertex differentiation is an
established method \cite{HammantV2}; the present contribution is the
explicit calculation and its conditional-Gaussian analytic control for
this carrier, not a general novelty claim about that method.

\section{Completing a column calculation without inverting column norms}
\label{sec:v3-frame}

The normalized witness of V2 is convenient at a regular internal point.
It is not a suitable coordinate system at every internal momentum: an
individual covariance column can vanish while the innovation subspace
has unchanged rank. For the following calculation it is better to retain
all covariance columns with their exact energy weights.

\begin{proposition}[Energy-weighted complete-column formula]
\label{prop:v3-frame}
Let $G_L,G_R\succeq0$ be finite Hermitian matrices and let
$H_L=\diag(h_{L,a})$, $H_R=\diag(h_{R,b})\succeq0$ satisfy
\[
 G_LH_LG_L=G_L,\qquad G_RH_RG_R=G_R.
\]
For an arbitrary compatible complex matrix $K$, put $M=G_LKG_R$. Then
\begin{equation}\label{eq:v3-frame}
 \boxed{\displaystyle
 \Tr(G_LKG_RK^*)
   =\sum_{a,b}h_{L,a}h_{R,b}|M_{ab}|^2.}
\end{equation}
Neither an inverse of $G_L,G_R$ nor a nonzero norm for each of their
columns is required. If these matrices are $C^2$ in a real parameter,
the identity may be differentiated twice as a finite identity, including
all derivatives of the energy weights.
\end{proposition}
\begin{proof}
The identity $GHG=G$ gives
\[
 G=\sum_a h_a(Ge_a)(Ge_a)^*.
\]
Insert this decomposition for each covariance into the trace and expand.
Each term is
$h_{L,a}h_{R,b}|e_a^*G_LKG_Re_b|^2$, proving
\cref{eq:v3-frame}. A zero energy weight or a zero covariance column
contributes zero, not an undefined normalized vector. Differentiation
is the ordinary product rule for a finite sum. Square roots of the
weights are unnecessary, even when a weight vanishes.
\end{proof}

This is an implementation of the complete Hilbert--Schmidt trace,
not a sum of separately normalized witness squares. Omitting energy
weights or freezing their derivatives would describe another quantity.
The independent verification below compares \cref{eq:v3-frame} with
direct matrix multiplication for both orderings and at the soft point.

\section{The sixteen-alias covariance and its full cubic bank}
\label{sec:v3-bank}

Let $k=(k_1,k_2)\in\T^2$ be the internal coarse momentum and
$q\in\R$ the external coarse momentum in direction one. Alias indices
are $\alpha=(r,s)\in\{0,1,2,3\}^2$, in lexicographic order. Set
\begin{align}
 u_r(k,q)&=\exp\!\left(\frac{\mathrm i(k_1+q+2\pi r)}4\right),\notag\\
 v_s(k)&=\exp\!\left(\frac{\mathrm i(k_2+2\pi s)}4\right),\notag\\
 \lambda_{rs}(k,q)&=4-u_r-u_r^{-1}-v_s-v_s^{-1},\label{eq:v3-phases}\\
 C_{rs}(k,q)&=\frac1{256\lambda_{rs}(k,q)},\qquad
 H_{rs}(k,q)=256\lambda_{rs}(k,q).\notag
\end{align}
The symbols $C,H$ below denote the diagonal matrices with these entries.
Initially the formula for $C$ is used where the fine eigenvalues are
nonzero. The innovations will be extended independently across their
soft point in \cref{sec:v3-soft}.

Write $D_4(z)=1+z+z^2+z^3$ and define the first and final observations
\begin{align}
 B_{ab,rs}(k,q)&=1_{\{r\equiv a,\ s\equiv b\ (2)\}}
                 \frac{(1+u_r)(1+v_s)}2,
                    &&a,b\in\{0,1\},\label{eq:v3-observations}\\
 W_{rs}(k,q)&=\frac{D_4(u_r)D_4(v_s)}4.\notag
\end{align}
Then
\begin{align}
 C^{[1]}&=CB^*(BCB^*)^{-1}BC,\qquad
 C^{[2]}=CW^*(WCW^*)^{-1}WC,\notag\\
 G_0&=C-C^{[1]},\qquad G_1=C^{[1]}-C^{[2]}.
                                                   \label{eq:v3-shells}
\end{align}
The factor $1/256$ is the finest covariance normalization in a two-step
alias class. The factor $1/2$ in $B$ follows by restricting the supplied
four-dimensional block symbol to $p_3=p_4=0$ and polarization $e_3$.
In particular, $\Pi(p)e_3=e_3$ throughout this carrier; no transverse
projection has been replaced by the identity on another polarization.
The final observation $W$ is the composition of the two blocks, since
$(1+z)(1+z^2)=D_4(z)$ in each active coordinate.

\begin{proposition}[Actual innovation ranks and constraints]
\label{prop:v3-shells}
Away from the soft fine point, \cref{eq:v3-shells} are precisely the first
two innovations on this carrier. Their ranks are $12$ and $3$, and
\begin{equation}\label{eq:v3-constraints}
 \begin{gathered}
 BG_0=0,\qquad WG_0=WG_1=0,\\
 G_iHG_i=G_i\quad(i=0,1),\qquad G_0HG_1=0.
 \end{gathered}
\end{equation}
They are positive semidefinite and Hermitian. The same statements extend
to the soft point by the construction in \cref{sec:v3-soft}.
\end{proposition}
\begin{proof}
The four rows of $B$ have disjoint supports and each is nonzero. The
nonzero row $W$ lies in their span by the composition identity. Therefore
the ranks of the first and final observation spaces are four and one.
Whitening the positive diagonal $C$ identifies their explained
covariances with nested orthogonal projections, as in
\cref{lem:v1-projectors}. Their differences have ranks $16-4=12$ and
$4-1=3$. Unwhitening gives positivity and
\cref{eq:v3-constraints}. The extension assertion follows from the
explicit bounded formulas below; it is not an evaluation of $C$ at a
zero eigenvalue.
\end{proof}

\subsection{A representative with no internal inverse}
For an alias coefficient vector $a$ define
\[
 h_a(x,y;k,q)=\sum_{r,s}a_{rs}u_r(k,q)^xv_s(k)^y.
\]
Use Bloch continuation in $x,y$ and periodic continuation in the third
coordinate $t$. On a macrocell put
\begin{equation}\label{eq:v3-tree}
 \begin{aligned}
 (T_ka)_1&=-t\Delta_1h_a,&
 (T_ka)_2&=-t\Delta_2h_a,\\
 (T_ka)_3&=4h_a\,1_{\{t=3\}},&
 (T_ka)_4&=0,
 \end{aligned}
 \qquad 0\le t<4.
\end{equation}
The notation $T_k$ includes the shift $k_1+q$ when used on the left leg.
It is obtained by subtracting the gradient of $\phi=t h_a$ from the
transverse fine field. The average of this potential is
\[
 \frac1{4^4}\sum_{x,y,t,x_4}t h_a(x,y)
      =\frac3{32}\sum_{x,y}h_a(x,y)=\frac38 Wa.
\]
Consequently, on either innovation range $Wa=0$, the residual cell-root
gauge correction vanishes identically. Formula \cref{eq:v3-tree} is
therefore the declared common tree representative on both ranges, with
zero coarse longitudinal component. Its extension to arbitrary alias
vectors is used only as a matrix carrier; contractions with $G_0,G_1$
remove those vectors outside the required ranges. No inverse of
$\lambda$ or of a coarse gradient appears in this representative.

The external background is $b_4(q)$ from \cref{eq:v2-treefield}, with
colour $E_3$ and transverse coarse polarization $e_2$. For second jets it
may equivalently be generated by
\[
 h_b(x;q)=\frac{e^{\mathrm i qx/4}}
                {\sum_{j=0}^3e^{\mathrm i qj/4}}
\]
\emph{before the same tree operation}. The code evaluates both the exact
rational harmonic profile and this replacement and compares all bank jets.

\subsection{A complete finite summation prescription}
Let $P_{\mu\nu}(k,q;x,y,t)$ be the $4$-by-$16$ matrix taking $a$ to the
four oriented plaquette entries of \cref{eq:v3-tree}, in the order
\[
 (F_\mu(s),F_\nu(s+e_\mu),-F_\mu(s+e_\nu),-F_\nu(s)).
\]
Let $b=(b_0,b_1,b_2,b_3)$ denote the corresponding entries of the external
background on this plaquette. Define the fixed antisymmetric matrix
$A_{ij}=1_{\{j>i\}}-1_{\{j<i\}}$, and
\begin{equation}\label{eq:v3-local}
 \ell_j(b)=\sum_i b_iA_{ij},\qquad
 T(b)_{ij}=\left(\sum_m b_m\right)A_{ij}-\ell_j(b)+\ell_i(b).
\end{equation}
This matrix is antisymmetric under ordinary transpose. The ordered
Wilson cubic identity \cref{eq:v2-cubic-local} says that its local value
is $a^{\mathsf T}T(b)c$ for fluctuation colours $E_1,E_2$.

The complete one-colour-pair bank on the carrier is
\begin{equation}\label{eq:v3-K}
 \boxed{\displaystyle
 K(k,q)=4\sum_{x,y,t=0}^3\ 
       \sum_{1\le\mu<\nu\le3}
 P_{\mu\nu}(k,q;x,y,t)^*
 T(b_{\mu\nu}(q;x,y,t))
 P_{\mu\nu}(k,0;x,y,t).}
\end{equation}
The factor four counts direction four. Plaquettes involving direction
four contribute zero because all fourth components vanish and the fields
are independent of that coordinate. The bar on the left Fourier leg is
included in the adjoint. All $256$ entries, rather than one chosen
contraction, are specified by this finite formula.

\section{An explicit removal of the internal soft pole}
\label{sec:v3-soft}

The trace is to be constructed from conditional innovations, not from
unrestricted fine covariances considered separately at zero momentum.
Their cancellation can be performed before numerical evaluation or
external differentiation.

\begin{lemma}[Deflated diagonal conditioning]
\label{lem:v3-deflation}
Let $D\succ0$ be a finite diagonal matrix, let $\beta$ be a nonzero scalar,
let $b$ be a row vector, and set
\[
 C(t)=\diag(t^{-1},D),\quad L=(\beta,b),\quad
 s=bDb^*,\quad d=|\beta|^2+t s.
\]
For $t>0$,
\begin{equation}\label{eq:v3-deflation}
 \boxed{\displaystyle
 C-C L^*(LCL^*)^{-1}LC
 =\begin{pmatrix}
 s/d&-\overline\beta\,bD/d\\
 -Db^*\beta/d&D-tDb^*bD/d
 \end{pmatrix}.}
\end{equation}
The right side is regular at $t=0$. There it is the covariance of the
quadratic energy $\diag(0,D^{-1})$ restricted to $Ly=0$, not a covariance
assigned to the unrestricted zero-energy coordinate.
\end{lemma}
\begin{proof}
One has $LCL^*=|\beta|^2/t+s=d/t$. Substitute this scalar inverse into
the left side. Its upper-left entry is
$t^{-1}-|\beta|^2/(td)=s/d$; the other entries give the displayed blocks.
At $t=0$, the constraint writes $y_0=-\beta^{-1}b y_h$. The remaining
energy is $y_h^*D^{-1}y_h$, so the covariance is the congruence of $D$
by $y_h\mapsto(-\beta^{-1}b y_h,y_h)$. This is exactly the displayed
limit. Its denominator is nonzero at zero, proving regularity, also for
analytic parameter-dependent $t,D,\beta,b$ in a sufficiently small chart.
\end{proof}

\begin{theorem}[Real-analytic innovation symbols over the reduced torus]
\label{thm:v3-internal}
The matrices $G_0,G_1$ in \cref{eq:v3-shells}, in alias charts related
by their natural permutations, have real-analytic extensions across
the entire internal torus $\T^2$. These extensions satisfy
\cref{eq:v3-constraints}, remain positive semidefinite, and have ranks
$12$ and $3$. The complete scalar bubbles defined below are real analytic
in $(k,q)$ for $k\in\T^2$ and sufficiently small $|q|$.
In particular, their fixed-order internal and external derivatives are
bounded on compact smaller parameter domains.
\end{theorem}
\begin{proof}
For real phases, $\lambda=|u-1|^2+|v-1|^2$. In a fundamental internal
chart the only fine zero eigenvalue is the alias $(r,s)=(0,0)$ at $k=0$;
other occurrences are its relabelled charts. All other diagonal
covariances are analytic and bounded near this point.

The first innovation $G_0$ is a direct sum of four conditionings, one on
each parity group. On the group containing the soft alias apply
\cref{lem:v3-deflation} with
$t=256\lambda_{00}$ and with the corresponding row of $B$.
Its soft-coordinate coefficient is $\beta=2$ at zero. The other parity
groups have no soft diagonal entry and use their ordinary conditioning
formula. Next form the total eliminated covariance
\[
 \Sigma=C-CW^*(WCW^*)^{-1}WC=G_0+G_1
\]
by the same deflation, now with the row $W$. Its soft-coordinate
coefficient is $\beta=4$ at zero. Therefore both $G_0$ and $\Sigma$
are analytic there, and so is $G_1=\Sigma-G_0$. This construction never
inverts $256\lambda_{00}$ at its zero.

There are no other singular conditioning denominators. Indeed, outside
the soft point $0<\lambda\le8$, hence $C_{rs}\ge1/2048$. On each parity
group,
\[
 \sum |(1+u_r)(1+v_s)/2|^2=4,
\]
so its scalar $BCB^*$ entry is at least $1/512$. The geometric-sum
identity $\sum_{r=0}^3|D_4(\xi\mathrm i^r)|^2=16$ for $|\xi|=1$
also gives $\sum_{r,s}|W_{rs}|^2=16$ and $WCW^*\ge1/128$ away from the
soft point. The preceding deflated charts replace these infinite
unrestricted expressions at the soft point itself.

Positivity and all polynomial identities in \cref{eq:v3-constraints}
pass to the finite limits. Away from zero,
$\Tr(HG_0)=12$ and $\Tr(HG_1)=3$, so these identities persist at zero.
For a positive semidefinite $G$ with $GHG=G$, the matrix
$G^{1/2}HG^{1/2}$ is the identity on $\Ran G$ and zero on its orthogonal
complement. Thus $\operatorname{rank}G=\Tr(HG)$ even when $H$ has a null direction.
The stated ranks persist.

Changing $k_a$ by $2\pi$ cyclically permutes the corresponding fine
aliases and permutes first-block rows. These are compatible changes of
chart for both $G_i$ and the vertex. The fine representative
\cref{eq:v3-tree} is polynomial in its phases and has no internal inverse.
The external harmonic profile is analytic by \cref{thm:v2-uniform}.
Consequently \cref{eq:v3-K} and all contracted scalar traces are analytic
on the described charts; their scalar values are periodic. Compactness
of $\T^2$ and a finite chart cover give the derivative bounds on a smaller
external interval. Equivalently, each chart's rational real-analytic
expressions extend to a small complex neighbourhood before the finite
cover is taken.
\end{proof}

\begin{scope}[What has, and has not, lost its soft pole]
The pole is absent from the \emph{two eliminated innovations} at this
fixed macrocell size. The unrestricted fine covariance and the retained
coarse sector are not assigned a massive limit. In particular the
massless Gaussian fixed point in the monograph is unchanged. The theorem
proves an actual internal-momentum regularity input for this reduced,
fixed-level bank; it is not a uniform estimate as the number of levels,
the full polarization space, or all four internal momenta are restored.
\end{scope}

\section{Complete trace jets at a regular point and at the soft point}
\label{sec:v3-traces}

For the matrices constructed above, use the ordered scalar traces
\begin{equation}\label{eq:v3-ordered}
 \begin{aligned}
 F_{01}(k,q)&=\Tr\bigl[G_0(k+qe_1)K(k,q)G_1(k)K(k,q)^*\bigr],\\
 F_{10}(k,q)&=\Tr\bigl[G_1(k+qe_1)K(k,q)G_0(k)K(k,q)^*\bigr],\\
 \mathcal C_{\rm mix}(k,q)&=F_{01}(k,q)+F_{10}(k,q).
 \end{aligned}
\end{equation}
Here $G_i(k+qe_1)$ means the left covariance in
\cref{eq:v3-phases,eq:v3-shells}; the shift is not applied twice.
Each ordered trace is a complete squared Hilbert--Schmidt norm on its
two innovation spaces. For the two charged colours, each ordered
colour trace is $2F_{ij}$. Thus at $q=0$ the Gaussian amplitude mixed debit
is $\mathcal C_{\rm mix}=2F_{01}$, in agreement with V2's normalization.
No extra monograph colour or extraction factor is included in
\cref{eq:v3-ordered}.

\begin{theorem}[Full sixteen-alias jets at the V2 internal phase]
\label{thm:v3-regular-card}
Let $k_*= (0,k_{*,2})$ have $e^{\mathrm i k_{*,2}/4}=(3+4\mathrm i)/5$,
with the same continuous Bloch convention as V2. Put
\begin{align}
 f_*&=\frac{41170821816340625}{47371097767929936},\notag\\
 a_{01}&=\frac{12761661882010890169796242633}
                   {106703436119379187613288924160},\notag\\
 a_{10}&=\frac{1683014474524685332524536681283856723}
                   {28529098119792187335933429769896000960}.
                                                     \label{eq:v3-cards}
\end{align}
The exact Taylor expansions are
\begin{equation}\label{eq:v3-Fjets}
 F_{01}(k_*,q)=f_*+a_{01}q^2+O(q^3),\qquad
 F_{10}(k_*,q)=f_*+a_{10}q^2+O(q^3).
\end{equation}
In particular,
\begin{equation}\label{eq:v3-pair-card}
 \boxed{\begin{aligned}
 \mathcal C_{\rm mix}(k_*,0)
    &=\frac{41170821816340625}{23685548883964968},\\
 \partial_q\mathcal C_{\rm mix}(k_*,0)&=0,\\
 \partial_q^2\mathcal C_{\rm mix}(k_*,0)
    &=\frac{81521218246908525891635437353359304511}
            {228232784958337498687467438159168007680}>0.
 \end{aligned}}
\end{equation}
The last value is approximately $0.3571845222$.
\end{theorem}
\begin{proof}[Exact finite calculation]
Work in $\mathbb Q(\mathrm i)[q]/(q^3)$ with ordinary Taylor coefficients,
complex conjugation fixing real $q$. In \cref{eq:v3-phases} set
$u_r=\mathrm i^r(1+\mathrm i q/4-q^2/32)$ and
$v_s=((3+4\mathrm i)/5)\mathrm i^s$. Form $C,B,W$ using
\cref{eq:v3-observations}; the four first-block denominators are scalar
because their supports are disjoint. Form the two matrices in
\cref{eq:v3-shells} through degree two. Compute every entry of $K$
by \cref{eq:v3-local,eq:v3-K}, using either the exact rational harmonic
profile or its proved second-jet replacement.

For a scalar or matrix series $X(q)=\sum_{j=0}^2X_jq^j$, all products use
$[XY]_n=\sum_{j=0}^nX_jY_{n-j}$. The only inverses needed in these
covariance formulas are scalar inverses with nonzero constant term;
their recursion is
\[
 y_0=x_0^{-1},\qquad
 y_n=-x_0^{-1}\sum_{j=1}^n x_jy_{n-j},\quad n=1,2.
\]
Taking the two finite traces gives constant coefficient $f_*$, zero
linear coefficient, and the quadratic coefficients in
\cref{eq:v3-cards}. These are reproducible rational identities from
explicit phases, covariance rows, and a $4\cdot64\cdot3$ plaquette sum.
The accompanying report contains all covariance and vertex coefficients,
not merely the final scalars. An independent expansion by
\cref{eq:v3-frame} gives exactly the same three coefficients for each
ordering. Adding them and multiplying the quadratic coefficient by two
gives \cref{eq:v3-pair-card}.
\end{proof}

\subsection{Two exact regressions and the sign change}
The $(0,0)$ covariance columns reproduce the complete V2 moving entry
\cref{eq:v2-rjet} and its normalization \cref{eq:v2-betajet}. At $q=0$,
restricting both covariance legs to the old $r=0$ alias carrier gives
\[
 2\Tr\bigl[G_0^{r=0}K G_1^{r=0}K^*\bigr]
                 =\frac{74253125}{51495754},
\]
exactly the full axial subblock response of V2. The larger carrier has
additional entries, rather than a change of that subblock's normalization.

For the same regular internal point, the following true second derivatives
have different signs:
\begin{equation}\label{eq:v3-sign}
 \beta''(0)=-\frac{8439475}{192167082}<0,
 \qquad
 \partial_q^2F_{01}(k_*,0)=2a_{01}>0.
\end{equation}
The first is V2's energy-normalized single entry; the second sums the
entire declared innovation bank. This is a populated example of why an
entry's second derivative cannot replace the trace derivative. Positivity
of a squared norm does not prescribe either derivative's sign.

\subsection{Where the derivative is spent}
Let $G(q)$ denote the moving left covariance, $D$ the fixed right
covariance, and write primes for true external derivatives at zero.
The derivative of $F=\Tr(GKDK^*)$ is
\begin{equation}\label{eq:v3-derivative-ledger}
 \begin{aligned}
 F''={}&\Tr(G''KDK^*)
       +4\Rea\Tr(G'K'DK^*)\\
      &+2\Tr(GK'DK'^*)+2\Rea\Tr(GK''DK^*).
 \end{aligned}
\end{equation}
This follows by applying the product rule twice and grouping adjoint
pairs; Hermiticity of $G,G',G''$ and $D$ is used. The exact checker
verifies the four terms separately. At $k_*$ their approximate values
are shown only for readability; the report stores exact fractions.
\begin{center}
\begin{tabular}{@{}lrr@{}}
\toprule
Term in \cref{eq:v3-derivative-ledger}&$F_{01}''$&$F_{10}''$\\
\midrule
Moving covariance, second derivative&$-0.0095847760$&$-0.0566989822$\\
Covariance--vertex cross derivative&$ 0.0005967012$&$ 0.0355189917$\\
Squared first vertex derivative&$ 0.3780163910$&$ 0.3780163910$\\
Second vertex derivative&$-0.1298296109$&$-0.2388505835$\\
\midrule
Total&$0.2391987052$&$0.1179858170$\\
\bottomrule
\end{tabular}
\end{center}
Freezing the moving covariance removes the first two terms and changes
both answers. Likewise, the two orderings cannot be identified pointwise
merely because their constant terms agree.

\begin{theorem}[A finite full-bank card at the internal soft point]
\label{thm:v3-soft-card}
At $k=0$, use the analytic innovation extensions in
\cref{thm:v3-internal}, not an inverse of the zero fine eigenvalue.
Then
\begin{align}
 F_{01}(0,q)&=\frac{19}{48}+\frac{3343}{27648}q^2+O(q^3),\notag\\
 F_{10}(0,q)&=\frac{19}{48}+\frac{1357}{6912}q^2+O(q^3),
                                                   \label{eq:v3-soft-cards}\\
 \boxed{\displaystyle
 \mathcal C_{\rm mix}(0,0)=\frac{19}{24},\qquad
 \partial_q\mathcal C_{\rm mix}(0,0)=0,\qquad
 \partial_q^2\mathcal C_{\rm mix}(0,0)=\frac{8771}{13824}.}
                                                   \label{eq:v3-soft-pair}
\end{align}
\end{theorem}
\begin{proof}
Set $u_r=\mathrm i^r e^{\mathrm i q/4}$ and $v_s=\mathrm i^s$ in the
same truncated ring. The fine diagonal energy entry $H_{00}$ has zero
constant coefficient, so the direct formula for $C_{00}$ is not used.
For each parity group form $G_0$ from \cref{eq:v3-deflation}, choosing
its first index as pivot; on the group containing $(0,0)$ the pivot
coefficient of $B$ is two at zero. Form $\Sigma$ by the same formula
with row $W$, whose soft pivot coefficient is four, and set
$G_1=\Sigma-G_0$. All remaining scalar inverse constant terms are
nonzero. The finite representative and the background require no
internal inverse and are evaluated by the same plaquette sum.

The two traces give \cref{eq:v3-soft-cards}. Their independent weighted
column sums also agree, even though $H_{00}(0)=0$ and the corresponding
column normalization cannot be used to define V2's normalized witness
there. Adding the quadratic coefficients and converting to the true
second derivative gives \cref{eq:v3-soft-pair}. The analytic theorem
identifies these ring computations with derivatives of the actual
conditional-innovation trace, rather than an arbitrary soft-point
assignment.
\end{proof}

\section{Reciprocity, both orderings, and the reduced momentum integral}
\label{sec:v3-reciprocity}

The explicit discrepancy between the two pointwise second derivatives
has a useful exact interpretation. Write $K(k,q)$ for the ordered
colour-$E_1,E_2$ bank with left internal momentum $k+qe_1$ and right
internal momentum $k$.

\begin{proposition}[Reciprocal Wilson bank]
\label{prop:v3-reciprocity}
For real momenta, in compatible alias charts,
\begin{equation}\label{eq:v3-Krecip}
 K(k,q)^*=-K(k+qe_1,-q).
\end{equation}
Consequently,
\begin{equation}\label{eq:v3-Frecip}
 \boxed{F_{10}(k,q)=F_{01}(k+qe_1,-q).}
\end{equation}
These identities extend through the conditional soft chart.
\end{proposition}
\begin{proof}
The background satisfies $b(-q)=\overline{b(q)}$, first for its harmonic
formula and then for its real-coefficient tree operation. The local
matrix \cref{eq:v3-local} is linear in $b$ and antisymmetric under
transpose, so
\[
 T(b(q))^*=-T(b(-q)).
\]
Take the adjoint of each term of \cref{eq:v3-K}. The left and right
fluctuation representatives are exchanged, and the external momentum
is reversed. This gives \cref{eq:v3-Krecip}, with the minus sign arising
from the ordered charged colours. Substituting it into the right side
of \cref{eq:v3-Frecip} and using cyclicity of the trace gives the left
side. Finite analytic extensions preserve the identities at soft points.
The verification independently rebuilds the reverse bank and checks
all $256$ identities through degree two at $k_*$.
\end{proof}

\begin{theorem}[Differentiated reduced-torus identity]
\label{thm:v3-integral}
Let $\dd\bar k=(2\pi)^{-2}\dd k_1\dd k_2$ and define
\[
 I_{ij}(q)=\int_{\T^2}F_{ij}(k,q)\,\dd\bar k,\qquad
 I_{\rm mix}(q)=I_{01}(q)+I_{10}(q).
\]
For the fixed sixteen-alias Wilson component constructed above, these
functions are twice differentiable on a neighbourhood of zero, their
derivatives may be passed through the integral, and
\begin{equation}\label{eq:v3-integral}
 \boxed{\displaystyle
 I_{10}(q)=I_{01}(-q),\qquad
 I_{\rm mix}(q)=I_{01}(q)+I_{01}(-q),\qquad
 I_{\rm mix}''(0)=2\int_{\T^2}\partial_q^2F_{01}(k,0)\,\dd\bar k.}
\end{equation}
No pointwise identity $\partial_q^2F_{10}(k,0)
=\partial_q^2F_{01}(k,0)$ is asserted.
\end{theorem}
\begin{proof}
The derivative bounds from \cref{thm:v3-internal} give an integrable
constant majorant on the compact internal torus. The fundamental theorem
of calculus followed by dominated convergence justifies both derivative
passages. Apply \cref{eq:v3-Frecip} and translate $k_1$ by $q$ in the
periodic integral; alias permutations at a chart boundary do not change
the scalar trace. This proves the first identity, and the other two
follow by addition and differentiation.
\end{proof}

There is also a local total-derivative form of the distinction between
the two orderings. Twice differentiating \cref{eq:v3-Frecip} gives
\begin{equation}\label{eq:v3-total-derivative}
 \partial_q^2F_{10}(k,0)-\partial_q^2F_{01}(k,0)
  =\partial_{k_1}^2F_{01}(k,0)
       -2\partial_{k_1}\partial_qF_{01}(k,0).
\end{equation}
The right side integrates to zero by periodicity and the proved
regularity. Thus a doubled single ordering is valid for this integrated
second derivative, but not for a single sampled internal-momentum card.
This refines rather than contradicts V1's requirement to retain both
orderings before an actual comparison has been proved.

The theorem does not evaluate the displayed integral or identify its
sign. A positive trace curvature at two internal points does not decide
an integral over all internal momenta. Moreover this reduced integral
is not the full four-dimensional principal response; no equality with
$b_{\rm YM}$ follows from \cref{eq:v3-integral}.

\section{Verification, new proof boundary, and append-only record}
\label{sec:v3-verification}

\paragraph{Executed checks.}
The new driver \path{verify_ym_shell_mixing_V3.py} recomputes the covariance
and vertex jets from the phases and the plaquette rule. It does not load
a cached bank. Its report has $19$ passing exact test families, including:
the $64$ independent ordered-quaternion local cubic coefficients; the
symbolic deflation identity; left and right constraints, energy identities
and ranks; direct/deflated covariance agreement; both independent full curl-Gram jets; all $256$ harmonic/plane
bank jet equalities; all $256$ reciprocal bank jet equalities; the two
complete trace jets; both complete-column checks; both four-term true
derivative decompositions; the V2 moving-entry and axial-subblock
regressions; alias-wrap covariance; four soft-point internal directions;
and the soft full-bank and weighted-frame cards. Reported decimal values
are only renderings of stored exact rationals.

The reusable modules are \path{ym_jet_ring.py} and
\path{ym_v3_wilson_bank.py}. The latter provides both direct and deflated
conditional-covariance evaluation, the complete finite plaquette bank,
and its trace and weighted-frame evaluations. The driver stores all
three Taylor coefficients of every covariance and vertex entry at both
certified points. V1 and V2 executables and their fresh reverification
reports are also retained in the bundle. The proofs of full-torus
analyticity and differentiated integration rest on the arguments above,
not on sampling the torus. No Lean formalization or numerical quadrature
certificate is claimed.

\paragraph{The acquisition achieved here.}
The moving scalar witness has been replaced by a complete finite
innovation-bank trace, and the entire internal soft chart of this
reduced symbol is controlled. The differentiated-integral passage is now
proved for that symbol, rather than left as an unspecified premise.
The exact reciprocal identity further identifies the total-derivative
difference between its two pointwise orderings. These are new inputs to
a larger bank assembly, not evaluations of a complete anomaly scalar.

\paragraph{The next substantive target.}
The useful next extension is to restore the other fine momentum aliases
and physical polarizations in the same covariance and vertex pipeline.
The sixteen-alias matrices, the reciprocal identity, and the soft-point
card supply stringent embedded regressions. On the larger carrier, the
scalar deflation argument must be replaced by a proved matrix-valued
conditional cancellation on a common gauge quotient; the present scalar
argument must not simply be declared sufficient. The complete compact
principal terms must then be supplied or compared by exact transport.
Only after those acquisitions would a four-dimensional integral and
cofinal shell-separation estimate address the monograph's remaining
Gaussian matching coefficient. A high-precision integral of this reduced
symbol alone would not remove those obligations.

\paragraph{Physical conclusion.}
No interacting continuum theory or positive physical Hamiltonian gap is
asserted in this version. The original local-repair conclusions and the
separate construction, noncollapse, and physical-time contraction
requirements are unchanged. In particular, removing a conditional
innovation's finite-level soft pole does not remove the retained
massless infrared channel of the Gaussian theory.

\paragraph{Version record.}
The V1 and V2 mathematical bodies, including their bibliographies and
end markers, are retained byte-for-byte. Only version displays in the
preamble are updated. This continuation is appended before the final
\verb|\end{document}|. The provenance report records the original V2
source hash, the preserved-body hash, the V3 source, the rendered PDF,
and the executable verification files. A later requested V4 is to append
after the marker below, without rewriting the existing mathematical body.

% END OF VERSION 3 MATHEMATICAL BODY.
% Append subsequent requested versions below this marker; retain V1--V3 above.

\clearpage
\section{V4: upstream reconstruction instead of further reduced response cards}
\label{sec:v4-context}

This continuation changes the acquisition order. Versions V1--V3 established
useful composition identities and finite Wilson-component regressions, but
increasing the dimension of the same partial bank would still leave its
connection to the complete pushed action unresolved. Moreover, the supplied
archive already contains more of the actual one-step calculation than that
route was using. In particular, f12 V7 constructs the charged Hessian with
its multiplier, first and second constraint tangents, both momentum transfers,
and the measure term, and states its four-dimensional integral
\src{F12}{\texttt{sec:f12v7}, \texttt{prop:f12v7-integrand}}. That construction,
including its unevaluated integral, is an inherited input, not work to repeat.

The upstream task has two parts. First, prove that the entire Gaussian tower
can be eliminated using the \emph{same} full four-dimensional parity fibre at
every level, without separately enlarging a fine-alias regression. Second,
specify a closed, measured jet transformation which transports the complete
classical action and one-loop density, rather than repeatedly substituting
the original fine Wilson vertices. The results below supply both finite
structural statements. The scale-uniform bounds on the transported nonlinear
vertices, and the physical infrared contraction, remain distinct tasks.

\paragraph{The genuinely new spectral input.}
The positive lattice sum, the covariance fixed point, and the one-step
parity-fibre gap are inherited. The new estimate is an explicit comparison
with the free transverse covariance, valid at \emph{every} level and every
four-dimensional momentum. It transfers the fixed-fibre floor to all levels
and gives a uniform inverse bound on a fixed $45$-coordinate one-colour
fibre. A Schur recursion then removes the internal soft-momentum obstruction
for the full finite-level kernel, without assigning a massive covariance to
a retained zero mode.

\paragraph{The genuinely new response implementation.}
The constrained-Hessian correction and measured Jacobian are already part of
the source. Here they are assembled into a bordered determinant and an
executable transformation of fourth action jets, third block jets, and
second log-density jets. Its finite composition law is proved and tested
against independent scalar-fibre calculations. This is a re-expression of
the actual pushforward, not another unmeasured Gaussian reference. Neither
Gaussian elimination nor stationary phase is claimed as a new general
method; lattice Gaussian renormalization and background-field methods have
an extensive prior literature \cite{BalabanV4,LWV2}.

\section{An all-level comparison on the complete four-dimensional tower}
\label{sec:v4-ellipticity}

Retain the covariance and Fourier conventions of the monograph. Put
\[
 q_\mu(k)=e^{\mathrm i k_\mu}-1,\qquad
 \lambda(k)=\sum_{\mu=1}^4|q_\mu(k)|^2,\qquad
 \Pi(k)=I-\frac{q(k)q(k)^*}{\lambda(k)}
\]
for $k\ne0$ modulo $2\pi$. For $L=2^j$, let
\[
 D_L(p)=\sum_{r=0}^{L-1}e^{\mathrm i rp},\qquad
 p_M=(k+2\pi M)/L,
\]
where $M$ ranges over one complete four-dimensional alias class. The exact
finite-level formula in \src{YM3}{\texttt{eq:gammaj}} is
\begin{equation}\label{eq:v4-gamma}
 G^{(j)}(k)=\Pi(k)\diag_\mu(\gamma^{(j)}_\mu(k))\Pi(k),\qquad
 \gamma^{(j)}_\mu(k)=L^{-12}\sum_M
 \frac{\prod_{\nu=1}^4|D_L(p_{M,\nu})|^2\,
                  |D_L(p_{M,\mu})|^2}{\lambda(p_M)}.
\end{equation}
All apparent zero-over-zero factors in a geometric sum are evaluated as
$D_L(0)=L$. No restriction to an axial internal momentum or to one link
polarization is made in this section.

\begin{lemma}[Exact alias moments]
\label{lem:v4-moments}
For real $t$ and $p_r=(t+2\pi r)/L$,
\begin{align}
 \sum_{r=0}^{L-1}|D_L(p_r)|^2&=L^2,\label{eq:v4-moment2}\\
 \sum_{r=0}^{L-1}|D_L(p_r)|^4
 &=\frac{L^2}{3}\{2L^2+1+(L^2-1)\cos t\}\le L^4.
                                                        \label{eq:v4-moment4}
\end{align}
Consequently the diagonal weights in \cref{eq:v4-gamma} satisfy
\begin{equation}\label{eq:v4-weight-mass}
 L^{-12}\sum_M\prod_\nu|D_L(p_{M,\nu})|^2
                         |D_L(p_{M,\mu})|^2\le L^{-2}.
\end{equation}
\end{lemma}
\begin{proof}
The Fourier coefficients of $|D_L(p)|^2$ are
$c_n=L-|n|$ for $|n|<L$ and zero otherwise. Summing over the aliases
annihilates every coefficient whose index is not divisible by $L$.
Only $c_0=L$ survives in the second moment, proving
\cref{eq:v4-moment2}. In the square of this polynomial the possible
surviving indices are $0,L,-L$, and
\[
 \sum_n c_n^2=\frac{2L^3+L}{3},\qquad
 \sum_n c_n c_{L-n}=\frac{L^3-L}{6}.
\]
These follow by summing squares and $a(L-a)$ for $1\le a<L$.
Multiplication by the alias count $L$ proves
\cref{eq:v4-moment4}, also at $L=1$. Factor the four-dimensional sum:
there is one fourth moment and three second moments. Their product is
at most $L^{10}$, proving \cref{eq:v4-weight-mass}.
\end{proof}

\begin{theorem}[Uniform relative ellipticity at every level]
\label{thm:v4-ellipticity}
Set
\begin{equation}\label{eq:v4-constants}
 m_0=(2/\pi)^{12},\qquad M_0=5.
\end{equation}
For every $j\ge0$ and $k\ne0$ modulo $2\pi$,
\begin{equation}\label{eq:v4-ellipticity}
 \boxed{\displaystyle
 m_0G^{(0)}(k)\preceq G^{(j)}(k)\preceq5G^{(0)}(k),\qquad
 G^{(0)}(k)=\frac{\Pi(k)}{\lambda(k)}.}
\end{equation}
The same comparison holds for the limiting covariance $G^*$. Thus,
with inverses on the transverse quotient,
\begin{equation}\label{eq:v4-kernel-comparison}
 \frac15 S^{(0)}(k)\preceq S^{(j)}(k)
       \preceq m_0^{-1}S^{(0)}(k),\qquad
 S^{(0)}=\lambda I-qq^*.
\end{equation}
The kernels extend continuously to $S^{(j)}(0)=0$, uniformly in $j$.
This is not an inverse assigned to the divergent covariance at zero.
\end{theorem}
\begin{proof}
Choose $k\in[-\pi,\pi]^4\setminus\{0\}$ and centered alias representatives.
The principal alias is $p_0=k/L$. The inequalities
\[
 |D_L(k_\nu/L)|\ge(2/\pi)L,\qquad
 \lambda(k/L)\le |k|^2/L^2,\qquad
 \lambda(k)\ge4|k|^2/\pi^2
\]
follow from the elementary sine inequalities on $[-\pi/2,\pi/2]$.
The first has its continuous meaning when $k_\nu=0$. Keeping only the
principal positive summand in \cref{eq:v4-gamma} gives
\[
 \gamma_\mu^{(j)}(k)
 \ge\frac{(2/\pi)^{10}}{|k|^2}
 \ge\frac{(2/\pi)^{12}}{\lambda(k)}.
\]

For the upper bound, $|D_L(p)|\le L$ and concavity of sine gives
$\lambda(k/L)\ge\lambda(k)/L^2$. The principal summand is therefore
at most $1/\lambda(k)$. For each nonprincipal centered alias at least
one fine component has magnitude at least $\pi/L$. Since its magnitude
is at most $\pi$, its energy is at least $4/L^2$. Applying
\cref{eq:v4-weight-mass} to the remaining weights gives
\[
 \gamma_\mu^{(j)}(k)\le\frac1{\lambda(k)}+\frac14
                      \le\frac5{\lambda(k)},
\]
because $\lambda(k)\le16$. For $L=1$ the nonprincipal sum is empty.
Compression by $\Pi$ proves \cref{eq:v4-ellipticity}. Pass to the known
pointwise covariance limit to obtain the comparison for $G^*$.
Inversion reverses positive inequalities on the common transverse
space and gives \cref{eq:v4-kernel-comparison}. In particular
$\|S^{(j)}(k)\|\le m_0^{-1}\lambda(k)\to0$ uniformly in $j$, proving
the asserted continuous extension.
\end{proof}

This bound is stronger in a different direction than a positive constant
floor for $G^*$: it compares the entire finite tower to the massless free
covariance, including its infrared order, and it holds before any
large-$j$ threshold. It does not claim optimal constants.

\begin{corollary}[Kernel convergence without a late-level positivity threshold]
\label{cor:v4-kernel-rate}
Use the inherited estimate
$\sup_{k\ne0}\|G^{(j)}(k)-G^*(k)\|\le C_G4^{-j}$,
where $C_G=4d_1/3$ in the monograph. Put $c_0=m_0/16=256/\pi^{12}$.
Then
\begin{equation}\label{eq:v4-kernel-rate}
 \sup_k\|S^{(j)}(k)-S^*(k)\|
                         \le c_0^{-2}C_G4^{-j}
\end{equation}
for every $j\ge0$, with both kernels zero at $k=0$.
\end{corollary}
\begin{proof}
By \cref{thm:v4-ellipticity}, both covariances are at least
$c_0\Pi$ since $\lambda\le16$. On $\Ran\Pi$ apply
$S^{(j)}-S^*=S^{(j)}(G^*-G^{(j)})S^*$ and bound each inverse by
$c_0^{-1}$. The continuous extensions give the assertion at zero.
\end{proof}

\section{The same forty-five-coordinate elimination at every level}
\label{sec:v4-fibre}

We now use the \emph{actual linearized parity block} of f12. Let
$\mathcal B=\{0,1\}^4$. Its $64$ fine edges contain the $15$ ordered-tree
edges whose parent clears the highest nonzero bit. Let $J$ embed the
remaining $49$ coordinates into all fine edges. The block row on this
tree slice is denoted $L(k)$ in this section; it must not be confused
with the dyadic length used in \cref{sec:v4-ellipticity}.

For a non-tree internal edge of orientation $\mu$, its coefficient in
$L_\mu$ is $(1+e^{\mathrm i k_\mu})/16$. For a boundary edge of that
orientation the coefficient is $1/8$. Other rows have coefficient zero.
These coefficients follow by counting its two occurrences in the sixteen
chronological paths. Choose the boundary pivot at $s=e_\mu$ in each
orientation. The inclusion $E(k):\C^{45}\to\C^{49}$ is the identity on
the $17$ non-tree internal and $28$ nonpivot boundary coordinates; the
pivot is minus eight times the remaining block row. Define
$R:\C^4\to\C^{49}$ by copying a retained link to its eight boundary
links and assigning zero to internal links. Then
\begin{equation}\label{eq:v4-ER}
 L(k)E(k)=0,\qquad L(k)R=I_4,\qquad
 \Ran E(k)=\ker L(k).
\end{equation}
The source already uses this independent fibre and right inverse
\src{F12}{major proved interfaces, actual block and independent fibre}.

\begin{lemma}[A bounded full-dimensional coordinate inclusion]
\label{lem:v4-E}
For every real $k$,
\begin{equation}\label{eq:v4-Ebound}
 I_{45}\preceq E(k)^*E(k)\preceq15I_{45}.
\end{equation}
\end{lemma}
\begin{proof}
The free rows are the identity. The four pivot rows have disjoint
supports. Their free internal coordinate counts are $7,6,4,0$, and
there are seven free boundary coordinates in each row. Every internal
pivot coefficient has modulus at most one, as does every boundary
coefficient. Thus each pivot row has squared norm at most $14$.
The squared norm of their direct sum is at most $14$, which proves
the upper bound after adding the free identity rows. The lower bound
is immediate from those rows.
\end{proof}

Let $U(k)$ be the unitary cell Fourier matrix, with
\[
 U_{(s,\mu),(\epsilon,\nu)}(k)
       =\tfrac14e^{\mathrm i(k/2+\pi\epsilon)\cdot s}\delta_{\mu\nu}.
\]
It includes every one of the sixteen aliases and all four link
components. Put
\begin{align}
 A_j(k)&=J^*U(k)\diag_\epsilon
            S^{(j)}(k/2+\pi\epsilon)U(k)^*J,\label{eq:v4-Aj}\\
 H_j(k)&=E(k)^*A_j(k)E(k).\label{eq:v4-Hj}
\end{align}
These are $49$- and $45$-dimensional one-colour matrices. At zero
background the Lie-algebra colour factor is the identity; it can be
tensored back in without changing the spectral comparison. The factor
$1/4$ in $U$ is essential: in unnormalized plane-wave amplitudes it
corresponds to the $1/16$ covariance factor in the tower recursion.

\paragraph{Explicit inherited input.}
In the canonical one-colour curl normalization, f12 records
\begin{equation}\label{eq:v4-base-floor}
 H_0(k)=E(k)^*d(k)^*d(k)E(k)\succeq\kappa_0I,
                     \qquad \kappa_0=24/2327.
\end{equation}
Here $d$ is restricted to the $49$ tree-slice coordinates. This global
base estimate is imported as a proved source result; the new finite
checks reconstruct its incidence but do not independently recertify
that numerical lower bound over the whole torus. A different common
normalization multiplies every $H_j$ and its floor by the same factor.

\begin{theorem}[Uniform Gaussian fibre control on a fixed carrier]
\label{thm:v4-fibre}
For every $j\ge0$ and every $k\in\T^4$,
\begin{equation}\label{eq:v4-Hcomparison}
 \boxed{\displaystyle
 \tfrac15H_0(k)\preceq H_j(k)\preceq m_0^{-1}H_0(k),\qquad
 H_j(k)\succeq\frac{24}{11635}I_{45}.}
\end{equation}
In particular,
\begin{equation}\label{eq:v4-Hinverse}
 m_0H_0(k)^{-1}\preceq H_j(k)^{-1}\preceq5H_0(k)^{-1},\qquad
 \|H_j(k)^{-1}\|\le\frac{11635}{24}.
\end{equation}
The same assertions hold at the Gaussian fixed point. Moreover,
\begin{align}
 \sup_k\|H_j-H_*\|&\le15c_0^{-2}C_G4^{-j},\label{eq:v4-Hrate}\\
 \sup_k\|H_j^{-1}-H_*^{-1}\|
     &\le(5/\kappa_0)^2\,15c_0^{-2}C_G4^{-j}.\notag
\end{align}
\end{theorem}
\begin{proof}
Apply \cref{eq:v4-kernel-comparison} to each of the sixteen fine
momenta. Unitary transport and compression by $JE$ preserve Loewner
order, proving the first comparison. At a fine soft momentum both
kernels in that comparison are their continuous zero extensions, so
no excluded alias is needed. The inherited base floor gives the
numerical bound. Inversion and passage to the limit give
\cref{eq:v4-Hinverse}. The rate follows from
\cref{cor:v4-kernel-rate}, $\|E\|^2\le15$, and the inverse identity
with the common inverse bound $5/\kappa_0$.
\end{proof}

Thus the zero-background Gaussian elimination need not use a growing
$4^{4j}$ fine-alias carrier to establish its invertibility. It can use
one $45$-coordinate fibre per current cell with the transported kernel
$S^{(j)}$. This does not say that the transported cubic and quartic
vertices equal the original Wilson ones, or that their position-space
range stays fixed.

\subsection{Full four-dimensional soft-momentum removal by a Schur recursion}

\begin{theorem}[Exact full-kernel recursion and finite-level analyticity]
\label{thm:v4-Schur}
For the actual block and normalization above,
\begin{equation}\label{eq:v4-Schur}
 \boxed{\displaystyle
 S^{(j+1)}(k)=R^*A_j(k)R
  -R^*A_j(k)E(k)H_j(k)^{-1}E(k)^*A_j(k)R.}
\end{equation}
Each finite-level $S^{(j)}$, $A_j$, $H_j$, and $H_j^{-1}$ is real
analytic on the entire four-dimensional Bloch torus, in its compatible
cell and alias frames. In particular, the inverse in
\cref{eq:v4-Schur} is nonsingular at every internal soft momentum.
The retained kernel still satisfies $S^{(j)}(0)=0$.
\end{theorem}
\begin{proof}
Every tree-slice field with block $c$ is uniquely $Ez+Rc$. Minimizing
its quadratic energy gives
\[
 z_*=-H_j^{-1}E^*A_jRc.
\]
Substitution gives the right side of \cref{eq:v4-Schur}. At $k\ne0$
this is the same Gaussian quotient minimization as the retained
covariance recursion: fine gauge changes map to coarse gauge changes
by the Ward identity. The energy of a coarse gradient is zero, and on
the transverse quotient its inverse is exactly the stated $G^{(j+1)}$.
The equality at $k=0$ follows by continuity, rather than by inverting
a coarse zero mode.

For analyticity, start with the Laurent-polynomial kernel
$S^{(0)}=\lambda I-qq^*$. If $S^{(j)}$ is real analytic, so are
$A_j$ and $H_j$ in the local alias frames. The uniform real lower bound
in \cref{thm:v4-fibre} implies that $H_j^{-1}$ is real analytic at
every real momentum. All factors on the right of
\cref{eq:v4-Schur} are therefore analytic, establishing the induction.
The alias changes at a $2\pi$ boundary simply permute the Fourier
columns and give the same cell operator. This proves the global
statement. It supplies no $j$-independent complex strip or derivative
bound: real-axis spectral bounds alone do not prove those stronger
assertions.
\end{proof}

\begin{scope}[The gap belongs to the eliminated Gaussian fibre]
The lower bound $24/11635$ is dimensionless, on a recentered Gaussian
fibre with the retained coarse data held fixed. It is not a gap of the
conditional compact diffusion ruled out in f13 V8, nor a gap of the
physical transfer Hamiltonian. The retained Schur kernel remains
infrared massless. In particular, no substitution of this constant
into $-\tau^{-1}\log q$ is justified.
\end{scope}

\section{A bordered determinant for the actual measured block}
\label{sec:v4-bordered}

The previous section concerns the quadratic principal kernel. A complete
background response also requires the \emph{nonlinear} constraint and
its measure. Work on one fixed finite-dimensional real tree chart,
with a smooth fine action $S(x)$, positive fine density $a(x)=e^{\ell(x)}$,
and a submersion $F:\R^n\to\R^m$, $0<m<n$. Assume
\[
 F(0)=0,\qquad DS(0)=0,
\]
and assume that $D^2S(0)$ is positive on $\ker DF(0)$. It need not be
invertible on the ambient fine space.

For $c$ near zero, let $(x(c),\eta(c))$ be the local stationary branch
satisfying the sign convention
\begin{equation}\label{eq:v4-stationary}
 \nabla S(x)+DF(x)^*\eta=0,\qquad F(x)=c.
\end{equation}
Define its Lagrangian Hessian and bordered matrix by
\begin{equation}\label{eq:v4-border}
 \mathcal A(c)=D^2S(x(c))+\sum_{\alpha=1}^m
                         \eta_\alpha(c)D^2F_\alpha(x(c)),\qquad
 \mathcal K(c)=
 \begin{pmatrix}\mathcal A(c)&DF(x(c))^*\\DF(x(c))&0\end{pmatrix}.
\end{equation}
With the source's alternative convention $\lambda=-\eta$, the first
block is $S''-\langle\lambda,F''\rangle$, exactly its constrained
fine Hessian. The multiplier correction is not optional.

\begin{lemma}[Bordered determinant and coarea]
\label{lem:v4-coarea}
Let $J_0$ have full row rank and let a real symmetric $A_0$ be positive
on $\ker J_0$. If $T$ is an orthonormal inclusion of $\ker J_0$, then
\begin{equation}\label{eq:v4-coarea}
 \det\begin{pmatrix}A_0&J_0^*\\J_0&0\end{pmatrix}
       =(-1)^m\det(J_0J_0^*)\det(T^*A_0T)\ne0.
\end{equation}
For the local stationary branch of \cref{eq:v4-stationary}, the one-loop
term of the pushforward density relative to \emph{Lebesgue} $dc$ is
\begin{equation}\label{eq:v4-Q}
 \boxed{\displaystyle
 \mathcal Q(c)=\tfrac12\log|\det\mathcal K(c)|-\ell(x(c)).}
\end{equation}
Its classical term is $\mathcal G(c)=S(x(c))$. Relative to a declared
coarse reference $w(c)\,dc$, add $\log w(c)$ to $\mathcal Q$.
\end{lemma}
\begin{proof}
Complete $T$ to an orthogonal basis $(T,N)$ of $\R^n$. In these
coordinates $J_0=(0,B)$ with $B=J_0N$ invertible. Expand the bordered
determinant first along its last $m$ rows and then the corresponding
columns. The remaining tangent block is $T^*A_0T$ and the two normal
factors have product $\det(B)^2=\det(J_0J_0^*)$. Their exchange gives
$(-1)^m$, proving \cref{eq:v4-coarea}.

Locally the pushforward integral can be written
\[
 Z_\beta(c)=\int\delta(F(x)-c)e^{-\beta S(x)}a(x)\,dx.
\]
Coarea places the factor $\det(DF\,DF^*)^{-1/2}$ on the level
surface. In orthonormal tangent coordinates at its strict stationary
minimum, the surface Hessian is $T^*\mathcal A T$. To see this,
differentiate $F(x(z))=c$ twice and use
$\nabla S=-DF^*\eta$ in the second chain rule for $S(x(z))$.
The Gaussian factor from the tangent integral is
$(2\pi/\beta)^{(n-m)/2}\det(T^*\mathcal A T)^{-1/2}$.
Combining it with coarea and \cref{eq:v4-coarea} gives
\[
 Z_\beta(c)\sim (2\pi/\beta)^{(n-m)/2}
 e^{-\beta\mathcal G(c)}
 \frac{a(x(c))}{\sqrt{|\det\mathcal K(c)|}},
\]
which proves \cref{eq:v4-Q}. This leading local coefficient can be
justified by a cutoff equal to one near the stationary point: rescale
tangent coordinates by $\beta^{-1/2}$, use the positive tangent Hessian
for domination on a smaller chart, and bound its fixed exterior by a
strict local action increment. This is a fixed-dimensional local
stationary-phase statement, not a global-minimum or cofinal remainder
claim. Division of the density by $w(c)$ gives the last assertion.
\end{proof}

The determinant contains the constrained Hessian \emph{and} the coarea
normalization. Adding a separate inverse-pivot Jacobian to
\cref{eq:v4-Q} after that same factor has already been represented by
coarea would double count it. Alternatively, the source's independent
chart $(z,c)$ gives $\tfrac12\log\det H_c-\ell_{\rm chart}$; the two
expressions are equal up to the fixed coordinate constant. Nonlinear
changes of the fine tree chart still require the usual transformation
of the fine density $a$.

\section{Closed fourth-action and second-density jet transport}
\label{sec:v4-jets}

All tensors in this section are evaluated at $x=0$. Write
\[
 A=S_2,\quad L=F_1,\quad T=S_3,\quad V=S_4,\quad B=F_2,\quad C=F_3,
 \qquad
 K=\begin{pmatrix}A&L^*\\L&0\end{pmatrix}.
\]
The subscripts denote \emph{true derivatives}, not homogeneous Taylor
coefficients. For example $B[a,\cdot]$ is an $m$-by-$n$ map and
$\eta\cdot B$ is a symmetric $n$-by-$n$ form. This use of $B,C$ is
local to the jet theorem; they are not covariance matrices.

For an arbitrary real retained direction $v\in\R^m$, define
\begin{align}
 \binom{a}{\eta_1}&=K^{-1}\binom{0}{v},\label{eq:v4-first-solve}\\
 w&=B[a,a],\qquad
 r=T[a,a,\cdot]+2B[a,\cdot]^*\eta_1,\qquad
 \zeta=\binom{r}{w},\notag\\
 \binom{b}{\eta_2}&=-K^{-1}\zeta.\label{eq:v4-second-solve}
\end{align}
Here the covector $T[a,a,\cdot]$ is identified with a vector by the
fixed real Euclidean coordinates. These solves use no inverse of $A$
and no inverse of a retained massless kernel.

\begin{theorem}[Complete finite jet compiler]
\label{thm:v4-compiler}
The derivatives of the stationary branch along $c=tv$ are
$x'(0)=a$, $x''(0)=b$, $\eta'(0)=\eta_1$, and
$\eta''(0)=\eta_2$. The outgoing classical action satisfies
\begin{align}
 \mathcal G_2[v,v]&=-\langle\eta_1,v\rangle,\label{eq:v4-Gjets}\\
 \mathcal G_3[v,v,v]&=T[a,a,a]+3\langle\eta_1,B[a,a]\rangle,\notag\\
 \mathcal G_4[v,v,v,v]&=V[a,a,a,a]
   +4\langle\eta_1,C[a,a,a]\rangle-3\langle\zeta,K^{-1}\zeta\rangle.
 \notag
\end{align}
Define the first and second derivatives of the \emph{whole border} by
\begin{align}
 K_1&=\begin{pmatrix}
 T[a,\cdot,\cdot]+\eta_1\cdot B&B[a,\cdot]^*\\
 B[a,\cdot]&0
 \end{pmatrix},\label{eq:v4-K1}\\
 K_2&=\begin{pmatrix}
 V[a,a,\cdot,\cdot]+T[b,\cdot,\cdot]
       +\eta_2\cdot B+2\eta_1\cdot C[a,\cdot,\cdot]
       &(C[a,a,\cdot]+B[b,\cdot])^*\\
 C[a,a,\cdot]+B[b,\cdot]&0
 \end{pmatrix}.\label{eq:v4-K2}
\end{align}
Then the complete one-loop derivatives relative to Lebesgue measure are
\begin{align}
 \mathcal Q_1[v]&=\tfrac12\Tr(K^{-1}K_1)-\ell_1[a],\label{eq:v4-Q1}\\
 \boxed{\displaystyle
 \mathcal Q_2[v,v]
  =\tfrac12\Tr(K^{-1}K_2-K^{-1}K_1K^{-1}K_1)
                -\ell_2[a,a]-\ell_1[b].}\label{eq:v4-Q2}
\end{align}
Thus $S_{\le4}$, $F_{\le3}$, and $\ell_{\le2}$ determine
$\mathcal G_{\le4}$ and $\mathcal Q_{\le2}$. No fourth block jet or
fifth action jet is needed at a stationary zero reference. Mixed
retained derivatives are recovered by polarization of the displayed
homogeneous polynomials.
\end{theorem}
\begin{proof}
At the reference point the stationary multiplier is zero: $DS(0)=0$
and injectivity of $L^*$ imply $\eta(0)=0$. The derivative of the
stationary equations \cref{eq:v4-stationary} is the invertible border
$K$, so the implicit function theorem gives the local branch.
Differentiating once gives \cref{eq:v4-first-solve}. Differentiating
again gives
\[
 Ab+L^*\eta_2=-T[a,a,\cdot]-2B[a,\cdot]^*\eta_1,
                  \qquad Lb=-B[a,a],
\]
which is \cref{eq:v4-second-solve}.

Since $\mathcal G'(t)=-\langle\eta(t),v\rangle$, the second derivative
is the first line of \cref{eq:v4-Gjets}. For the remaining derivatives
write $x(tv)=ta+t^2b/2+t^3h/6+O(t^4)$. The constraint gives
\[
 La=v,\quad Lb=-B[a,a],\quad
 Lh=-3B[a,b]-C[a,a,a].
\]
Together with $Aa=-L^*\eta_1$, the cubic term of $S(x(tv))$ is
$T[a,a,a]+3\langle\eta_1,B[a,a]\rangle$ as a true derivative.
Its fourth derivative is
\[
 V[a,a,a,a]+4\langle\eta_1,C[a,a,a]\rangle
             +6\langle r,b\rangle+3\langle b,Ab\rangle.
\]
The second solve implies
$\langle r,b\rangle=-\langle b,Ab\rangle+\langle\eta_2,w\rangle$.
Also
\[
 \langle\zeta,K^{-1}\zeta\rangle
  =\left\langle\binom b{\eta_2},K\binom b{\eta_2}\right\rangle
  =\langle b,Ab\rangle-2\langle\eta_2,w\rangle.
\]
These two identities prove the quartic formula.

Differentiate \cref{eq:v4-border} along the branch. The first
Lagrangian-Hessian derivative is $T[a]+\eta_1\cdot B$; the second is
$V[a,a]+T[b]+\eta_2\cdot B+2\eta_1\cdot C[a]$.
The two Jacobian derivatives are $B[a,\cdot]$ and
$C[a,a,\cdot]+B[b,\cdot]$. This gives \cref{eq:v4-K1,eq:v4-K2}.
A term involving $F_4$ is multiplied by $\eta(0)=0$ and is absent.
On a small branch the nonzero determinant of $\mathcal K$ has fixed
sign, so ordinary differentiation of $\log|\det\mathcal K|$ gives
\[
 (\log|\det\mathcal K|)''
   =\Tr(K^{-1}K_2-K^{-1}K_1K^{-1}K_1).
\]
The second chain rule for $\ell(x(tv))$ is
$\ell_2[a,a]+\ell_1[b]$. Apply \cref{eq:v4-Q} to finish the proof.
\end{proof}

The border is indefinite; its inverse is not a covariance of additional
physical degrees of freedom. Its trace formula is a determinant identity.
Its use retains the multiplier and measure derivatives which are absent
from a bare $\Tr(\Gamma R\Gamma R)$ Wilson component. In the source's
Cartan case, oddness can make $b$ and $\eta_2$ vanish; such a simplification
is made only after proving the symmetry, as f12 V3 does. The general
compiler does not impose it by default.

\subsection{Why the nonlinear block cannot be inferred from Gaussian data}

\begin{proposition}[Unchanged harmonic branch, arbitrary missing one-loop curvature]
\label{prop:v4-nonidentifiability}
Let
\[
 S(x,y)=\tfrac12(ax^2+y^2),\qquad a>0,\qquad
 F_\kappa(x,y)=x(1+\kappa y^2/2),\qquad \ell=0.
\]
On a sufficiently small chart and for sufficiently small retained $c$,
the local minimizing branch is $(x,y)=(c,0)$ for every real $\kappa$.
All these models have the same fine action, the same $DF(0)$ and
$D^2F(0)$, the same zero-background fibre covariance, and the same
entire harmonic branch. Nevertheless their complete local one-loop
terms obey
\begin{equation}\label{eq:v4-counter}
 \boxed{\displaystyle
 \mathcal Q_\kappa(c)-\mathcal Q_\kappa(0)
   =\tfrac12\log(1-a\kappa c^2),\qquad
 \mathcal Q_\kappa''(0)=-a\kappa.}
\end{equation}
\end{proposition}
\begin{proof}
The constraint solves $x=c/(1+\kappa y^2/2)$. At $y=0$ the reduced action
is stationary and its fibre Hessian is $1-a\kappa c^2>0$ near zero.
Continuity and the implicit function theorem identify this strict local
branch. Along it $DF=(1,0)$, so the coarea Jacobian is one. The reduced
fine density there is also one. Formula \cref{eq:v4-Q} gives
\cref{eq:v4-counter}. The nonlinear part of $F_\kappa$ starts at degree
three, proving all the equal-data assertions.
\end{proof}

This finite analytic example is \emph{not} a Yang--Mills block and does
not claim freedom to change the supplied compact rule. It proves the
relevant nonidentifiability statement: even exact harmonic information
and complete fine-action vertices do not determine the pushed one-loop
response without the nonlinear block data. For the actual programme
those data are fixed by the root-path convention in
\src{F12}{major proved interfaces}. They are inputs to transport, not
adjustable counterterms.

\section{Nonlinear composition without resetting the action or density}
\label{sec:v4-composition}

For iteration use a fixed Lebesgue chart convention. Associate to a block
$F$ the finite jet operation
\begin{equation}\label{eq:v4-map}
 \mathfrak R_F:\ (S_{\le4},\ell_{\le2})
       \longmapsto(\mathcal G_{\le4},\ell_{+,\le2}),
 \qquad \ell_+=-\mathcal Q.
\end{equation}
The reference action is centered, while its field-independent vacuum
constant and $(2\pi/\beta)$ factors are recorded separately. Coarse Haar
terms may instead be carried relative to Haar, but converting conventions
must retain the same density factor once, not discard or duplicate it.

\begin{theorem}[Finite one-loop jet composition]
\label{thm:v4-composition}
Let $F_1:\R^n\to\R^r$ and $F_2:\R^r\to\R^m$ be smooth submersions
with compatible local stationary branches. Assume that the eliminated
Hessians for the two stages, and equivalently for their common local
composition, are strictly positive on their constraint kernels. Then
\begin{equation}\label{eq:v4-composition}
 \boxed{\displaystyle
 \mathfrak R_{F_2}\bigl(\mathfrak R_{F_1}(S_{\le4},\ell_{\le2})\bigr)
   =\mathfrak R_{F_2\circ F_1}(S_{\le4},\ell_{\le2})}
\end{equation}
for the fourth classical-action and second one-loop-density jets, with
field-independent normalizing constants accounted for consistently.
The composite block needs only the third jets of $F_1,F_2$.
\end{theorem}
\begin{proof}
In common local coordinates, exact pushforward gives
\[
 \int dy\,\delta(F_2(y)-c)
       \int dx\,\delta(F_1(x)-y)a(x)e^{-\beta S(x)}
 =\int dx\,\delta(F_2(F_1(x))-c)a(x)e^{-\beta S(x)}.
\]
This identity includes the coarea factors; they are not inserted as
independent approximations. Classical minimization on the compatible
branches is associative: first minimize at fixed $y$, then minimize
the resulting action under $F_2(y)=c$. The tangent quadratic Hessian
of the joint minimum is a block matrix whose first fibre block is
positive and whose Schur complement is the second eliminated Hessian.
Its Gaussian determinant is the product of those two determinants.
The ordinary change-of-variables formula supplies the product of the
same two coarea factors. Thus the leading stationary-phase density
computed in two stages equals the one-step density of the composite.
Equivalently, apply \cref{lem:v4-coarea} in each set of local coordinates
to this determinant factorization. There is no contribution from the
order-$\beta^{-1}$ term to the order-$\beta^0$ coefficient at a fixed
number of stages.

By \cref{thm:v4-compiler}, the first output action and density through
the specified orders need no higher fine data. The same is true at
the second stage. The chain rule through degree three determines
$(F_2\circ F_1)_{\le3}$ from the two third block jets. Therefore
truncating to the stated jets between stages loses no term in the
claimed output. This proves \cref{eq:v4-composition}. The argument is
finite-stage; no uniform estimate on a growing number of discarded
higher-loop remainders is asserted.
\end{proof}

This is the nonlinear version of the accounting discipline needed by V1.
Mixed Gaussian returns, constraint tangents, and measure changes are
carried by the transformed action and density. They must not additionally
be charged as discarded original-shell diagrams. In particular, after
one step the input to the next is $(\mathcal G,\ell_+)$, not the original
Wilson cubic/quartic action with a freshly normalized unit amplitude.

\subsection{The full finite-level principal germ is now a defined calculation}

\begin{corollary}[A fixed-fibre construction of the complete finite-level germ]
\label{cor:v4-full-germ}
Start from the supplied small identity chart of the compact Wilson law,
its actual root-path block, and its fine Haar density. Retain the same
chart and density conventions and iterate \cref{eq:v4-map} a fixed finite
number of times on its local stationary branch. The Gaussian quadratic
part is the stated tower, and every elimination uses the full
$45(\dim\mathfrak g)$ fibre with the lower bound of
\cref{thm:v4-fibre} in the declared normalization. The fourth tree-action
and second one-loop-density jets are determined by the primitive
$S_{\le4},F_{\le3},\ell_{\le2}$ and the explicit bordered solves.
No independent large fine-alias vertex prescription is required.

For translation-invariant primitive jets with their finite-range Fourier
symbols, at every fixed number of stages the resulting momentum vertices
and one-loop integrands are real analytic on their real internal
four-dimensional tori near zero external momentum. Their fixed-order
external derivatives can therefore be passed under the finite-level
internal momentum integrals. The constants in these analytic statements
are allowed to depend on the number of stages.
\end{corollary}
\begin{proof}
At zero field, $F_1$ is the given block-average map and the classical
quadratic minimization is exactly \cref{thm:v4-Schur}. Induction therefore
identifies the quadratic kernels with $S^{(j)}$. The fixed-fibre floor
proves the required invertibility at every stage on the identity branch.
The other entries are generated by \cref{thm:v4-compiler} and the measured
composition theorem. The density of a relative Haar presentation is
converted as specified above, so its local Jacobian is not lost.

For the final assertion, the initial finite-range vertices are analytic
trigonometric functions of their incident momenta. At any fixed stage,
the compiler uses sums, products, contractions, and inverses of the
current fibre or bordered matrices. On the tree slice $L$ has the fixed
right inverse $R$, and the bordered matrix is invertible because its
fibre block is positive. Its real-momentum inverse is analytic, just as
in the proof of \cref{thm:v4-Schur}. These statements hold also at sums
of incident momenta, including zero. A finite chart cover of the compact
momentum tori gives a common analytic neighbourhood at that fixed stage.
Finite tree contractions preserve analyticity. The one-loop traces have
only a finite number of momentum integrals at that stage; on this compact
domain their derivatives have constant integrable majorants. The
fundamental theorem of calculus and dominated convergence justify the
differentiated integrals and complete the induction. This constructs
finite-level perturbative principal coefficients, not an interacting
infinite-volume measure or a scale-uniform analytic neighbourhood.
\end{proof}

The root-path logarithm must be expanded with its actual distinguished
path and chronological products. Its already known second tangent cannot
be replaced by zero because a corner regression happens to vanish. The
source's $19/48$ one-step measure coefficient is likewise retained in its
own normalization, not reused as a universal density at every level.

\section{Executed checks, changed acquisition order, and version record}
\label{sec:v4-verification}

\subsection{An independent nonlinear composition certificate}
For reproducibility, the complete jet compiler is tested on a genuinely
nonlinear two-stage finite model. This model verifies the implementation;
it is not presented as Yang--Mills data. Its fine action and log amplitude
are
\begin{align}
 S(x,y,z)={}&\tfrac12(3x^2+2y^2+5z^2)
  +\tfrac13xyz+\tfrac15x^2z+\tfrac17yz^2
  +\tfrac1{11}(x^2+y^2+z^2)^2,\label{eq:v4-check-model}\\
 \ell(x,y,z)={}&x/7-y/9+z/13+xz/8-yz/10+x^2/17+z^2/19.\notag
\end{align}
The first block is
$F_1=(x+z^2/3,\,y+xz/4)$ and the second is
$F_2(u,v)=u+v^2/5+uv/7+uv^2/9$. Direct composition and two-stage
transport both give the exact output jets
\begin{align}
 \mathcal G(c)&=\tfrac32c^2+\frac{11047}{269500}c^4+O(c^5),
                                                       \label{eq:v4-check-output}\\
 \ell_+(c)-\ell_+(0)&=\frac9{14}c+
                       \frac{2663743}{21441420}c^2+O(c^3).\notag
\end{align}
The two stage border determinants at zero are $5$ and $-2$, while the
direct determinant is $-10$; their absolute factors agree.

The first calculation is independently checked in the triangular chart
$x=u-z^2/3$, $y=v-xz/4$, whose eliminated coordinate is $z$ and whose
coordinate Jacobian is one. The second uses
\[
 u=\frac{c-v^2/5}{1+v/7+v^2/9},
\]
and its logarithmic Jacobian is \emph{retained} as
$-\log(1+v/7+v^2/9)$. In these checks the stationary scalar series is
solved directly from the derivative of the reduced action, not from the
bordered formulas. They agree with \cref{eq:v4-check-output} exactly.
The driver also verifies the nonidentifiability example, an ambient
singular Hessian with a nonsingular border, and the independence from
unused fifth-action, fourth-block, and third-density terms.

\subsection{Full-dimensional geometry and covariance checks}
The actual parity geometry is reconstructed from its sites and oriented
plaquettes. Exact checks verify all $96$ curl rows, the scalar Ward map,
$LE=0$, $LR=I$, and the independent $45$ coordinate rows. The geometric
moment identities are additionally checked by exact Fourier coefficients
and complex-rational alias sums. Three full four-dimensional corner
calculations independently compare the $45$-dimensional Schur solve with
the Gaussian alias covariance. At the soft corner the retained kernel
is exactly zero without inverting a zero fine mode. At phases
$(-1,1,1,1)$ it is $\diag(0,8,8,8)$; at phases $(-1,-1,-1,1)$ it is
\[
 \begin{pmatrix}
 32&-16&-16&0\\-16&32&-16&0\\-16&-16&32&0\\0&0&0&24
 \end{pmatrix}.
\]
These are regression tests of the normalization and recursion, not new
response cards and not substitutes for the analytic bounds.

The new exact drivers report $30$, $3$, and $6$ passing test families,
respectively. A separately labelled floating-point diagnostic evaluates
$28$ full $45$-dimensional fibre matrices at levels zero through three,
including a soft point and a nearby point, and compares their Schur
outputs with the closed tower. These numerical samples are not global
certificates; the all-momentum, all-level theorem is the proof in
\cref{sec:v4-ellipticity,sec:v4-fibre}. The inherited rational base floor
is used openly rather than purportedly certified by these samples.

\subsection{The acquisition to pursue, and what not to repeat}
The work now has a complete, composable finite-level principal object on
a fixed full fibre. The missing quantitative task is to control the
\emph{transported} third and fourth action jets and second density jets
as the number of scales grows, in norms that include their momentum
arguments, Ward-related sums, and loop traces. The uniform inverse bound
above pays for Gaussian fibre inversions; it does not pay for arbitrary
vertex growth. The internal derivative passage has been proved for
each fixed stage, but a common strip and cofinal bounds still need proof.

Accordingly, another enlarged \emph{bare Wilson-component} trace or a
higher-precision integral of V3's reduced symbol is not selected as the
next main acquisition. The next calculation should evaluate and estimate
$\mathfrak R_F$ on the already specified compact primitive jets, retaining
the entire transformed packet. The first nontrivial two-stage comparison
must use that packet on both routes, not reset its cubic/quartic vertices
or Haar amplitude. Its purpose is to acquire quantitative scale control
of the complete response, not to establish another abstract ownership
identity. The source's existing one-step actual-vertex formula is the
baseline for that implementation.

\begin{scope}[Remaining physical obligations]
The all-level estimates here concern the Gaussian principal tower.
The measured jet composition is exact at the stated formal orders and
at a fixed number of local stationary steps. It does not prove a common
small-field neighbourhood for infinitely many nonlinear steps, control
all large fields, identify the complete limiting Yang--Mills coefficient,
construct the interacting continuum theory, or establish a positive
physical Hamiltonian gap. The local-vacuum normalization, noncollapse,
and calibrated physical-time contraction of the closure monograph remain
separate. The reason for these boundaries is the specific missing
nonlinear and infrared estimates, not the historical status of the problem.
\end{scope}

\paragraph{Append-only record.}
The complete V1--V3 mathematical body, including all earlier end markers,
is retained byte-for-byte. Only the version display in the preamble is
changed. The new code, exact reports, numerical diagnostics, and hashes
are supplied separately. Later V5 work is to append after the marker
below, with any needed correction stated explicitly rather than silently
changing a prior proof.

\begin{thebibliography}{99}
\bibitem[B84]{BalabanV4}
T.~Balaban, \emph{Propagators and renormalization transformations for
lattice gauge theories. I}, Communications in Mathematical Physics
\textbf{95} (1984), 17--40. DOI: 10.1007/BF01215753.
Cited for historical Gaussian-renormalization context only. No external
estimate is imported as a missing bound for the present block rule.
\end{thebibliography}

% END OF VERSION 4 MATHEMATICAL BODY.
% Append subsequent requested versions below this marker; retain V1--V4 above.

\clearpage
\section{V5: a common analytic domain for the full Gaussian tower}
\label{sec:v5-context}

The next missing linear estimate in V4 is not another real-axis spectral
comparison.  It is a common complex-momentum domain, with summable spatial
kernels, for every level of the actual four-dimensional Gaussian tower.
Without that estimate the uniform fibre inverse in
\cref{thm:v4-fibre} cannot legitimately be used as a uniform momentum-jet
bound in the nonlinear transformation.  We prove that estimate here,
including at the complex continuation of the soft momentum.

\paragraph{Nonduplication.}
The supplied f12 V5 already proves a complex strip for the \emph{initial}
finite-range fibre, by its explicit Laurent stencil
\src{F12}{\texttt{prop:f12v5-strip}}.  V4 proves finite-level analyticity
and an all-level real inverse bound, but explicitly does not provide a
common strip.  Neither assertion is relabelled as new.  This version uses
the exact all-level alias formula to control the spatially non-finite-range
kernels generated after arbitrarily many Gaussian steps.  The resulting
strip, exponential locality, and locality-norm convergence are uniform
in the Gaussian level.  A further calculation determines the universal
fourth \emph{momentum} jet of the quadratic action at every level.
Momentum order and field order are distinct throughout.

\paragraph{Inputs and scope.}
The positive alias formula, the real comparisons, the real convergence
rate, and the inherited base fibre floor are precisely the inputs of V4.
The nonlinear measured transformation of V4 is retained unchanged.  The
new argument does not assume that the transported nonlinear vertices
are uniformly bounded.  Gaussian locality as a renormalization input is
classical context \cite{BalabanV4}; the proof below concerns the specific
block-average tower in the supplied programme and imports no external
locality theorem or identification of a Yang--Mills coefficient.

\section{Complex alias estimates before taking a gauge quotient}
\label{sec:v5-alias}

For complex $z\in\mathbb C^4$ set
\[
 q_\mu(z)=e^{\mathrm i z_\mu}-1,\qquad
 q^\sharp_\mu(z)=e^{-\mathrm i z_\mu}-1,\qquad
 \lambda(z)=\sum_\mu q_\mu(z)q^\sharp_\mu(z).
\]
The symbol $q^\sharp$ is a row when used in a matrix product.  For a
matrix germ $F$, the holomorphic adjoint is
$F^\sharp(z)=\overline{F(\overline z)}^{\mathsf T}$, not a
pointwise conjugate transpose at complex $z$.  It agrees with the adjoint
on real momenta.  In particular $\lambda(z)$ is a holomorphic function,
not $\sum|q_\mu(z)|^2$ off the real torus.

Write $L=2^j$, $p_M=(z+2\pi M)/L$, and
\[
 D_L(p)=\sum_{a=0}^{L-1}e^{\mathrm iap},\qquad
 d_{L,M,\mu}(z)=D_L(p_{M,\mu})D_L(-p_{M,\mu}).
\]
The meromorphic continuation of the diagonal covariance formula is
\begin{equation}\label{eq:v5-gamma}
 \gamma_{L,\mu}(z)=L^{-12}\sum_{M\bmod L}
       \frac{\prod_\nu d_{L,M,\nu}(z)\,d_{L,M,\mu}(z)}{\lambda(p_M)}.
\end{equation}
On the real torus this is the diagonal factor in
\cref{eq:v4-gamma}.  We will not evaluate its principal pole separately
at $z=0$.

\begin{lemma}[Uniform complex weight mass]
\label{lem:v5-weight}
For $z=x+\mathrm iy$ with $\|y\|_\infty\le a$,
\begin{equation}\label{eq:v5-weight}
 L^{-12}\sum_{M\bmod L}
       \left|\prod_\nu d_{L,M,\nu}(z)\,d_{L,M,\mu}(z)\right|
                    \le e^{10a}L^{-2}.
\end{equation}
This holds for all real $x$, all dyadic $L$, and every component $\mu$.
\end{lemma}
\begin{proof}
For one coordinate, root-of-unity orthogonality gives
\[
 \sum_{m=0}^{L-1}\left|D_L((x+\mathrm iy+2\pi m)/L)\right|^2
       =L\sum_{r=0}^{L-1}e^{-2ry/L}\le L^2e^{2a}.
\]
The same bound applies with the argument negated.  Cauchy--Schwarz yields
$\sum_m|d_{L,m}|\le L^2e^{2a}$.  Also
$\max_m|d_{L,m}|\le L^2e^{2a}$ by the geometric sum itself, whence
$\sum_m|d_{L,m}|^2\le L^4e^{4a}$.  Factoring the four-coordinate sum
uses the last estimate once and the preceding estimate three times.
Their product, followed by $L^{-12}$, is
\cref{eq:v5-weight}.  These absolute-value estimates do not change the
holomorphic definition of the factors.
\end{proof}

Define, in a neighbourhood of zero,
\begin{align}
 f_{L,\mu}(z)&=L^{-2}D_L(z_\mu/L)D_L(-z_\mu/L),\notag\\
 P_L(z)&=\prod_{\nu=1}^4f_{L,\nu}(z),\qquad
 s_L(z)=L^2\lambda(z/L),\label{eq:v5-soft-data}\\
 b_{L,\mu}(z)&=L^{-12}\sum_{M\ne0}
       \frac{\prod_\nu d_{L,M,\nu}(z)\,d_{L,M,\mu}(z)}{\lambda(p_M)}.
       \notag
\end{align}
For $L=1$ the last sum is zero.  Thus
\begin{equation}\label{eq:v5-split}
             \gamma_{L,\mu}=\frac{P_Lf_{L,\mu}}{s_L}+b_{L,\mu}.
\end{equation}
The normalization in the principal summand is important: its numerator
is $P_Lf_{L,\mu}$, not an additional power of $L$.

\begin{lemma}[A common soft chart for the nonprincipal sum]
\label{lem:v5-tail}
On the polydisc $\|z\|_\infty<r_0=1/8$, every $b_{L,\mu}$ is
holomorphic and
\begin{equation}\label{eq:v5-tail}
              |b_{L,\mu}(z)|\le e^{5/4}/3<2
\end{equation}
uniformly in $L$.  Each $b_{L,\mu}$ is even in each coordinate,
vanishes at zero, and is $O(\|z\|_\infty^2)$ uniformly on a smaller
polydisc.  The functions $f_{L,\mu}$ are even entire functions of their
coordinate, with $f_{L,\mu}(0)=1$, and
\begin{equation}\label{eq:v5-fclose}
 |f_{L,\mu}(z)-1|\le \tfrac12|z_\mu|^2e^{|z_\mu|}.
\end{equation}
\end{lemma}
\begin{proof}
For $L\ge2$, a nonprincipal alias has at least one real fine component
whose distance to $2\pi\mathbb Z$ is at least $(2\pi-r_0)/L$,
interpreted modulo $2\pi$.  In particular its real energy is at least
$4/L^2$, by $\sin(\pi/(2L))\ge1/L$.  For $p=u+\mathrm iv$,
\[
 \operatorname{Re}\lambda(p)
 \ge\lambda(u)-2\sum_\nu(\cosh v_\nu-1)
 \ge \lambda(u)-4r_0^2\cosh(r_0)/L^2.
\]
Here $|v_\nu|\le r_0/L$.  Since $4r_0^2\cosh(r_0)<1$, every
nonprincipal denominator has real part at least $3/L^2$.  It therefore
has no zero in the polydisc.  Applying \cref{lem:v5-weight} to the
nonprincipal terms gives \cref{eq:v5-tail}.

At zero, a nonprincipal alias has at least one nontrivial root of unity,
so its corresponding $D_L$ vanishes.  Every summand defining $b$ is
therefore zero.  Replacing any $z_\nu$ by $-z_\nu$ relabels the
corresponding alias by its negative; the principal alias remains
principal.  This proves the stated evenness.  The uniform holomorphic
bound and the absence of constant and linear coefficients give the
uniform quadratic order on a smaller polydisc by Cauchy's formula.
Finally
\[
 f_{L,\mu}(z)=L^{-2}\sum_{r,t=0}^{L-1}
             \cos((r-t)z_\mu/L).
\]
Use $|(r-t)/L|\le1$ and
$|\cos w-1|\le |w|^2e^{|w|}/2$ to obtain
\cref{eq:v5-fclose}.
\end{proof}

\section{Exact gauge cancellation and a level-independent complex strip}
\label{sec:v5-strip}

Let
\begin{align}
 \alpha_{L,\mu}&=P_L f_{L,\mu},\qquad
 e_{L,\mu}=\alpha_{L,\mu}+s_Lb_{L,\mu},\notag\\
 h_L&=P_L^{-1}-\sum_{\mu=1}^4
       \frac{q_\mu q^\sharp_\mu b_{L,\mu}}
                         {\alpha_{L,\mu}e_{L,\mu}},\qquad
 D_e=\diag_\mu(e_{L,\mu}^{-1}).\label{eq:v5-eh}
\end{align}
These are defined without dividing by $s_L$.  In particular a complex
null direction of the quadratic form $\sum z_\mu^2$ is not excluded.

\begin{proposition}[Holomorphic soft-kernel formula]
\label{prop:v5-cancel}
On the common polydisc $\|z\|_\infty<r_1=1/64$,
\begin{equation}\label{eq:v5-softkernel}
 \boxed{\displaystyle
          S_L(z)=s_L(z)D_e(z)
                    -D_e(z)q(z)h_L(z)^{-1}q^\sharp(z)D_e(z)}
\end{equation}
is holomorphic, agrees with $S^{(j)}$ at real nonzero momentum, and
satisfies $S_Lq=0=q^\sharp S_L$.  On this polydisc,
\begin{equation}\label{eq:v5-soft-bounds}
 |e_{L,\mu}-1|<\tfrac14,\qquad |h_L-1|<\tfrac14,
       \qquad \|S_L(z)\|\le64\|z\|_\infty^2
\end{equation}
uniformly in $L$.
\end{proposition}
\begin{proof}
For real nonzero $z$, if $D=\diag(\gamma_{L,\mu})$, the inverse of
$\Pi D\Pi$ on the transverse space, extended by zero on $q$, is
\begin{equation}\label{eq:v5-quotient-inverse}
 D^{-1}-D^{-1}q(q^\sharp D^{-1}q)^{-1}q^\sharp D^{-1}.
\end{equation}
Indeed it annihilates $q$, and multiplication by $\Pi D\Pi$ gives
$\Pi$ on that space.  The key exact identity is
\begin{equation}\label{eq:v5-ward-factor}
 \sum_\mu\frac{q_\mu q^\sharp_\mu}{\alpha_{L,\mu}}
       =\frac{s_L}{P_L},\qquad
       q^\sharp D_eq=s_Lh_L.
\end{equation}
The first equality follows componentwise from
$q_\mu q^\sharp_\mu/f_{L,\mu}
 =L^2(2-2\cos(z_\mu/L))$.  This geometric-sum identity extends
analytically through coordinate zeros.  For the second use
\[
 e_{L,\mu}^{-1}
       =\alpha_{L,\mu}^{-1}
              -\frac{s_Lb_{L,\mu}}{\alpha_{L,\mu}e_{L,\mu}}.
\]
Since $D^{-1}=s_LD_e$, substitution in
\cref{eq:v5-quotient-inverse} cancels the entire soft scalar and yields
\cref{eq:v5-softkernel} on its original domain.  Multiplication by $q$
or $q^\sharp$ then proves both Ward identities using
\cref{eq:v5-ward-factor}.

For completeness, the following conservative bounds make the common
soft radius explicit.  Put $t=\|z\|_\infty\le r_1$.  By
\cref{eq:v5-fclose}, $|f_\mu-1|\le t^2$,
$|P_L-1|\le8t^2$, and $|\alpha_{L,\mu}-1|\le10t^2$.
Also $|s_L|\le8t^2$ and
$\sum_\mu|q_\mu q^\sharp_\mu|\le8t^2$.  Thus
\[
 |e_{L,\mu}-1|\le26t^2<\tfrac14.
\]
Since $|\alpha_\mu|,|e_\mu|\ge3/4$ and $|P_L^{-1}-1|\le16t^2$,
\[
 |h_L-1|\le16t^2+\frac{256}{9}t^2<45t^2<\tfrac14.
\]
Finally $\|q\|,\|q^\sharp\|\le4t$.  The first term of
\cref{eq:v5-softkernel} is at most $32t^2/3$, and the second at most
$1024t^2/27$.  Their sum is below $64t^2$.
Every denominator in the displayed formula is therefore uniformly
nonzero on the stated polydisc, including where $s_L=0$.  This proves
holomorphy and the estimates.
\end{proof}

\begin{theorem}[A common complex strip for the complete Gaussian kernels]
\label{thm:v5-strip}
There are constants $\delta_S>0$ and $M_S<\infty$, independent of
$j\ge0$, such that $S^{(j)}$ extends to a periodic holomorphic matrix
on
\[
 \mathcal T_{\delta_S}
       =\{z\in\mathbb C^4:\|\operatorname{Im}z\|_\infty<\delta_S\},
 \qquad
       \sup_{j,z\in\mathcal T_{\delta_S}}\|S^{(j)}(z)\|\le M_S.
\]
Its restriction near each soft point is \cref{eq:v5-softkernel}.
The limiting kernel has the same type of bounded holomorphic extension
on every narrower strip.  Consequently all fixed-order momentum
 derivatives of the kernels have bounds independent of $j$.
\end{theorem}
\begin{proof}
The soft neighbourhood is already uniform.  It remains to control a
closed real region at fixed positive distance $\varepsilon$ from
$2\pi\mathbb Z^4$, with $\varepsilon$ smaller than that soft
neighbourhood.  In a centered real chart, every alias satisfies
\[
 \lambda((x+2\pi M)/L)
                    \ge\frac{4}{\pi^2L^2}|x|^2.
\]
The reason is that in each coordinate its distance to $2\pi\mathbb Z$
is at least $|x_\mu|/L$, followed by the elementary sine bound.
For $z=x+\mathrm iy$ with $|x|\ge\varepsilon$ and
$\|y\|_\infty\le a$, the real part of its denominator is therefore
at least $c\varepsilon^2/L^2$, uniformly in $M,L$, when
$a\le c'\varepsilon$ is small enough.  This follows from the same
$O(a^2/L^2)$ estimate used in \cref{lem:v5-tail}.
Combining it with \cref{lem:v5-weight} gives a bound
\[
                    |\gamma_{L,\mu}(z)|\le C_\varepsilon
\]
on a common complex collar.  Taking a slightly larger real region and
a slightly wider collar first, Cauchy's estimate also gives uniform
first derivatives on a smaller collar.

The coarse $\lambda(z)$ is nonzero on a sufficiently thin such collar,
so $\Pi(z)=I-q(z)q^\sharp(z)/\lambda(z)$ is holomorphic and bounded
there.  Put
\[
        B_L(z)=\Pi(z)\diag(\gamma_{L,\mu}(z))\Pi(z)+I-\Pi(z).
\]
On the real region, V4's lower bound for the diagonal factors gives
\[
        B_L(x)\succeq\min\{1,m_0/16\}I=(m_0/16)I,
                      \qquad m_0=(2/\pi)^{12}.
\]
The preceding uniform holomorphic bounds give a uniform derivative
bound $C_1$ for $B_L$ in a smaller collar.  Thus
$\|B_L(x+\mathrm iy)-B_L(x)\|\le4C_1\|y\|_\infty$.
Choose a further common width so this is at most $m_0/32$.
A Neumann series then inverts every $B_L(z)$ there with bound
$32/m_0$.  The expression
$\Pi(z)B_L(z)^{-1}\Pi(z)$ is the required uniformly bounded kernel
on this collar.

It agrees on real overlap with \cref{eq:v5-softkernel}.  Holomorphic
uniqueness, applied successively to the four variables on an open real
set, makes the continuations agree on smaller overlapping charts.
Alias permutations give the same matrices at periodic boundaries.
A finite real cover, whose radii and bounds above are independent of
$L$, produces the asserted periodic strip and uniform bound.

For the limit, choose $0<\sigma<\delta_S$.  Periodic contour shifts
give Fourier coefficient bounds $M_Se^{-\sigma|n|_1}$.
The inherited uniform real convergence makes each Fourier coefficient
converge.  On any narrower strip these coefficients form an absolutely
and uniformly convergent series, so their limit is holomorphic and
agrees with the known real $S^*$.  Cauchy's formula on polydiscs inside
a common strip gives, for every multi-index $\alpha$, a bound of the
form $C\alpha! r^{-|\alpha|}$, independent of $j$.
\end{proof}

The theorem proves existence of a common positive strip width, with a
constructive sequence of estimates.  It does not claim a numerically
optimized or interval-enclosed value for $\delta_S$.  Its proof uses
the explicit complex alias weights and the Ward cancellation; it is
not an inference from real ellipticity alone.

\section{The exact all-level fourth momentum correction}
\label{sec:v5-quartic}

Let $D_{1/2}(z)=\diag(e^{\mathrm iz_\mu/2})$ and set
\[
 \overline S_L(z)=D_{1/2}(z)^{-1}S_L(z)D_{1/2}(z),\qquad
 r_\mu(z)=2\sin(z_\mu/2).
\]
This local midpoint frame removes edge phases: $q=\mathrm iD_{1/2}r$
and $q^\sharp=-\mathrm ir^{\mathsf T}D_{1/2}^{-1}$.
For the bare kernel,
$\overline S_1=\lambda I-rr^{\mathsf T}$.
Write $Z_2=\sum_\mu z_\mu^2$, $Z_4=\sum_\mu z_\mu^4$, and define
\begin{equation}\label{eq:v5-Q4}
 \mathcal P_{\mu\nu}(z)
 =\delta_{\mu\nu}\{Z_4+Z_2^2+Z_2z_\mu^2\}
                 -z_\mu z_\nu\{Z_2+z_\mu^2+z_\nu^2\}.
\end{equation}
This is a homogeneous fourth-degree momentum matrix.

\begin{theorem}[Universal fourth momentum jet and uniform remainder]
\label{thm:v5-quartic}
For all dyadic $L=2^j$,
\begin{equation}\label{eq:v5-quartic}
 \boxed{\displaystyle
  \overline S_L(z)=\overline S_1(z)
       +\frac{1-L^{-2}}{12}\mathcal P(z)
       +O(\|z\|_\infty^6).}
\end{equation}
The remainder bound holds on a fixed smaller polydisc with a constant
independent of $L$.  In particular every midpoint-frame momentum jet
through degree three equals the bare one, and the fourth correction
converges with the exact coefficient difference $L^{-2}/12$.
The limiting kernel obeys the same expansion with coefficient $1/12$.
For real $z$, $\mathcal P(z)\succeq0$ and $\mathcal P(z)z=0$.
\end{theorem}
\begin{proof}
Put $a_L=(1-L^{-2})/12$.  Geometric-sum expansion, with uniform
remainders by the preceding holomorphic bounds, gives
\[
 f_{L,\mu}=1-a_Lz_\mu^2+O(\|z\|^4),\qquad
 P_L=1-a_LZ_2+O(\|z\|^4),\qquad
 s_L=Z_2-\frac{Z_4}{12L^2}+O(\|z\|^6).
\]
Here and below the norm in local remainders is $\|\cdot\|_\infty$.
By \cref{lem:v5-tail}, $b_\mu=O(\|z\|^2)$.  Therefore
\[
 e_\mu^{-1}=1+a_L(Z_2+z_\mu^2)+O(\|z\|^4),\qquad
 h_L^{-1}=1-a_LZ_2+O(\|z\|^4).
\]
Insert these expansions in the midpoint form of
\cref{eq:v5-softkernel}:
\[
 (\overline S_L)_{\mu\nu}
 =\delta_{\mu\nu}\frac{s_L}{e_\mu}
                 -\frac{r_\mu r_\nu}{e_\mu e_\nu h_L}.
\]
Use $s_L-\lambda=a_LZ_4+O(\|z\|^6)$ and
$r_\mu=z_\mu+O(z_\mu^3)$.  The fourth-degree difference is exactly
$a_L\mathcal P_{\mu\nu}$.  All scalar factors are even in each
coordinate and $r$ is odd, so the matrix has no odd total degree in
this frame.  The common holomorphic bound and Cauchy estimates give a
uniform sixth-order remainder after subtracting the computed Taylor
polynomial.  Passage to the limiting kernel is justified by the uniform
convergence on a smaller polydisc proved above.

Direct multiplication shows $\mathcal P(z)z=0$.  For real $z$ and a
real $v$ with $v\cdot z=0$,
\[
 v^{\mathsf T}\mathcal P(z)v
       =(Z_4+Z_2^2)|v|^2+Z_2\sum_\mu z_\mu^2v_\mu^2\ge0.
\]
The real symmetric matrix has $z$ in its kernel, so decomposing any
vector into a multiple of $z$ and its orthogonal complement proves
positive semidefiniteness.  The case $z=0$ is immediate.
\end{proof}

For an axial real momentum $z=te_1$ and a transverse component,
\[
 (\overline S_L-\overline S_1)_{22}
             =\frac{1-L^{-2}}6t^4+O(t^6),
\]
which agrees with the exact all-level axial inverse in V2.  The new
formula retains all four momentum components and all link polarizations.
It is a statement about the \emph{quadratic Gaussian action}, not the
third or fourth field derivatives of the nonlinear action.  Its positive
quartic correction is not a mass: the leading retained energy remains
quadratic in momentum and zero at the origin.

\section{Uniformly local fibres, harmonic lifts, and measured borders}
\label{sec:v5-packet}

Retain the actual parity matrices $J,U,E,R,L(k)$ and $A_j,H_j$ of
\cref{sec:v4-fibre}.  For complex momentum all adjoints in their
formulas are holomorphic adjoints.  Define
\begin{align}
 \mathscr C_j&=E H_j^{-1}E^\sharp,\qquad
 \mathscr T_j=R-\mathscr C_j A_jR,\label{eq:v5-CT}\\
 \mathscr B_j&=\begin{pmatrix}A_j&L^\sharp\\L&0\end{pmatrix}.
 \notag
\end{align}
Here $\mathscr C_j$ is the conditional covariance on the 49-coordinate
tree carrier and $\mathscr T_j$ is its retained harmonic lift.  Neither
is an unrestricted retained massless covariance.

\begin{theorem}[Uniform complex control of the full fixed-fibre packet]
\label{thm:v5-packet}
There exist $\delta>0$ and $M<\infty$, independent of $j$, such that
all matrices
\begin{equation}\label{eq:v5-packet}
  S^{(j)},\quad A_j,\quad H_j,\quad H_j^{-1},\quad
  \mathscr C_j,\quad\mathscr T_j,\quad
  \mathscr B_j^{-1}
\end{equation}
are periodic holomorphic symbols on $\mathcal T_\delta$, with norm at
most $M$.  In a possibly narrower common strip,
\begin{equation}\label{eq:v5-complex-H}
             \|H_j(z)^{-1}\|\le \frac{2}{\kappa},
                          \qquad \kappa=24/11635.
\end{equation}
The measured border has the exact inverse
\begin{equation}\label{eq:v5-border-inverse}
 \boxed{\displaystyle
   \mathscr B_j^{-1}
    =\begin{pmatrix}
         \mathscr C_j&\mathscr T_j\\
         \mathscr T_j^\sharp&-S^{(j+1)}
      \end{pmatrix}.}
\end{equation}
All assertions include the real soft momenta and arbitrary fixed colour
multiplicity.
\end{theorem}
\begin{proof}
The alias substitution $k\mapsto k/2+\pi\epsilon$ halves the imaginary
part.  The finitely many entries of $U,J,E,R,L$ are exponentials or
fixed Laurent polynomials, bounded on each fixed strip.  Thus
\cref{thm:v5-strip} supplies a common bounded holomorphic strip for
$A_j$ and $H_j$, with no growth in $j$.  Cauchy's estimate on a narrower
strip supplies a uniform first-derivative bound $C_H$.  V4 gives
$H_j(x)\succeq\kappa I$ for real $x$.  For
$\|y\|_\infty\le\kappa/(8C_H)$, narrowed to remain in that strip,
\[
       \|H_j(x+\mathrm iy)-H_j(x)\|\le\kappa/2.
\]
Factoring by $H_j(x)$ and using a Neumann series proves
\cref{eq:v5-complex-H}.  Products in \cref{eq:v5-CT} are now uniformly
bounded and holomorphic.

To prove the border identity on real momenta, solve
\[
             A_jx+L^*\eta=u,\qquad Lx=v.
\]
Write $x=Rv+Ez$.  Multiplication of the first equation by $E^*$ gives
\[
 z=H_j^{-1}E^*(u-A_jRv),\qquad
 x=\mathscr C_ju+\mathscr T_jv.
\]
Since $R^*L^*=I$, the multiplier is
\[
 \eta=R^*(u-A_jx)
       =\mathscr T_j^*u-R^*A_j\mathscr T_jv
       =\mathscr T_j^*u-S^{(j+1)}v.
\]
The last equality is V4's exact Schur identity.  This proves
\cref{eq:v5-border-inverse} without inverting $A_j$.  The formula
continues holomorphically, and its bounded factors give the remaining
claim.  The periodic cell representation follows from alias
permutations exactly as in V4.  Tensoring with a fixed colour identity
changes the bookkeeping dimension but not the inverse argument.
\end{proof}

The border remains indefinite; it is not a covariance of extra physical
variables.  Its formula nevertheless shows that the zero-background
linear solves needed by the \emph{measured} V4 transformation have a
common analytic domain.  This is stronger than a scalar bound on
$\|H_j^{-1}\|$ and does not require an ambient inverse of the retained
massless kernel.

\section{Exponential spatial locality and summable convergence}
\label{sec:v5-locality}

For a matrix-valued translation-invariant kernel on the current cell
lattice, write
\[
 F(k)=\sum_{n\in\mathbb Z^4}\widehat F(n)e^{\mathrm ik\cdot n},
 \qquad
 \|F\|_a=\sum_{n\in\mathbb Z^4}
                      e^{a|n|_1}\|\widehat F(n)\|.
\]
The norm on finite matrices is the operator norm.  This weighted
convolution norm is submultiplicative because
$|n|_1\le|m|_1+|n-m|_1$.  Positions here are in \emph{current rescaled
cell units}.  A statement uniform in those units is not a physical
mass statement at fixed microscopic distance.

\begin{theorem}[Uniform locality and convergence in weighted kernel norms]
\label{thm:v5-locality}
For each family $F_j$ in \cref{eq:v5-packet}, there are $\sigma>0$,
$M_F,C_F<\infty$, independent of $j$, such that
\begin{equation}\label{eq:v5-locality}
             \|\widehat F_j(n)\|\le M_Fe^{-\sigma|n|_1}.
\end{equation}
The limiting symbol $F_*$ exists.  For any $0<\theta<1$ and
$0<a<(1-\theta)\sigma$,
\begin{equation}\label{eq:v5-local-convergence}
 \boxed{\displaystyle
       \|F_j-F_*\|_a\le C_{a,\theta,F}\,4^{-\theta j}.}
\end{equation}
The same estimate holds with an additional factor $|n|_1^r$ in the
weighted sum, for each fixed integer $r\ge0$.  In particular all fixed
momentum derivatives converge geometrically and their differences are
summable over $j$.  Periodization onto any finite cell torus preserves
uniform unweighted row-sum bounds and the corresponding bounds with
periodic distance.
\end{theorem}
\begin{proof}
Shift each of the four Fourier contours in the sign of $-n_\mu$, by
any fixed width $\sigma<\delta$ from \cref{thm:v5-packet}.  Periodicity
cancels the end faces and holomorphy permits the shifts.  The resulting
factor is $e^{-\sigma|n|_1}$, proving \cref{eq:v5-locality}.

For $S^{(j)},H_j,H_j^{-1}$ the uniform real convergence rate
$C_F4^{-j}$ is inherited from V4.  The bounded products defining
$A_j,\mathscr C_j,\mathscr T_j$ and the displayed border inverse give
the same real rate for all other families, using
$XY-X_*Y_*=(X-X_*)Y+X_*(Y-Y_*)$.  Their Fourier coefficients therefore
converge to those of the stated real limits.  The uniform exponential
bound constructs their holomorphic limits on a smaller strip, as in
\cref{thm:v5-strip}.

Let $\epsilon_j=C_F4^{-j}$.  A coefficient difference satisfies both
the real integral bound and the complex contour bound:
\[
 \|\widehat F_j(n)-\widehat F_*(n)\|
       \le\min\{\epsilon_j,2M_Fe^{-\sigma|n|_1}\}
       \le\epsilon_j^\theta(2M_F)^{1-\theta}
                              e^{-(1-\theta)\sigma|n|_1}.
\]
Summation gives the explicit bound
\begin{equation}\label{eq:v5-sumconstant}
 \|F_j-F_*\|_a\le C_F^\theta(2M_F)^{1-\theta}4^{-\theta j}
       \left(\frac{1+e^{-b}}{1-e^{-b}}\right)^4,
                         \qquad b=(1-\theta)\sigma-a>0.
\end{equation}
Inserting $|n|_1^r$ leaves a finite exponential sum and proves the
higher-derivative assertion.  Since $4^{-\theta}<1$, the errors are
summable in $j$.

For a finite torus, periodization sums $\widehat F(n+N\ell)$ over
integer wraps.  Its periodic distance is no greater than
$|n+N\ell|_1$.  Summing absolute values therefore costs no more than
the corresponding infinite-lattice weighted sum.  This proves the
volume-uniform assertions without an additional compactness hypothesis.
\end{proof}

For example, choosing $\theta=1/2$ and $a<\sigma/2$ gives a fully
summable $2^{-j}$ locality-norm error.  The smaller rate compared with
the real-axis $4^{-j}$ is the explicit price of this interpolation;
no unproved derivative-level $4^{-j}$ rate is asserted.

\begin{corollary}[A volume-independent perturbation radius for the border]
\label{cor:v5-perturb}
Fix $a$ for which
$\sup_j\|\mathscr B_j^{-1}\|_a\le C_a$.  For a perturbation $V$ in
the same convolution algebra with
$\|V\|_a\le(2C_a)^{-1}$,
\begin{equation}\label{eq:v5-perturb}
 \|(\mathscr B_j+V)^{-1}\|_a\le2C_a,\qquad
 \|(\mathscr B_j+V)^{-1}-\mathscr B_j^{-1}\|_a
                                      \le2C_a^2\|V\|_a.
\end{equation}
The same statement holds for matrices indexed by finite cell volumes,
using the weighted maximum of row and column sums, whenever the
perturbation is measured in that algebra and the unperturbed bound is
converted to its fixed-component norm.
\end{corollary}
\begin{proof}
Factor $\mathscr B_j+V=\mathscr B_j(I+\mathscr B_j^{-1}V)$.
The submultiplicative norm of $\mathscr B_j^{-1}V$ is at most $1/2$, so the
Neumann series converges in the same algebra and gives both estimates.
For the finite-volume row/column norm, the triangle inequality for
periodic distance proves submultiplicativity.  Fixed component dimension
only changes the common conversion constant.
\end{proof}

The corollary is a radius in \emph{operator perturbation norm}, not yet
a uniform physical field radius.  Turning a field bound into this norm
requires the very nonlinear vertex estimates still to be acquired.
It does, however, prove that no additional Gaussian spatial inverse
loss is hidden in that step.

\section{What is now available to the nonlinear continuation}
\label{sec:v5-next}

The level-uniform linear analytic inputs to V4's measured transformation
are now populated for the actual full tower.  Its conditional covariance,
harmonic lift, fibre inverse, and measured border inverse have common
momentum domains, exponential spatial bounds, and summably convergent
coefficients.  These conclusions do not rely on retaining only a
polarization subspace, enlarging a finite response card, or treating a
pointwise covariance column as a global coordinate system.

For a fixed finite expression made from these matrices, fixed finite-range
vertex kernels, their translates $k+p$, and finitely many derivatives,
the product rule and Cauchy bounds now give common bounds whenever all
complex momentum arguments stay in the common strip.  Normalized traces
satisfy $|\Tr A|\le d\|A\|$ at the fixed component dimension $d$;
there is no hidden multiplication by the number of cells in a trace per
cell.  Products of a fixed number of such factors also inherit summable
convergence from \cref{thm:v5-locality}.  This is the usable consequence
for a fixed specified diagram, not a bound for a growing unclassified
family of nonlinear diagrams.

\paragraph{The remaining nonlinear quantity is not covered by those words.}
After exact blocking, the cubic and quartic action tensors are themselves
transported objects, not fixed primitive kernels.  In V4's notation,
\[
 \mathcal G_3[v^3]=S_3[a^3]+3\langle\eta_1,F_2[a,a]\rangle,
\]
\[
 \mathcal G_4[v^4]=S_4[a^4]+4\langle\eta_1,F_3[a^3]\rangle
                                -3\langle\zeta,K^{-1}\zeta\rangle.
\]
This version controls the linear maps generating $a,\eta_1$ and the
inverse $K^{-1}$ in spatial and momentum norms.  It does not bound the
size of the evolving $S_3,S_4$, their source/colour contractions, or the
second log-density tensor under arbitrarily many nonlinear iterations.
An inequality with an unspecified vertex bound inserted on its right
side would merely restate that obligation and is not called a closure
result here.

The next decisive acquisition is therefore a bound on the \emph{actual
complete nonlinear update} in an anchored locality norm, after the
retained gauge constraints and the marginal field-strength normalization
are separated.  It must use the fixed root-path second and third jets
and the inherited full one-step measured response, not reset the Wilson
vertices or the Haar amplitude.  A useful first calculation would bound
the bilinear forcing $\zeta$ and the resulting quartic exchange term in
that norm, with the factors of two from cell rescaling visible.  The
linear inverses and their spatial tails no longer need to be left as
unknowns in that calculation.  Mere boundedness of those inverses still
does not imply that the nonlinear renormalization map contracts.

\begin{scope}[Physical boundary after V5]
The new all-level theorem is about the Gaussian principal tower and its
linear stationary solves.  The retained kernel remains zero at zero
momentum, and its leading order remains the massless curl kernel.
Exponential locality of an \emph{action kernel} or a conditional
innovation is not exponential decay of unrestricted physical
correlations.  This version does not identify the complete interacting
renormalization coefficient, construct the interacting continuum,
prove a common nonlinear small-field neighbourhood, or establish the
fixed-physical-time contraction required for the Yang--Mills mass gap.
The remaining boundary is the specified nonlinear and physical infrared
estimate, not a historical objection to attempting the problem.
\end{scope}

\section{Verification and append-only record}
\label{sec:v5-verification}

The new checker \path{verify_ym_uniform_locality_V5.py} uses rational,
Gaussian-rational, and symbolic polynomial arithmetic.  Its formal jets
store ordinary coefficients through degree six, not numerical derivatives.
It reconstructs the finite alias sums, separates the principal alias,
and evaluates the regular soft formula without ever forming its divergent
covariance.  The test directions include a complex direction with
$\sum_\mu z_\mu^2=0$; holomorphic reversal of phases, not numerical
complex conjugation, is used there.

The executed report contains thirteen passing test families.  They
include seven actual dyadic/directional germs, all sixteen entries of
the fourth-momentum formula, the two Ward identities as exact series,
midpoint symmetry and vanishing odd coefficients, the quadratic order
of the nonprincipal tails, the exact soft-denominator factorization,
the transverse positive-form identity for $\mathcal P$, and independent
full $53$-by-$53$ border-inverse checks at a soft and a nonsoft actual
parity momentum.  The output contains the exact fourth and sixth
coefficients of the checked kernels.  The symbolic quartic formula is
proved above for all levels and directions; the finite cases are
implementation checks, not an inference from sampling.

The soft radius $1/64$ is justified analytically in
\cref{prop:v5-cancel}.  The global strip width and locality constants
are existence constants supplied by the proof, not purported numerical
enclosures.  No finite list of complex samples is used to prove a
uniform strip, and no Lean formalization is claimed.  Prior executable
materials and their original reports are retained separately; the
provenance report identifies any additional fresh reverifications.

The entire V1--V4 mathematical body, including its end markers and
bibliographies, is retained byte-for-byte.  Only the displayed version
in the preamble changes.  This appendix is inserted before the final
\verb|\end{document}|.  A subsequent version is to append after the
marker below, stating any necessary correction explicitly rather than
silently editing an earlier proof.

% END OF VERSION 5 MATHEMATICAL BODY.
% Append subsequent requested versions below this marker; retain V1--V5 above.

\clearpage
\section{V6: normalized remainders for nonlinear transport}
\label{sec:v6-context}

The new comparison manuscripts suggest replacing a coefficient-by-coefficient
remainder campaign by a differentiated, normalized shell estimate.  This is
adopted, with two qualifications which determine the actual theorem needed.
First, the supplied SM--Einstein estimate controls one retained-field derivative,
whereas a density Hessian requires two and a fourth-action estimate requires
four.  Second, absolute moments of an extensive centered interaction need not
be uniform in volume even for independent sites.  The appropriate replacement
is a \emph{connected, locally anchored} estimate, before taking absolute values.

The Gaussian inputs in V5 are retained, not reproved: its conditional covariance,
harmonic lift and measured border have uniform spatial bounds.  The new work
constructs an exact measured interpolation, specializes the inherited normalized
score hierarchy to its third strength derivative, and proves a rooted locality
estimate through four retained derivatives.  Five-entry connected moments suffice
for the Hessian remainder and seven suffice for the fourth derivative.  An
additional theorem reduces those connected bounds to a source-restricted covariance
estimate for products and local fourteenth moments.  Thus a separate all-order
analytic source expansion is not required.  A correlated nonlinear comparison
over the actual V5 Gaussian fibre family and a stable quartic product family
populate the estimate on their whole real interpolation intervals.

\paragraph{Source audit and nonduplication.}
The first-derivative shell identity and its one-sided stable-tilt qualification
are inherited from \cite[Section ``Differentiated shell remainders in the
comparison norm'']{SM6}.  The all-order normalized score differentiation rule is
already in the YM closure monograph, including its warning that an identity is
not a norm estimate \cite[Section ``Higher normalized responses'']{YM4V6}.
Neither is advertised below as a new general identity.  The new acquisition is
an exact measured-law implementation and a volume-independent finite-order
locality estimate with an explicit number of marks, reference-score terms,
and transport factors.

The five newly supplied manuscripts were structurally inspected and their
relevant arguments read: normalized shells and full-face rigidity in
\cite{SM6}, source shorts in \cite{GT76V6}, commutant stability in
\cite{DualV6}, residual/leakage forcing in \cite{ClassV6}, and normalized
pressure reconstruction in \cite{LorV6}.  This is a dependency audit, not a
line-by-line recertification of those manuscripts.  The Library copy of
\cite{YM4V6} was also read where it incorporates f13 V17.  It records genuine
active-source repair and still excludes the general transition-cutoff,
high-occupancy and other noncentral configurations from that certificate.
The standalone V17 proof archive was not independently retrieved or rerun.
The present version therefore makes no new coverage claim for that line.

Clustering estimates for unbounded systems are established constructive tools;
\cite{APSV6} is external methodological context.  No estimate from that paper
is imported as a proved property of the present compact YM law.  The finite
lemmas below are proved here; their application hypotheses remain visible.

\section{The exact residual must come from the measured law}
\label{sec:v6-exact}

The finite jet transformation in V4 determines a fourth classical jet and a
second one-loop jet.  It does not determine the full non-Gaussian interaction.
To avoid mistaking a formal polynomial for an exact density, start from the
actual independent fibre chart of the compact block.  In a fixed finite volume
write its measured density as
\[
                 e^{-f(y,c)/\epsilon^2+\ell(y,c)}\,dy,
                 \qquad \epsilon>0.
\]
Here $c$ is the retained field, $y$ the independent fibre, and $\ell$ includes
its \emph{actual} coordinate-Haar and inverse-pivot Jacobian.  Let $y_*(c)$ be
one specified smooth local minimum, put
\[
 G(c)=f(y_*(c),c),\qquad H(c)=f_{yy}(y_*(c),c)\succ0,
 \qquad \ell_*(c)=\ell(y_*(c),c),
\]
and let $d$ be the real fibre dimension.  The hypotheses here are local on a
common finite retained domain; their uniformity under infinitely many nonlinear
steps is not inferred from the Gaussian theorem.

Choose a nonnegative cutoff $\chi_\epsilon(\eta)$, independent of $c$, with
positive Gaussian mass, such that the displayed chart is defined on its support
under $y=y_*(c)+\epsilon\eta$.  A compactly supported cutoff makes this a
well-defined local contribution.  The uncut choice is permissible only when
that exact chart and integral are justified globally.  Define
\begin{align}
 A(c,\eta)&=\tfrac12\langle\eta,H(c)\eta\rangle,\qquad
 N_\chi(c)=\int \chi_\epsilon(\eta)e^{-A(c,\eta)}\,d\eta,\notag\\
 d\nu_c(\eta)&=N_\chi(c)^{-1}\chi_\epsilon(\eta)e^{-A(c,\eta)}\,d\eta,
                                                        \label{eq:v6-reference}\\
 Z_\epsilon(c,\eta)&=\epsilon^{-1}\left\{
 \frac{f(y_*(c)+\epsilon\eta,c)-G(c)}{\epsilon^2}
 -\tfrac12\langle\eta,H(c)\eta\rangle\right.\notag\\
 &\hspace{6em}\left.-\ell(y_*(c)+\epsilon\eta,c)+\ell_*(c)\right\}.\notag
\end{align}

\begin{proposition}[Exact measured local factorization]
\label{prop:v6-exact}
The local contribution defined by the cutoff obeys
\begin{equation}\label{eq:v6-exact}
 \boxed{\displaystyle
 I_\chi(c)=\epsilon^d e^{-G(c)/\epsilon^2+\ell_*(c)}N_\chi(c)
             \E_{\nu_c}e^{-\epsilon Z_\epsilon(c,\eta)}.}
\end{equation}
For fixed $\epsilon$, the normalized interpolation
\begin{equation}\label{eq:v6-tilt}
 W_u(c)=-\log\E_{\nu_c}e^{-uZ_\epsilon(c,\eta)},\qquad
 d\mu_{u,c}=\frac{e^{-uZ_\epsilon(c,\eta)}d\nu_c}
                    {\E_{\nu_c}e^{-uZ_\epsilon(c,\eta)}},\quad 0\le u\le\epsilon
\end{equation}
has the exact local shell at its endpoint.  The parameter $\epsilon$ in
$Z_\epsilon$ is held fixed when differentiating $u$.
\end{proposition}
\begin{proof}
Substitute $y=y_*(c)+\epsilon\eta$ in the fibre integral; its Jacobian is
$\epsilon^d$.  The definition of $Z_\epsilon$ splits its exponent exactly
into the constant, quadratic and residual terms in \cref{eq:v6-exact}.
Division by $N_\chi(c)$ proves the normalized identity.  No Taylor truncation
has entered it.
\end{proof}

For an uncut Gaussian reference,
$N_1=(2\pi)^{d/2}\det H(c)^{-1/2}$.  More generally, if
$p_\chi(c)=\E_{N(0,H(c)^{-1})}\chi_\epsilon$, then
\begin{equation}\label{eq:v6-cutoff}
 -\log I_\chi=\frac{G}{\epsilon^2}
 -d\log\epsilon-\frac d2\log(2\pi)
 +\frac12\log\det H-\ell_*-\log p_\chi+W_\epsilon.
\end{equation}
Thus V4's measured one-loop determinant is recovered in its stated convention,
and the cutoff acceptance and its retained derivatives are not dropped.
A partition into such local contributions and their complement must be made in
the original integral.  An estimate on $I_\chi$ alone does not control that
complement.  A cutoff depending explicitly on $c$ would introduce additional
score derivatives and is not covered by suppressing them in
\cref{eq:v6-reference}.

For fixed finite jets the residual begins as
\begin{equation}\label{eq:v6-leadingZ}
 Z_\epsilon=\frac16 f_{yyy}(y_*,c)[\eta^3]-\ell_y(y_*,c)[\eta]
 +\epsilon\left\{\frac1{24}f_{yyyy}(y_*,c)[\eta^4]
                   -\frac12\ell_{yy}(y_*,c)[\eta^2]\right\}+O(\epsilon^2).
\end{equation}
This is a local expansion, not a definition of the exact $Z_\epsilon$.
It must not be extended as a naked cubic interaction on an unbounded Gaussian
space: its exponential integral need not exist.  Nor does a bound on the
Gaussian inverse control the retained derivatives of every term in this exact
residual.  The proposed moment norm measures those derivatives under the law;
it does not declare them absent.

\begin{example}[Finite compiler data do not determine the exact shell]
\label{ex:v6-not-exact}
Take $F(y,c)=c$, $\ell=0$ and
$f_\lambda(y,c)=y^2/2+\lambda c^2y^4$, $\lambda\ge0$.  All these actions have
the same fourth jet at $(0,0)$, the same block, the same minimizing branch
$y_*=0$, and $H(c)=1$.  V4's stated classical and one-loop output jets are
therefore identical.  The exact uncut shell nevertheless has
\[
 W_\epsilon(c)=-\log\E_{N(0,1)}e^{-\lambda\epsilon^2c^2\eta^4},
 \qquad W_\epsilon''(0)=6\lambda\epsilon^2.
\]
Indeed $Z_\epsilon=\lambda\epsilon c^2\eta^4$ and $\E\eta^4=3$.
The omitted term is of higher order in the fine action and of higher loop
order in this example.  This is not a YM counterexample; it proves why the
exact interpolation must be constructed from the full measured law rather
than recovered from the compiler's finite output.
\end{example}

\section{The required normalized hierarchy has five or seven entries}
\label{sec:v6-hierarchy}

We give the result for a moving reference, since recentering the actual shell
normally gives precisely that situation.  Let
\[
 d\nu_x=N(x)^{-1}e^{-A(x,\eta)}\,dm(\eta),\qquad
 W_u(x)=-\log\E_{\nu_x}e^{-uZ(x,\eta)}.
\]
The common underlying measure $m$ can include the fixed cutoff.  All indicated
finite products and derivatives are assumed integrable, uniformly on compact
subdomains in $x$ and $0\le u\le u_0$.  In particular assume the dominated
passages needed through three strength derivatives and four retained
 derivatives.  No two-sided exponential moment of a quartic source is assumed.

Write $\kappa_{u,x}$ for the finite joint cumulants, defined by
\begin{equation}\label{eq:v6-cumulant}
 \kappa(Y_1,\ldots,Y_m)
 =\sum_{\pi\in\Pi([m])}(-1)^{|\pi|-1}(|\pi|-1)!
                  \prod_{B\in\pi}\E\prod_{i\in B}Y_i.
\end{equation}
Cumulants of order at least two vanish when any argument is constant.
For a set $I$ of labelled commuting retained derivatives, put
$Z_I=\partial_I Z$, $Z_\varnothing=Z$, and
\begin{equation}\label{eq:v6-score}
                         P_I=A_I+uZ_I\qquad(I\ne\varnothing).
\end{equation}
Only the nonconstant part of $P_I$ is relevant.  In
\cref{eq:v6-reference}, $A_I=\frac12\eta^T(\partial_IH)\eta$; it is not
legitimate to differentiate that moving Gaussian as though $A_I=0$.
A transport to a fixed reference is an alternative, provided its entire
Radon--Nikodym derivative is included once.

\begin{lemma}[Three interaction roots in the inherited score hierarchy]
\label{lem:v6-marks}
One has $\partial_u^3W_u=\kappa_{u,x}(Z,Z,Z)$, and
\begin{equation}\label{eq:v6-marks}
 \boxed{\displaystyle
 \partial_I\kappa(Z,Z,Z)=
 \sum_{I_1\sqcup I_2\sqcup I_3\sqcup J=I}
 \ \sum_{\pi\in\Pi(J)}(-1)^{|\pi|}
       \kappa\bigl(Z_{I_1},Z_{I_2},Z_{I_3},(P_B)_{B\in\pi}\bigr).}
\end{equation}
The first three slots are labelled, whereas the blocks in $\pi$ are
unordered.  For $r=|I|$, the number of terms with $k$ additional score slots is
\begin{equation}\label{eq:v6-count}
 N_{r,k}=\sum_{s=k}^{r}\binom rs3^{r-s}\left\{\begin{matrix}s\\k\end{matrix}\right\},
\end{equation}
where braces denote a Stirling number of the second kind.  The rows for
$r=0,1,2,3,4$ are respectively
\[
 (1),\quad(3,1),\quad(9,7,1),\quad(27,37,12,1),
                    \quad(81,175,97,18,1).
\]
\end{lemma}
\begin{proof}
Strength differentiation of the normalized expectation gives
$W_u'=\E_uZ$, $W_u''=-\kappa_u(Z,Z)$ and $W_u'''=\kappa_u(Z,Z,Z)$.
The normalized score rule in \cite{YM4V6} gives
\[
 \partial_i\kappa(Y_1,\ldots,Y_m)
 =\sum_a\kappa(Y_1,\ldots,\partial_iY_a,\ldots,Y_m)
                  -\kappa(Y_1,\ldots,Y_m,P_i).
\]
It follows directly by differentiating \cref{eq:v6-cumulant} and the
normalized expectation.  Bounded truncations followed by the assumed
dominated limits justify the identity without an analytic moment-generating
function.  A new derivative either joins one of the three labelled roots,
joins an existing score block, or creates a new singleton score block.
Induction gives \cref{eq:v6-marks}.  Choose the $s$ marks assigned to score
blocks, partition them into $k$ nonempty blocks, and assign the remaining
marks to three roots.  This proves \cref{eq:v6-count}.
\end{proof}

In particular the exact Hessian identity is
\begin{align}\label{eq:v6-C2}
 \partial_a\partial_b\kappa(Z,Z,Z)={}&
 3\kappa(Z_{ab},Z,Z)+6\kappa(Z_a,Z_b,Z)\notag\\
 &-3\kappa(Z_a,Z,Z,P_b)-3\kappa(Z_b,Z,Z,P_a)
       -\kappa(Z,Z,Z,P_{ab})\notag\\
 &+\kappa(Z,Z,Z,P_a,P_b).
\end{align}
For a fixed reference $P_I=uZ_I$.  The first derivative then reduces to
$3\kappa(Z_a,Z,Z)-u\kappa(Z,Z,Z,Z_a)$, which is exactly the SM--Einstein
centered-moment expression after its disconnected products are combined.
The fifth cumulant in \cref{eq:v6-C2}, and the seventh cumulant in the fourth
derivative, cannot in general be replaced by pair covariances of the
individual insertions alone.

\begin{proposition}[Differentiated integral remainder]
\label{prop:v6-remainder}
For $0\le\epsilon\le u_0$, let
\[
 R_{3,\epsilon}=W_\epsilon-\epsilon\E_{\nu_x}Z
                       +\frac{\epsilon^2}{2}\operatorname{Var}_{\nu_x}Z.
\]
For $|I|\le4$,
\begin{equation}\label{eq:v6-integral}
 \partial_I R_{3,\epsilon}
 =\frac12\int_0^\epsilon(\epsilon-u)^2
                  \partial_I\kappa_{u,x}(Z,Z,Z)\,du.
\end{equation}
This is an exact real-interpolation formula, including moving-reference scores.
\end{proposition}
\begin{proof}
Apply Taylor's theorem in $u$ at zero with the third-derivative integral
remainder, then the stipulated dominated retained derivatives.  The integral
of the nonnegative kernel is $\epsilon^3/6$.
\end{proof}

The cumulant order is bounded here, but the interaction is not truncated to a
polynomial.  For \cref{eq:v6-tilt}, every formula holds with the fixed exact
$Z_\epsilon$.  Higher exact action tails can contribute to its first two
cumulants as well as to this remainder.

\section{Local connected bounds remove the volume loss}
\label{sec:v6-locality}

Let $\Lambda$ be a finite subset or torus of the four-dimensional cell lattice.
Finite field components may be included in a site index.  Write $d$ for its
graph distance and $\tau(i_1,\ldots,i_r)$ for the minimum length of a connected
graph joining the indicated external sites; repeated sites are permitted.
For a retained function $F$ define the anchored derivative seminorm
\begin{equation}\label{eq:v6-norm}
 \|F\|_{a,r}=\sup_{x\in\mathcal U}\sup_{i_1\in\Lambda}
 \sum_{i_2,\ldots,i_r\in\Lambda}
 e^{a\tau(i_1,\ldots,i_r)}
             |\partial_{i_1}\cdots\partial_{i_r}F(x)|,
 \qquad 1\le r\le4.
\end{equation}
For $r=1$ the sum has one term and $\tau=0$.
Vacuum constants are not measured by these seminorms and are retained
separately.  No global extensive random variable is inserted into an absolute
moment as a shortcut for \cref{eq:v6-norm}.

Suppose the \emph{exact} interaction and the differentiated reference admit
absolutely convergent local decompositions
\[
 Z=\sum_y z_y,\qquad P_I=\sum_y p_{y,I}.
\]
Finite-volume sums always have their literal meaning; a useful uniform
application must establish locality of these decompositions.  Assume the
following bound for every term of \cref{eq:v6-marks}.  Its $m=3+k$ local
insertions have centers $y_1,\ldots,y_m$, with derivative labels $I_1,\ldots,I_m$
partitioning all $r$ external labels.  If $O_{y_b,I_b}$ denotes the indicated
$z$ or $p$ insertion, suppose
\begin{equation}\label{eq:v6-treehyp}
 \left|\kappa_{u,x}(O_{y_1,I_1},\ldots,O_{y_m,I_m})\right|
 \le C_{m,r}\sum_{T\in\mathcal T_m}
       \prod_{\{b,c\}\in T}h(y_b,y_c)
       \prod_{b=1}^m\prod_{i\in I_b}w(y_b,i),
\end{equation}
uniformly in the full interval $0\le u\le u_0$, the retained domain and the
finite volume.  Here $h$ is nonnegative and symmetric, $w$ is nonnegative, and
\begin{align}\label{eq:v6-rows}
 J_a&=\sup_y\sum_z e^{ad(y,z)}h(y,z)<\infty,\notag\\
 D_a&=\max\left\{\sup_y\sum_i e^{ad(y,i)}w(y,i),
                 \sup_i\sum_y e^{ad(y,i)}w(y,i)\right\}<\infty.
\end{align}
Put $\mathsf K_{m,r}=C_{m,r}m^{m-2}J_a^{m-1}D_a^r$.

\begin{theorem}[Anchored differentiated-shell bound through order four]
\label{thm:v6-rooted}
Under the exact interpolation, derivative and connected-locality hypotheses,
\begin{equation}\label{eq:v6-rooted}
 \boxed{\displaystyle
 \|R_{3,\epsilon}\|_{a,r}
 \le\frac{\epsilon^3}{6}\sum_{k=0}^rN_{r,k}\mathsf K_{3+k,r},
                       \qquad 1\le r\le4.}
\end{equation}
The bound has no factor $|\Lambda|$.  If the reference is fixed and
\cref{eq:v6-treehyp} is instead imposed on the local $z$ derivatives only,
then
\begin{equation}\label{eq:v6-fixed}
 \|R_{3,\epsilon}\|_{a,r}
 \le\frac{\epsilon^3}{6}\sum_{k=0}^rN_{r,k}u_0^k\mathsf K^{Z}_{3+k,r}.
\end{equation}
In particular the budgets in this fixed-reference convention are
\begin{align}\label{eq:v6-twofour}
 \mathcal M_2={}&9\mathsf K^Z_{3,2}+7u_0\mathsf K^Z_{4,2}
                            +u_0^2\mathsf K^Z_{5,2},\notag\\
 \mathcal M_4={}&81\mathsf K^Z_{3,4}+175u_0\mathsf K^Z_{4,4}
 +97u_0^2\mathsf K^Z_{5,4}+18u_0^3\mathsf K^Z_{6,4}
                            +u_0^4\mathsf K^Z_{7,4}.
\end{align}
\end{theorem}
\begin{proof}
Expand each cumulant in \cref{eq:v6-marks} into its local insertions.
For one of its spanning trees, attach every external derivative site to the
center of the insertion carrying that mark.  This yields a connected graph
which joins all external sites.  Consequently
\[
 \tau(i_1,\ldots,i_r)
 \le\sum_{\{b,c\}\in T}d(y_b,y_c)
          +\sum_b\sum_{i\in I_b}d(y_b,i).
\]
Multiply \cref{eq:v6-treehyp} by the external exponential weight and absorb
it into the edge factors.  Root this augmented tree at the fixed external
site $i_1$.  Summing its other external leaves and then its internal leaves
costs at most $D_a^rJ_a^{m-1}$.  Both orientations of an attachment edge are
covered by the two bounds in $D_a$.  There are $m^{m-2}$ labelled internal
spanning trees.  This proves the bound $\mathsf K_{m,r}$ for one marked
pattern.  Sum the $N_{r,k}$ patterns and use \cref{eq:v6-integral}.
When $A_I=0$, each new score slot contributes $u$, giving
\cref{eq:v6-fixed}.  The explicit rows follow from \cref{eq:v6-count}.
Absolute convergence justifies the same argument for countable local
expansions; otherwise the asserted theorem is read for the finite sums.
\end{proof}

\subsection{A completely explicit reason to keep the connected cancellation}

\begin{example}[The imported absolute-moment budget can grow on a product law]
\label{ex:v6-volume}
Let the reference variables $\eta_y$ be independent standard real Gaussians,
and take
\[
             Z(x,\eta)=\sum_{y=1}^N(1+x_y)\eta_y^2.
\]
At $x=0$ the tilted law remains a product Gaussian of variance
$v=(1+2u)^{-1}$.  With $X=Z-\E_uZ$ and
$D_i=\eta_i^2-\E_u\eta_i^2$, one has
\[
 \E X^2=2Nv^2,\quad \E XD_i=2v^2,\quad
 \E X^2D_i=8v^3,\quad \E X^3D_i=48v^4+12Nv^4.
\]
The exact normalized derivative is
\begin{equation}\label{eq:v6-volcancel}
 3\E X^2D_i+3u\E X^2\E XD_i-u\E X^3D_i
                         =24v^4,
\end{equation}
independent of $N$.  The positive term
$3uC_{20}C_{11,i}$ in the absolute-moment estimate alone is at least
$12uNv^4$ for $u>0$.  The product cancellation which removes that factor
is precisely the fourth connected cumulant.
\end{example}
\begin{proof}
Under the tilt the $\eta_y^2$ are independent squared Gaussians.  Their
second, third and fourth cumulants are $2v^2,8v^3,48v^4$.
Independence gives the displayed mixed moments by the moment--cumulant
partition formula.  Substitution gives \cref{eq:v6-volcancel} because
$v^{-1}=1+2u$.  Absolute expectations dominate the two nonnegative signed
moments in the asserted lower bound.
\end{proof}

This does not invalidate the source theorem.  Its centered-moment hypotheses
are explicit, and its stable domination result is stated there as a
finite-regulator estimate.  It shows why those hypotheses are not the most
useful volume-uniform quantities for the present continuation.

\section{Seven connected entries from covariance responses and local moments}
\label{sec:v6-covariance}

It is not necessary to postulate seven-point clustering as an indivisible new
input.  At this fixed finite order it can be obtained from covariance estimates
for \emph{products of local insertions}, provided their local moments are also
controlled.  This is stronger than checking pair covariances of the individual
insertions but weaker in scope than demanding a physical Hamiltonian gap.

Let $b_m=\sum_{\pi\in\Pi([m])}(|\pi|-1)!$; for $m=1,\ldots,7$ these are
$1,2,6,26,150,1082,9366$.

\begin{lemma}[Product-covariance factorization controls a connected moment]
\label{lem:v6-cut}
Let $2\le m\le7$.  Suppose $\|O_i\|_{L^{2m}(\mu)}\le a_i$, and for every
pair of disjoint nonempty subsets $B,C\subseteq[m]$,
\begin{equation}\label{eq:v6-productmix}
 \left|\Cov_\mu\left(\prod_{i\in B}O_i,\prod_{i\in C}O_i\right)\right|
 \le K|B||C|\left(\prod_{i\in B\cup C}a_i\right)
                         e^{-\nu d(y_B,y_C)}.
\end{equation}
Here $d(y_B,y_C)=\min_{i\in B,j\in C}d(y_i,y_j)$.  Then
\begin{equation}\label{eq:v6-cutbound}
 |\kappa_\mu(O_1,\ldots,O_m)|
 \le K m^2 b_m\left(\prod_i a_i\right)
 \sum_{T\in\mathcal T_m}\prod_{\{i,j\}\in T}
                                e^{-\nu d(y_i,y_j)/(m-1)}.
\end{equation}
\end{lemma}
\begin{proof}
Choose a minimal spanning tree of the terminal centers in their metric, and
remove a longest edge of length $d_*$.  Its two nonempty terminal classes
$B_*,C_*$ have all cross distances at least $d_*$: a shorter cross edge
would replace that edge and contradict minimality.  Also
$d_*\ge\tau_{\rm MST}/(m-1)$.

Write $M_J=\E\prod_{i\in J}O_i$, with $M_\varnothing=1$, and replace it by
$\widetilde M_J=M_{J\cap B_*}M_{J\cap C_*}$.  These are the moments of
a law in which the two groups, with their original within-group joint laws,
are independent.  Its full $m$th cumulant is zero.  If $J$ crosses the cut,
\cref{eq:v6-productmix} bounds $M_J-\widetilde M_J$ by
$K|J\cap B_*||J\cap C_*|e^{-\nu d_*}\prod_{i\in J}a_i$.
Every other moment factor is at most the product of its $a_i$ by H\"older.
Telescoping the products belonging to one partition therefore costs at most
$Km^2 e^{-\nu d_*}\prod_i a_i$.  Sum its absolute cumulant coefficient
$(|\pi|-1)!$ over all partitions.  This gives
$Km^2b_m e^{-\nu\tau_{\rm MST}/(m-1)}\prod_i a_i$.
The exponential is the product on one minimal spanning tree and is no larger
than the sum over all trees in \cref{eq:v6-cutbound}.  Coincident centers cause
no problem; a zero longest edge simply gives the unweighted bound.
\end{proof}

\begin{theorem}[A source-restricted covariance and fourteenth-moment certificate]
\label{thm:v6-covcompiler}
On each exact tilted law $\mu_{u,x}$, assume for the products appearing below
that
\begin{equation}\label{eq:v6-covinput}
 |\Cov(F,G)|\le C_R\sum_{v,w}e^{-\nu d(v,w)}
                          \|\nabla_{\eta_v}F\|_2\|\nabla_{\eta_w}G\|_2.
\end{equation}
Assume each admitted local marked insertion $O_{y,I}$ obeys
\begin{align}\label{eq:v6-mominput}
 \|O_{y,I}\|_{14}&\le A\prod_{i\in I}w(y,i),\notag\\
 \|\nabla_{\eta_v}O_{y,I}\|_{14}
 &\le A\prod_{i\in I}w(y,i)\,1_{\{d(v,y)\le R\}}.
\end{align}
All constants are uniform in $x,u$, scale and volume.  Let $B_R$ bound the
number of scalar detail coordinates in a radius-$R$ ball.  Then
\cref{eq:v6-treehyp} holds, simultaneously for $3\le m\le7$ and $r\le4$, with
\begin{equation}\label{eq:v6-Cconstant}
 h(y,z)=e^{-\nu d(y,z)/6},\qquad
 C_{m,r}=C_RB_R^2e^{2\nu R}\,m^2b_m A^m.
\end{equation}
Consequently \cref{thm:v6-rooted} applies for every $0<a<\nu/6$ for which
$D_a<\infty$.  On the cubic cell lattice $J_a$ is finite uniformly in volume.
\end{theorem}
\begin{proof}
For a product $F=\prod_{i\in B}O_i$ with $|B|\le7$, the product rule and
H\"older with exponent $2|B|\le14$ give
\[
 \|\nabla_{\eta_v}F\|_2
 \le\left(\prod_{i\in B}a_i\right)
                         \sum_{i\in B}1_{\{d(v,y_i)\le R\}},
 \quad a_i=A\prod_{j\in I_i}w(y_i,j).
\]
Insert this and the analogous bound for $G$ into \cref{eq:v6-covinput}.
For each pair of source balls the sum over $v,w$ is at most
$B_R^2 e^{2\nu R}e^{-\nu d(y_i,y_j)}$.  Thus
\cref{eq:v6-productmix} holds with $K=C_RB_R^2e^{2\nu R}$.
Apply \cref{lem:v6-cut}.  Since $m-1\le6$, replacing its exponential by
$e^{-\nu d/6}$ weakens the bound and proves \cref{eq:v6-Cconstant}.
For $\mathbb Z^4$,
\[
 \sum_{n\in\mathbb Z^4}e^{-t|n|_1}
                         =\left(\frac{1+e^{-t}}{1-e^{-t}}\right)^4<\infty
 \quad(t>0),
\]
and a finite torus has no larger periodic-distance row sum.  This proves the
last assertion.  The assumptions explicitly include all products required
by the moment factorization, not only the individual observables.
\end{proof}

Finite range in \cref{eq:v6-mominput} can be replaced by a nonnegative
localization kernel $b(y,v)$.  The same proof applies whenever its actual
convolution with the covariance-response kernel obeys
\[
 \sum_{v,w}b(y,v)e^{-\nu d(v,w)}b(z,w)
                      \le C_b e^{-\nu' d(y,z)}.
\]
For example, if $\sup_y\sum_v e^{\nu'd(y,v)}b(y,v)\le B$ and the same
bound holds with the roles of source and detail reversed, then the left side
is at most $B^2e^{-\nu'd(y,z)}$ for $0<\nu'\le\nu$, by the triangle
inequality.  This includes absolutely summable spatial tails; the needed
marked-insertion norms still have to be bounded.

\subsection{What V5 supplies, and the precise input it does not supply}

\begin{proposition}[Gaussian reference endpoint from the existing locality stock]
\label{prop:v6-Gauss}
For a real Gaussian detail field of covariance $C$, one has
\begin{equation}\label{eq:v6-Gcov}
 |\Cov(F,G)|\le\sum_{v,w}|C_{vw}|
                           \|\partial_vF\|_2\|\partial_wG\|_2
\end{equation}
for smooth observables with square-integrable derivatives.  Therefore V5's
uniformly exponentially local covariance supplies \cref{eq:v6-covinput} at
$u=0$ for its uncut Gaussian reference.  Local polynomial insertions of bounded
degree, bounded local coefficient norms and bounded support have the local
fourteenth-moment stock in \cref{eq:v6-mominput}, uniformly in that Gaussian
family.
\end{proposition}
\begin{proof}
Let $(\eta_t,\eta_t')$ be jointly Gaussian with marginal covariance $C$ and
cross covariance $tC$, $0\le t\le1$.  Differentiating its Gaussian expectation
by integration by parts gives
\[
 \Cov(F,G)=\int_0^1\sum_{v,w} C_{vw}
             \E[\partial_vF(\eta_t)\partial_wG(\eta_t')]\,dt.
\]
Approximation covers a degenerate supported Gaussian.  Cauchy--Schwarz uses
only the fixed marginal laws and proves \cref{eq:v6-Gcov}.  For a polynomial
of degree at most $p$, expand its monomials, use H\"older and the scalar
Gaussian moments up to degree $14p$.  V5 bounds every coordinate variance;
a fixed coefficient sum and support therefore give a uniform constant.
Apply the same argument to its detail derivative, of degree at most $p-1$.
A deterministic nonzero mean requires its own uniformly bounded local stock.
\end{proof}

This proposition does not extend the Gaussian covariance-response inequality
to $u>0$.  A cutoff also changes that reference and needs its own argument.
The hard interacting input is now explicit: a \emph{tilted} response kernel
on the specified product sources, together with their local fourteenth
moments and derivative norms.  The SM--Einstein $M$-matrix comparison is one
possible source of such a response estimate when its conditional Poincar\'e,
incidence and domain hypotheses hold for the same tilted law.  Source-short
bounds may be useful on this restricted source family, but they do not
identify an auxiliary source short with this covariance response by notation.

The existing four-mark polymer theorem is another conditional route.  Its
connected-cluster argument extends to any fixed number of marks, with the
number of source radii and its inclusion--exclusion factor retained.  Here it
would have to be available through seven marks \emph{uniformly along the real
tilt}.  At $u=0$ an unbounded quartic insertion has no two-sided exponential
source neighbourhood, so that complex-source hypothesis cannot simply be
asserted for the uncut Gaussian law.  The finite-moment argument above avoids
that unnecessary assumption.  No new four-mark theorem is counted as an
acquisition in this version.

\subsection{A correlated full-tilt certificate over the actual Gaussian tower}

The source covariance condition can be populated beyond a product measure.
This next statement uses the actual independent Gaussian fibre kernels of V5,
with a declared bounded nonlinear perturbation.  The perturbation is a comparison
interaction, not the compact Yang--Mills residual in \cref{eq:v6-reference}.
Its purpose is to show that the Gaussian locality theorem can support an
all-level, correlated, non-Gaussian differentiated remainder once a genuine
Hessian bound is available.

\begin{lemma}[Local covariance response under diagonal Hessian perturbations]
\label{lem:v6-Witten}
On a finite set of real detail coordinates let $H=H^T\succeq\kappa I$, with
\[
 \max\left\{\sup_i\sum_j e^{a d(i,j)}|H_{ij}|,
             \sup_j\sum_i e^{a d(i,j)}|H_{ij}|\right\}\le M_a,
                  \qquad a>0.
\]
Let $U(\eta)=\eta^TH\eta/2+V(\eta)$ be smooth, with $V''(\eta)$ diagonal,
bounded, and $V''(\eta)\succeq-bI$ for $0\le b<\kappa$.  Put
$c=\kappa-b>0$, and assume the usual integration-by-parts closures on
$d\mu\propto e^{-U}d\eta$; smooth bounded $V$ and all its derivatives suffice.
For any
\begin{equation}\label{eq:v6-Witten-radius}
 0<\sigma\le\min\{a/2,\,ac/(4M_a)\},
\end{equation}
one has
\begin{equation}\label{eq:v6-Witten-bound}
 \boxed{\displaystyle
 |\Cov_\mu(F,G)|\le\frac2c\sum_{i,j}e^{-\sigma d(i,j)}
                      \|\partial_iF\|_{L^2(\mu)}
                      \|\partial_jG\|_{L^2(\mu)}.}
\end{equation}
The constants do not depend on the number of coordinates.
\end{lemma}
\begin{proof}
Let $\mathcal A=-\Delta+\nabla U\cdot\nabla$ be the nonnegative scalar
reversible operator in $L^2(\mu)$.  Integration by parts and its differentiated
identity give, on a smooth core,
\[
 \|\mathcal A f\|_2^2
   =\|D^2f\|_2^2+\E\langle\nabla f,U''\nabla f\rangle
   \ge c\langle f,\mathcal A f\rangle.
\]
The closure and spectral theorem give spectrum in $\{0\}\cup[c,\infty)$;
its kernel consists of constants, since zero Dirichlet energy forces every
derivative to vanish on the connected support.  Thus the centered Poisson
problem is solvable.  The differentiated Poisson equation is
\[
 \nabla\mathcal A f=\mathcal W\nabla f,\qquad
 \mathcal W=\mathcal A\otimes I+H+V''\succeq cI.
\]
It follows, first on the core and then by closure, that
\[
        \Cov(F,G)=\langle\nabla F,\mathcal W^{-1}\nabla G\rangle.
\]
This is an inverse on vector-valued functions of the same law, not an
auxiliary time evolution substituted for the physical Hamiltonian.

For a fixed coordinate $j_0$ conjugate the spatial component index by
$Q_i=e^{\sigma d(i,j_0)}$.  This commutes with $\mathcal A\otimes I$ and
with the diagonal multiplication $V''$.  The change in $H$ has matrix entries
bounded by $|H_{ij}|(e^{\sigma d(i,j)}-1)$.  Its row and column sums are at most
\[
 M_a\sup_{t\ge0}(e^{\sigma t}-1)e^{-at}
 \le M_a\frac{\sigma}{a-\sigma}\le\frac{2M_a\sigma}{a}\le c/2.
\]
The second estimate uses $e^{\sigma t}-1\le\sigma t e^{\sigma t}$ and
$t e^{-(a-\sigma)t}\le(a-\sigma)^{-1}$.  The row/column bound controls the
$L^2(\mu)$ vector operator norm as well.  The resolvent Neumann series therefore
gives $\|Q\mathcal W^{-1}Q^{-1}\|\le2/c$.  At each finite volume $Q$ is a
bounded invertible component matrix; its possibly large condition number does
not enter this last bound.  Choosing $j_0=j$ yields
\[
       \| (\mathcal W^{-1})_{ij}\|_{L^2(\mu)\to L^2(\mu)}
                                     \le(2/c)e^{-\sigma d(i,j)}.
\]
Insert these block bounds into the covariance identity.  For the bounded
smooth perturbations used below the integration identities and closures follow
by Gaussian-confining cutoff approximation; all bounds persist as the cutoff
is removed.  This proves \cref{eq:v6-Witten-bound}.
\end{proof}

\begin{theorem}[A correlated nonlinear remainder uniform over the V5 fibre family]
\label{thm:v6-correlated}
Let $H_j$ be the actual real independent-fibre kernel of V5 on any finite
periodic cell volume, with fixed component dimension and its common weighted
row/column bound $M_a$.  Use its inherited lower bound
$\kappa=24/11635$ in the declared one-colour normalization.  For a retained
scalar parameter $x_i$ per detail coordinate, define the \emph{comparison}
interaction
\begin{equation}\label{eq:v6-cosmodel}
 Z_j(x,\eta)=\sum_i a_i(x_i)(1-\cos\eta_i),\qquad
 \nu_j=N(0,H_j^{-1}),
\end{equation}
where $0\le a_i\le A_0$ and all derivatives through order four are bounded
uniformly in $i,j$.  If $u_0 A_0\le\kappa/2$, then the normalized laws
$e^{-uZ_j}d\nu_j/\E_{\nu_j}e^{-uZ_j}$, $0\le u\le u_0$, satisfy the hypotheses
of \cref{thm:v6-covcompiler}, with constants uniform in level and volume.
Consequently their exact third remainders obey
\begin{equation}\label{eq:v6-cosbound}
                    \|R_{3,j,\epsilon}\|_{a',r}\le C_r\epsilon^3,
           \qquad r=1,2,3,4,\quad 0\le\epsilon\le u_0,
\end{equation}
for one $a'>0$ independent of the level and volume.
\end{theorem}
\begin{proof}
The exact tilted potential is
$U=\eta^TH_j\eta/2+u\sum_i a_i(x_i)(1-\cos\eta_i)$.
Its Hessian perturbation is the diagonal matrix
$u a_i(x_i)\cos\eta_i\ge-u_0A_0$, so it has uniform convexity at least
$\kappa/2$.  \Cref{lem:v6-Witten} applies with $c\ge\kappa/2$ and, for example,
\[
 \sigma=\min\{a/2,a\kappa/(8M_a)\},\qquad C_R=4/\kappa.
\]
Every marked interaction insertion and its detail derivative is bounded
pointwise by a fixed multiple of the maximum coefficient derivative, and is
supported on its one coordinate.  Thus all local fourteenth-moment bounds
hold for the actual tilted law without any change-of-density comparison.
The reference is fixed in $x$, so all score insertions are $u$ times the
corresponding marked interaction.  Set $w(y,i)=1_{y=i}$; the source-support sums now have exactly one
coordinate each.  Apply \cref{thm:v6-covcompiler,thm:v6-rooted} with $0<a'<\sigma/6$.
The constants $C_r$ are obtained from \cref{eq:v6-Cconstant,eq:v6-fixed};
they may be large, but have no level or volume dependence.
\end{proof}

This example is correlated whenever the fibre covariance is not diagonal.
It uses a nonlinear interaction on the actual Gaussian fibre family, not an
independent-site surrogate for its covariance.  It still does not identify
\cref{eq:v6-cosmodel} with the completed Wilson residual.  In particular the
uniform Hessian inequality just proved for the bounded cosine interaction is
not inferred for the actual compact fibre, whose hidden metastability must
not be contradicted.  The subsequent quartic example addresses a different
issue: unbounded insertions with no two-sided exponential source moment.


\section{A fully populated stable non-Gaussian test of the remainder theorem}
\label{sec:v6-stable}

The preceding implication is not restricted to Gaussian tilted laws.
The following family realizes all real-tilt estimates without an absolute
exponential moment of a quartic polynomial or a factor growing with volume.
It is a controlled comparison model, not a relabelling of the compact YM law.

At each site let
\begin{equation}\label{eq:v6-quartic}
 d\nu_x(\eta)\propto\prod_y e^{-h(x_y)\eta_y^2/2}\,d\eta_y,
 \qquad Z(x,\eta)=\sum_y a(x_y)\eta_y^4,
\end{equation}
where $h\ge h_0>0$, $a\ge0$, and $h,a$ have bounded derivatives through order
four on the declared retained interval.  The normalization of $\nu_x$ is
retained; its local reference scores are
$P_{y,I}=\frac12h^{(|I|)}(x_y)\eta_y^2+u a^{(|I|)}(x_y)\eta_y^4$
when all derivative labels are at $y$, and zero otherwise.

\begin{theorem}[Volume-uniform real-tilt fourth-derivative remainder]
\label{thm:v6-quartic}
For \cref{eq:v6-quartic}, $0\le u\le u_0$, the constants in
\cref{thm:v6-rooted} are finite uniformly in the number of sites.
More explicitly, choose $B\ge1$ so that every needed coefficient of a
$z$ or $p$ insertion is bounded by $B$ in the estimate
$|O_y|\le B(1+|\eta_y|^4)$.  Set
\begin{equation}\label{eq:v6-quarticconstant}
 C_m^{\rm prod}=b_m B^m 2^{m-1}
                      \left\{1+(4m-1)!!\,h_0^{-2m}\right\},\quad 3\le m\le7.
\end{equation}
Then, for $1\le r\le4$,
\begin{equation}\label{eq:v6-quarticbound}
 \boxed{\displaystyle
 \|R_{3,\epsilon}\|_{a,r}
 \le\frac{\epsilon^3}{6}\sum_{k=0}^rN_{r,k}C_{3+k}^{\rm prod}}
\end{equation}
for any $a\ge0$.  The moving reference is included.  The reference quartic
observable nevertheless satisfies $\E e^{t\eta^4}=\infty$ for every $t>0$.
\end{theorem}
\begin{proof}
The tilted measure is the product of the one-site measures
$e^{-h\eta^2/2-u a\eta^4}$ with their own normalizers.  For every increasing
nonnegative function $g(|\eta|)$, its expectation under this tilt is at most
its Gaussian expectation at the same $h$.  Indeed the multiplying weight
$e^{-u a|\eta|^4}$ is decreasing, and
\[
 2\Cov(g(X),w(X))
     =\E[(g(X)-g(X'))(w(X)-w(X'))]\le0
\]
for independent copies $X,X'$ of the Gaussian absolute value.  Divide the
covariance inequality by $\E w>0$.  In particular
\[
 \E_{u,x}|\eta_y|^{2n}\le(2n-1)!!h_0^{-n}.
\]
All derivative moments needed for dominated real interpolation are finite;
on compact retained intervals this also follows directly from the uniform
Gaussian lower confinement and a one-site normalizer lower bound.

A connected cumulant of insertions from two independent site classes vanishes.
When all entries belong to one site, \cref{eq:v6-cumulant}, H\"older and
$|O_y|\le B(1+|\eta_y|^4)$ give
\[
 |\kappa(O_1,\ldots,O_m)|
 \le b_m B^m\E(1+|\eta_y|^4)^m
 \le C_m^{\rm prod}.
\]
Every nonzero retained derivative also has all its labels at that same site.
Thus the rooted spatial sums have a single term, with no Cayley factor or
volume cost.  Apply \cref{eq:v6-marks,eq:v6-integral} to prove
\cref{eq:v6-quarticbound}.  Finally $t\eta^4-h\eta^2/2$ tends to positive
infinity quartically on either tail, proving the last assertion.
\end{proof}

For the fixed reference $h=1$ and $a(x)=e^x$ on a bounded interval,
\[
 \kappa_0(\eta^4,\eta^4,\eta^4)
       =10395-3(105)(3)+2(3)^3=9504.
\]
The fourth retained derivative of the local third cumulant at $u=x=0$ is
$81\cdot9504=769824$.  The finite checker verifies these cumulants through
seven entries and the general marked derivatives against an independently
normalized finite-state computation.  The proof of the entire real-tilt
bound is the monotonicity argument, not a finite list of sampled strengths.

\section{Cofinal use: subtraction, physical weights and transport amplification}
\label{sec:v6-cofinal}

The remainder theorem supplies one error in a complete shell comparison.  It
does not make the first two strength coefficients summable, fix a marginal
coupling, or construct an invariant local subtraction.

\begin{proposition}[Summable remainders and their actual rate requirement]
\label{prop:v6-summable}
Let $\delta_j=\epsilon_j^3\mathcal M_j/6$ be one of the proved rooted
remainder bounds after all reference conventions have been matched.  If
$\epsilon_j\le C(j+j_0)^{-1/2}$ and
$\mathcal M_j\le M(j+j_0)^p$ with $p<1/2$, then
\[
 \sum_{j\ge J}\delta_j
 \le\frac{C^3M}{6(1/2-p)}(J+j_0-1)^{p-1/2}
\]
whenever $J+j_0>1$.  In particular bounded budgets give a tail of order
$(J+j_0)^{-1/2}$.  These are conditional rate implications; the stated
coupling and moment rates are not established for the nonlinear YM trajectory.
\end{proposition}
\begin{proof}
The summand is bounded by $(C^3M/6)(j+j_0)^{-3/2+p}$.
Comparison with its decreasing integral gives the displayed estimate.
\end{proof}

\begin{example}[A zero third remainder does not control a running action]
\label{ex:v6-marginal}
Let $Z_j(x,\eta)=x^4$ be deterministic, and $\epsilon_j=(j+1)^{-1/2}$.
Every connected cumulant of order at least two, and hence every third
remainder, is zero.  Nevertheless
$\sum_{j<N}W_{j,\epsilon_j}(x)=x^4\sum_{j<N}\epsilon_j$ has an unbounded
fourth derivative.  This elementary example shows why matched first and
second coefficients and the retained local trajectory are independent inputs.
\end{example}

The localized comparison seminorm in \cite{SM6} uses the physical weights
$w_{S,j}(i)=e^{-m_*d_j^{\rm phys}(S,i)}$ and $L^2$ norms in a stated retained
law.  When its derivative coordinates agree with those in \cref{eq:v6-norm},
\begin{equation}\label{eq:v6-physical}
 \|R_{3,j}\|_{\mathcal B_{S,j}}
 \le\mathcal N_{S,j}\|R_{3,j}\|_{a,1},\qquad
                   \mathcal N_{S,j}=\sum_iw_{S,j}(i).
\end{equation}
This follows by the uniform retained-field bound and the triangle inequality.
The source norm is not silently replaced by a normalized per-cell norm.  On a
mesh of spacing $h_j$, the unweighted site count in a fixed physical region
can grow as $h_j^{-4}$.  A physical cell-volume weight, canonical field
rescaling, or source-specific cancellation must be declared if used to
remove that factor.  Equation \cref{eq:v6-physical} alone does not do so.
Different derivative coordinates require their additional conversion norm.

\begin{lemma}[Subsequent nonlinear steps can amplify an otherwise summable error]
\label{lem:v6-amplification}
Suppose exact and approximate normalized steps have local defect at most
$\delta_j$ in their declared common norms, and step $j$ has Lipschitz bound
$L_j$.  Their trajectory errors obey
\[
 d_N\le\left(\prod_{\ell=0}^{N-1}L_\ell\right)d_0
           +\sum_{j=0}^{N-1}\delta_j\prod_{\ell=j+1}^{N-1}L_\ell.
\]
If $L_j\le1+b_j$, $b_j\ge0$, and $\sum_jb_j<\infty$, the amplification products are
uniformly bounded by $e^{\sum_jb_j}$.  A genuine adjacent-cutoff comparison
then has a summable remainder contribution whenever its physically converted
$\delta_j$ are summable.  Other comparison terms are still required.
\end{lemma}
\begin{proof}
The one-step inequality is $d_{j+1}\le L_jd_j+\delta_j$.
Induct and use $1+b\le e^b$.  With an identification into one common complete
space, summable bounds for \emph{all} actual adjacent differences give the
Cauchy conclusion.  That last hypothesis is not implied by a bound for only
this one remainder term.
\end{proof}

For example $d_{j+1}=2d_j+2^{-(j+1)}$, $d_0=0$, gives
$d_N=(2^N-2^{-N})/3$, despite summable raw defects.  Thus an unstable
relevant coordinate cannot be ignored in the error transport.  The leakage
and adjacent-cutoff method in \cite{ClassV6} is a useful model for such a
complete comparison, not an already populated YM stability estimate.

\section{The secondary transfers and their exact scope}
\label{sec:v6-secondary}

\subsection{Full-face information must be compared in a common gauge frame}

The rigid full-face identity in \cite{SM6} is a proved geometric statement
about affine Euclidean rigid fields.  It does not by itself identify a
negative-curvature or heat-bath repair direction of the completed YM action.
There is a legitimate small-logarithm comparison which makes part of the
suggested transfer precise.

\begin{lemma}[A chordal comparison for affine Lie-algebra face data]
\label{lem:v6-face}
Let $\xi_\pm(y)$ be skew-Hermitian matrix fields on a flat face $f$, with
$\|\xi_\pm(y)\|_{\rm op}\le r<\log2$.  Suppose their difference is affine,
$\xi_+(y)-\xi_-(y)=J+Ty$, with the centered face coordinates of side $\ell$.
Then
\begin{align}\label{eq:v6-face}
 &(2-e^r)^2\left\{\ell^{d-1}\|J\|_\HS^2
                     +\frac{\ell^{d+1}}{12}\|TP_\mu\|_\HS^2\right\}\notag\\
 &\quad\le\int_f\|e^{\xi_+(y)}-e^{\xi_-(y)}\|_\HS^2\,dS\notag\\
 &\quad\le e^{2r}\left\{\ell^{d-1}\|J\|_\HS^2
                     +\frac{\ell^{d+1}}{12}\|TP_\mu\|_\HS^2\right\}.
\end{align}
The comparison requires one common matrix frame on the face.
\end{lemma}
\begin{proof}
For $X,Y$ of operator norm at most $r$,
$X^n-Y^n=\sum_{k=0}^{n-1}X^k(X-Y)Y^{n-1-k}$ gives
\[
 \|e^X-e^Y-(X-Y)\|_\HS\le(e^r-1)\|X-Y\|_\HS,
 \quad \|e^X-e^Y\|_\HS\le e^r\|X-Y\|_\HS.
\]
The reverse triangle inequality gives the lower bound.
Integrate the squared affine difference.  Its cross term vanishes and its
second moment is $\ell^{d+1}P_\mu/12$, as in the source identity.
\end{proof}

A large mismatch between \emph{unidentified} cell frames is not a gauge-invariant
energy diagnostic.  Indeed links $U_{xy}=g_x^{-1}g_y$ have identity holonomy on
every contractible plaquette for arbitrarily varying $g_x$, whereas the raw
frame differences $\|g_x-g_y\|$ can be large.  Their Wilson action is a minimum,
so that raw mismatch cannot imply a lowering action direction.  This observation
is made before gauge fixing; it does not claim such arbitrary frames survive
inside the source's fixed independent tree coordinates.  A usable extension
must specify parallel transport or a common tree, retain actual pivot and
Jacobian derivatives, and prove the action comparison in those variables.
\Cref{lem:v6-face} supplies none of the missing global coverage implication.

The source's V17 active-motion coefficient $k/(8-k)$ and its curvature cost
remain the actual available certificate, as summarized in \cite{YM4V6}.
A centre-value test and a full-face functional may be compared on a specified
small-logarithm family, but the unrestricted transition-cutoff and high-occupancy
complement remains outside this version.

\subsection{Source shorts and Ward protection are distinct tasks}

The source-short transport theorem in \cite{GT76V6} correctly retains the
Feshbach returning-memory debit.  It can preserve a chosen source reserve
without a global form gap.  For the present transfer its useful target is the
\emph{actual covariance response on the finite-product insertion family} in
\cref{eq:v6-covinput}.  Its form operator, source map and original clock must
first be identified with that response.  A bound on a fixed unrelated source
bank is not a uniform bound for every local observable needed in the reflected
continuum theory.  No new source-short theorem is reproved here.

For Ward compatibility the exact normalized interpolation has a simpler
algebraic test, before any commutant spectral argument is invoked.

\begin{proposition}[Equivariant exact interpolation preserves its Ward identity]
\label{prop:v6-Ward}
Suppose a retained symmetry $g$ transports the reference by
$(T_g)_*\nu_x=\nu_{gx}$ and the exact correction satisfies
$Z(gx,T_g\eta)=Z(x,\eta)$, including cutoff and measured density conventions.
Then $W_u(gx)=W_u(x)$ for every admitted $u$, and its first two strength
coefficients and exact remainder have the same invariance.  At a stationary
retained background, differentiating the infinitesimal identity puts the
symmetry tangent in the Hessian kernel.  No assertion that this kernel has
no additional directions follows.
\end{proposition}
\begin{proof}
Change variables by $T_g$ in the normalized expectation.  Differentiate in
$u$ and subtract the resulting invariant terms.  For an infinitesimal field
$X(x)$ the identity is $DW(x)X(x)=0$.  Its derivative is
$D^2W(v,X)+DW(DX[v])=0$; the last term vanishes at a stationary background.
\end{proof}

The commutant Laplacian in \cite{DualV6} controls a represented algebra's
commutant, which is not automatically the gauge tangent of this action.
Its finite-anchor perturbation estimate is useful once the generators and
margin are specified.  Its compact-Mosco theorem separately assumes the
absence of extra zero directions in the limiting commutant.  That hypothesis
must not be used as though convergence alone had established the desired
kernel exclusion.  These spectral tools remain useful insurance, not a
substitute for exact Ward transport or for a physical mass gap.

\section{The changed acquisition order and the remaining nonlinear estimate}
\label{sec:v6-next}

The adopted direction is now
\[
 \begin{gathered}
 \text{full measured local shell and reference scores}\\
 \longrightarrow\ \text{rooted product-source covariance and local moments}\\
 \longrightarrow\ C^2/C^4\text{ normalized remainder bounds}\\
 \longrightarrow\ \text{physically normalized, stable adjacent-cutoff comparison}.
 \end{gathered}
\]
This is a quantitative transfer of the proposed method, not a claim that the
last two lines have already been populated on the interacting YM family.
The exact law and cutoff identity, the derivative count, the cancellation of
the extensive disconnected stock, and the reduction of seven connected
entries to covariance and fourteenth moments are proved.  The correlated
cosine comparison supplies an all-level real-tilt estimate over the actual
Gaussian fibre family under its declared diagonal Hessian bound.  The stable
quartic family verifies that the resulting assumptions can hold on a genuinely
non-Gaussian real tilt without impossible quartic exponential moments.

For the actual compact program, the next input to acquire is
\cref{eq:v6-covinput,eq:v6-mominput} on the \emph{same tilted reduced law},
including its complete measured local score insertions.  V5 supplies the
Gaussian $u=0$ spatial response, not the entire interval.  A small-field
conditional response estimate may begin the argument; its common domain,
cutoff acceptance and complement must remain in the comparison.  Replacing
that task by a false unrestricted hidden-diffusion gap would undo the source's
metastability result, and is not selected.

Nor can the complete low-order recursion be discarded.  The first two shell
coefficients, the marginal field-strength normalization, and the Ward-compatible
local projection must be matched, while transport amplification and physical
source weights are recorded as in \cref{sec:v6-cofinal}.  Once those inputs
are verified, the new remainder bound is summable under the stated rates;
before then it is a precise sufficient estimate with an identified interacting
response obligation.  No physical Hamiltonian mass gap, interacting continuum
construction, or new unrestricted interface coverage is asserted.

\section{Executed checks and append-only provenance}
\label{sec:v6-verification}

The new script \path{verify_ym_normalized_shell_V6.py} executes exact rational,
symbolic and Gaussian-moment checks.  Its normalized-derivative calculation
constructs and divides the finite density series directly in labelled
square-free jet variables.  It does not use the marked cumulant formula to
produce both sides of the test.  The latter side is independently assembled
from all moment partitions, including moving-reference scores.  Distinct and
repeated retained directions through order four are checked.

The report contains nineteen passing test families: the complete mark census;
ten normalized fixed/moving-reference jet tests; the extensive product-Gaussian
cancellation; quartic cumulants through seven entries; connected Wick cycles
against full moment partitions; a rooted spatial tree sum; the exact residual
not determined by the compiler jets; stable quadratic remainder derivatives;
a rational noncommuting pure-gauge plaquette; and transport amplification with
summable raw errors.  The finite-state and product-law examples are explicitly
verification models, not claimed compact YM measurements.  No finite set of
tests proves the analytic covariance hypotheses for the interacting tilt.

The V5 source, its complete mathematical body and all preceding version markers
are preserved without modification.  Only version displays in the preamble
are advanced.  The new report, code, source comparisons and hashes are bundled
with the prior V5 materials.  Prior reports remain archival unless a separate
fresh execution is recorded.  No Lean formalization or independent rerun of
the supplied V17 certificate chain is claimed.  Future V7 work is to append
after the marker below; any correction to a preceding proof must be stated
explicitly rather than silently replacing it.

\begin{thebibliography}{99}
\bibitem[SM6]{SM6}
\emph{SM--Einstein Continuum Monograph}, supplied updated manuscript,
\path{SM_Einstein_Continuum_Monograph_Updated(3)(5).tex}.
Relevant labels: \texttt{thm:localized-shell},
\texttt{prop:stable-polynomial-tilt}, \texttt{prop:full-face-rigid},
\texttt{thm:covariance}.
\bibitem[GT76]{GT76V6}
\emph{Renewal Grand Tensor Monograph}, V076, supplied manuscript,
\path{Renewal_Grand_Tensor_Monograph_V076(5).tex}.
Relevant section: \texttt{sec:GT-source-short-transport}.
\bibitem[YM4]{YM4V6}
\emph{Renewal Yang--Mills Closure Monograph}, V4, Library snapshot,
\path{Renewal_Yang_Mills_Closure_Monograph_V4.tex}, version 1 of the retrieved
Library file.  Relevant labels: \texttt{thm:connected-score-hierarchy},
\texttt{thm:four-mark-clustering}, \texttt{thm:active-repair}.
This source summarizes the standalone f13 V17 line; its certificate executables
were not independently supplied in this continuation.
\bibitem[DU6]{DualV6}
\emph{Spacetime--Gauge Commutant Duality: Finite Rigidity and Cofinal Stability},
supplied manuscript,
\path{Spacetime_Gauge_Commutant_Duality_CMP_Analytic_Strengthened(6).tex}.
Relevant labels: \texttt{thm:finite-anchor}, \texttt{thm:cofinal-duality}.
\bibitem[CE6]{ClassV6}
\emph{Renewal Geometry Classical Einstein--SM}, supplied manuscript,
\path{Renewal_Geometry_Classical_Einstein_SM_CMP_Submission(6).tex}.
Relevant labels: \texttt{thm:native-source}, \texttt{thm:native-closure}.
\bibitem[LS6]{LorV6}
\emph{Renewal Lorentzian Spacetime}, supplied revised manuscript,
\path{Renewal_Lorentzian_Spacetime_CMP_Revised (1)(3).tex}.
Relevant label: \texttt{thm:supp-reward-pressure}.
\bibitem[APS09]{APSV6}
A. Abdesselam, A. Procacci and B. Scoppola,
\emph{Clustering bounds on n-point correlations for unbounded spin systems},
Journal of Statistical Physics \textbf{136} (2009), 405--452,
arXiv:0901.4756, DOI: 10.1007/s10955-009-9789-y.
\url{https://arxiv.org/abs/0901.4756}.
External methodological context only; no estimate is imported as a proved
property of the present Yang--Mills law.
\end{thebibliography}

% END OF VERSION 6 MATHEMATICAL BODY.
% Append subsequent requested versions below this marker; retain V1--V6 above.


\clearpage
\section{V7: a uniform measured two-loop expansion on the actual local box}
\label{sec:v7-context}

The next step is not another nonlinear comparison interaction.  The actual
small-field Wilson fibre already has more analytic structure than the
comparison examples in V6 use.  The f12 archive constructs a local convex
extension, proves moments at every fixed order and spatial covariance decay,
and returns from that extension to the original fixed box through two
retained derivatives.  Its final original-box estimate is
$O(\beta^{-1/2})$ in a localized Hessian norm
\src{F12}{V10--V12, \texttt{thm:f12v12-local}}.  Its opening summary also
records the usual fixed-volume expansion through the next Laplace coefficient,
with an $O_{C^2}(\beta^{-2})$ remainder at that fixed volume.  Neither the
finite-volume coefficient nor the existing local convexity is a new
acquisition here.

The new conclusion is the \emph{volume-uniform fourth-retained-derivative
version} of that measured expansion on the same original local box:
\[
 \widehat W_\beta^{\rm loc}
 =\beta G+\frac{n_\Lambda}{2}\log(\beta/2\pi)
       +\widehat Q+\beta^{-1}\mathcal B_1+\mathcal R_\beta,
 \qquad \|\mathcal R_\beta\|_{a,r}\le C_r\beta^{-2},\quad 1\le r\le4.
\]
Here $n_\Lambda$ is the real number of eliminated coordinates.  The norm is
V6's anchored derivative norm, not a norm divided by volume; only the
zero-derivative pressure is divided by $n_\Lambda$.  All Haar and inverse-pivot
terms remain in the measured coefficient.  A strength-integrable five-entry
wall estimate extends the source's two-derivative transfer to four derivatives.

The proof goes upstream from V6's generic $O(\epsilon^3)$ remainder to the
\emph{actual signed small-noise family}.  Its complete normalizer is even in
$\epsilon$, irrespective of whether the fine action is even in its field.
Eight-entry connected bounds, rather than seven, then give a fourth-strength
Taylor remainder $O(\epsilon^4)$.  The bounds are populated from the inherited
local Wilson chart, not postulated on the unrestricted compact law.

Parity in Laplace's method, Gaussian Wick contraction, and local stationary
expansions are not claimed as new general methods; see, for example,
\cite{KatsevichV7} for a different high-dimensional Laplace setting.  The
additional result here is the spatially anchored, volume-independent estimate
and its differentiated transfer to this programme's measured fixed-box law.
No bound from that external reference is used to supply a Yang--Mills input.

\section{The actual local inputs and their higher marked derivatives}
\label{sec:v7-inputs}

Let $\Lambda$ be a finite periodic cell set in four dimensions, with a fixed
number of field components per cell and periodic graph distance $d$.  Work on
the fixed independent-fibre chart of the supplied root-path block.  In its
original coordinates the local integral is
\begin{equation}\label{eq:v7-local-integral}
 I_\beta^{\rm loc}(c)=\int_{\mathcal D}
                 e^{-\beta f(z,c)+\ell(z,c)}\,dz,
 \qquad \mathcal D=[-r,r]^{n_\Lambda}.
\end{equation}
The box is independent of $c$ and $\beta$.  It is the small box from f12
V11--V12, not the full compact fibre.  The logarithmic density $\ell$ includes
the actual fine Haar and inverse-pivot coordinate Jacobian.  Set
\[
 z_*=z_*(c),\quad G(c)=f(z_*,c),\quad H_c=f_{zz}(z_*,c),\quad
 \ell_*(c)=\ell(z_*,c),\quad C_c=H_c^{-1}.
\]
For a coarse reference $j_c(c)\,dc$, use
\begin{equation}\label{eq:v7-normalizations}
 \widehat W_\beta^{\rm loc}=-\log I_\beta^{\rm loc}+\log j_c,
 \qquad \widehat Q=\tfrac12\log\det H_c-\ell_*+\log j_c.
\end{equation}
The reference conversion is included once.  Changing a constant coordinate
normalization changes only the separately recorded vacuum constant.

\paragraph{Precisely inherited inputs.}
The following properties are used with their original small-chart scope:
\begin{enumerate}[label=\textup{(L\arabic*)},leftmargin=3em]
\item $f=\sum_y f_y$ and $\ell=\sum_y\ell_y$ are sums of analytic local
stencils of bounded support and incidence.  On a smaller common chart their
local derivatives through order twelve are uniformly bounded.  This finite
order is a sufficient, nonoptimal choice; it follows from the fixed analytic
stencils, and is not inferred from V4's fourth-jet compiler.
\item The specified classical minimum satisfies
$\|z_*\|_\infty\le r/2$.  Its fibre Hessian has a common positive lower bound,
finite interaction range, bounded row sums, and an exponentially decaying
inverse.  The actual measured potential $\beta f-\ell$ is uniformly strongly
convex on $\mathcal D$ for sufficiently large $\beta$.
\item The centred local cutoff construction of f12 V10 gives a whole-space
convex extension agreeing with the actual law near $z_*$.  On the original
box its interpolation with the actual potential has the same interior
measured minimum and volume-uniform local tails.  For the extension with a
fixed-coordinate convex wall, f12 V12 supplies moments, marginal wall-layer
bounds, and the wall-sensitive covariance resolvent, uniformly in wall
strength.
\end{enumerate}
These are the source results in \src{F12}{V10--V12}, specifically the
convex interpolation and centred moment propositions, the same-box
interpolation and local-tail propositions, and the concentration,
wall-layer and wall-sensitive resolvent propositions.  Their original
labels are recorded for precise comparison:
\begin{quote}\small
\nolinkurl{prop:f12v10-convex}, \nolinkurl{prop:f12v10-moments};\\
\nolinkurl{prop:f12v11-interpolation}, \nolinkurl{prop:f12v11-tail};\\
\nolinkurl{prop:f12v12-concentration}, \nolinkurl{prop:f12v12-layer};\\
\nolinkurl{prop:f12v12-resolvent}.
\end{quote}
Their constants are not numerically recertified here.  The source already
states the moment and layer estimates at every fixed moment order; only its
retained-derivative conclusions stop at order two.  The proof below extends
that derivative argument explicitly.

Write $\tau(i,I)$ for the length of a shortest spatial tree joining the
fibre site $i$ and the retained sites labelled by $I$.  Repeated labels are
allowed.  The same notation includes local stencil centres as terminals.

\begin{lemma}[Four marked responses of the moving minimum]
\label{lem:v7-centre}
On a smaller fixed retained chart, for $1\le |I|\le4$,
\begin{equation}\label{eq:v7-centre}
 |\partial_I z_{*,i}(c)|\le A_{|I|}e^{-\nu_{|I|}\tau(i,I)}.
\end{equation}
The same tree estimate holds for the marked coefficients of any local
centred stencil required below, anchored at that stencil's centre.  All
constants are independent of $\Lambda$.
\end{lemma}
\begin{proof}
Differentiate $f_z(z_*(c),c)=0$.  At first order the response is
$-H_c^{-1}f_{zc_i}$, whose inverse kernel and local source are inherited.
At order $s$, separate the unique term $H_c\partial_Iz_*$.  Every other term
is a derivative of one local $f$ stencil contracted with derivatives of $z_*$
of strictly smaller orders; the derivative labels are distributed among
those factors.  This is the finite multivariate chain rule.  There are only
finitely many partitions for $s\le4$.

Inductively insert the lower-order response bounds and multiply by
$-H_c^{-1}$.  A term is represented by a connected tree of inverse links and
local stencil incidences joining its output to all marked sites.  Its total
link length, up to a fixed finite support allowance, is at least $\tau(i,I)$.
Keep half the decay of every link for the terminal tree, and use the other
half to sum its internal vertices.  Exponential row sums are uniform on
four-dimensional periodic lattices.  The number of vertices and local
incidences is bounded at each of the four orders, so a common smaller
positive exponent and finite constants result.  This proves the induction.
The identical chain rule for a centred local coefficient proves the last
statement.  Derivatives through the finite order specified in (L1) suffice.
No property of a growing family of transported nonlinear stencils is inferred.
\end{proof}

For comparison with V6's product attachment weights, note that
\[
 \tau(y,I)\ge |I|^{-1}\sum_{i\in I}d(y,i).
\]
Consequently a tree bound through four marks implies a product bound with
$w(y,i)=C e^{-\nu d(y,i)/4}$ after reducing $\nu$.  The two weighted row sums
required in \cref{eq:v6-rows} are then finite uniformly in volume.

\section{Signed-noise interpolation of the complete measured extension}
\label{sec:v7-noise}

Keep the actual local Taylor remainders from f12 V10.  In centred variables
$v=z-z_*(c)$ write
\begin{equation}\label{eq:v7-centred}
 f(z_*+v,c)-G(c)=\tfrac12v^TH_cv+\sum_y r_y(v,c),\qquad
 \ell(z_*+v,c)-\ell_* =\sum_y d_y(v,c).
\end{equation}
Each $r_y$ is the remainder after subtracting its local action Taylor
polynomial through degree two, and each $d_y$ is the local density difference.
The sum of the subtracted linear action terms is zero by stationarity.
Choose the inherited smooth local cutoff
$\chi_y(v)=\prod_{i\in S_y}\chi(v_i/\delta)$ and set
\[
 \widehat r_y=\chi_y r_y,\qquad \widehat d_y=\chi_y d_y.
\]
These functions are extended by zero outside their fixed local supports in
field space.  Fix $\delta>0$ small enough for the source's convex extension
and mean-absorption estimates.  It does not depend on volume or noise.

For a \emph{signed} real parameter $t$ define on all $\eta\in\R^{n_\Lambda}$
\begin{equation}\label{eq:v7-U}
 U_t(c,\eta)=\tfrac12\eta^TH_c\eta+
       \sum_y\left\{\frac{\widehat r_y(t\eta,c)}{t^2}
                         -\widehat d_y(t\eta,c)\right\}.
\end{equation}
The value at $t=0$ is its smooth limit
$U_0=\eta^TH_c\eta/2$.  Put
\begin{equation}\label{eq:v7-W}
 \mathcal W(t,c)=-\log\frac{\int e^{-U_t(c,\eta)}\,d\eta}
                                  {\int e^{-U_0(c,\eta)}\,d\eta}.
\end{equation}
This is not the $u$ derivative of V6 with $Z_\epsilon$ frozen.  Changing $t$
differentiates the complete scaled residual and the measured density.  The
reflection argument below is for this family as defined, not a parity
assumption imposed on an isolated third strength cumulant.

Only $t=\beta^{-1/2}>0$ is used to identify the physical local integral.
Negative $t$ is a proof parameter, not negative inverse coupling.  In
particular, no logarithm of a negative $t$ or signed Gaussian volume factor
is introduced.

\begin{lemma}[A uniformly controlled signed family]
\label{lem:v7-noise-stock}
There are $t_0,m>0$, independent of volume and the small retained field, such
that for $|t|\le t_0$ the potential \cref{eq:v7-U} is uniformly strongly convex,
$D_\eta^2U_t\succeq mI$, with a common finite-range off-diagonal Hessian bound.
Its normalized laws satisfy
\begin{equation}\label{eq:v7-noise-moments}
 \sup_{t,c,i}\E_{t,c}|\eta_i|^p\le C_p
 \quad\text{for every fixed }p<\infty,
\end{equation}
and a componentwise spatial covariance estimate of the form
\cref{eq:v6-covinput}, with common constants.

Every local insertion $\partial_t^a\partial_I U_{t,y}$ needed for four
noise derivatives and at most four retained derivatives, and its first
$\eta$ derivative, has a uniform local $L^{16}$ bound with the attachment
weights of \cref{lem:v7-centre}.  These statements apply to the actual
cutoff extension of (L3), not to a substituted cosine or quartic model.
\end{lemma}
\begin{proof}
Cutoff Taylor estimates give
\[
 \|D^2\widehat r\|_{\rm row}\le C\delta,\qquad
 \|D^2\widehat d\|_{\rm row}\le C_\delta.
\]
The Hessian of \cref{eq:v7-U} is
$H_c+D^2\widehat r(t\eta)-t^2D^2\widehat d(t\eta)$.
Choose $C\delta\le h/4$ and $t_0^2C_\delta\le h/4$, where
$H_c\succeq hI$.  This proves the lower bound $m=h/2$; bounded stencil
incidence proves the common range and off-diagonal bound.

For clarity, moments are bounded about the \emph{centred origin}, not merely
about an uncontrolled mean.  Uniform convexity also holds after any linear
tilt.  The Poincare estimate for those tilts yields a sub-Gaussian coordinate
law about its mean with variance parameter $m^{-1}$.  Integration by parts in
\cref{eq:v7-U} gives, for $t\ne0$,
\[
 H_c\E\eta=-t^{-1}\E\nabla\widehat r(t\eta)
                                  +t\E\nabla\widehat d(t\eta).
\]
The inherited inverse row bound and
$|\nabla\widehat r_y(v)|\le C\delta|v_{S_y}|$ imply, for
$M=\max_i|\E\eta_i|$,
\[
 M\le KC\delta(M+m^{-1/2})+KC_\delta|t|.
\]
Choose the same fixed $\delta$ so $KC\delta\le1/2$.  This bounds $M$
independently of dimension.  Combining with the sub-Gaussian estimate proves
\cref{eq:v7-noise-moments}.  Using the sharper quadratic gradient bound then
also gives $M\le C|t|$.  At $t=0$ the assertion is Gaussian; negative $t$
follows by reflection or the same estimates.

The covariance proof is the source's differentiated Poisson argument.
The vector operator is $\mathcal L_t\otimes I+D^2U_t\succeq mI$.
Conjugation on its spatial component index by $e^{\sigma d(i,i_0)}$ commutes
with the differential part and with the diagonal Hessian.  The remaining
error is bounded by $K(e^{\sigma R}-1)$ using the uniform finite interaction
range $R$.  Choose it at most $m/2$ and invert by a Neumann series.  The
resulting inverse blocks are bounded by $(2/m)e^{-\sigma d(i,j)}$.
Insertion into the same-law covariance representation gives
\cref{eq:v6-covinput}.  This proof allows field-dependent off-diagonal entries;
their common bound, not diagonal structure of the perturbation, is used.
Whole-space quadratic confinement justifies closure and integration by parts.

To check the differentiated insertion bounds without an inverse power of $t$,
use the identity
\[
 \frac{r_y(t\eta,c)}{t^2}
 =\int_0^1(1-s)\bigl[f_{y,zz}(z_*+st\eta,c)
                         -f_{y,zz}(z_*,c)\bigr][\eta,\eta] \,ds.
\]
For $a\ge1$ its $a$th noise derivative has the local factor
$s^a f_{y,z^{a+2}}[\eta^{a+2}]$ under the integral.
When derivatives hit $\chi_y(t\eta)$ they give further local powers of $\eta$;
on that support $|t\eta_{S_y}|$ is bounded.  Taylor's vanishing remainder
therefore removes every apparent negative noise power.  The resulting
insertions and their first detail derivatives have polynomial growth of
bounded degree, which can conservatively be bounded by
$C(1+|\eta_{S_y}|^7)$ through the orders used here.  Density terms are treated
the same way.  Retained derivatives differentiate $z_*$ and local coefficients,
with the weights of \cref{lem:v7-centre}.  The finite local moment bound at
order $112$ pays for the stated $L^{16}$ norms.  This intentionally generous
moment order is available from \cref{eq:v7-noise-moments}; no exponential
moment of a quartic insertion is assumed.
\end{proof}

\begin{proposition}[Exact reflection of signed noise]
\label{prop:v7-even}
For all admitted $t,c$,
\begin{equation}\label{eq:v7-even}
 U_{-t}(c,\eta)=U_t(c,-\eta),\qquad
 \mathcal W(-t,c)=\mathcal W(t,c).
\end{equation}
In particular $\mathcal W(0,c)=0$ and its first and third noise derivatives
vanish at zero, including after four retained derivatives.
\end{proposition}
\begin{proof}
Every nonquadratic argument in \cref{eq:v7-U} is $t\eta$, and its sole explicit
denominator is $t^2$.  The quadratic reference is even.  The change of
variables $\eta\mapsto-\eta$ in \cref{eq:v7-W} proves the assertion.
Local polynomial domination and quadratic confinement justify the required
finite derivatives.  No evenness of $f$, $\ell$, or their centred Taylor
coefficients is required.  In particular, odd cubic action and odd measured
linear density insertions have not been deleted from the integrand.
\end{proof}

\section{Eight connected entries give a fourth-noise bound without volume loss}
\label{sec:v7-connected}

The general cumulant identities used next are inherited normalized
identities, not a replacement by moments of an extensive random variable.
For a finite set $A$ of labelled derivatives of the complete potential,
\begin{equation}\label{eq:v7-loghierarchy}
 \partial_A\left[-\log\int e^{-U}\right]
 =\sum_{\pi\in\Pi(A)}(-1)^{|\pi|-1}
                  \kappa\bigl((\partial_B U)_{B\in\pi}\bigr).
\end{equation}
All derivatives in a block commute; equal physical directions may have
distinct labels.  Formula \cref{eq:v7-loghierarchy} follows by one derivative
of the log normalizer and repeated use of the normalized score rule already
recorded in V6 and \cite{YM4V6}.

For $\partial_t^4\partial_I$ there are four labelled strength directions and
$r=|I|$ retained marks.  Hence the largest connected order is $4+r\le8$.
The partition counts for $r=0,1,2,3,4$ are
\[
                         15,\quad52,\quad203,\quad877,\quad4140.
\]
These are counts of terms before any cancellations or identical-direction
coefficients are combined.  They are finite-order constants, not bounds for
an unrestricted analytic series.

\begin{lemma}[The finite-product argument at eight entries]
\label{lem:v7-eight}
The product-covariance implication of \cref{lem:v6-cut,thm:v6-covcompiler}
extends to $m\le8$ with $L^{16}$ insertion and gradient bounds.  If the
covariance response decays at rate $\sigma$, its resulting connected tree
bound may be taken at rate $\sigma/7$, before the external attachment weights
are summed.
\end{lemma}
\begin{proof}
For a product of $b\le8$ factors its $L^2$ gradient is bounded by the product
rule and Holder with exponent $2b\le16$.  Apply the covariance estimate to
any two nonempty disjoint groups.  In the connected moment formula, cut a
longest edge of a minimal spanning tree on the $m$ insertion centres and
replace mixed moments by moments with those two groups independent.
The replaced cumulant vanishes.  Telescoping each moment product gives the
same finite bound as in \cref{lem:v6-cut}, with
$b_m=\sum_{\pi\in\Pi([m])}(|\pi|-1)!$ and the factor $m^2b_m$.
The longest-edge distance is at least the tree length divided by $m-1\le7$.
This gives the stated decay.  The proof uses no additional property specific
to $m\le7$; the change of moment exponent and decay loss are kept explicitly.
\end{proof}

\begin{theorem}[Uniform mixed fourth-noise and fourth-retained derivatives]
\label{thm:v7-fourthnoise}
For the measured local extension above, there exist $a>0$ and finite $C_r$,
independent of volume, $t\in[-t_0,t_0]$, and the small retained field, such that
\begin{equation}\label{eq:v7-fourthnoise}
 \|\partial_t^4\mathcal W(t,\cdot)\|_{a,r}\le C_r,
                      \qquad 1\le r\le4,
 \qquad
 \frac{\sup_c|\partial_t^4\mathcal W(t,c)|}{n_\Lambda}\le C_0.
\end{equation}
The same assertion, with its own constants, holds for the second noise
derivative.  Here and below the anchored seminorm is exactly
\cref{eq:v6-norm}.
\end{theorem}
\begin{proof}
The reference normalizer in \cref{eq:v7-W} has no $t$ dependence, so it drops
out of every expression under consideration.  Apply
\cref{eq:v7-loghierarchy} and expand each insertion into its local stencils.
\Cref{lem:v7-noise-stock} gives the actual tilted covariance response and
local $L^{16}$ values and gradients, throughout the signed interval.
\Cref{lem:v7-eight} therefore bounds every connected term by a sum of
spanning-tree exponentials, with the attachment weights of
\cref{lem:v7-centre}.  For a cumulant with one entry, use its direct local
moment bound instead.

When $r\ge1$, root the insertion tree at the first retained site and attach
all other external sites to the insertions carrying their labels.  The total
edge length is at least the external tree length in
$\|\cdot\|_{a,r}$.  Choose $a$ smaller than both available decay rates and
sum leaves successively.  Spatial row sums and attachment row and column
sums are uniformly finite.  This is the rooted summation argument of
\cref{thm:v6-rooted}, now on at most eight insertions.  The finite partition
count completes the bound.  Without retained marks choose one insertion
centre as a root; its sum costs at most a fixed multiple of $n_\Lambda$,
while all remaining tree sums are bounded.  This proves the pressure estimate.
The second-strength case has at most $2+r\le6$ entries and is covered by the
same argument.  Finite-dimensional dominated differentiation justifies the
identities first; the rooted estimates are what remove the volume loss from
the final constants.
\end{proof}

\section{The complete measured two-loop coefficient and parity improvement}
\label{sec:v7-coefficient}

All tensors in this section are evaluated at the actual moving classical
minimum, in the same independent fibre coordinates.  Define the local
polynomial insertions
\begin{equation}\label{eq:v7-PQ}
 P_c(\eta)=\frac16 f_{zzz}(z_*,c)[\eta^3]-\ell_z(z_*,c)[\eta],
 \qquad
 Q_c(\eta)=\frac1{24}f_{zzzz}(z_*,c)[\eta^4]
                              -\frac12\ell_{zz}(z_*,c)[\eta^2].
\end{equation}
Let $\gamma_c=N(0,C_c)$, and put
\begin{equation}\label{eq:v7-B1}
 \boxed{\displaystyle
 \mathcal B_1(c)=\E_{\gamma_c}Q_c
                         -\tfrac12\E_{\gamma_c}P_c^2.}
\end{equation}
Although $P_c,Q_c$ are extensive sums, their indicated expectation combination
is interpreted by connected local contractions.  The mean of $P_c$ is zero
by Gaussian parity; no mean of a retained-dependent reference is frozen when
this formula is differentiated in $c$.

\begin{proposition}[Measured Wick formula with the tadpole shift completed]
\label{prop:v7-Wick}
Write $T=f_{zzz}$, $V=f_{zzzz}$, $p=\ell_z$, $L_2=\ell_{zz}$, and
$\vartheta_i=\sum_{a,b}T_{abi}(C_c)_{ab}$.  Summing all repeated fibre indices,
\begin{equation}\label{eq:v7-Wick}
 \boxed{\begin{aligned}
 \mathcal B_1={}&\frac18V_{abcd}C_{ab}C_{cd}
                -\frac12(L_2)_{ab}C_{ab}\\
 &-\frac1{12}T_{abc}T_{def}C_{ad}C_{be}C_{cf}\\
 &-\frac12\left(p-\frac12\vartheta\right)^T
                  C\left(p-\frac12\vartheta\right).
 \end{aligned}}
\end{equation}
This is a reformulation of the finite-volume Laplace coefficient, not a
new beta-function coefficient or an assertion of a sign for the complete
Yang--Mills response.
\end{proposition}
\begin{proof}
The four-leg Gaussian moment has three pairings; the symmetry of $V$ gives
the factor $3/24=1/8$.  The square of the cubic has six pairings joining all
three legs across its two vertices, and nine pairings with an internal pair
at each vertex and one cross pair.  With its prefactor $-1/72$, these are the
third displayed term and $-\vartheta^TC\vartheta/8$.  The cubic--linear
cross term has three pairings and contributes $\vartheta^TCp/2$.
The density-linear square gives $-p^TCp/2$.  Complete those last three
terms to the displayed square, and retain the density Hessian term.
All Wick contractions use the same $C=H_c^{-1}$.  The actual pivot and Haar
parts of $p,L_2$ are thus required, not optional addenda.
\end{proof}

\begin{theorem}[Uniform fourth-retained-derivative expansion of the extension]
\label{thm:v7-extension}
For $|t|\le t_0$,
\begin{equation}\label{eq:v7-extension}
 \boxed{\displaystyle
 \|\mathcal W(t,\cdot)-t^2\mathcal B_1\|_{a,r}
                        \le\frac{C_r}{24}|t|^4,\qquad 1\le r\le4.}
\end{equation}
The pressure difference divided by $n_\Lambda$ satisfies the analogous
bound.  The coefficient $\mathcal B_1$ itself has uniform anchored
seminorms through order four and an extensive zeroth-order bound.
\end{theorem}
\begin{proof}
At fixed $\eta$, the cutoff is identically one in a neighbourhood of $t=0$.
Therefore \cref{eq:v7-U} has its exact first two noise coefficients
\[
 U_t=\tfrac12\eta^TH_c\eta+tP_c(\eta)+t^2Q_c(\eta)+O_\eta(t^3).
\]
Differentiate the \emph{normalized} integral: $\mathcal W''(0)=
2\E Q_c-\operatorname{Var}(P_c)=2\mathcal B_1$.  The polynomial moment
bounds justify this identification after retained differentiation.
\Cref{prop:v7-even} gives $\mathcal W(0)=\mathcal W'(0)=
\mathcal W'''(0)=0$.  Taylor's formula is consequently
\begin{equation}\label{eq:v7-taylor4}
 \mathcal W(t,c)-t^2\mathcal B_1(c)
 =\frac16\int_0^t(t-s)^3\partial_s^4\mathcal W(s,c)\,ds.
\end{equation}
For negative $t$ use evenness.  Apply \cref{thm:v7-fourthnoise} and
$\int_0^{|t|}(|t|-s)^3/6\,ds=|t|^4/24$.
The coefficient bound follows from the same theorem for the second noise
derivative.  It could equivalently be proved by local Gaussian contractions
in \cref{eq:v7-Wick}.  An uncut cubic exponential has never been integrated:
the exact globally convex extension, not its cubic truncation, defines every
intermediate law.
\end{proof}

\section{A fifth-entry hard-wall estimate closes the fourth-derivative transfer}
\label{sec:v7-wall}

The preceding theorem still concerns an extension.  Returning it to the actual
fixed box is a separate step.  This section keeps the source's exact
partition-ratio construction and proves the higher derivative bound that is
needed here.  There is no assertion that all coordinates stay in the chart
simultaneously with probability tending to one uniformly in volume.

In the original $z$ coordinates let
\[
 E_\beta(z,c)=\beta G-\ell_*+
       \frac\beta2(z-z_*)^TH_c(z-z_*)+
       \beta\widehat r(z-z_*,c)-\widehat d(z-z_*,c).
\]
Define the fixed wall and its interpolation by
\begin{equation}\label{eq:v7-walllaw}
 \psi_i(z)=\tfrac12(|z_i|-r)_+^2,\quad
 U_{\lambda,\beta}=E_\beta+\lambda\sum_i\psi_i,\quad
 F_{\lambda,\beta}=-\log\int e^{-U_{\lambda,\beta}}\,dz.
\end{equation}
The wall is independent of $c$; no derivative of its $C^{1,1}$ Hessian in
$c$ is introduced.  Use the inherited bounds, for every fixed $q\ge1$,
\begin{align}\label{eq:v7-wallmoments}
 \|\psi_i\|_q&\le C_qe^{-b_q\beta}(\beta+\lambda)^{-1-1/(2q)},\\
 \|\psi_i'\|_q&\le C_qe^{-b_q\beta}(\beta+\lambda)^{-1/2-1/(2q)}.\notag
\end{align}
For a vector source supported in component $i$ and vanishing off
$\{|z_i|>r\}$, the covariance resolvent has the additional factor
$(m\beta+\lambda)^{-1/2}$ on that source.  Its other side has the ordinary
$\beta^{-1/2}$ factor, and both have the source's spatial exponential.
The dependence of a source component's coefficient on other field variables
does not remove this component-support property
\src{F12}{\texttt{prop:f12v12-layer}, \texttt{prop:f12v12-resolvent}}.

\begin{lemma}[An integrable connected wall insertion through five entries]
\label{lem:v7-wallcluster}
Let $2\le m\le5$, and let $X_1,\ldots,X_{m-1}$ be local score insertions
with values and first fibre gradients bounded in $L^{2m}$ by $A_1,\ldots,A_{m-1}$,
uniformly in $\lambda$.  Their support radii are bounded.  On the exact law
\cref{eq:v7-walllaw},
\begin{equation}\label{eq:v7-wallcluster}
 \left|\kappa(X_1,\ldots,X_{m-1},\psi_i)\right|
 \le C_m e^{-\omega_m\beta}(1+\lambda)^{-1-1/(4m)}
     \left(\prod_{j=1}^{m-1}A_j\right)
     e^{-\nu_m\tau_{\rm MST}(y_1,\ldots,y_{m-1},i)}.
\end{equation}
A smaller common positive exponent and the power $21/20$ work for every
$m\le5$.  The constants are independent of total volume and wall strength.
\end{lemma}
\begin{proof}
First consider a covariance of two products, one of which contains the unique
wall factor.  Differentiating that product gives two types of terms.
If a score is differentiated, the undifferentiated wall factor is bounded by
its $L^{2m}$ norm in \cref{eq:v7-wallmoments}, yielding the power
$1+1/(4m)$ with the ordinary covariance inverse.  If the wall is
differentiated, its norm gives the power $1/2+1/(4m)$ and the
wall-supported inverse contributes another $1/2$.  The product rule and
Holder with $2m$ control all remaining factors, including their gradients.
The support radius changes only a spatial constant.  Thus both terms have
integrable strength decay $1+1/(4m)$ and the distance exponential between
the two groups.  Factors $\beta^{-1/2}$ or $\beta^{-1}$ can be weakened
since $\beta\ge1$.

To pass from products to a connected moment, use the longest-edge cut from
\cref{lem:v6-cut}.  In the moment-partition formula replace mixed moments
by the product of their two cut-group moments.  The resulting full cumulant
is zero.  Telescope each partition's products.  If the changed moment
contains the wall, use the covariance bound just proved.  If it does
\emph{not} contain the wall, another moment factor in that same product
contains it.  Holder and the first line of \cref{eq:v7-wallmoments} supply
exactly the missing exponential rarity and strength power; the changed
score-only moment supplies the cross-distance decay.  This second case is
necessary and cannot be dropped in a higher-order cumulant argument.

All other moment factors are bounded by their $A_j$.  The partition sum is
finite, and the longest cut edge is at least
$\tau_{\rm MST}/(m-1)$.  This proves \cref{eq:v7-wallcluster} after reducing
the spatial rate.  As $m\le5$, $1+1/(4m)\ge21/20$.  The $m=1$ expectation
uses the stronger $L^1$ wall power $3/2$ directly.  No derivative of the
wall's discontinuous second derivative is used.
\end{proof}

\begin{theorem}[Exponentially small hard-wall ratio through four retained derivatives]
\label{thm:v7-wall}
For the actual extension and original fixed box set
\[
 B_\beta(c)=-\log\mu^{\rm ext}_{\beta,c}(\mathcal D).
\]
Then, on the same smaller retained chart, for some $a,b>0$,
\begin{equation}\label{eq:v7-wall-bound}
 \boxed{\displaystyle
 \|B_\beta\|_{a,r}\le C_re^{-b\beta},\quad 1\le r\le4,
 \qquad 0\le n_\Lambda^{-1}\sup_c B_\beta(c)\le C_0e^{-b\beta}.}
\end{equation}
\end{theorem}
\begin{proof}
At fixed $z$, all retained derivatives of $U_{\lambda,\beta}$ equal those
of $E_\beta$.  Decompose them into their actual local centred stencils.
\Cref{lem:v7-centre} supplies the four-mark attachment weights.
Each needed local score and its first fibre derivative has $L^{10}$ norm at
most $C\beta$ times those weights, uniformly in $\lambda$.  Outside the
cutoff support the extension is quadratic in $z-z_*$, and the inherited
all-order local wall moments bound that growth.  Inside it, bounded local
analytic derivatives suffice.  Higher $c$ derivatives do not increase the
integration-variable support.

For a labelled retained set $I$, $1\le |I|=r\le4$, normalized differentiation
gives the exact fixed-coordinate identity
\begin{equation}\label{eq:v7-wallderivative}
 \partial_\lambda\partial_I F_{\lambda,\beta}
 =\sum_i\sum_{\pi\in\Pi(I)}(-1)^{|\pi|}
                \kappa\bigl(\psi_i,(E_{\beta,B})_{B\in\pi}\bigr).
\end{equation}
There are at most $r+1\le5$ entries, one of which is a wall insertion.
Apply \cref{lem:v7-wallcluster}, expand the other entries locally, and root
at the first external retained site.  The weighted sums are the same finite
spatial tree sums used earlier.  Uniformly in $\lambda$,
\[
 \|\partial_\lambda F_{\lambda,\beta}\|_{a,r}
       \le C_r\beta^r e^{-b'\beta}(1+\lambda)^{-21/20}.
\]
For $r=0$ divide its extensive expectation sum by $n_\Lambda$ and use the
$L^1$ wall moment.  The strength integral is finite and explicit:
\[
                    \int_0^\infty(1+\lambda)^{-21/20}\,d\lambda=20.
\]
Integrate and absorb the fixed powers of $\beta$ into a smaller positive
exponent.

At each fixed finite volume, $e^{-\lambda\sum_i\psi_i}$ decreases to
$1_{\mathcal D}$.  On compact smaller retained charts, derivatives of the
integrand through order four are bounded by a polynomial times a confining
Gaussian, independently of $\lambda$.  Thus the partition function and its
first four derivatives converge.  The limiting partition function is
positive, so the logarithmic differences converge as well, with
$F_{\infty,\beta}-F_{0,\beta}=B_\beta$.  This identifies the differentiated
hard-wall limit rather than treating a uniform bound in $\lambda$ as
integrable.  The inherited $C^{1,1}$ wall resolvent can be approximated by
convex smooth walls at fixed $\lambda$ as in f12; no additional wall derivative
is required in this extension.
\end{proof}

\begin{proposition}[The actual and extended potentials on the same original box]
\label{prop:v7-samebox}
Let
\[
 A_\beta=\beta f-\ell,\qquad
 \mathcal E_\beta=-\log\int_{\mathcal D}e^{-A_\beta}
                             +\log\int_{\mathcal D}e^{-E_\beta}.
\]
Then $\|\mathcal E_\beta\|_{a,r}\le C_re^{-b\beta}$ for $1\le r\le4$,
with the analogous pressure-per-coordinate bound.
\end{proposition}
\begin{proof}
Use exactly the source's fixed-domain interpolation
$E_\beta+s(A_\beta-E_\beta)$, $0\le s\le1$.  Its local defect is
\[
 \Delta_y=(1-\chi_y(z-z_*))\{\beta r_y(z-z_*,c)-d_y(z-z_*,c)\}.
\]
By the inherited common interior minimum, convexity and local tails, the
event on which this defect is nonzero has probability at most
$Ce^{-b\beta}$, uniformly in interpolation strength and volume.
At fixed $z$, every derivative through four retained marks and one detail
gradient vanishes on the same inner open set.  Its magnitude on the fixed
box is at most $C\beta$ times the attachment weights, including derivatives
of the moving $z_*$.  Therefore its $L^{10}$ norm is bounded by
$C\beta e^{-b'\beta}$ times those weights.

Apply the normalized expectation hierarchy to
$\partial_s[-\log\int_{\mathcal D}e^{-E_\beta-s\Delta}]=\E_s\Delta$.
Every term through four retained derivatives has one rare defect insertion
and at most four local score insertions.  The inherited boxed covariance
bound applies to products of their gradients.  The longest-edge cumulant
argument gives a connected spatial tree with the rare factor retained.
If the changed moment in that argument omits the defect, its remaining
partition factor contains the defect and supplies the same rarity; this is
the finite-box analogue of the two cases in \cref{lem:v7-wallcluster}.
Root the tree at a retained mark and integrate over $s\in[0,1]$.
Fixed powers of $\beta$ are absorbed by reducing the rarity exponent.
The zero-derivative sum is extensive and is divided by $n_\Lambda$.
All differentiations use the original fixed box, so no moving boundary
term is suppressed.  This extends, rather than repeats, the two-derivative
statement of f12 V11.
\end{proof}

\section{The actual fixed-box theorem and a compact Wilson calibration}
\label{sec:v7-main}

\begin{theorem}[Volume-uniform measured local two-loop expansion in $C^4$]
\label{thm:v7-main}
For the actual local compact-block integral \cref{eq:v7-local-integral} with
the inherited inputs (L1)--(L3), there exist a fixed smaller retained chart,
$a,b>0$, $\beta_0<\infty$, and $C_r<\infty$, independent of finite volume,
such that
\begin{equation}\label{eq:v7-main-expansion}
 \boxed{\displaystyle
 \widehat W_\beta^{\rm loc}(c)
 =\beta G(c)+\frac{n_\Lambda}{2}\log(\beta/2\pi)
          +\widehat Q(c)+\frac{\mathcal B_1(c)}{\beta}
          +\mathcal R_\beta(c),}
\end{equation}
where $\mathcal B_1$ is \cref{eq:v7-B1}, and
\begin{equation}\label{eq:v7-main-bound}
 \boxed{\displaystyle
 \|\mathcal R_\beta\|_{a,r}\le C_r\beta^{-2},\quad 1\le r\le4,
 \qquad \frac{\sup_c|\mathcal R_\beta(c)|}{n_\Lambda}\le C_0\beta^{-2}.}
\end{equation}
The coefficient $\mathcal B_1$ has uniform anchored seminorms through four
retained derivatives.  The constants may depend on the original local
chart, group, and bounded primitive stencil derivatives, but not on
$\beta$ or volume.
\end{theorem}
\begin{proof}
Put $t=\beta^{-1/2}>0$.  Rescaling the exact whole-space extension gives
\[
 \widehat W_\beta^{\rm ext}
 =\beta G+\frac{n_\Lambda}{2}\log(\beta/2\pi)
                                      +\widehat Q+\mathcal W(t,c).
\]
This identity uses the actual measured density, not its expansion.  The
original fixed-box law is related to it by the source's exact ratios,
\begin{equation}\label{eq:v7-exact-transfer}
 \widehat W_\beta^{\rm loc}
            =\widehat W_\beta^{\rm ext}+B_\beta+\mathcal E_\beta.
\end{equation}
Coarse Haar conversion cancels in each ratio.  Apply
\cref{thm:v7-extension,thm:v7-wall,prop:v7-samebox}; the first error is
$O(t^4)=O(\beta^{-2})$, and the other two are exponentially small through
four retained derivatives.  Their pressure-per-coordinate versions give
the zeroth-order bound.  The coefficient bound is already proved in
\cref{thm:v7-extension}.  Thus the auxiliary extension has disappeared from
the conclusion.  It has not been identified with the compact complement.
\end{proof}

In the source's convention $\widehat W=\cdots+\widehat Q-\beta^{-1}A_1+\cdots$,
one has $\mathcal B_1=-A_1$.  The result is not the discovery of a previously
unknown fixed-volume $A_1$.  It upgrades that coefficient to a
volume-uniform, localized fourth-derivative expansion with a controlled
original-box remainder.  In particular, the measure derivatives in
\cref{eq:v7-Wick} must not be replaced by the primitive Haar derivatives alone
when the actual inverse-pivot density is also present.

\begin{proposition}[A sharp actual $\mathrm{SU}(2)$ one-plaquette calibration]
\label{prop:v7-SU2}
Use the one-plaquette normalization of \cref{eq:v2-Wilson-action}, with
positive scalar coefficient $a=e^c$:
\[
 f(Y,c)=a\bigl(1-\tfrac12\Tr Y\bigr),\qquad Y\in\mathrm{SU}(2).
\]
For its full normalized Haar integral, the negative logarithm after subtracting
its Gaussian term and field-independent normalization has the expansion
\begin{equation}\label{eq:v7-SU2}
 \frac{3e^{-c}}{8\beta}+\frac{3e^{-2c}}{16\beta^2}
                                              +O_{C^4}(\beta^{-3})
\end{equation}
on each bounded $c$ interval.  Consequently the fourth $c$ derivative of the
remainder after the first displayed term equals
$3\beta^{-2}+O(\beta^{-3})$ at $c=0$.  The power in
\cref{eq:v7-main-bound} therefore cannot generally be improved without
subtracting another coefficient.
\end{proposition}
\begin{proof}
Write $Y=\exp(y\cdot E)$ and $\rho=|y|$.  In these coordinates
$f=a(1-\cos\rho)$ and normalized Haar measure has, up to its fixed constant,
the density $(\sin\rho/\rho)^2\,d^3y$.  This follows by dividing the spherical
volume element $\sin^2\rho\,d\rho\,d\Omega$ of the unit three-sphere by
$\rho^2d\rho\,d\Omega$.  Hence
\[
 f=\tfrac a2\rho^2-\tfrac a{24}\rho^4+\tfrac a{720}\rho^6+O(\rho^8),
 \qquad \ell=-\rho^2/3-\rho^4/90+O(\rho^6).
\]
For a three-dimensional Gaussian of covariance $a^{-1}I$, let $R=|\eta|^2$.
The noise-square and noise-fourth potential coefficients are
\[
 Q=-aR^2/24+R/3,\qquad Q_4=aR^3/720+R^2/90.
\]
Its radial moments are
$\E R^k=(3\cdot5\cdots(2k+1))a^{-k}$.
They give $\E Q=3/(8a)$ and
$\E Q_4-\operatorname{Var}(Q)/2=3/(16a^2)$.
Equivalently, expansion of the \emph{density} itself gives coefficients
$-3/(8a)$ and $-15/(128a^2)$ before taking its negative logarithm;
the two calculations agree.

This finite-dimensional expansion may be justified with a fixed neighbourhood
of the identity and Gaussian rescaling; analytic Taylor remainders and a
smaller Gaussian confinement bound dominate the integrands through the next
finite order.  Outside that neighbourhood $1-\cos\rho$ has a strictly
positive lower bound on the compact group, so the discarded part and its
first four $c$ derivatives are exponentially small, uniformly on bounded
$c$ intervals.  This proves \cref{eq:v7-SU2}.  Finally
$\partial_c^4(3e^{-2c}/16)|_{c=0}=3$.

This is an actual compact one-plaquette calibration.  It is not the
many-link root-path constrained block, not a value of its $\mathcal B_1$,
and not a computation of a Yang--Mills renormalization coefficient.
\end{proof}

\section{What this removes from the next RG estimate}
\label{sec:v7-next}

V6's sought local fourth-derivative conclusion is now obtained on the \emph{actual
fixed local Wilson box}: the exact measured response has a complete
localized two-loop approximation with an error uniform through four retained
field derivatives.  The signed-noise interpolation, the high local moments, the eight connected
entries, and the hard-wall passage are accounted for.  This route does not
assert that every hypothesis of V6's different fixed-residual, linear-strength
interpolation has thereby been established on the full compact law.
The initial Gaussian inverse has not been reproved, and the conclusion is
not another bounded-interaction comparison model.

At a translation-invariant retained background, the anchored spatial bounds
also control fixed momentum derivatives of the derivative tensors.  More
precisely, after fixing one external site, a Fourier derivative multiplies a
kernel coefficient by a polynomial in external separations.  Its degree-$s$
absolute value is bounded by $C_s\tau^s$.  For $a'<a$,
\[
                 \tau^s e^{a'\tau}\le C_{s,a-a'}e^{a\tau}.
\]
Thus every fixed number of momentum derivatives of the remainder's
$r$-field vertex, $1\le r\le4$, obeys the same $O(\beta^{-2})$ bound in a
smaller spatial weight.  This is a consequence of the new \emph{nonlinear
local} derivative bound, not an identification of field order with momentum
order.  It does not supply a Ward identity for a Cartesian cutoff that fails
to preserve the full gauge symmetry; V6's exact equivariance requirement
remains in force.

If a separately controlled scale family has uniform copies of (L1)--(L3),
$\beta_j\ge b(j+j_0)$, compatible physically normalized source norms and
bounded subsequent error transport, the contribution of this remainder has
\begin{equation}\label{eq:v7-tail}
 \sum_{j\ge J}\|\mathcal R_{\beta_j}\|_{a,r}
      \le\frac{C_r}{b^2(J+j_0-1)},\qquad J+j_0>1.
\end{equation}
This is an integral bound for the inverse-square series, not a proof of the
stated YM running-coupling or stability hypotheses.  The coefficient
$\mathcal B_1/\beta_j$ generally is not summable; it belongs to the retained
running action/density and must be matched rather than discarded.  V6's
physical-site weights and amplification factors still apply.

\paragraph{The genuinely remaining common-scale input.}
The local chart and stencil derivatives used here are verified in the source
for the primitive small-field block.  After many exact nonlinear steps the
stencils are transported, not identical primitive Wilson functions.  V5's
all-level Gaussian locality does not by itself imply uniform bounds for
those nonlinear derivatives or a common compact small-field domain.  The
next substantive target is to propagate the local hypotheses on the actual
renormalized packet, with an invariant local subtraction and the
$\mathcal B_1$ density retained, rather than to prove another generic
moment identity.  Extending this theorem to a growing nonlinear family
without that stock would simply conceal the previous bottleneck.

\paragraph{The compact complement remains a different problem.}
Equation \cref{eq:v7-exact-transfer} returns to the \emph{original local box}
only.  It does not compare that box with all compact configurations, other
conditional wells, transition cutoffs or high active-source occupancy.  None
of the convex estimates is applied to the unrestricted hidden diffusion
whose metastability is already proved.  No continuum measure, physical
noncollapse, or fixed-physical-time transfer contraction follows from
\cref{thm:v7-main}.  The $\beta^{-1}$ coefficient here is a local Laplace
coefficient, not the two-loop beta function or a physical mass.

\section{Exact checks and append-only record}
\label{sec:v7-verification}

The executable \path{verify_ym_measured_two_loop_V7.py} contains twenty-five
passing exact test families.  Its independent bivariate Taylor ring stores
ordinary coefficients through fourth noise and fourth retained order.  It
constructs a finite symmetric normalized density with a moving reference,
then compares its direct log series with all labelled cumulant partitions
through eight entries.  The counts include $4140$ partitions at the largest
mixed derivative; no finite-difference approximation is used.

Other checks evaluate the complete measured coefficient on a correlated
three-coordinate rational Gaussian, independently expand its next formal
normalizer coefficient, verify the tadpole square in
\cref{eq:v7-Wick}, and test a nonorthogonal common-coordinate change.
An independent nonlinear one-dimensional coordinate change of an exactly
Gaussian integral gives a zero correction only when its logarithmic
Jacobian is kept.  Omitting that Jacobian produces a nonzero result.
The actual $\mathrm{SU}(2)$ primitive is checked both by its logarithmic Haar
series and by direct multiplication of its density series, including the
sharp fourth-retained-derivative coefficient.  Further rational tests record
the integrable wall powers, their $20$ strength integral, the $1/24$ Taylor
constant, and the conditional inverse-square tail estimate.

These calculations check algebra, normalization and finite calibrations.
They do not numerically verify the all-volume analytic theorem, evaluate the
full constrained YM block's $\mathcal B_1$, or recertify the older large
computational chains.  The analytic conclusions rest on the proofs above
and the expressly inherited local inputs.  Prior reports are retained as
archival reports unless a fresh execution is separately identified.

The entire mathematical body of V1--V6, including its bibliographies and
end markers, is preserved byte-for-byte.  Only the version display in the
preamble is advanced.  The bundle records the preserved-body hash, input
source hashes, new code and report, and the output provenance.  Future work
is to append after the following marker; a correction to any earlier claim
must be identified explicitly rather than silently rewriting its proof.

\begin{thebibliography}{99}
\bibitem[K24]{KatsevichV7}
A.~Katsevich, \emph{The Laplace asymptotic expansion in high dimensions},
arXiv:2406.12706 (2024).  \url{https://arxiv.org/abs/2406.12706}.
External context only.  The dimension-dependent hypotheses and estimates
of that work are not substituted for the localized lattice estimates here.
\end{thebibliography}

% END OF VERSION 7 MATHEMATICAL BODY.
% Append subsequent requested versions below this marker; retain V1--V7 above.

\clearpage
\section{V8: closing the field--loop packet before attempting scale induction}
\label{sec:v8-context}

The preceding version gives a uniform $C^4$ local remainder after the complete
$\beta^{-1}$ Laplace coefficient has been retained. It does not say that fourth
field jets of every coefficient form an iteratable packet. In fact they do not.
A sixth-degree term in an intermediate one-loop density can feed a fourth-degree
two-loop coefficient at the next elimination. This version resolves that finite
closure issue before attempting a nonlinear scale bound.

Write $h=\beta^{-1}$. The three negative-log coefficient functions are denoted
$S,Q,B$, so the exponent is $S/h+Q+hB$; in the original primitive Lebesgue
chart $Q=-\ell$. These labels are loop orders zero, one and two of the effective
action $S+hQ+h^2B$. They are not momentum orders or coefficients of the gauge
beta function. We prove that the joint field--loop packet
\[
                J^8S,\qquad J^6Q,\qquad J^4B
\]
is closed under the measured finite two-loop pushforward, using the actual
block through order seven. Closure means determination of these finite Taylor
coefficients and composition of their maps. It does not mean that their norms
are bounded through infinitely many nonlinear steps, or that finite jets
determine an exact nonperturbative integral.

\paragraph{Nonduplication.}
The one-loop compiler and its composition are inherited from V4; the measured
Laplace coefficient and actual-box remainder are inherited from V7. The new
results are the sharp two-loop derivative budget, its invariant joint
filtration, the feed from the transported one-loop density, and an executable
higher-order oracle for the literal primitive Wilson/root-path/Haar packet.
An actual $\mathrm{SU}(2)$ compact convolution supplies a noncommutative
two-stage calibration with a fully evaluated retained background and a direct
independent four-link check. Its arbitrary-length coefficient formulas also
expose a scale loss which finite composition alone cannot remove. This compact
convolution is not identified with the four-dimensional parity block.

Perturbative composition is a general method, not a novelty claim about
Gaussian integration or Feynman diagrams; \cite{CMRV8} provides external
perturbative-gluing context. No theorem from that reference is used below to
supply a missing estimate for the present regulator. The mathematical
statements here follow from the specified finite local integrals and charts.

\section{Measured two-loop pushforward, including the incoming correction}
\label{sec:v8-map}

Let $F:\R^n\to\R^m$ be an $h$-independent smooth submersion near zero, with
$F(0)=0$, and let $DS(0)=0$. Assume that $D^2S(0)$ is positive on
$\ker DF(0)$. Work on the specified local stationary branch and choose a
regular chart
\[
 X:(y,c)\longmapsto x,\qquad F(X(y,c))=c,
\]
with $y\in\R^{n-m}$. The chart is used on a neighbourhood where its Jacobian
does not vanish. A compatible localization equal to one near the branch is
understood. All conclusions in this section are finite-regulator asymptotic
coefficient identities; they need neither a global minimum classification nor
an ambient inverse of a retained massless Hessian.

For the incoming density
\[
                 \exp[-S(x)/h-Q(x)-hB(x)]\,dx
\]
define the chart data
\begin{equation}\label{eq:v8-chart}
 f=S\circ X,\qquad q=Q\circ X-\log|\det DX|,\qquad b=B\circ X.
\end{equation}
In particular, the coarea/inverse-pivot factor appears in $q$ exactly once.
It is not permissible to use $Q\circ X$ here and omit the chart determinant,
nor to add the same determinant again after a bordered representation of it.
Let $y_*(c)$ be the local classical minimum of $f(\cdot,c)$ and put
\[
 H=f_{yy}(y_*,c),\qquad C=H^{-1},\qquad G(c)=f(y_*,c).
\]
All fibre derivatives in the next statement are evaluated at $(y_*(c),c)$.

\begin{theorem}[The measured two-loop step]
\label{thm:v8-step}
With the field-independent Gaussian factor $(2\pi h)^{(n-m)/2}$ separated,
the local pushforward has negative logarithm
\begin{equation}\label{eq:v8-step}
              h^{-1}G(c)+Q_+(c)+hB_+(c)+O(h^2),
\end{equation}
where
\begin{align}
 Q_+&=q(y_*,c)+\tfrac12\log\det H,\label{eq:v8-Qplus}\\
 B_+&=b(y_*,c)+\E_C\!\left[\tfrac1{24}f_{yyyy}[\eta^4]
                                      +\tfrac12q_{yy}[\eta^2]\right]
       -\tfrac12\E_C\!\left[\tfrac16f_{yyy}[\eta^3]+q_y[\eta]\right]^2.
                                                        \label{eq:v8-Bplus}
\end{align}
Here $\E_C$ is the centered Gaussian expectation of covariance $C$. If
$T=f_{yyy}$, $V=f_{yyyy}$ and $\vartheta_i=T_{abi}C_{ab}$, then equivalently
\begin{equation}\label{eq:v8-contracted}
\boxed{\begin{aligned}
 B_+={}&b(y_*,c)+\tfrac18 V_{abcd}C_{ab}C_{cd}
                    +\tfrac12 q_{ab}C_{ab}\\
       &-\tfrac1{12}T_{abc}T_{def}C_{ad}C_{be}C_{cf}
        -\tfrac12(q_y+\tfrac12\vartheta)^TC(q_y+\tfrac12\vartheta).
\end{aligned}}
\end{equation}
The formulas are independent of the chosen regular fibre chart, provided
\cref{eq:v8-chart} is transformed with its full Jacobian.
\end{theorem}
\begin{proof}
Set $y=y_*+\sqrt h\,\eta$. The exponent after the classical and stationary
amplitude factors are removed is
\[
 \tfrac12\eta^TH\eta
 +\sqrt h\left(\tfrac16 f_{yyy}[\eta^3]+q_y[\eta]\right)
 +h\left(\tfrac1{24}f_{yyyy}[\eta^4]+\tfrac12q_{yy}[\eta^2]+b(y_*,c)\right)
 +O_\eta(h^{3/2}).
\]
The $\sqrt h$ term is odd. Expand the normalized exponential through $h$,
then take its negative logarithm. This gives \cref{eq:v8-Qplus,eq:v8-Bplus}.
The pairings in \cref{prop:v7-Wick}, with $\ell=-q$, give
\cref{eq:v8-contracted}; in particular the sign of the density-gradient
shift is positive inside the last square in this convention.

At any fixed regulator and on a compact smaller retained domain, the local
positive Hessian and ordinary smooth Laplace estimates justify these
coefficient identities with a sufficiently high finite differentiability
order. Alternatively they are identities of formal stationary expansions.
Neither interpretation asserts a regulator-uniform remainder here.
Two regular fibre charts describe the same localized integral by the exact
change-of-variables formula. Uniqueness of its asymptotic coefficients proves
chart independence. Changes of localization away from the same strict local
branch leave these coefficients unchanged. A cutoff comparison at growing
volume still requires the separate estimates of V7 or another proved bound.
\end{proof}

The incoming $B$ is evaluated at the classical minimum at this order; its fibre
derivatives first enter higher-loop corrections. Its retained derivatives are,
of course, differentiated along that minimum. The first-stage $Q$, in contrast,
contributes its fibre gradient and Hessian at the second stage. Replacing it by
a unit amplitude loses the last two types of density terms in
\cref{eq:v8-contracted}. An order-$h$ shift of the measured saddle is already
accounted for by the negative square and must not be charged a second time.

An $h$-independent coarse reference $w(c)\,dc$ adds $\log w$ to $Q_+$ and does
not alter $B_+$. Before another step in Lebesgue convention, that conversion
must be undone. An $h$-dependent reference or a change of the running parameter
$h$ has its own coefficient conversion and is not absorbed into notation.

\begin{theorem}[Finite measured composition through two loops]
\label{thm:v8-composition}
For two $h$-independent submersions $F_1,F_2$ on compatible stationary branches,
assume the first and second eliminated Hessians are positive. Denote the map
in \cref{thm:v8-step} by $\mathfrak R^{[2]}_F$. Then, with all normalizing
constants and intermediate density factors retained,
\begin{equation}\label{eq:v8-composition}
 \boxed{\mathfrak R^{[2]}_{F_2}\circ\mathfrak R^{[2]}_{F_1}
                         =\mathfrak R^{[2]}_{F_2\circ F_1}.}
\end{equation}
This equality concerns the classical, one-loop and two-loop coefficient
germs, not equality of arbitrarily chosen nested cutoff domains.
\end{theorem}
\begin{proof}
For compatible localizations, the exact pushforward identity is
\[
 \int du\,\delta(F_2(u)-c)
      \int dx\,\delta(F_1(x)-u)e^{-S/h-Q-hB}
 =\int dx\,\delta(F_2(F_1(x))-c)e^{-S/h-Q-hB}.
\]
One may choose the localization jointly before applying this identity.
Classical minimization is associative on the common branch. Positivity of the
first eliminated Hessian and of the second Schur complement gives positivity
of the combined tangent Hessian. The determinant factors and coordinate
Jacobians therefore give the same leading coefficient, as in V4.
Expand the inner integral through relative order $h$ on a compact small
neighbourhood of the outer saddle. Its remainder is $O(h^2)$ there, with the
finite number of parameter derivatives needed in the outer Laplace estimate.
Its $hB_1$ term contributes $hB_1$ evaluated at the outer classical minimum;
its $Q_1$ term participates in the outer coefficient exactly as in
\cref{eq:v8-Bplus}. Terms of relative order $h^2$ cannot change the coefficient
of $h$. Thus both sides have the same three coefficient germs by uniqueness
of their finite asymptotic expansions. Formal Gaussian expansion gives the
same identity without an asymptotic remainder assertion. A different local
cutoff equal to one near the same branch changes no coefficient, but this
observation is not a uniform compact-complement estimate.
\end{proof}

\section{The closed joint field--loop filtration and its sharp degree budget}
\label{sec:v8-filtration}

Write $J^d$ for the Taylor polynomial through total field degree $d$ at the
stationary zero reference. Total field degree counts both retained and
eliminated fields before a step. It is not the polynomial degree in external
momentum. Constants, including $h$-dependent vacuum coefficients, are recorded
rather than silently reset.

\begin{theorem}[A closed $(8,6,4)$ packet for a fourth-field two-loop target]
\label{thm:v8-budget}
Under the local hypotheses of \cref{thm:v8-step}, the input
\begin{equation}\label{eq:v8-packet}
 \boxed{\bigl(J^8S,\ J^6Q,\ J^4B;\ J^7F\bigr)}
\end{equation}
determines
\[
                         J^8G,\qquad J^6Q_+,\qquad J^4B_+.
\]
This packet is closed under any fixed finite number of measured two-loop
steps. In particular, truncating the first output to these three respective
degrees loses no coefficient in the same packet at the next step. The
composite block needs only the seventh jets of its constituent blocks.
\end{theorem}
\begin{proof}
Put $L=DF(0)$, choose a right inverse $R$ and a basis inclusion $E$ of
$\ker L$, and construct a chart by
\[
               X(y,c)=Ey+Rv(y,c),\qquad F(X(y,c))=c.
\]
The matrix $[E,R]$ is invertible. At every degree $d\ge2$, the homogeneous
degree-$d$ equation for $v$ has identity linear coefficient because $LR=I$.
Its right side depends only on lower chart coefficients and $J^dF$.
Induction therefore determines $J^7X$ from $J^7F$. Since $DS(0)=0$, a degree
at least eight change in $X$ first changes $S\circ X$ in degree at least nine.
Thus $J^8 f$ is determined. The determinant of $DX$ through degree six uses
only $J^7X$, so $J^6q$ is determined; also $J^4b$ is determined.

The stationary equation $f_y(y_*(c),c)=0$ has invertible fibre linearization.
Its homogeneous recursion determines $J^7y_*$ from $J^8f$. In $G=f(y_*,c)$,
the uncomputed degree-eight saddle coefficient is multiplied by the vanishing
first derivative at the origin and cannot enter degree eight. Thus $J^8G$ is
known. The degree-six germ of $H(c)=f_{yy}(y_*,c)$ and its determinant requires
only $J^8f$ and the already determined saddle coefficients. It gives
$J^6Q_+$ with $J^6q$.

In \cref{eq:v8-contracted}, four retained derivatives of $f_{yyyy}$ use at
most $J^8f$. Those of $q_{yy}$ use at most $J^6q$, and those of $b(y_*,c)$
use at most $J^4b$. The other terms use no greater field orders. Derivatives
of $H^{-1}$ are computed by differentiating $HH^{-1}=I$ and introduce no
new primitive derivatives. This proves the one-step determination.
The same reasoning applies to the output packet, whose classical first
derivative at zero still vanishes. The chain rule determines
$J^7(F_2\circ F_1)$ from $J^7F_1,J^7F_2$. Combine with
\cref{thm:v8-composition} to prove finite-stage closure.
\end{proof}

There is a general diagrammatic reason for the descending field degrees. Give
a field degree one and $h$ degree two. For a coefficient of loop order
$\ell$, its joint degree is $d+2\ell$ if its field degree is $d$.

\begin{proposition}[Joint degree cannot be lowered by connected Gaussian elimination]
\label{prop:v8-degree}
After the quadratic completion and removal of vacuum constants from the
interaction expansion, a connected graph with vertices of field degrees
$d_v$ and input loop orders $\ell_v$, $E$ internal Gaussian contractions,
$V$ vertices and $r$ external fields has output loop order
\[
       L=E-V+1+\sum_v\ell_v,
\]
and obeys the exact identity
\begin{equation}\label{eq:v8-degree}
 \boxed{r+2L=2+\sum_v(d_v+2\ell_v-2).}
\end{equation}
Every nonconstant interaction vertex has $d_v+2\ell_v\ge3$. Consequently a
vertex of joint degree greater than $N$ cannot affect an output coefficient
of joint degree at most $N$. For $N=8$ and output loop order at most two,
each non-vacuum graph has at most six interaction vertices and seven internal
contractions.
\end{proposition}
\begin{proof}
A Gaussian contraction contributes a factor $h$ and removes two field legs.
A vertex of loop order $\ell_v$ in the exponential contributes
$h^{\ell_v-1}$. Taking $-h$ times the connected logarithm supplies the final
factor $h$. This proves the expression for $L$. Since
$r=\sum_vd_v-2E$, adding $2L$ gives \cref{eq:v8-degree}.
The classical quadratic terms are in the completed Gaussian, and the
classical linear term vanishes at the reference. Nonconstant one-loop
vertices have degree at least one, and higher-loop nonconstant vertices
have still greater joint degree. All the summands on the right are therefore
positive integers. If one vertex has joint degree exceeding $N$, its summand
alone is at least $N-1$, giving output joint degree greater than $N$.
For a graph of degree at most eight, $V\le6$. At loop order $L\le2$,
$E=L+V-1-\sum_v\ell_v\le V+1\le7$. A field-independent vertex contributes
only its own vacuum term to the connected logarithm and is handled separately.
\end{proof}

The underlying pairing identity is not new combinatorics. Its relevance here
is that it supplies a finite, invariant packet for the measured update. It
prevents an apparent endless requirement to raise field order after each
fixed-loop composition. It supplies no bound on the size of the finitely many
retained tensors as the number of steps grows.

\subsection{Three sharp missing-input tests and one generated intermediate term}

\begin{proposition}[The three degrees in the packet cannot generally be lowered]
\label{prop:v8-sharp}
For a fourth retained derivative of $B_+$, the action degree eight, density
degree six and block degree seven in \cref{eq:v8-packet} are each necessary
in general. More precisely, the following one-fibre local models have
identical lower input jets as their parameter is varied, but the displayed
$B_+$ changes:
\begin{center}\small
\begin{tabular}{@{}p{.61\linewidth}p{.25\linewidth}@{}}
\toprule
Changed primitive & Two-loop output\\
\midrule
$S(c,y)=\frac12(c^2+y^2)+\lambda c^4y^4$, $Q=B=0$, projection block
& $B_+=3\lambda c^4$\\[3pt]
$S(c,y)=\frac12(c^2+y^2)$, $\ell=\lambda c^4y^2$, $B=0$, projection block
& $B_+=-\lambda c^4$\\[3pt]
$S(x,y)=\frac12(x^2+y^2)$, $Q=B=0$, $F=x+\kappa x^3y^4$
& $B_+=-3\kappa c^4$\\
\bottomrule
\end{tabular}
\end{center}
The block example is read on its regular small identity chart. The density
example is integrated for sufficiently small $h$ on a bounded retained domain.
\end{proposition}
\begin{proof}
For the first row the fibre minimum is zero, $C=1$, and
$\E(\lambda c^4\eta^4)=3\lambda c^4$. The perturbation begins in total
field degree eight. In the second row, $q=-\ell$ and
$\frac12\E q_{yy}\eta^2=-\lambda c^4$; the perturbation begins in degree six.
For the last row solve the constraint near zero:
\[
 x=c-\kappa c^3y^4+O(y^8),\qquad
 f=\tfrac12(c^2+y^2)-\kappa c^4y^4+O(y^8).
\]
Its chart density is $|\partial x/\partial c|$, so
$q=\log(1+3\kappa x^2y^4)$ begins at fourth fibre order. It has no first or
second fibre derivative at $y=0$ and makes no contribution at this loop
order. The action quartic gives $-3\kappa c^4$. The block perturbation begins
in total degree seven and does not change lower block jets. Each minimum is
strict on a small chart; no global statement about the nonlinear block is
required for this test.
\end{proof}

\begin{theorem}[A generated sixth-order density feeds the next fourth-order coefficient]
\label{thm:v8-feed}
Take the stable local action, with $\lambda\ge0$,
\[
 S(c,y,z)=\tfrac12(c^2+y^2+z^2)+\lambda c^4y^2z^2,
 \qquad Q=B=0,
\]
and first eliminate $y$, then $z$. The first elimination is exactly Gaussian:
\[
 G_1(c,z)=\tfrac12(c^2+z^2),\quad
 Q_1(c,z)=\tfrac12\log(1+2\lambda c^4z^2),\quad B_1=0.
\]
Its degree-six term $\lambda c^4z^2$ produces
\begin{equation}\label{eq:v8-feed}
             \boxed{B_2(c)=\lambda c^4.}
\end{equation}
Truncating $Q_1$ through total degree four instead gives zero and loses the
exact fourth retained derivative $24\lambda h$ of the final negative-log
correction at $c=0$.
\end{theorem}
\begin{proof}
Perform the first one-dimensional Gaussian integral explicitly. At the
second classical minimum $z=0$, the second fibre derivative of $Q_1$ is
$2\lambda c^4$. In \cref{eq:v8-Bplus} all other terms vanish, leaving
$\lambda c^4$. Independently, eliminating $(y,z)$ together and rescaling by
$\sqrt h$ gives the normalized expectation
\[
              \E\exp[-h\lambda c^4Y^2Z^2],
              \qquad Y,Z\text{ independent }N(0,1).
\]
Its negative logarithm is
$h\lambda c^4-4h^2\lambda^2c^8+O(h^3c^{12})$ as a local asymptotic series.
The coefficient of $c^4$ is exactly $h\lambda$ for every $h>0$, by direct
fourth differentiation at zero. This proves both the missing term and the
normalization. The example concerns a derivative-closure implication, not a
counterexample to the specified Yang--Mills action or to V7's theorem.
\end{proof}

\section{The literal primitive packet can be generated through the required orders}
\label{sec:v8-primitive}

The new higher derivative inventory need not be filled by an invented block.
Use the source's tree-gauge parity cell and chronological two-link paths. The
fine action per cell is
\[
 S_{\rm W}=\sum_{s\in\{0,1\}^4}\sum_{\mu<\nu}
 \left(1-\tfrac12\Tr[
 e^{A_\mu(s)}e^{A_\nu(s+e_\mu)}
 e^{-A_\mu(s+e_\nu)}e^{-A_\nu(s)}]\right).
\]
For a coarse edge the literal primitive block is
\begin{equation}\label{eq:v8-root}
 F_\mu(A)=\log\left[P_0\exp\left(
        \tfrac1{16}\sum_{s\in\{0,1\}^4}\log(P_0^{-1}P_s)\right)\right],
 \quad P_s=e^{A_\mu(s)}e^{A_\mu(s+e_\mu)}.
\end{equation}
The reference path starts at the cell root. These are the inherited
identity-chart definitions, not a replacement by an associative group mean.
Neighbour-cell inputs must remain distinct unless a periodic restriction is
explicitly imposed.

For $A=a\cdot E\in\mathfrak{su}(2)$, the primitive fine Haar logarithm is
\begin{equation}\label{eq:v8-Haar}
 \ell_{\rm Haar}(A)
 =-\frac{|a|^2}{3}-\frac{|a|^4}{90}-\frac{2|a|^6}{2835}
                                      +O(|a|^8).
\end{equation}
It is summed over the 49 non-tree link groups in each cell. In the primitive
Lebesgue presentation $Q=-\ell_{\rm Haar}$ and the incoming $B$ is zero.
The nonlinear block's inverse-pivot factor is supplied by the chart Jacobian
in \cref{eq:v8-chart} or by the equivalent measured border; it is not omitted
and is not added to the primitive Haar density twice.

\begin{proposition}[An explicit complex germ and a high-order primitive oracle]
\label{prop:v8-oracle}
For every local fine Lie input with operator norm at most $r_0=1/64$, the
matrix logarithms in \cref{eq:v8-root} have a common analytic identity branch.
On the associated complex polydisc,
\begin{equation}\label{eq:v8-root-bound}
                 \|F_\mu\|_{\rm op}<\frac18.
\end{equation}
In particular, for a direction of local maximum operator norm at most one,
\[
            \|D^kF_\mu(0)[A^k]\|\le k!\,8^{-1}64^k.
\]
The ordered quaternion power-series algorithm in
\path{ym_primitive_packet_V8.py} computes $J^8S_{\rm W}$, $J^7F$ and
$J^6\ell_{\rm Haar}$ exactly on arbitrary rational directions. Its unperiodized
mode keeps the full neighbour-cell stencil. Homogeneous polarization of these
directional values determines every mixed tensor in the stated primitive
packet.
\end{proposition}
\begin{proof}
For matrices $\|A_i\|\le r$, a product of $p$ exponentials satisfies
$\|\prod_i e^{A_i}-I\|\le e^{pr}-1$. Thus
$\|P_0^{-1}P_s-I\|\le e^{4r}-1<1$ on the stated disc. The convergent logarithm
series bounds each inner logarithm by $m_r=-\log(2-e^{4r})$.
The average of the sixteen logarithms has the same bound. The outer product
has distance from identity at most $e^{2r+m_r}-1$; hence its logarithm is
bounded by
\[
                 -\log\left(2-\frac{e^{2r}}{2-e^{4r}}\right).
\]
At $r=1/64$ this is less than $1/8$. For example,
$e^{1/32}\le32/31$, $e^{1/16}\le16/15$ show that the expression is at most
$-\log(194/217)$. Since
$e^{-1/8}\le1-1/8+1/128=113/128<194/217$, the asserted strict bound follows.
Cauchy's one-variable estimate on the direction disc proves the derivative
bound. Mixed-direction bounds follow by using discs of radius $r_0/k$ in
$k$ labelled directions. These constants concern one fixed primitive stencil,
not the transported nonlinear functions at later scales.

A quaternion jet is a list of rational quaternion coefficients modulo $t^9$.
Ordered multiplication is convolution of those lists. The finite series for
$e^X$ and $\log(I+X)$, truncated at the required degree, are exact because
$X(0)=0$. They preserve noncommutative order. Apply them to the literal
plaquette and path words, then take the scalar trace for the action and the
three imaginary components for the block. Formula \cref{eq:v8-Haar} supplies
the density germ. Its analytic meaning follows from the nonzero sine ratio
near zero; $|(\sin\sqrt z)/\sqrt z-1|$ is bounded by the convergent even
power series at $|z|\le r_0^2$. There is no singular square-root branch in
that even analytic function.

For a homogeneous degree-$k$ polynomial $p_k(A)=D^kf(0)[A^k]/k!$, inclusion--
exclusion gives
\[
 D^kf(0)[A_1,\ldots,A_k]
 =\sum_{J\subseteq[k]}(-1)^{k-|J|}p_k\!\left(\sum_{j\in J}A_j\right).
\]
Only monomials using all $k$ labelled directions survive, proving the claim
about mixed tensors. The exponentials and logarithms in the program are
computed, not fitted to sampled values. A finite list of evaluations checks
the implementation; the formal-series argument is its general justification.
\end{proof}

For the fully specified rational periodic direction in the report, the action
coefficients at degrees two and eight are respectively
\[
                     \frac{3451}{256},\qquad
                     -\frac{63407775911}{86586540687360}.
\]
All four seventh-degree block vectors and the full unperiodized stencil output
are also stored, together with the input direction on every edge. Independent
checks recover the fine curl quadratic form, the exact block derivative,
commuting-path linearity, and global conjugation covariance. Changing the
reference path is tested separately, rather than assuming its high-order jets
are identical. These are primitive data and regressions, not an evaluation of
the full two-stage four-dimensional $B_+$ tensor.

\section{What V5's spatial estimates pay for in the enlarged finite packet}
\label{sec:v8-spatial}

The joint-degree closure is compatible with anchored spatial norms, but does
not by itself create a scale contraction. This distinction can be made without
leaving the role of the linear inverses unspecified.

For a tensor kernel $T(i_1,\ldots,i_d)$ let its anchored exponentially weighted
row norm be the supremum, over any one distinguished slot, of the sum of
absolute values over the remaining slots with weight
$e^{a\tau(i_1,\ldots,i_d)}$. Use the maximum over slots when the tensor is not
symmetric. For a linear kernel use weighted row and column sums. Component
counts are fixed. The primitive translation-covariant stencil norms are finite;
V5 bounds the Gaussian inverse and harmonic-lift norms on the full parity
fibre.

\begin{proposition}[No additional volume loss in a fixed joint-degree step]
\label{prop:v8-norm}
Suppose the input tensor norms through the degrees in \cref{eq:v8-packet} are
finite, with maximum $M$, and the required linear chart, harmonic and fibre
inverse kernels have norm at most $K$. Then the non-vacuum output packet has
anchored norm bounded by a finite polynomial $\mathcal P_8(M,K)$, with constants
independent of cell volume. Its vacuum coefficients have an analogous bound
per cell. The polynomial is independent of the number of prior steps used to
produce the input, but its arguments are their \emph{actual} input norms.
\end{proposition}
\begin{proof}
The homogeneous chart recursion in \cref{thm:v8-budget} involves only finitely
many products, inverse linear maps and contractions through degree seven.
Its coefficients therefore have a finite polynomial bound in $M,K$ in the
same weighted tensor algebra. Finite determinant and logarithm series give
the same assertion for the transformed density through degree six.

After quadratic completion, each output coefficient is a sum over finitely
many connected contractions. By \cref{prop:v8-degree} there are at most six
interaction vertices and seven internal Gaussian links in the packet. Fix one
external site. Every term's vertices and links form a connected graph joining
all remaining external sites. The external tree distance is at most the sum
of its link lengths and vertex-support tree lengths. Its exponential weight
is therefore absorbed into the factors' weights. Choose a spanning tree of the
contraction graph. Summing vertices away from the anchored external site uses
the maximum anchored tensor norms and the two orientations of linear row/column
norms. The remaining cycle links are bounded by their weighted supremum, which
is no greater than the corresponding row norm. This leaves a product of input
bounds and no free volume sum. With no external site, choose one vertex centre
as an anchor and sum it once; that costs one volume factor, appropriate for the
vacuum coefficient per cell. The finite graph and chart-recursion counts give
$\mathcal P_8$. They are coefficientwise finite calculations, not an all-order
convergent diagram expansion.
\end{proof}

A polynomial majorant need not map a fixed ball into itself. Thus this
proposition cannot justify replacing $M$ by one uniform number along the
nonlinear trajectory. Its contribution is narrower: the derivative inventory
is finite and invariant, and V5 pays for its linear contractions and spatial
sums. The actual remaining estimate is stability of this enlarged measured
packet after the Ward-compatible relevant/marginal projection and physical
rescaling. The Cartesian local cutoff used in V7 is not declared gauge
invariant by this argument. Nor do these finite Taylor norms bound its
higher-order analytic tails needed to reapply the full remainder theorem.

\section{An actual compact two-stage calibration with the complete density retained}
\label{sec:v8-SU2}

This section evaluates a compact model exactly enough to test the new
composition law, including a retained background. It is not a claim that
one-dimensional product blocking replaces the interacting four-dimensional
parity block. On $\mathrm{SU}(2)$ set
\[
 w_h(U)=\exp[-h^{-1}(1-\tfrac12\Tr U)],\qquad
 K_{n,h}=\underbrace{w_h*\cdots*w_h}_{n\text{ factors}},
\]
where convolution uses normalized Haar measure. This is the actual constrained
compact product integral with product of its $n$ links equal to the retained
$W$. Write $W=\exp(\theta E_3)$ on a small real interval about zero. For each
fixed $n$ its local Laplace expansion has the form
\begin{equation}\label{eq:v8-Kn}
 K_{n,h}(W)=c_H^{\,n-1}(2\pi h)^{3(n-1)/2}
       \exp[-G_n(\theta)/h-Q_n(\theta)-hB_n(\theta)+O(h^2)],
\end{equation}
with $c_H=1/(2\pi^2)$, the identity exponential-chart Haar constant.
The $h$-dependent normalization in this display is retained in composition.

\begin{theorem}[Two compact Wilson factors and a fully measured second step]
\label{thm:v8-SU2-two}
For $|\theta|$ sufficiently small, let $C=\cos(\theta/4)$ and
$c=\cos(\theta/2)=2C^2-1$. The two-link coefficients are
\[
 G_2=2(1-\cos(\theta/2)),\quad
 Q_2=\tfrac32\log(2\cos(\theta/2)),\quad
 B_2=\frac{3}{16\cos(\theta/2)}.
\]
Composing this \emph{complete} two-link density with itself gives
\begin{equation}\label{eq:v8-four}
 \boxed{\begin{aligned}
 G_4(\theta)&=4(1-C),\\
 Q_4(\theta)&=3\log2+\log c+\tfrac52\log C,\\
 B_4(\theta)&=\frac{7C^2-4}{32C^3}.
 \end{aligned}}
\end{equation}
In particular,
\begin{equation}\label{eq:v8-fourjets}
 B_4(\theta)=\frac3{32}-\frac5{1024}\theta^2
                      -\frac{97}{196608}\theta^4+O(\theta^6),
 \qquad B_4^{(4)}(0)=-\frac{97}{8192}.
\end{equation}
For this theorem the $O(h^2)$ and its first four retained derivatives are
understood at the fixed number of four factors.
\end{theorem}
\begin{proof}
Identify $\mathrm{SU}(2)$ with the unit quaternions. Rotation invariance of
Haar measure makes
\[
 K_{2,h}(W)=e^{-2/h}\int_{\mathrm{SU}(2)}
                 \exp[(2\cos(\theta/2)/h)\Re V]\,dV.
\]
Its unique small-background minimum is $U=W^{1/2}$. The action near that
minimum is $2-2\cos(\theta/2)\cos|y|$ and the Haar density is
$c_H(\sin|y|/|y|)^2$. Gaussian radial integration, equivalently the primitive
calibration in \cref{prop:v7-SU2}, gives the displayed $G_2,Q_2,B_2$.

For the second step use the Haar coordinate
$U=W^{1/4}\exp(y\cdot E)W^{1/4}$, put $R=|y|^2$, and choose its third axis
along $E_3$. The two scalar arguments of the two-link function are
\[
 u_\pm=c\cos\sqrt R\ \pm\ \sin(\theta/2)y_3
                                \frac{\sin\sqrt R}{\sqrt R},
 \qquad A_\pm=\sqrt{(1+u_\pm)/2}.
\]
The second-stage classical action is $4-2(A_++A_-)$, with minimum at $y=0$.
Its Hessian is
\[
                       H=\operatorname{diag}(c/C,c/C,C).
\]
It is positive on a common small retained interval. Its cubic term vanishes
by $y\mapsto-y$. The quartic term has the form
$\alpha R^2+\zeta y_3^2R+\gamma y_3^4$, where, writing
$s^2=\sin^2(\theta/2)$,
\begin{align*}
 \alpha&=-\frac{c}{24C}+\frac{c^2}{32C^3},\\
 \zeta&=-\frac{s^2}{24C^3}+\frac{3cs^2}{64C^5},\qquad
 \gamma=\frac{5s^4}{512C^7}.
\end{align*}
These follow by the fourth-degree binomial expansion of $A_\pm$, including
both signs. The log density at this step includes both transported one-loop
terms and one Haar factor. Its quadratic part is
\[
 \ell_2^{\rm hom}=l_RR+l_3y_3^2,\quad
 l_R=\frac{3c}{8C^2}-\frac13,\quad l_3=\frac{3s^2}{16C^4}.
\]
Indeed the transported part is $-\tfrac32\log(A_+A_-)$ up to its constant,
and the Haar quadratic part is $-R/3$. The incoming two-loop coefficient,
evaluated at the minimum, is $3/(8C)$.

Let $a=C/c$ and $b=1/C$ be the transverse and longitudinal Gaussian variances.
Then
\[
 \E R=2a+b,\quad \E R^2=8a^2+4ab+3b^2,\quad
 \E(y_3^2R)=2ab+3b^2,\quad \E y_3^4=3b^2.
\]
Formula \cref{eq:v8-Bplus} gives
\[
 B_4=\frac3{8C}+\alpha(8a^2+4ab+3b^2)
             +\zeta(2ab+3b^2)+3\gamma b^2-l_R(2a+b)-l_3b.
\]
Substitute $s^2=4C^2(1-C^2)$ and $c=2C^2-1$. The result simplifies to
$(7C^2-4)/(32C^3)$. The stationary one-loop terms are
$3\log(2C)+\tfrac12\log\det H$, giving $Q_4$ in \cref{eq:v8-four}.
Expansion at zero gives \cref{eq:v8-fourjets}.

For fixed $n=2,4$, a sufficiently small background has the unique global
minimum described by the shortest equally divided group geodesic. One direct
argument is that an action no larger than the displayed candidate has every
link angle below $\pi/2$. On that interval $1-\cos t$ is strictly convex.
The triangle inequality for the group distance followed by Jensen gives
$\sum(1-\cos t_i)\ge n(1-\cos(\theta/n))$, with equality only on the equally
divided short geodesic. Off a fixed neighbourhood of this unique minimum the
compact action has a positive gap, uniformly on a smaller retained interval.
The local analytic Laplace estimates and this compact gap justify the stated
finite-order differentiated expansion. No claim uniform in growing $n$ is
made for that remainder.
\end{proof}

At zero background the incoming term is $3/8$, whereas the second integration's
new Wick correction is $-9/32$. Their sum is $3/32$. Resetting the incoming
$B_2$ to zero would give $-9/32$ rather than $3/32$, even with the same
classical action and covariance. Resetting the intermediate $Q_2$ would
introduce a different error. This is an actual measured compact composition
check, not a choice of counterterms to force agreement.

\subsection{Independent direct four-link verification}

There is an independent identity-background calculation with no intermediate
coefficient inserted. Choose three independent group links and set the fourth
to their inverse product. In their exponential coordinates
$x_1,x_2,x_3\in\R^3$, the zero-field Gaussian covariance is
\[
               (I_3-\boldsymbol 1\boldsymbol 1^T/4)\otimes I_3.
\]
The cubic action is $P=x_1\cdot(x_2\times x_3)$. Writing $R_i=|x_i|^2$, its
quartic part is
\[
 S_4^{\rm hom}=-\frac1{12}\sum_iR_i^2
  -\sum_{i<j}\left[\frac14R_iR_j+
       (x_i\cdot x_j)\left(\frac{R_i+R_j}{6}+\frac{R_k}{2}\right)\right],
 \qquad\{i,j,k\}=\{1,2,3\}.
\]
The Haar contribution to the order-$h$ potential is $\sum_iR_i/3$.
Direct nine-dimensional Wick contraction gives
\begin{equation}\label{eq:v8-direct}
 \E S_4^{\rm hom}=-\frac{45}{32},\qquad
 \E P^2=\frac32,\qquad \E\sum_iR_i/3=\frac94,
 \qquad -\frac{45}{32}-\frac34+\frac94=\frac3{32}.
\end{equation}
The quartic word follows by multiplying three quaternion exponentials and
using invariance of the scalar part under inversion. The cubic-square moment
is also $6\det(I_3-\boldsymbol1\boldsymbol1^T/4)=3/2$.
Thus \cref{eq:v8-direct} checks the staged answer without reusing its
square-root action expansion or its intermediate density.

\section{An all-length compact coefficient exposes the remaining scale issue}
\label{sec:v8-length}

Exact finite composition does not imply bounded unrescaled coefficients. The
same compact product family makes this distinction calculable at arbitrary
length, not just at two sampled stages.

\begin{theorem}[Classical and one-loop germs at every length, and the flat two-loop coefficient]
\label{thm:v8-alln}
For every integer $n\ge2$, the coefficient functions of \cref{eq:v8-Kn} obey
\begin{align}
 G_n(\theta)&=n(1-\cos(\theta/n)),\label{eq:v8-Gn}\\
 Q_n(\theta)&=\tfrac12\log n+\tfrac{n-1}{2}\log\cos(\theta/n)
             +\log\frac{\sin\theta}{\sin(\theta/n)},\label{eq:v8-Qn}\\
 B_n(0)&=\boxed{-\frac{(n-1)(n-5)}{8n}}.\label{eq:v8-Bn}
\end{align}
The logarithms have their continuous real germs at zero; in particular
$Q_n(0)=\tfrac32\log n$. These are identities of finite-length coefficients.
They do not assert a remainder estimate uniform in $n$.
\end{theorem}
\begin{proof}
Replace independent links by their successive partial products
$q_1,\ldots,q_{n-1}\in S^3$, with $q_0=1,q_n=W$. Haar invariance gives exactly
the product Haar measure in these variables. The action is
\[
                         \sum_{i=0}^{n-1}(1-q_i\cdot q_{i+1}).
\]
The stationary short path has $q_i=\exp(i\theta E_3/n)$, proving
\cref{eq:v8-Gn}. With $\delta=\theta/n$, use exponential tangent coordinates
at each $q_i$ and a parallel frame along the great circle. The longitudinal
Hessian is $\cos\delta$ times the Dirichlet tridiagonal matrix $D$ with
entries $2,-1,-1$. Each of the two transverse Hessians is the tridiagonal
matrix $T$ with entries $2\cos\delta,-1,-1$. This follows by expanding
$q_i\cdot q_{i+1}$ to second order; the longitudinal tangent dot product is
$\cos\delta$, whereas transverse parallel tangents have dot product one.
The linear terms cancel at every interior vertex.

The determinant recursions give $\det D=n$ and
$\det T=\sin(n\delta)/\sin\delta$. Thus
\[
            \det H=n(\cos\delta)^{n-1}
                          \left(\frac{\sin(n\delta)}{\sin\delta}\right)^2.
\]
The Haar density at the origin of each tangent chart is $c_H$, independent of
the retained background, and already extracted in \cref{eq:v8-Kn}.
Half the log determinant proves \cref{eq:v8-Qn}. For $|\theta|<\pi$ this
stationary Hessian is positive: the smallest transverse eigenvalue is
$2\cos(\theta/n)-2\cos(\pi/n)>0$, and the longitudinal block is positive.
Only a smaller common real neighbourhood is needed for the global-minimum
argument in the preceding proof.

At $\theta=0$ all cubic terms vanish in these partial-product coordinates.
The scalar covariance is
\[
                C_{ij}=\min(i,j)-ij/n,\qquad 1\le i,j<n,
\]
with three independent colour components. Let $a=C_{ii}$, $b=C_{i+1,i+1}$,
$c=C_{i,i+1}$, putting endpoint covariances equal to zero. For one edge the
quartic action expectation is
\[
 -\frac58(a^2+b^2)-\frac94ab-\frac32c^2+\frac52c(a+b).
\]
Indeed its quartic action is
$-(|x|^4+|y|^4)/24-|x|^2|y|^2/4+
(x\cdot y)(|x|^2+|y|^2)/6$, and the three-component Gaussian moments are
$15a^2$, $15b^2$, $9ab+6c^2$, and $15c(a+b)$.
The Haar contribution is $\sum_{i=1}^{n-1}C_{ii}$. Substitution of
\[
 a=i(n-i)/n,\quad b=(i+1)(n-i-1)/n,\quad c=i(n-i-1)/n
\]
and summation over $0\le i<n$ gives
\[
 \sum_{i=0}^{n-1}\left[-\tfrac58(a^2+b^2)-\tfrac94ab-\tfrac32c^2
                     +\tfrac52c(a+b)\right]
       +\sum_{i=1}^{n-1}\frac{i(n-i)}n
                =-\frac{(n-1)(n-5)}{8n}.
\]
This is a polynomial sum in $i$, evaluable by the sums of powers through four.
It proves \cref{eq:v8-Bn}. At $n=4$ it independently reproduces
\cref{eq:v8-direct}, now in partial-product rather than independent-link
coordinates; the Haar factors are retained in both.
\end{proof}

The coefficient $B_n(0)=-n/8+3/4-5/(8n)$ grows linearly in magnitude. It is a
vacuum coefficient in this example, not a fourth-field or physical mass
instability. Nevertheless it disproves any rule that would infer bounded
unrescaled loop coefficients solely from correct finite composition.

There is also a positive, explicitly rescaled comparison. On each complex
disc $|\theta|<\rho<\pi$, with derivatives on smaller discs,
\begin{equation}\label{eq:v8-chainlimits}
 \begin{aligned}
 nG_n(\theta)&=\tfrac12\theta^2+O_\rho(n^{-2}),\\
 Q_n(\theta)-Q_n(0)&=\log(\sin\theta/\theta)+O_\rho(n^{-1}),\\
 B_n(0)/n&=-\tfrac18+\frac3{4n}-\frac5{8n^2}.
 \end{aligned}
\end{equation}
The first statement follows from the cosine series. For the second write its
left side as
\[
 \frac{n-1}{2}\log\cos(\theta/n)
 +\log(\sin\theta/\theta)
 -\log\left(\frac{\sin(\theta/n)}{\theta/n}\right).
\]
The first and last functions are respectively $O_\rho(n^{-1})$ and
$O_\rho(n^{-2})$ on the common disc. Cauchy estimates give the corresponding
bounds for every fixed derivative on a smaller disc. For dyadic $n=2^j$ these
coefficient-germ errors are summable. The natural chain fluctuation parameter
is $h_{\rm eff}=nh$; changing to it also changes the separately retained
Gaussian normalizations. No uniform statement for the full integral at
$nh$ fixed follows from these coefficient formulas. This one-dimensional
calibration does not establish the four-dimensional running coupling,
nonlinear Ward subtraction, or the compact interface estimates.

\section{The next scale estimate and the exact verification boundary}
\label{sec:v8-next}

The immediate induction target is now a specified packet, not three
undifferentiated words ``cubic, quartic, density''. For the two-loop
fourth-field coefficient one should propagate
\[
 (J^8S_j,J^6Q_j,J^4B_j)
\]
with the measured root-path block through order seven, the full spatial and
momentum norm, and its Ward-compatible marginal projection. V5 supplies the
linear stationary maps and inverses. This version supplies the invariant
finite derivative inventory, a primitive high-order evaluation procedure,
and the exact measured two-stage rule. A later application cannot reset
$Q_j$ or $B_j$, and cannot truncate each coefficient to field order four.

This does not prove boundedness of the packet under iteration. In particular,
\cref{prop:v8-norm}'s polynomial majorant is not a contraction, and the compact
chain shows why scale factors must actually be measured. Nor is the
$(8,6,4)$ packet a sufficient replacement for V7's stronger local analytic and
higher-derivative hypotheses on the \emph{whole action} when proving its exact
$O(\beta^{-2})$ remainder. Formal coefficient closure and uniform remainder
closure are separate. An analytic-tail estimate or a separately propagated
higher-order norm remains necessary for the latter. The compact complement,
other wells and transition-cutoff configurations remain outside the original
local-box theorem.

\paragraph{Executed verification.}
The new exact driver \path{verify_ym_two_loop_transport_V8.py} reports 27
passing exact test families. They include the sharp action/density/block examples, the generated intermediate density feed,
an independent direct normalizer series, nonlinear Jacobian cancellation,
the fully retained two-step $\mathrm{SU}(2)$ coefficient through its fourth
background derivative, independent quaternion square-root and nine-dimensional
Wick calculations, arbitrary-length Hessian determinants and the closed flat
coefficient sum. It also rebuilds the actual periodic and unperiodized
root-path primitive jets, checks the curl normalization, commuting inputs and
rational global conjugations, and exports the reconstructible directional
packet. These checks concern finite algebra and explicitly stated compact
calibrations. They do not evaluate the full four-dimensional two-stage packet,
certify a scale-invariant nonlinear ball, or prove a physical mass gap.

The V7 proof and its local assumptions are inherited, not independently
recertified by these algebra checks. Its 25 exact test families were freshly
rerun and passed; older reports in the retained archive are not described as
fresh executions. The new general identities rest
on their proofs, not on counting passing sample cases. No Lean verification
is asserted.

\paragraph{Append-only record.}
The complete V1--V7 mathematical body and all its end markers are retained
byte-for-byte. Only preamble version displays change. The bundle records the
preserved-body hash and the new source, executable and report hashes. Future
V9 work is to append below the next marker. Any correction to an earlier
proof must be stated explicitly rather than silently changing that proof.

\begin{thebibliography}{99}
\bibitem[CMR15]{CMRV8}
A.~S. Cattaneo, P. Mnev and N. Reshetikhin,
\emph{Perturbative quantum gauge theories on manifolds with boundary},
arXiv:1507.01221 (2015). \url{https://arxiv.org/abs/1507.01221}.
External perturbative-gluing context only; no cofinal estimate for the present
Yang--Mills block is imported from this reference.
\end{thebibliography}

% END OF VERSION 8 MATHEMATICAL BODY.
% Append subsequent requested versions below this marker; retain V1--V8 above.


\clearpage
\section{V9: critical-energy control before nonlinear constraint transport}
\label{sec:v9-context}

V8 closes the finite field--loop inventory, but its polynomial majorant does
not provide an invariant norm ball. Raising that inventory again would not
address this problem. This version tests a different upstream statement:
can the \emph{actual fine Wilson action}, before its nonlinear coarse
observable is imposed, be controlled in a norm independent of the number of
fine subdivisions? The answer is affirmative on a small global critical-energy
ball. It yields a uniform analytic classical minimizing branch for the
\emph{linear quotient block} underlying the supplied Gaussian tower, and a
uniform bound on the Wilson part of the quartic exchange.

The distinction in the preceding sentence is essential. The Wilson action and
the Gaussian block-average family used below are the supplied ones. The
constraint in the nonlinear minimization theorem is deliberately their
\emph{linearized quotient constraint}, not the nonlinear root-path constraint.
No Haar or inverse-pivot integral is evaluated by this classical theorem. A
last proposition isolates the additional uniform nonlinear-observable estimate
which would transfer the classical conclusion. A thin-line example shows why
control of all long logarithmic paths cannot simply be inferred from a small
critical-energy norm.

\paragraph{Nonduplication and external inputs.}
V4's all-level kernel comparison and the explicit harmonic minimizer are
inherited. V5's spatially local one-step inverses and V7--V8's measured quantum
coefficients are unchanged. The source audit found no corresponding
all-refinement critical-energy comparison in the inspected YM passages; this is
not a claim that small-background variational methods are new. Balaban's work
on lattice background minima \cite{BalabanV9} is relevant prior context. Its
averaging and regularity hypotheses are not asserted for the present root rule
by citation. The only standard analytic ingredient used below is the ordinary
four-dimensional Sobolev inequality, with its finite-periodic-lattice consequence
proved explicitly. No optimal Sobolev constant is needed; see
\cite{TalentiV9} for the classical Euclidean inequality.

\section{The full harmonic lift has the canonical amplitude scale}
\label{sec:v9-harmonic}

Let $L=2^j$, let $k\in[-\pi,\pi]^4\setminus\{0\}$, and choose centered aliases
$p_M=(k+2\pi M)/L\in[-\pi,\pi]^4$. Write
\[
 D_L(p)=\sum_{s=0}^{L-1}e^{\mathrm i sp},\qquad
 W_{M,\mu}=L^{-4}D_L(p_{M,\mu})\prod_{\nu=1}^4D_L(p_{M,\nu}).
\]
The diagonal matrix $W_M$ is exactly the composition of the source's dyadic
block symbols. For a transverse coarse polarization $c$, the actual quadratic
minimizer in the fine transverse representative is
\begin{equation}\label{eq:v9-lift}
 a_M=\mathcal M_{L,M}(k)c,
 \qquad \mathcal M_{L,M}=L^{-4}\frac{\Pi(p_M)}{\lambda(p_M)}
                                  W_M^*S^{(j)}(k).
\end{equation}
All matrices here use real momenta. This section does not claim that the
transverse representative, on arbitrary polarizations, has a regular complex
extension through its gauge singularity.

Put $m_0=(2/\pi)^{12}$, as in V4, and for $s>0$ define
\begin{equation}\label{eq:v9-As}
 A_s=\sum_{M\in\mathbb Z^4\setminus\{0\}}
       \frac{1}{|M|_2^s}\prod_{\nu=1}^4\frac1{1+|M_\nu|}.
\end{equation}

\begin{lemma}[A summable real alias envelope]
\label{lem:v9-alias}
The constants $A_1,A_2$ are finite. For a centered nonprincipal alias,
\[
 |D_L(p_{M,\nu})|/L\le\frac2{1+|M_\nu|},\quad
 |D_L(p_{M,\nu})|/L\le1,
 \qquad \lambda(p_M)\ge4|M|_2^2/L^2.
\]
\end{lemma}
\begin{proof}
For $M_\nu\ne0$, the sine formula and $|\sin t|\ge2|t|/\pi$ on
$|t|\le\pi/2$ give
\[
 \frac{|D_L(p_{M,\nu})|}{L}
 \le\frac{\pi}{|k_\nu+2\pi M_\nu|}
 \le\frac1{2|M_\nu|-1}\le\frac2{1+|M_\nu|}.
\]
For $M_\nu=0$ use the geometric-sum bound $|D_L|\le L$.
Similarly $\lambda(p_M)\ge4|k+2\pi M|^2/(\pi^2L^2)$; every nonzero
component of $M$ has $|k_\nu+2\pi M_\nu|\ge\pi|M_\nu|$.
To prove summability, group $M$ by $2^r\le|M|_\infty<2^{r+1}$. The product
weight summed over this group is at most
\[
 \left(1+2\sum_{m=1}^{2^{r+1}}\frac1{1+m}\right)^4
 \le [1+2\log(2^{r+1}+1)]^4.
\]
Multiplying by $2^{-sr}$ and summing over $r\ge0$ proves the assertion for
every $s>0$. Boundary ties only select one representative of an alias and do
not change the estimates.
\end{proof}

\begin{theorem}[All-level amplitude and one-difference bounds]
\label{thm:v9-amplitude}
For the full four-dimensional harmonic lift,
\begin{equation}\label{eq:v9-amplitude}
 \boxed{\begin{aligned}
 L\sum_M\|\mathcal M_{L,M}(k)\|&\le m_0^{-1}(1+64A_2),\\
 L^2\sum_M\sqrt{\lambda(p_M)}\|\mathcal M_{L,M}(k)\|
                         &\le m_0^{-1}(2\pi+128A_1).
 \end{aligned}}
\end{equation}
The bounds are independent of $j$ and all four components of real $k$.
Consequently the reconstructed fine field satisfies
\[
 \sup_x|a(x)|\le C|c|/L,\qquad
 \sup_{x,\nu}|a(x+e_\nu)-a(x)|\le C|c|/L^2.
\]
At $k=0$ the constant harmonic lift $a=c/L$ has the same bounds. This last
choice does not identify every direction-dependent transverse projector limit.
\end{theorem}
\begin{proof}
V4 gives $\|S^{(j)}(k)\|\le m_0^{-1}\lambda(k)$. At the principal alias,
$\|W_0\|\le L$ and $\lambda(k/L)\ge\lambda(k)/L^2$, so
$\|\mathcal M_{L,0}\|\le m_0^{-1}/L$. Also
$\sqrt{\lambda(k/L)}\le |k|/L\le2\pi/L$.
For a nonprincipal alias, drop the extra directional geometric-sum factor
using $|D_L|/L\le1$. The lemma yields
\[
 \|W_M\|\le16L\prod_\nu(1+|M_\nu|)^{-1}.
\]
Since $\|\Pi(p_M)\|=1$ and $\lambda(k)\le16$, substitution into
\cref{eq:v9-lift} gives
\[
 L\|\mathcal M_{L,M}\|\le
 64m_0^{-1}|M|_2^{-2}\prod_\nu(1+|M_\nu|)^{-1},
\]
and, with one square-root energy factor, the bound with coefficient $128$
and $|M|_2^{-1}$. Summation proves \cref{eq:v9-amplitude}.
A forward difference multiplies each Fourier term by $e^{\mathrm i p_\nu}-1$,
whose modulus is at most $\sqrt\lambda$. The triangle inequality gives the
pointwise consequences. Finite superpositions have the same estimates with
$|c|$ replaced by the sum of their coarse Fourier amplitudes.
\end{proof}

These are long-lift estimates in a smooth transverse representative. They do
not follow merely by multiplying a one-step bound $j$ times, which could
introduce a spurious exponential loss. Conversely, an absolute alias estimate
with two or more differences needs more cancellation than the envelope used
here; no such extension is asserted by dropping the energy denominator.

\section{The critical-energy norm and its uniformly coercive detail space}
\label{sec:v9-critical}

On the periodic lattice $(\mathbb Z/N\mathbb Z)^4$ use unnormalized counting
norms. A link field takes values in $\mathfrak{su}(2)\simeq\mathbb R^3$ with
the quaternion coefficient norm; in this convention the quadratic Wilson
action is $\frac12\|dA\|_2^2$. The forward difference is $\nabla_\mu$,
$dA$ is the plaquette curl, and $d^*$ is its usual scalar-to-link adjoint
divergence. The discrete Hodge identity is
\begin{equation}\label{eq:v9-Hodge}
 \sum_{\mu,\nu}\|\nabla_\mu A_\nu\|_2^2
                       =\|dA\|_2^2+\|d^*A\|_2^2.
\end{equation}
Define the linear Coulomb space by $d^*A=0$, allowing its constant modes, and
put
\begin{equation}\label{eq:v9-Enorm}
                       \|A\|_{\mathcal E}=\|dA\|_2+\|A\|_4.
\end{equation}
This is a global norm, not a local-max-norm chart or V6's anchored vertex norm.

\begin{lemma}[A lattice Sobolev constant independent of side length]
\label{lem:v9-Sobolev}
There is a fixed $C_S$ such that every mean-zero scalar lattice field satisfies
\[
                       \|u\|_4\le C_S\|\nabla u\|_2.
\]
The same statement holds componentwise. Therefore, on mean-zero Coulomb link
fields,
\begin{equation}\label{eq:v9-Sobolev}
                  \|v\|_{\mathcal E}\le(1+C_S)\|dv\|_2,
\end{equation}
with a changed fixed component constant, independent of $N$.
\end{lemma}
\begin{proof}
Periodically interpolate $u$ by the tensor product of linear hat functions on
unit cubes. On one cube, equivalence of norms on the fixed space of multilinear
polynomials gives a constant lower bound for its integral fourth norm by the
sum of fourth powers of its vertex values. This norm equivalence is valid also
for complex or vector coefficients. Summing cubes gives
$\|u\|_{\ell^4}\le C\|Iu\|_{L^4}$.
The gradient of the interpolant is a convex combination of edge differences
in each coordinate. Squaring, integrating and summing gives
$\|\nabla Iu\|_{L^2}\le C\|\nabla u\|_{\ell^2}$.
The integral mean of the interpolant equals the discrete mean.

The ordinary Sobolev--Poincare inequality on a unit four-torus is
$\|v\|_{L^4}\le C\|\nabla v\|_{L^2}$ for mean-zero $v$. It follows from
the Euclidean Sobolev inequality, a fixed finite partition of unity, and
Poincare's inequality for the zero mean; the Euclidean inequality is the
standard analytic input cited above. Scaling the torus from side one to side
$N$ multiplies both displayed norms by $N$, so its constant is unchanged.
Apply it to $Iu$ to prove the scalar lattice statement. Fourier expansion or
summation by parts proves \cref{eq:v9-Hodge}; for Coulomb fields it gives the
last assertion. The finite component count changes only $C_S$.
\end{proof}

Let the coarse torus have side $N$ and the fine torus side $NL$, with $L=2^j$.
Let $B_L$ be the supplied composed \emph{linear} link average, and let
$\mathcal Q_L=\Pi_c B_L$ on fine Coulomb fields. At nonzero coarse momenta
$\Pi_c$ is the transverse projector; at zero it is the identity, so constant
retained holonomies are not deleted. On coarse Coulomb fields set
\begin{equation}\label{eq:v9-Ynorm}
                  \|c\|_{\mathcal Y_N}=\|dc\|_2+N|\overline c|,
\end{equation}
where $\overline c$ is the four-component spatial mean. Let
$\mathcal H_{L,N}c$ be the fine harmonic lift of \cref{eq:v9-lift}, with
constant mean $\overline c/L$, and define
\[
 \mathcal V_{L,N}=\{v:d^*v=0,\ \mathcal Q_Lv=0\},
                \qquad \|v\|_V=\|dv\|_2.
\]

\begin{proposition}[Uniform critical splitting for the actual Gaussian quotient]
\label{prop:v9-splitting}
Every $v\in\mathcal V_{L,N}$ has zero mean, and $\|\cdot\|_V$ is a norm.
The identities
\begin{equation}\label{eq:v9-energy}
 \mathcal Q_L\mathcal H_{L,N}c=c,\qquad
 \langle d\mathcal H_{L,N}c,dv\rangle=0,
 \qquad
 \|d\mathcal H_{L,N}c\|_2^2=\langle c,S^{(j)}c\rangle
\end{equation}
hold in the inherited Fourier normalization. There is a constant $K_0$
independent of $L,N$ such that
\begin{equation}\label{eq:v9-splitting}
 \boxed{\|\mathcal H_{L,N}c\|_{\mathcal E}\le K_0\|c\|_{\mathcal Y_N},
 \qquad \|\mathcal Q_LA\|_{\mathcal Y_N}\le K_0\|A\|_{\mathcal E}.}
\end{equation}
The maps $(c,v)\mapsto\mathcal H_{L,N}c+v$ and its inverse are uniformly
bounded between $\mathcal Y_N\oplus\mathcal V_{L,N}$ and the fine Coulomb
space in these norms.
\end{proposition}
\begin{proof}
At coarse zero momentum every nonconstant alias has a vanishing geometric-sum
factor in $B_L$, while its constant coefficient is $L$. Thus
$\overline{\mathcal Q_LA}=L\overline A$, and a detail has zero mean. A
zero-curl Coulomb field on the torus is constant by \cref{eq:v9-Hodge}; this
proves positivity of the detail norm.
The quadratic variational identities in \cref{eq:v9-energy} are exactly the
source's harmonic minimization, including its $L^4$ fine-alias factor. V4 gives
\[
 \tfrac15\langle c,S^{(0)}c\rangle\le
 \langle c,S^{(j)}c\rangle\le m_0^{-1}\langle c,S^{(0)}c\rangle.
\]
Use the upper bound and the lattice Sobolev estimate on the mean-zero part of
the lift. The constant part has fourth norm at most a fixed component factor
times $(NL)|\overline c|/L=N|\overline c|$. This proves the first bound.
For $c=\mathcal Q_LA$, minimization gives
$\langle c,S^{(j)}c\rangle\le\|dA\|_2^2$. Its lower comparison controls
$\|dc\|_2$, and Holder gives $(NL)|\overline A|\le C\|A\|_4$. This proves
the second bound. Finally $v=A-\mathcal H_{L,N}\mathcal Q_LA$ and
\cref{eq:v9-Sobolev} prove the direct-sum assertion.
\end{proof}

The constants do not deteriorate with the number of cells. The \emph{class of
allowed fields} is nevertheless stronger than a volume-independent maximum
bound: for example, a fixed nonzero constant coarse field has
$\|c\|_{\mathcal Y_N}=N|c|$. Small global critical norm is not a substitute
for a thermodynamic small-field domain.

\section{The Wilson nonlinear force is uniformly small in this norm}
\label{sec:v9-Wilson}

On the full fine lattice use the actual primitive action
\[
 \mathcal S(A)=\sum_{x,\mu<\nu}
 \left(1-\tfrac12\Tr[e^{A_\mu(x)}e^{A_\nu(x+e_\mu)}
                      e^{-A_\mu(x+e_\nu)}e^{-A_\nu(x)}]\right),
\]
with $A_\mu\in\mathfrak{su}(2)$. Its complexification is the same trace
formula, not a nonholomorphic real-part operation. Put
$\mathcal N(A)=\mathcal S(A)-\frac12\|dA\|_{\rm bil}^2$, where the suffix
indicates the bilinear complex extension when needed.

\begin{lemma}[Primitive Wilson analytic estimates in critical energy]
\label{lem:v9-nonlinearity}
There exist $r_*,C>0$, independent of lattice size, such that for
$\|A\|_{\mathcal E}\le r_*$,
\begin{equation}\label{eq:v9-nonlinearity}
 \boxed{\begin{aligned}
 |\mathcal N(A)|&\le C\|A\|_{\mathcal E}^3,\\
 \|D\mathcal N(A)\|_{\mathcal E^*}&\le C\|A\|_{\mathcal E}^2,\\
 \|D^2\mathcal N(A)\|_{\mathcal E\to\mathcal E^*}
                                      &\le C\|A\|_{\mathcal E}.
 \end{aligned}}
\end{equation}
The action has a holomorphic extension to this complex norm ball with common
bounds. Its third derivative at zero obeys
\begin{equation}\label{eq:v9-cubic}
 |D^3\mathcal S(0)[A,B,C]|
 \le C\sum_{\rm cyc}\|dA\|_2\|B\|_4\|C\|_4.
\end{equation}
For $m\ge4$ the primitive homogeneous coefficient has a bound
$|\mathcal S_m(A)|\le C_0 8^m\|A\|_4^m/m!$, with a fixed incidence constant.
\end{lemma}
\begin{proof}
For a plaquette let $X_0,\ldots,X_3$ be its oriented Lie inputs. The quadratic
term is half the squared coefficient norm of $\sum X_i$. The cubic term is
\[
 -\frac14\Tr\left[(\sum_iX_i)\sum_{i<j}[X_i,X_j]\right].
\]
This is the same ordered-exponential identity used in V2, not a changed
Wilson vertex. On complex $\mathfrak{sl}_2$ the trace of the cube of one
traceless matrix vanishes, so the identity also gives its holomorphic germ.
The first factor is the actual plaquette curl. Holder with exponents $2,4,4$
and the fixed plaquette incidence proves \cref{eq:v9-cubic} by polarization.
This derivative factor is crucial: a raw $\ell^3$ estimate would have an
uncontrolled volume factor.

At degree $m\ge4$, each exponential word is a product of $m$ oriented fields
with coefficient $\prod_i(n_i!)^{-1}$, $\sum_i n_i=m$. Its trace is bounded by
the product of fixed matrix norms. Sum its locations using Holder with exponent
$m$; unnormalized counting norms satisfy $\|A\|_m\le\|A\|_4$. Shifts and
orientations have bounded incidence, independent of volume. The sum of word
coefficients is $4^m/m!$; the fixed conversion from complex quaternion
coefficient norm to matrix norm is absorbed into $8^m$, giving the stated
diagonal bound. These estimates
hold for complex directions. Polarization gives locally normally convergent
multilinear series, with at most a fixed exponential factor in $m$. Therefore
its first two derivative series converge on a common ball. Together with the
cubic estimate they give \cref{eq:v9-nonlinearity}. The quadratic term is
already controlled by $\|dA\|_2^2$, proving the remaining assertion.
\end{proof}

\section{A uniform nonlinear classical branch with the linear quotient constraint}
\label{sec:v9-minimizer}

For this section and only this comparison, impose $\mathcal Q_LA=c$ and the
linear Coulomb condition. Thus every allowed field is $A=a+v$ with
$a=\mathcal H_{L,N}c$ and $v\in\mathcal V_{L,N}$. The constraint is linear
in Lie coordinates. It is not $F_{\rm root}(A)=c$.

\begin{theorem}[Uniform analytic Wilson minimum in a global critical ball]
\label{thm:v9-minimizer}
There exist $\delta,C>0$ independent of $L,N$ such that for
$\|c\|_{\mathcal Y_N}<\delta$ there is a real-analytic stationary point
\begin{equation}\label{eq:v9-minimizer}
 A_{L,N}(c)=\mathcal H_{L,N}c+v_{L,N}(c),\qquad
                    \|v_{L,N}(c)\|_V\le C\|c\|_{\mathcal Y_N}^2.
\end{equation}
It is the unique stationary point in a common small detail-energy ball and
the strict minimum on that ball. Its fibre Hessian satisfies
\begin{equation}\label{eq:v9-Hessian}
 \tfrac12\|dw\|_2^2\le D^2\mathcal S(A_{L,N}(c))[w,w]
                               \le\tfrac32\|dw\|_2^2,
                  \qquad w\in\mathcal V_{L,N}.
\end{equation}
There is a holomorphic branch on a common complex retained ball. Consequently,
on a smaller common retained ball, for each fixed $m$, its fine-field
output derivatives (in $\mathcal E$) and those of
\[
                         \mathcal G_{L,N}(c)=\mathcal S(A_{L,N}(c))
\]
are bounded as multilinear maps in $\mathcal Y_N$ by constants independent of
$L,N$. In particular this holds through field degree eight, with no further
increase of the required radius as $L$ grows.
\end{theorem}
\begin{proof}
Let $\mathcal J:V\to V^*$ be the energy Riesz map,
$\langle v,\mathcal Jw\rangle=\langle dv,dw\rangle$. It and its inverse
have norm one in the stated energy norms. Let $i:V\to\mathcal E$ be the
inclusion, with $\|i\|\le1+C_S$. Orthogonality in \cref{eq:v9-energy}
reduces stationarity to the exact fixed-point equation
\begin{equation}\label{eq:v9-FP}
                 v=-\mathcal J^{-1}i^*D\mathcal N(a+iv).
\end{equation}
The bounds in \cref{lem:v9-nonlinearity} and
$\|a\|_{\mathcal E}\le K_0\|c\|_{\mathcal Y_N}$ show that its right side
has norm at most $C_1(K_0\|c\|+(1+C_S)\|v\|_V)^2$ and Lipschitz constant
at most $C_2(K_0\|c\|+(1+C_S)\|v\|_V)$. Choose a fixed detail-ball radius
and then $\delta$ small enough that the map sends that ball into itself with
Lipschitz constant at most $1/2$. Alternatively, on the smaller ball of radius
$C_3\|c\|^2$ the same inequalities give invariance, proving the sharper bound
in \cref{eq:v9-minimizer}. Both choices use only common constants.

The contraction theorem gives existence and uniqueness. The Hessian estimate
follows from $D^2\mathcal S=\mathcal J+i^*D^2\mathcal N i$ on the fibre by
making the last perturbation norm at most $1/2$. Every line segment in the
small detail ball stays in the same critical neighbourhood, so this Hessian
bound gives strict convexity there and proves the minimum assertion.
For complex fields, use the complexified bilinear energy Riesz map with its
induced operator norm. The same analytic majorants prove uniform contraction
on a smaller complex ball. Picard iteration converges uniformly there to a
holomorphic branch. Cauchy estimates on complex lines, followed by polarization,
give common multilinear derivative bounds at each fixed order. The action is
uniformly bounded by $C\|c\|^2$ on this branch, which gives the corresponding
classical coefficient bounds. No inverse bound for a fine $\ell^2$ Laplacian
has been used; the potentially soft inverse is measured in energy-dual norm.
\end{proof}

This is a genuine nonlinear statement for the complete fine Wilson action,
not a truncation of its cubic exponential. It is also a \emph{comparison
problem}: the affine linear quotient fibre is not the compact nonlinear
root-path fibre. Neither a loop determinant nor a microscopic-to-continuum
probability law is supplied. Uniqueness is only on the declared small ball,
not among all compact critical points or other conditional wells.

\section{The Wilson part of the quartic exchange is uniformly bounded}
\label{sec:v9-exchange}

Let $T=D^3\mathcal S(0)$ and $V_4=D^4\mathcal S(0)$. For
$a=\mathcal H_{L,N}c$ define the actual Wilson bilinear forcing on the detail
space by $g_a(w)=T[a,a,w]$.

\begin{theorem}[Energy-dual forcing and quartic exchange]
\label{thm:v9-exchange}
For all $L,N$,
\begin{equation}\label{eq:v9-exchange-bound}
 \boxed{\|g_a\|_{V^*}\le C\|c\|_{\mathcal Y_N}^2,
 \qquad 0\le\langle g_a,\mathcal J^{-1}g_a\rangle
                                      \le C^2\|c\|_{\mathcal Y_N}^4.}
\end{equation}
The quartic coefficient of the nonlinear minimum in the preceding theorem is
\begin{equation}\label{eq:v9-quartic}
 \boxed{D^4\mathcal G_{L,N}(0)[c^4]
       =V_4[a^4]-3\langle g_a,\mathcal J^{-1}g_a\rangle.}
\end{equation}
Both terms have uniform fourth-degree bounds. The inner product in the middle
is nonnegative for real data, but the complete quartic coefficient is not
assigned a sign by this fact.
\end{theorem}
\begin{proof}
By \cref{eq:v9-cubic},
\[
 |T[a,a,w]|\le C\left(\|da\|_2\|a\|_4\|w\|_4
                              +\|a\|_4^2\|dw\|_2\right).
\]
The critical splitting and Sobolev estimate give the first bound. The Riesz
identity $\langle g,\mathcal J^{-1}g\rangle=\|g\|_{V^*}^2$ gives the second.
Differentiate stationarity at $c=tc_0$: $A'(0)=a$ and
$v''(0)=-\mathcal J^{-1}g_a$. Expanding $\mathcal S(ta+t^2v''(0)/2+\cdots)$
through $t^4$, the terms with $A'''(0)$ vanish by harmonic orthogonality.
The remaining terms are $V_4[a^4]+6g_a(v''(0))+3\|dv''(0)\|_2^2$, giving
\cref{eq:v9-quartic}. The primitive fourth-degree bound controls the direct
term. This reproduces the general exchange identity in the \emph{linear}
constraint case while adding its all-refinement bound for the Wilson forcing.
\end{proof}

In V4's measured root-constraint notation, the forcing also contains
$2F_2[a,\cdot]^*\eta_1$, its constraint component $F_2[a,a]$, and the direct
term $4\langle\eta_1,F_3[a^3]\rangle$. Those are \emph{not} set to zero in
the root program by \cref{thm:v9-exchange}; they vanish here because this
section's constraint is linear. The estimate controls the Wilson component of
that source, not the entire root-constraint border or its quantum response.

\section{An exact all-level noncommuting background of the comparison problem}
\label{sec:v9-constant}

Take one coarse periodic cell, a fine torus of side $L\ge2$, and the retained
constant components $c_1=xE_1$, $c_2=yE_2$, $c_3=c_4=0$. Its harmonic lift is
the constant fine field $a_1=xE_1/L$, $a_2=yE_2/L$. These two components do
not commute when $xy\ne0$.

\begin{proposition}[Exact constant-connection minimum and irrelevant corrections]
\label{prop:v9-constant}
On the small branch of \cref{thm:v9-minimizer}, this constant lift is the
\emph{exact} nonlinear minimizer under the linear quotient constraint. Its
action is
\begin{equation}\label{eq:v9-constant}
 \boxed{\mathcal G_{L,1}(x,y)=2L^4\sin^2(x/L)\sin^2(y/L).}
\end{equation}
Hence its fourth field term is $2x^2y^2$, and through joint field degree eight,
\begin{align}\label{eq:v9-constantjets}
 \mathcal G_{L,1}(x,y)={}&2x^2y^2
 -\frac{2}{3L^2}(x^4y^2+x^2y^4)\notag\\
 &+\frac1{L^4}\left[\frac4{45}(x^6y^2+x^2y^6)
                               +\frac29x^4y^4\right]+O_{\rho}(L^{-6})
\end{align}
on each fixed smaller complex $(x,y)$ polydisc. The remainder includes only
field terms of degree at least ten. It and every fixed retained derivative
are $O_\rho(L^{-6})$ on still smaller polydiscs. Thus the displayed irrelevant
corrections decay exactly as $L^{-2}$ and $L^{-4}$ on this noncommuting family.
\end{proposition}
\begin{proof}
At a constant fine connection, translation invariance makes every component
of the action gradient independent of position. Every detail has zero mean,
so this gradient annihilates the whole detail space, not only constant trial
variations. The uniqueness theorem therefore identifies the constant field
with its local nonlinear minimum. In particular its cubic forcing on details
vanishes, consistently with \cref{eq:v9-quartic}.

The quaternion commutator of $e^{xE_1/L}$ and $e^{yE_2/L}$ has scalar part
$1-2\sin^2(x/L)\sin^2(y/L)$. Only the $1,2$ plaquettes contribute, and there
are $L^4$ of them. This proves \cref{eq:v9-constant}. Multiply the ordinary
sine-square series to obtain \cref{eq:v9-constantjets}. The analytic tails
are uniformly bounded by the next powers of $1/L$, and Cauchy estimates give
the derivative statement. The result is a noncommuting classical calibration,
not a determination of the constrained root block's two-loop coefficient.
\end{proof}

\section{Why the coordinate and logarithmic-chart issues cannot be skipped}
\label{sec:v9-obstructions}

\begin{proposition}[The ordered tree does not preserve the fine amplitude scale]
\label{prop:v9-tree}
For a single constant commuting fine component $a_2=cE/L$ on a side-$L$
macrocell, the linear ordered-tree representative of V2 has
\[
 b_2=cE\,1_{\{s_2=L-1\}},\qquad b_1=b_3=b_4=0.
\]
Its pointwise amplitude is $|c|$, rather than $|c|/L$, and
\begin{equation}\label{eq:v9-tree}
                   \|a\|_4=|c|,\qquad \|b\|_4=L^{3/4}|c|.
\end{equation}
Both configurations have identity plaquette holonomy when implemented as the
corresponding exactly gauge-related link fields. Thus this growing coordinate
norm is not an energy or mass instability.
\end{proposition}
\begin{proof}
At $q=0$, V2's harmonic profile is $1/L$ and its tree formula gives the displayed
boundary field. There are $L^4$ nonzero links in the original component and
$L^3$ in the boundary component. Count their fourth powers to obtain
\cref{eq:v9-tree}. The exact commuting tree transformation concentrates the
same wrapping holonomy on the boundary; every contractible plaquette remains
identity. No noncommuting finite vertex is moved between these representatives
without its higher coordinate terms.
\end{proof}

More generally, if an exact gauge change of the curve $e^{tA}$ has logarithmic
link coordinates
$B(t)=tb_1+t^2b_2/2+t^3b_3/6+\cdots$, gauge invariance and $D\mathcal S(0)=0$
give the exact fourth-order identity
\begin{equation}\label{eq:v9-gaugejets}
 D^4\mathcal S(0)[A^4]
 =D^4\mathcal S(0)[b_1^4]+6D^3\mathcal S(0)[b_1,b_1,b_2]
        +3D^2\mathcal S(0)[b_2,b_2]+4D^2\mathcal S(0)[b_1,b_3].
\end{equation}
It is the ordinary fourth chain rule. The cubic formula similarly contains
$3D^2\mathcal S(0)[b_1,b_2]$. A bound on only the first term in
\cref{eq:v9-gaugejets} can therefore be a coordinate artefact; the measured
transport must retain all terms and its density Jacobian.

\begin{proposition}[Small critical energy does not control every long logarithmic path]
\label{prop:v9-line}
On a fine torus of side $L\ge3$, put
\[
 A_1(x)=\frac\pi L E_3\,1_{\{x_2=x_3=x_4=0\}},\qquad A_2=A_3=A_4=0.
\]
This is a Coulomb field and
\begin{equation}\label{eq:v9-line}
 \boxed{\|dA\|_2^2=\frac{6\pi^2}{L},\qquad
 \|A\|_4=\pi L^{-3/4},\qquad
 \mathcal S(A)=6L(1-\cos(\pi/L))\longrightarrow0.}
\end{equation}
Nevertheless the group product on the distinguished wrapping line is $-I$.
Thus no fixed identity-logarithm radius for \emph{all} paths of length $L$
follows from $\|A\|_{\mathcal E}\to0$.
\end{proposition}
\begin{proof}
The nonzero component is independent of $x_1$, so its divergence is zero.
For each of the three transverse orientations exactly $2L$ plaquettes have
nonzero curl, of magnitude $\pi/L$, and their holonomies are
$e^{\pm\pi E_3/L}$. Count these plaquettes and the $L$ nonzero links to obtain
\cref{eq:v9-line}. Along the wrapping line the factors commute, and their
product is $e^{\pi E_3}=-I$. Its distance from the identity does not shrink.
This argument is a path-control obstruction, not a claim that the exact
\emph{nested} root average equals this single path or necessarily exits its
chart. That average may suppress the example; such suppression must be proved
from its actual rule.
\end{proof}

The two examples separate different issues. Tree concentration is a
representative-dependent norm loss even on a flat configuration. The thin
line is a genuine nontrivial holonomy with small four-dimensional energy.
Neither can be resolved by declaring a critical Sobolev estimate to be a
uniform logarithmic-chart theorem.

\section{A precise upstream transfer target for the classical root problem}
\label{sec:v9-transfer}

The classical Wilson action now has a common nonlinear bound in the smooth
critical coordinates. The unknown quantity can therefore be moved to the
coarse observable, without pretending it is already linear.

\begin{criterion}[A uniform nonlinear-observable bound suffices for the classical part]
\label{crit:v9-transfer}
Suppose the \emph{actual composed nonlinear coarse observable}, after an
explicit compatible source and gauge identification, is a holomorphic map
$\mathcal F_L$ from a common fine critical ball to $\mathcal Y_N$, with
\[
 \mathcal F_L(0)=0,\quad D\mathcal F_L(0)=\mathcal Q_L,\quad
 \mathcal F_L(A)=\mathcal Q_LA+\mathcal Z_L(A),
 \qquad
 \sup_{L,N}\sup_{\|A\|_{\mathcal E}<r}
                           \|\mathcal Z_L(A)\|_{\mathcal Y_N}\le M r^2.
\]
Here $r>0$ and $M<\infty$ are fixed, the maps are real on real fields, and the
source/gauge identification is part of the hypothesis. Then, on a smaller
common retained ball, the true constraint $\mathcal F_L(A)=c$ admits a
uniform analytic local classical Wilson minimum and uniformly bounded
classical field jets through degree eight (indeed through each fixed degree).
The constants depend on $M,r$, not on $L,N$.
\end{criterion}
\begin{proof}
Uniform complex Cauchy estimates give uniform derivatives of $\mathcal Z_L$
on a smaller ball. Its value and first derivative vanish at zero, so these
bounds imply $\|\mathcal Z_L(A)\|\le C\|A\|^2$ and
$\|D\mathcal Z_L(A)\|\le C\|A\|$ there. In the bounded splitting of
\cref{prop:v9-splitting}, solve
\[
                  s=c-\mathcal Z_L(\mathcal H_{L,N}s+v).
\]
The same uniform contraction argument gives a chart
$A(c,v)=\mathcal H_{L,N}s(c,v)+v$ of the nonlinear constraint, with
$A(c,v)=\mathcal H_{L,N}c+v+O((\|c\|+\|v\|_V)^2)$ and uniform analytic
bounds. Its quadratic Wilson action is
$\frac12\langle c,S^{(j)}c\rangle+\frac12\|v\|_V^2$ by harmonic
orthogonality. The higher chart terms and \cref{lem:v9-nonlinearity} give a
uniform cubic analytic remainder in $(c,v)$; the quadratic coarse term has a
uniform bilinear bound in $\mathcal Y_N$ by V4. Apply the proof of
\cref{thm:v9-minimizer} to the chart action, now with fibre variable $v$.
It gives the uniform analytic minimum and its derivative estimates. All
nonlinear constraint terms participate in that chart action; none is set to
zero. This proves the conditional implication, not its stated coarse-map
hypothesis for the root rule.
\end{proof}

This target is narrower than assuming bounds for an arbitrarily growing
family of transported classical vertices: the primitive Wilson nonlinear
estimates are now supplied, and the proposed remaining bound is on one
\emph{exact} nonlinear coarse observable in a specified norm. It may still
require a more regular norm or controlled path averages, as
\cref{prop:v9-line} makes clear. The one-step primitive analytic radius in V8
alone does not bound an arbitrarily long composition. Conversely, a failure
to prove this sufficient critical-norm criterion would not disprove a
construction in another suitable local norm.

Nothing in this criterion bounds a quantum log determinant, the transported
$Q_j$ or $B_j$, or an exact remainder on all compact fields. Those requirements
retain their measured conventions from V7--V8. Nor does a small global energy
ball replace their volume-uniform local-field domains or the independent
interface globalization theorem.

\section{Verification and the next acquisition}
\label{sec:v9-verification}

The accompanying checker reports fourteen passing exact test families and
sixteen separately labelled floating-point diagnostics. It separates exact
algebra from numerical diagnostics.
Its exact tests reconstruct the full discrete Hodge identity, the primitive
Wilson cubic against ordered quaternion multiplication, the constant
noncommuting action through degree eight, the vanishing of its full
mean-zero directional response, the thin-line plaquette incidence and norms,
the gauge-curve chain rule, and the quartic exchange formula by an independent
stationary-series solve. A separate floating-point driver reconstructs the
full harmonic alias map at several dyadic levels and genuine four-dimensional
momenta, checking the quotient constraint, harmonic energy, and amplitude
scales. Those samples are diagnostics; \cref{thm:v9-amplitude} is proved by
the summable alias envelope, not by their extrema. The analytic contraction
and lattice Sobolev constants are existence constants, not claimed numerical
enclosures. All twenty-seven V8 exact test families were freshly rerun and passed; that
report and its executable dependencies are retained in the new bundle. Older
reports remain archival unless a fresh execution is explicitly recorded.

\paragraph{What has advanced.}
The full harmonic map has canonical fine amplitude and one-difference scaling.
The Wilson cubic force and its quartic exchange have all-refinement energy-dual
bounds. The complete nonlinear fine action, not a low-order truncation, has a
common classical analytic minimizing branch on the declared linear quotient
fibre. These statements control the classical Wilson self-interaction in a
specific norm and identify an exact remaining nonlinear-observable transfer.
The noncommuting constant family supplies a genuine all-level nonlinear
calibration rather than another isolated loop card.

\paragraph{What should happen next.}
Test the actual nested root-path observable in the smooth critical coordinates,
including its second and third derivatives and its compatible gauge/source
identification. The thin-line configuration is a useful falsification test
for any proposed path-norm shortcut. A successful estimate must control the
observable itself or pass to a stronger local regularity norm; it cannot infer
that control from the Gaussian inverse alone. The $Q_j,B_j$ density packet,
its Ward-compatible subtractions, analytic tails, and compact complement stay
in the program. No four-dimensional quantum measure or physical Hamiltonian
mass gap is asserted by this classical comparison.

\paragraph{Append-only record.}
The entire V1--V8 mathematical body and its end markers are preserved
byte-for-byte. Only preamble version displays are advanced. Source, report,
code and preserved-body hashes are recorded. A later V10 is to append after
the V9 marker, with any correction stated explicitly rather than silently
rewriting a preceding proof.

\begin{thebibliography}{99}
\bibitem[T76]{TalentiV9}
G.~Talenti, \emph{Best constant in Sobolev inequality}, Annali di Matematica
Pura ed Applicata \textbf{110} (1976), 353--372.
DOI: \nolinkurl{10.1007/BF02418013}.
\bibitem[B85]{BalabanV9}
T.~Balaban, \emph{The variational problem and background fields in
renormalization group method for lattice gauge theories}, Communications in
Mathematical Physics \textbf{102} (1985), 277--309.
DOI: \nolinkurl{10.1007/BF01229381}.
\end{thebibliography}

% END OF VERSION 9 MATHEMATICAL BODY.
% Append subsequent requested versions below this marker; retain V1--V9 above.


\clearpage
\section{V10: the exact root observable in a controlled coarse gauge}
\label{sec:v10-context}

V9 obtains a uniform classical Wilson minimum with a linear constraint. Its
last criterion does not yet establish a comparable bound for the actual root
observable. We address that observable before adding another conditional
nonlinear-RG majorant. Two concrete results are obtained. On arbitrary
four-dimensional lattice fields with a scaled small link amplitude, an exact
coarse gauge correction gives a common holomorphic domain for every nested
root step, with explicit amplitude and one-difference estimates. On homogeneous
noncommuting connections the same construction gives a convergent nonlinear
source germ, including all its finite derivatives and an exact cubic term.
The original root rule, distinguished path and ordered trees are retained.

These results have different quantifiers from V9's desired critical-energy
criterion. The general domain here is $L\|A\|_\infty<r$, not a fixed ball in
$\|dA\|_2+\|A\|_4$. The homogeneous cofinal limit does not assert convergence
for arbitrary changing inhomogeneous inputs. Neither result supplies a
conditional integral, a fluctuation measure, or a new global repair floor.
They do populate an actual nonlinear source map rather than replacing it by
its linearization.

\paragraph{Source audit and nonduplication.}
The primitive root-path rule and ordered tree are inherited from f12 and V8;
V8's directional quaternion arithmetic is retained as an elementary utility.
V9's thin-line example and source-transfer criterion are the specific audit
and acquisition targets. The source inspection distinguishes these results
from the earlier harmonic lift, isolated Wilson response cards, and local
Laplace remainder. No independent verification of the complete old certificate
chain is asserted. Holomorphic linearization near an expanding fixed point is
standard background \cite{AbateV10}; its elementary special-case proof is
included, rather than importing a theorem as a missing bound for this block.
The new input to that argument is the exact root map, its common gauge
identification, and its nonzero nonlinear coefficients.

\section{Exact ordered-tree blocking and its compensating coarse gauge}
\label{sec:v10-block}

On an even periodic lattice write a fine site as $x=2r+s$, with
$s\in\{0,1\}^4$. A link $U_{x,\mu}$ is oriented from $x$ to $x+e_\mu$.
Let $t_{r,s}(U)$ be the chronological product on the ordered path from $2r$
to $2r+s$, moving first in direction one, then two, three and four. This is
the tree whose parent clears the highest active bit; $t_{r,0}=I$.
Tree gauge transforms a fine link into
\[
 \widetilde U_{x,\mu}=t_{r,s}U_{x,\mu}t_{r',s'}^{-1},
 \qquad x+e_\mu=2r'+s'.
\]
The sixteen chronological two-link products on a coarse edge consequently are
\begin{equation}\label{eq:v10-path}
 P_{r,s,\mu}
 =t_{r,s}U_{2r+s,\mu}U_{2r+s+e_\mu,\mu}
                         t_{r+e_\mu,s}^{-1}.
\end{equation}
In particular $P_{r,0,\mu}$ is the actual root path. On the common identity
logarithm branch set
\begin{equation}\label{eq:v10-K}
 K_{r,\mu}=P_{r,0,\mu}\exp\left\{
   \frac1{16}\sum_s\log(P_{r,0,\mu}^{-1}P_{r,s,\mu})\right\}.
\end{equation}
This is precisely the source's rule applied after its tree gauge. It is not an
untransported average of matrices whose endpoints lie at different sites.

Define an actual coarse-site group and a representative of the same output by
\begin{equation}\label{eq:v10-balance}
 \phi_r(A)=\frac1{16}\sum_s\log t_{r,s}(e^A),\quad
 h_r(A)=e^{-\phi_r(A)},\quad
 \widehat K_{r,\mu}=h_rK_{r,\mu}h_{r+e_\mu}^{-1},\quad
 \mathfrak T(A)=\log\widehat K.
\end{equation}
The correction is a coarse gauge transformation, not a modification of any
of the sixteen weights or of the distinguished path. Its dependence on the
fine field is kept visible.

\begin{proposition}[Exact orbit equivalence, including after iteration]
\label{prop:v10-orbit}
For real compact fields in the admitted branches, $K(U)$ and $\widehat K(U)$
have the same coarse gauge orbit. If
$U^g_{x,\mu}=g_xU_{x,\mu}g_{x+e_\mu}^{-1}$, then
\begin{equation}\label{eq:v10-equivariance}
 K(U^g)_{r,\mu}=g_{2r}K(U)_{r,\mu}g_{2r+2e_\mu}^{-1}.
\end{equation}
Repeated use of \cref{eq:v10-balance} at intermediate levels therefore leaves
the orbit of the literal nested root output unchanged. For any input measure
on the admitted set, the two pushforwards to the coarse orbit space are equal.
\end{proposition}
\begin{proof}
Chronological products telescope under a fine gauge change, so
$t^g_{r,s}=g_{2r}t_{r,s}g_{2r+s}^{-1}$. Thus each path in
\cref{eq:v10-path} has left endpoint factor $g_{2r}$ and right endpoint factor
$g_{2r+2e_\mu}^{-1}$. The relative matrix $P_0^{-1}P_s$ is conjugated at its
right endpoint. Its identity logarithm is conjugated by the same matrix.
Exponentiation and multiplication by $P_0$ give \cref{eq:v10-equivariance}.

At stage $s$, suppose the balanced field differs from the literal field by
coarse-site matrices $G^{(s)}_x$. At the next stage its unbalanced block differs
by $G^{(s)}_{2r}$ at the new roots, by \cref{eq:v10-equivariance}. The new
balanced correction changes that gauge matrix to
$G^{(s+1)}_r=h^{(s+1)}_rG^{(s)}_{2r}$. This induction proves the iterated
statement. Every gauge-invariant measurable observable has identical values
on the two outputs, pointwise in the input. Integration proves equality of
the quotient pushforwards. On real fields unitary conjugation preserves the
admitted relative-logarithm branch. Complex statements below concern its
holomorphic germ and this same algebraic identity.
\end{proof}

This does not license dropping a Jacobian in an arbitrary gauge-fixed density.
In a spatially varying fibre, $h_r(A)$ generally depends on eliminated data, not
only on the coarse source. Equality of orbit observables is exact without a
change of integration variables, but it is not an assertion that the two
representative-valued probability densities are equal. An actual source chart
or coarea computation still has to retain its measured convention.

\section{A volume-independent analytic stencil with the correct linear part}
\label{sec:v10-local}

Use the maximum matrix operator norm of all fine Lie inputs. Complexification
uses the same matrix words and holomorphic inverses, not conjugate transposes
of complex parameters. All local supports below have a fixed finite number
of links, independently of surrounding volume.

\begin{lemma}[A common primitive domain and a quadratic nonlinear remainder]
\label{lem:v10-local}
The map $\mathfrak T$ is holomorphic on the local polydisc
$\|A\|_\infty<r_0=1/4096$, with $\|\mathfrak T(A)\|_\infty<1/64$.
Its derivative at zero is the supplied linear sixteen-path block
\begin{equation}\label{eq:v10-B}
 (BA)_{r,\mu}=\frac1{16}\sum_s
             (A_{2r+s,\mu}+A_{2r+s+e_\mu,\mu}).
\end{equation}
On a smaller common ball there is a constant $C<\infty$, independent of
volume, such that
\begin{equation}\label{eq:v10-localbound}
 \begin{split}
 \mathfrak T(A)&=BA+E(A),\\
 \|BA\|_\infty&\le2\|A\|_\infty,\\
 \|E(A)\|_\infty&\le C\|A\|_\infty^2,\\
 \|DE(A)\|_{\infty\to\infty}&\le C\|A\|_\infty.
 \end{split}
\end{equation}
\end{lemma}
\begin{proof}
A tree has at most four links, a transported path in \cref{eq:v10-path} at
most ten, and $P_0^{-1}P_s$ at most twelve. For matrices of norm at most $r$,
$\|\prod_{i=1}^p e^{\pm A_i}-I\|\le e^{pr}-1$. Hence their logarithms have
convergent identity series whenever the displayed distances are below one.
Put $m_p(r)=-\log(2-e^{pr})$. A tree logarithm has norm at most $m_4(r)$,
and the inner mean in \cref{eq:v10-K} at most $m_{12}(r)$.
The full balanced product has distance from identity at most
\[
 \exp(2m_4(r)+2r+m_{12}(r))-1.
\]
At $r=1/4096$ this is less than one and its logarithm has norm less than
$1/64$. An entirely rational check uses $e^{pr}\le(1-pr)^{-1}$ for
$p=2,4,12$, followed by $-\log(1-x)\le x/(1-x)$: it gives the bound
$1469588221/270598387462<1/64$.

Let $\phi^{(1)}_{r,s}$ be the sum of the Lie inputs along the ordered tree.
The first-order term of \cref{eq:v10-path} is
$\phi^{(1)}_{r,s}+A_{2r+s,\mu}+A_{2r+s+e_\mu,\mu}
-\phi^{(1)}_{r+e_\mu,s}$. At first order \cref{eq:v10-K} averages these
sixteen terms. The gauge factors in \cref{eq:v10-balance} subtract exactly
the mean at $r$ and add the mean at $r+e_\mu$. The result is \cref{eq:v10-B}.
Its row coefficients are nonnegative and sum to two. Analytic Cauchy estimates
on the fixed local polydisc, and $E(0)=DE(0)=0$, give
\cref{eq:v10-localbound} on a smaller ball. Taking the maximum over translated
stencils introduces no volume factor. These bounds allow inputs on neighbouring
cells to be different; they are not proved by a constant-field restriction.
\end{proof}

\section{Every nested root step stays in a common logarithmic domain}
\label{sec:v10-scale}

Let the fine torus have side $NL$, where $L=2^n$. Successively apply the same
balanced map $\mathfrak T$ at the $n$ dyadic stages, and write
\[
 A^{(0)}=A,\quad A^{(s+1)}=\mathfrak T(A^{(s)}),\quad
 \widehat{\mathcal F}_{L,N}(A)=A^{(n)},\quad B_L=B^{\circ n}.
\]
The domain and codomain lattice sizes change at each stage. The mean tree
correction is applied on that stage's actual lattice, not a tree of length $L$
in one step.

\begin{theorem}[Uniform nonlinear observable bound in scaled amplitude]
\label{thm:v10-amplitude}
There are $r,C_1>0$ independent of $L,N$ such that for
$\rho=L\|A\|_\infty<r$, every balanced root step is defined and
\begin{equation}\label{eq:v10-amplitude}
 \boxed{\begin{aligned}
 \|A^{(s)}\|_\infty&\le 2\rho\,2^{s-n},\\
 \|\widehat{\mathcal F}_{L,N}(A)\|_\infty&\le2\rho,\\
 \|\widehat{\mathcal F}_{L,N}(A)-B_LA\|_\infty&\le C_1\rho^2.
 \end{aligned}}
\end{equation}
The map is holomorphic on this common scaled ball. Each fixed field derivative
has a bound independent of $L,N$ as a multilinear map from the norm
$\|A\|_{X_L^0}=L\|A\|_\infty$ to coarse $\ell^\infty$, on a smaller ball.
On real fields its output represents the actual nested root gauge orbit.
\end{theorem}
\begin{proof}
Let $a_s=\|A^{(s)}\|_\infty$ and $x_s=2^{n-s}a_s$. By
\cref{eq:v10-localbound}, while in the primitive domain,
\[
 a_{s+1}\le2a_s+Ca_s^2,\qquad
 x_{s+1}\le x_s+\frac C2\,2^{s-n}x_s^2.
\]
Choose $2r$ inside that domain and $2Cr\le1/2$. Bootstrap $x_s\le2\rho$:
its accumulated excess over $x_0=\rho$ is at most
$2C\rho^2\sum_{s<n}2^{s-n}\le2C\rho^2\le\rho/2$.
This stricter bound closes the bootstrap and ensures that no intermediate
step leaves the primitive analytic ball.

Let $e_s=\|A^{(s)}-B^{\circ s}A\|_\infty$. Then
$e_{s+1}\le2e_s+Ca_s^2$, $e_0=0$, and therefore
\[
 e_n\le C\sum_{s<n}2^{n-1-s}(2\rho\,2^{s-n})^2
       \le2C\rho^2.
\]
This proves \cref{eq:v10-amplitude}. Finite composition is holomorphic, and
these estimates hold on the same complex scaled ball. Cauchy estimates on
smaller complex direction polydiscs followed by polarization give the fixed-order
multilinear bounds with dimension-independent constants at each fixed order.
Only the output maximum and direction norms enter that argument. Exact orbit
identification is \cref{prop:v10-orbit}.
\end{proof}

For a field on a periodic lattice put
$\|\nabla A\|_\infty=\max_\nu\|\tau_{e_\nu}A-A\|_\infty$.

\begin{theorem}[The same exact map preserves one scaled spatial difference]
\label{thm:v10-difference}
In the preceding domain, set $d_0=\|\nabla A\|_\infty$. With a common $C_2$,
\begin{equation}\label{eq:v10-difference}
 \boxed{\begin{aligned}
 \|\nabla\widehat{\mathcal F}_{L,N}(A)\|_\infty
       &\le e^{C_2\rho}L^2d_0,\\
 \|\nabla(\widehat{\mathcal F}_{L,N}(A)-B_LA)\|_\infty
       &\le C_2\rho e^{C_2\rho}L^2d_0.
 \end{aligned}}
\end{equation}
In particular the actual source representative is uniformly bounded in coarse
$C^1$ for inputs with $L\|A\|_\infty+L^2\|\nabla A\|_\infty<r$.
Its nonlinear error is quadratic in this stronger norm.
\end{theorem}
\begin{proof}
Translation of a coarse stencil by one cell translates all its fine inputs by
two sites. Hence $\nabla_c\mathfrak T(A)$ is the difference of the same local
map on $\tau_{2e_\nu}A$ and $A$. The linear part is bounded by
$2\|\tau_{2e_\nu}A-A\|_\infty\le4\|\nabla A\|_\infty$.
The derivative bound for $E$ gives
\[
 d_{s+1}\le(4+2Ca_s)d_s.
\]
After dividing by $4^{s+1}$, multiply the factors
$1+Ca_s/2\le e^{Ca_s/2}$. Their exponent is at most $C\rho$ because
$\sum_{s<n}a_s\le2\rho$. This proves the first inequality.
For the derivative of the cumulative nonlinear error, its contribution created
at stage $s$ is at most $2Ca_sd_s$, and the later linear blocks amplify one
difference by at most $4^{n-1-s}$. Thus the sum is bounded by
\[
 \sum_{s<n}4^{n-1-s}2Ca_s e^{C\rho}4^sd_0
 \le C\rho e^{C\rho}L^2d_0.
\]
Enlarge $C_2$ to cover both bounds. No commutative assumption has entered this
argument.
\end{proof}

\begin{scope}[This is a stronger norm, not a completed critical-energy transfer]
The bounds hold for arbitrary spatially varying four-dimensional inputs, but
require $L\|A\|_\infty$ to be small. A fixed small critical-energy ball need
not satisfy this condition. Nor does a coarse maximum or first-difference
bound furnish the anchored high-field tensor norms used for the quantum packet.
The compensation provides a representative of the same coarse orbit, not a
proved uniform nonlinear Coulomb slice in $\mathcal Y_N$. Therefore
\cref{crit:v9-transfer} has not been established on its stated domain.
\end{scope}

\section{The thin-line diagnostic is suppressed by the actual nested average}
\label{sec:v10-line}

\begin{proposition}[Exact abelian cancellation and line attenuation]
\label{prop:v10-abelian}
For fields in one fixed abelian Lie subalgebra, whenever the displayed short
relative logarithms use their continuous identity values,
\begin{equation}\label{eq:v10-abelian}
                    \mathfrak T(A)=BA
\end{equation}
exactly, not merely through a finite field order. On a side-$L=2^n$ fine torus,
for V9's field
\[
 A_1(x)=\frac{\theta}{L}E_3\,1_{\{x_2=x_3=x_4=0\}},\qquad A_2=A_3=A_4=0,
\]
the final one-cell balanced root output is
\begin{equation}\label{eq:v10-line}
                       \boxed{c_1=\theta E_3/L^3,\qquad c_2=c_3=c_4=0.}
\end{equation}
For each fixed real $\theta$, all the nested short logarithms are admitted for
sufficiently large dyadic $L$. The literal output has the same gauge orbit.
\end{proposition}
\begin{proof}
All products and logarithms add in the fixed abelian subalgebra. The raw root
logarithm is the average of the transported path sums, which is
$BA+\overline\phi_r-\overline\phi_{r+e_\mu}$. The factors in
\cref{eq:v10-balance} cancel these last terms exactly, proving
\cref{eq:v10-abelian}.

After $s$ dyadic averages, the field is still independent of its longitudinal
coordinate and supported on one transverse line of the lattice of side
$L/2^s$. Its link amplitude is $\theta/(L4^s)$: the two longitudinal links
supply a factor two and the three transverse averages a factor $1/8$.
Inducting through $n$ stages proves \cref{eq:v10-line}. Its maximum amplitude
never exceeds $|\theta|/L$. Choose this below the fixed primitive radius;
each remaining short word then stays in the identity branch. Although the
wrapping product at $\theta=\pi$ is $-I$, no single unaveraged path of length
$L$ is logarithmically evaluated in the nested construction. This is exactly
the distinction left open by V9's diagnostic, rather than a rejection of that
diagnostic.
\end{proof}

This resolves that specific proposed falsification test positively. It does
not show that every small-critical-energy configuration is suppressed by the
same abelian mechanism. Noncommuting path dispersion remains in the source.

\section{The homogeneous root map and its first genuine nonlinear term}
\label{sec:v10-homogeneous}

For a translation-invariant fine connection let $A=(A_1,\ldots,A_4)$,
$U_\mu=e^{A_\mu}$, and
\[
 t_s=e^{s_1A_1}e^{s_2A_2}e^{s_3A_3}e^{s_4A_4}.
\]
The exact path formula simplifies to
\begin{equation}\label{eq:v10-hompaths}
 P_{s,\mu}=t_se^{2A_\mu}t_s^{-1},\qquad P_{0,\mu}=e^{2A_\mu}.
\end{equation}
Write $M(A)=\log K(A)$ for the raw homogeneous map and
$T(A)=\mathfrak T(A)$ for the balanced homogeneous map. Put
$S_A=\sum_\nu A_\nu$ and $C_A=\sum_{\nu<\rho}[A_\nu,A_\rho]$.
This $C_A$ is local notation for a Lie commutator sum, not the quadratic Casimir.

\begin{theorem}[The actual homogeneous cubic, with the tree rotation removed]
\label{thm:v10-cubic}
Near zero the two exact maps have expansions
\begin{align}\label{eq:v10-rawcubic}
 M_\mu(A)={}&2A_\mu+[S_A,A_\mu]
 +\frac14[S_A,[S_A,A_\mu]]\notag\\
 &+\frac14\sum_\nu[A_\nu,[A_\nu,A_\mu]]
 +\frac14[C_A,A_\mu]+O(\|A\|^4).
\end{align}
\begin{equation}
 \boxed{T_\mu(A)=2A_\mu+
             \frac14\sum_\nu[A_\nu,[A_\nu,A_\mu]]+O(\|A\|^4).}
 \label{eq:v10-cubic}
\end{equation}
In particular $T$ has no quadratic term. Its cubic is not in general a common
gauge rotation or zero.
\end{theorem}
\begin{proof}
Let $X_s=\log P_{s,\mu}$ and $R=2A_\mu$. Then $X_s=R+O(\|A\|^2)$, exactly
by \cref{eq:v10-hompaths}. The map
$\Psi_R(D)=\log(e^{-R}e^{R+D})$ has an invertible linear derivative in $D$
near $R=0$. Its terms nonlinear in $D$ vanish at $R=0$, since $\Psi_0(D)=D$.
Consequently averaging $\Psi_R(D_s)$ and applying its inverse changes the
average of the $D_s$ by $O(\|R\|\max_s\|D_s\|^2)$.
Here this is $O(\|A\|^5)$. It follows that $M_\mu$ agrees with
$16^{-1}\sum_sX_s$ through degree four. This does not declare them identical
at higher degree.

Expansion of $X_s=2\operatorname{Ad}(t_s)A_\mu$ through degree three gives
\[
 2A_\mu+2\sum_\nu s_\nu[A_\nu,A_\mu]
 +\sum_\nu s_\nu[A_\nu,[A_\nu,A_\mu]]
 +2\sum_{\nu<\rho}s_\nu s_\rho[A_\nu,[A_\rho,A_\mu]].
\]
Average $s_\nu$ and $s_\nu s_\rho$, whose means are $1/2$ and $1/4$.
The cubic becomes
$\frac14[S_A,[S_A,A_\mu]]+
\frac14\sum_\nu[A_\nu,[A_\nu,A_\mu]]+\frac14[C_A,A_\mu]$
by the Jacobi identity. This proves \cref{eq:v10-rawcubic}.

The actual tree-log mean has expansion
\[
 \phi(A)=\frac1{16}\sum_s\log t_s
       =\frac12S_A+\frac18C_A+O(\|A\|^3).
\]
For homogeneous fields the two endpoint gauge factors in
\cref{eq:v10-balance} coincide, so $T_\mu=e^{-\operatorname{ad}\phi}M_\mu$.
Its quadratic subtracts $[S_A,A_\mu]$. At cubic degree the remaining terms
from this conjugation are
$-\frac14[S_A,[S_A,A_\mu]]-\frac14[C_A,A_\mu]$.
They cancel precisely the first and last cubic terms in
\cref{eq:v10-rawcubic}, giving \cref{eq:v10-cubic}. For perpendicular
$A_1=xE_1$, $A_2=yE_2$, the cubic components are $-xy^2E_1$ and
$-x^2yE_2$, which change the component lengths. Thus they are not merely a
simultaneous infinitesimal rotation.
\end{proof}

For clarity, the general root-log mean has the familiar cubic term
\[
 \log\!\left(e^{X_0}\exp\bigl[\overline{\log(e^{-X_0}e^{X_s})}\bigr]\right)
 =\overline X+\frac1{12}\overline{[X_s-\overline X,[X_0,X_s-\overline X]]}
       +O(X^4).
\]
In \cref{eq:v10-hompaths} the dispersion starts in degree two, which postpones
this particular correction. The proof above justifies that postponement
without substituting a generic degree-four remainder for a fifth-order claim.
The exact checker finds a nonzero fifth-order difference from the log average;
it keeps the literal reference in every higher-order computation.

\section{A cofinal nonlinear source germ for the actual noncommuting root orbit}
\label{sec:v10-limit}

For $L=2^n$ define homogeneous maps on a fixed twelve-dimensional Lie carrier by
\begin{equation}\label{eq:v10-Phi}
 \Phi_n(c)=T^{\circ n}(c/2^n),\qquad
 \Xi_n(c)=M^{\circ n}(c/2^n).
\end{equation}
The raw and balanced outputs are globally conjugate for each $n$, by
\cref{prop:v10-orbit}. A homogeneous gauge correction is one simultaneous
conjugation of all four retained groups.

\begin{theorem}[Common analytic limit with summable dyadic error]
\label{thm:v10-limit}
There is $r_h>0$ such that the $\Phi_n$ are holomorphic on $\|c\|<r_h$ and
converge uniformly on every smaller ball to a holomorphic local diffeomorphism
$\Phi$ with $\Phi(0)=0$, $D\Phi(0)=I$. More precisely,
\begin{equation}\label{eq:v10-limitbound}
 \boxed{\|\Phi_n(c)-\Phi(c)\|\le C\|c\|^3 4^{-n},\qquad
       T(\Phi(c))=\Phi(2c)}
\end{equation}
where the second identity is read where both sides are in the common domain.
Every fixed derivative, and the inverse maps on a smaller common image,
converges at rate $O(4^{-n})$. The raw maps $\Xi_n$ also converge on a smaller
ball, at rate $O(2^{-n})$; their real limits and the balanced limit are in the
same global conjugacy orbit.
\end{theorem}
\begin{proof}
By \cref{thm:v10-cubic}, write $T(x)=2x+R(x)$ with
$\|R(x)\|\le C\|x\|^3$ and $\|DR(x)\|\le C\|x\|^2$ on a fixed
complex ball. An induction as in \cref{thm:v10-amplitude} gives
\[
 \|T^{\circ s}(c/2^n)\|\le2\|c\|2^{s-n},\qquad s\le n,
\]
for a sufficiently small $r_h$. For trajectories in that domain the derivative
product through $n$ steps is bounded by
\[
 \prod_{s<n}(2+C\|T^{\circ s}(x)\|^2)
       \le 2^n\exp\!\left(C'\sum_{s<n}\|c\|^2 4^{s-n}\right)
       \le C''2^n.
\]
The same bound holds on the segment connecting the two initial arguments
below, by slightly enlarging the trajectory constant.
Now compare
\[
 \Phi_{n+1}(c)=T^{\circ n}\bigl(T(c/2^{n+1})\bigr),\qquad
 \Phi_n(c)=T^{\circ n}(c/2^n).
\]
Their initial arguments differ by $R(c/2^{n+1})$, of size at most
$C\|c\|^3 2^{-3n-3}$. The derivative product therefore gives
$\|\Phi_{n+1}-\Phi_n\|\le C\|c\|^3 4^{-n}$.
Summing the geometric tail proves uniform convergence and
\cref{eq:v10-limitbound}. Uniform limits of holomorphic maps are holomorphic;
Cauchy estimates on smaller balls prove convergence of every fixed derivative.

The exact relation $T(\Phi_n(c))=\Phi_{n+1}(2c)$ passes to the limit.
Each map differs from the identity by $O(\|c\|^3)$ with common derivative
bounds. On a smaller ball the equation $x=y-(\Phi_n(x)-x)$ is a uniform
contraction. This gives inverse maps on a common ball and their uniform
analytic bounds; comparison of the fixed-point equations and Cauchy estimates
gives the same convergence rate for their derivatives.

For $M(x)=2x+O(\|x\|^2)$ use the same proof with a quadratic remainder.
The derivative product is still at most $C2^n$, while the initial discrepancy
is $O(\|c\|^2 4^{-n})$, yielding a summable $O(2^{-n})$ difference. On real
fields the conjugating groups at every finite $n$ lie in compact $SU(2)$.
A subsequence of them converges; taking limits proves orbit equivalence of
$\Xi$ and $\Phi$. No global linearization away from the identity branch is
claimed.
\end{proof}

\begin{corollary}[The limiting and finite cubic source tensors]
\label{cor:v10-sourcecubic}
For every $n\ge0$,
\begin{equation}\label{eq:v10-sourcecubic}
 \boxed{\Phi_{n,\mu}(c)
 =c_\mu+\frac{1-4^{-n}}{24}
       \sum_\nu[c_\nu,[c_\nu,c_\mu]]+O(\|c\|^4).}
\end{equation}
The fourth-order remainder is uniform on a common smaller complex ball. The
limiting formula has coefficient $1/24$. All third mixed derivatives are
obtained by polarization of this cubic. In particular
\[
 \Phi_{\mu=1}(xE_1,yE_2,0,0)
       =xE_1-\frac16xy^2E_1+O(\|(x,y)\|^4),
\]
with the analogous second component $yE_2-x^2yE_2/6$ for the limit.
\end{corollary}
\begin{proof}
Let $T_3$ be the cubic in \cref{eq:v10-cubic}. A term created at stage $s$
of $T^{\circ n}(c/2^n)$ has degree-three coefficient
$2^{n-1-s}T_3(2^{s-n}c)$. Summing $s=0,\ldots,n-1$ gives
$(1-4^{-n})T_3(c)/6$. This is \cref{eq:v10-sourcecubic}.
Uniform analyticity bounds the rest uniformly. Alternatively the homogeneous
conjugacy equation gives $(2^3-2)\Phi_3=T_3$. The $SU(2)$ bracket
$[a\cdot E,b\cdot E]=2(a\times b)\cdot E$ gives the final components.
\end{proof}

\subsection{Exact higher coefficients without a fitted scale limit}
For a rational direction $c(t)=t(E_1,E_2,0,0)$ the limit's components through
fifth degree are
\begin{align}\label{eq:v10-jetdata}
 \Phi_{\mu=1}(c(t))={}&tE_1-\tfrac16t^3E_1
                 +\tfrac5{84}t^4E_3+\tfrac3{40}t^5E_1+O(t^6),\\
 \Phi_{\mu=2}(c(t))={}&tE_2-\tfrac16t^3E_2
                 +\tfrac1{84}t^4E_3+\tfrac{13}{120}t^5E_2+O(t^6).
 \notag
\end{align}
The other two components vanish. The report includes all coefficients through
order seven. They are computed recursively from the literal quaternion map:
if $\Phi(t)=\sum_{k\ge1}p_kt^k$, $p_1=c(1)$, then
\begin{equation}\label{eq:v10-recursion}
 p_k=\frac{[t^k]T(\sum_{a<k}p_at^a)}{2^k-2},\qquad k\ge2.
\end{equation}
The denominators are nonzero. The analytic theorem identifies this unique
formal solution with the actual limiting germ. The independent check
$T(\Phi(t))=\Phi(2t)\pmod{t^8}$ uses the full reconstructed path words.
Oddness is not imposed: the ordered tree can leave even terms at higher order.
The cubic's permutation-symmetric expression does not imply full permutation
symmetry of every representative coefficient.

\section{Source reparametrization has a measured cost even after a pure gauge is removed}
\label{sec:v10-density}

There are two distinct transformations here. A common gauge rotation changes
the representative but not its invariant component lengths. The limiting
source map $\Phi$ changes those lengths at cubic order. Their Jacobians must
not be identified.

\begin{proposition}[Homogeneous gauge balance is volume preserving]
\label{prop:v10-gaugevolume}
On a sufficiently small homogeneous source chart, the map from raw output
$B=M(A)$ to balanced output $T(A)$ preserves the twelve-dimensional Lebesgue
volume and the product exponential-coordinate Haar density. It does not follow
that the full spatially varying, detail-dependent balancing operation has the
same representative-density property.
\end{proposition}
\begin{proof}
Since $DM(0)=2I$, the homogeneous map has an analytic local inverse. Put
$P(B)=\phi(M^{-1}(B))$. The maps $M$, $M^{-1}$ and $\phi$ are equivariant
under simultaneous global conjugation. Thus $P$ is equivariant and the source
change is $B_\mu\mapsto\operatorname{Ad}(e^{-P(B)})B_\mu$.

For any equivariant Lie-valued $P$, the vector field
$X_\mu(B)=-[P(B),B_\mu]$ keeps $P(B)$ fixed along its trajectories: the
infinitesimal equivariance identity gives $DP[X]=-[P,P]=0$.
Its time-one flow is exactly the preceding source change. If $Y_a$ are the
infinitesimal simultaneous adjoint rotations in an orthonormal Lie basis,
$X=-\sum_aP^aY_a$. Each $Y_a$ has zero divergence, and equivariance gives
$Y_aP^b=f_{ac}^{\ b}P^c$. Therefore
$\operatorname{div}X=-\sum_{a,c}f_{ac}^{\ a}P^c=0$, using the zero trace of
adjoint matrices. The flow preserves Lebesgue volume. It preserves the
invariant norm of every $B_\mu$, so the product Haar density
$\prod_\mu(\sin|B_\mu|/|B_\mu|)^2$ is unchanged as well.
This proof uses the homogeneous inverse $M^{-1}$; it is not an inversion of a
fine-to-coarse map with eliminated coordinates.
\end{proof}

\begin{proposition}[Nontrivial Jacobian of the actual limiting source germ]
\label{prop:v10-sourcejac}
In the twelve real homogeneous Lie coordinates, the limiting map satisfies
\begin{equation}\label{eq:v10-sourcejac}
 \boxed{\log\det D\Phi(c)=-\sum_{\mu=1}^4|c_\mu|^2+O(\|c\|^3).}
\end{equation}
For the finite map $\Phi_n$, its quadratic coefficient is multiplied by
$1-4^{-n}$. This is a source-coordinate Jacobian, not a quantum determinant
or a mass counterterm.
\end{proposition}
\begin{proof}
There is no quadratic term in $\Phi$. Hence the degree-two logarithmic
Jacobian is the divergence of its cubic. In $\mathfrak{su}(2)$,
$\operatorname{Tr}_{\R^3}\operatorname{ad}_a^2=-8|a|^2$.
The summand with $\nu=\mu$ in
$\sum_\nu[c_\nu,[c_\nu,c_\mu]]$ is identically zero. For each of the
other three values of $\mu$ the derivative in that component has trace
$-8|c_\nu|^2$. Division by 24 gives
$-\sum_\nu|c_\nu|^2$. The finite coefficient follows from
\cref{eq:v10-sourcecubic}. Determinants are positive near the identity
Jacobian, so the real logarithm has this germ. The higher-order remainder is
uniform on a smaller common ball by the analytic theorem.
\end{proof}

For example, under $v=\Phi(c)$ a density relative to $dc$ transforms with
$|\det D\Phi(c)|^{-1}$. Its negative logarithm acquires
$+\log|\det D\Phi(c)|$ evaluated at $c=\Phi^{-1}(v)$.
The sign and this coordinate conversion must be retained when such a
reparametrization is used for a measured homogeneous source calculation.
It is not supplied by the unit-Jacobian gauge rotation alone.

\subsection{A source-coordinate sextic term which does not decay away}
Let $a_n(v)=\Phi_n^{-1}(v)$ and put the constant fine connection
$A_\mu=a_{n,\mu}(v)/L$ on a side-$L=2^n$ torus. It has actual balanced
one-cell root output $v$ exactly. Evaluate its Wilson action, without claiming
stationarity against arbitrary nonhomogeneous constrained variations:
\begin{equation}\label{eq:v10-testaction}
 E_n^{\rm hom}(v)=L^4\sum_{\mu<\nu}
 \left(1-\tfrac12\Tr[e^{a_{n,\mu}/L}e^{a_{n,\nu}/L}
                     e^{-a_{n,\mu}/L}e^{-a_{n,\nu}/L}]\right).
\end{equation}

\begin{proposition}[Uniform homogeneous action limit in the actual source coordinate]
\label{prop:v10-testaction}
The functions in \cref{eq:v10-testaction} have a common analytic domain and
converge, with every fixed derivative on smaller domains, at rate $O(L^{-2})$
to
\[
 E_\infty^{\rm hom}(v)
      =\frac12\sum_{\mu<\nu}|[a_\mu(v),a_\nu(v)]|_{\rm bil}^2,
                 \qquad a(v)=\Phi^{-1}(v).
\]
For $v=(xE_1,yE_2,0,0)$ the homogeneous field terms through degree six are
\begin{equation}\label{eq:v10-sextic}
 \boxed{E_n^{\rm hom}(v)
 =2x^2y^2+\frac23(1-2L^{-2})(x^4y^2+x^2y^4)+O(\|(x,y)\|^7).}
\end{equation}
The last field-degree remainder is uniform in $n$. The sextic coefficient
therefore need not vanish merely because the primitive constant-connection
lattice correction in V9 is $O(L^{-2})$.
\end{proposition}
\begin{proof}
For two constant $SU(2)$ Lie vectors $a,b$, the scalar commutator formula is
\[
 1-\tfrac12\Tr[e^{a/L}e^{b/L}e^{-a/L}e^{-b/L}]
 =2\sin^2(|a|/L)\sin^2(|b|/L)
            \frac{|a|^2|b|^2-(a\cdot b)^2}{|a|^2|b|^2}.
\]
The sine quotients have entire even series, so the apparent divisions are
removable on the complex germ. Multiplying by $L^4$ gives
$|[a,b]|_{\rm bil}^2/2+O(L^{-2})$ uniformly on bounded complex inputs.
Use the inverse-map convergence of \cref{thm:v10-limit} to obtain the claimed
limit and derivatives.

For the displayed two-component source, inversion of the cubic gives
$a_1=xE_1+(1-L^{-2})xy^2E_1/6+O(v^4)$ and
$a_2=yE_2+(1-L^{-2})x^2yE_2/6+O(v^4)$. Insert this into V9's constant
Wilson action. The quartic term produces the sextic correction
$\frac23(1-L^{-2})(x^4y^2+x^2y^4)$; the primitive sine correction contributes
$-\frac{2}{3L^2}(x^4y^2+x^2y^4)$. Their sum is \cref{eq:v10-sextic}.
The common analytic domain bounds all higher terms uniformly.
\end{proof}

This is an admissible homogeneous evaluation, not a proof that the constant
field is the minimizer of the actual nonlinear root fibre. Unlike V9's linear
constraint, the nonlinear constraint differential can couple to inhomogeneous
details at a constant background. The surviving sextic term here is explicitly
a nonlinear source-coordinate effect. It is not evidence of a new physical
relevant interaction or a changed Yang--Mills beta function.

\section{What the next upstream estimate must add}
\label{sec:v10-next}

The literal root observable now has a controlled all-level representative on
$L\|A\|_\infty<r$, with a one-difference bound on its natural stronger domain.
The thin-line test is explicitly attenuated, and homogeneous noncommuting data
have a cofinal nonlinear source germ with summable derivative errors. The
second and third source derivatives are no longer unspecified on that family;
the source is provably not linear there. These are actual root-map results,
not another replacement block or bare Wilson-loop card.

The principal transfer gap is now sharply localized. To use V9's classical
critical-energy minimum, one must either show that the relevant minimizers and
constraint corrections also satisfy
\[
 L\|A\|_\infty+L^2\|\nabla A\|_\infty\le C\|c\|,
\]
or establish a different nonlinear root estimate on a larger energy domain.
V9 bounds its correction in energy by $O(\|c\|^2)$; that does not give
$\|v\|_\infty=O(\|c\|^2/L)$. Promoting it to this sharper local estimate
without a scale-aware elliptic argument would reintroduce the missing premise.
A bound for a full critical ball is not inferred from the successful commuting
thin-line calculation either.

A subsequent fluctuation integral has an additional obstacle: its typical
fine details are not deterministic links of size $r/L$ at all ultraviolet
scales. The actual measured $(8,6,4)$ coefficient packet, the high-order analytic
tail, Ward-compatible subtraction, and compact complement remain separate.
The coarse gauge identification proves equality of quotient observables, but
a chosen gauge-fixed source density must still be computed in its own chart.
The homogeneous source Jacobian above illustrates a genuine conversion term.

Accordingly, the next substantive calculation should combine the newly
controlled root stencil with a scale-aware local elliptic bound for the
\emph{actual constrained background}, not merely iterate a larger list of
conditional moment identities. Neither a four-dimensional continuum law nor
a physical Hamiltonian gap is asserted here. No claim of new coverage of the
transition-cutoff or high-active-source interface configurations is made.

\section{Executed exact checks and append-only record}
\label{sec:v10-verification}

The new driver passes 21 exact test families. All 14 V9 exact test families
were freshly rerun and passed. The new modules reconstruct both the spatial
root map and its homogeneous restriction directly from ordered quaternion words. The homogeneous limit is
computed by \cref{eq:v10-recursion}, not by floating-point extrapolation.
The verification report distinguishes these finite identities from the analytic
all-level theorems, which rest on the proofs above.

Checks include direct multiplication of tree-gauged links versus the collapsed
path formula; the full spatial linearization on a side-four fine torus; exact
gauge covariance and the two-stage orbit identity; abelian cancellation with
spatially varying input; the cubic formula on two noncommuting direction
families; the literal fifth-order root/log-average discrepancy; the limit
functional equation through order seven; five finite dyadic cubic factors;
source invariants under balancing; the infinitesimal homogeneous gauge Jacobian;
the source Jacobian and sextic source-coordinate terms; and the analytic-radius
rational majorant. Coefficients in the report are exact rational numbers.
No floating-point sample is used to establish a uniform domain or convergence.

The entire V1--V9 mathematical body, its citations and end markers are retained
byte-for-byte. Only the preamble's version display is advanced. The bundle
contains the new appendix, source, executable checks, reports, assembly script
and input/output hashes. Fresh reruns of old checks are identified separately;
other old reports are archival, not described as current executions. Future V11
work is to append after the following marker. Any correction must be stated
explicitly, not silently substituted into an earlier proof.

\begin{thebibliography}{99}
\bibitem[A03]{AbateV10}
M.~Abate, \emph{Discrete local holomorphic dynamics}, arXiv:math/0310089 (2003).
\url{https://arxiv.org/abs/math/0310089}.
External context for local linearization only. The common-domain and dyadic-rate
proof for the present root map is given in \cref{thm:v10-limit}; no estimate for
a Yang--Mills regulator is imported from this survey.
\end{thebibliography}

% END OF VERSION 10 MATHEMATICAL BODY.
% Append subsequent requested versions below this marker; retain V1--V10 above.


\clearpage
\section{V11: the first actual root-constrained relaxation, rather than another trial action}
\label{sec:v11-context}

V10 controls the literal nested root observable and evaluates a homogeneous
connection with a prescribed source. It explicitly does not identify that
connection with the minimum against inhomogeneous constrained variations.
The next scale-aware estimate should therefore begin by finding the first
variation which the nonlinear root constraint actually permits. This version
computes it at every dyadic refinement, on a one-cell coarse torus while
retaining \emph{all} inhomogeneous fine Coulomb details.

The result is not another linear-constraint comparison. The first nonzero
detail displacement of the actual root-constrained classical minimum is of
field degree four. Its profile is an explicitly solved periodic Poisson
problem and has the fine scales $L^{-1}$ in amplitude and $L^{-2}$ in one
spatial difference. The induced decrease of the minimized action is of degree
eight, with an exact coefficient tending to a nonzero limit. Consequently,
V10's homogeneous evaluation gives the correct classical germ through degree
seven, but misses a definite eighth-order term for every noncommuting
homogeneous source. This resolves a real distinction which matters to V8's
$J^8S$ packet.

\paragraph{Scope and inherited inputs.}
The balanced root map, its exact orbit equivalence, its homogeneous maps
$T,\Phi_n$, and their cofinal analytic limit are inherited from V10. The Wilson
action, Coulomb detail space, and constant-connection formula are inherited
from V9. They are not reproved as new acquisitions. The new input is the
\emph{mixed} root Hessian with one homogeneous and one arbitrary detail leg,
not the already known Hessian on two homogeneous legs. Its exact first spatial
moment gives a complete solvable leading Euler source.

At each fixed $L$, ordinary local implicit-function arguments construct the
constrained stationary branch. Its higher-order remainders below may depend
on $L$. The results uniform in $L$ are the explicitly displayed coefficients,
their scaled profiles, and their convergence through degree eight. No common
finite-amplitude domain for the full constrained minimum is inferred from a
bound on its first nonzero displacement coefficient. Classical background-field
minimization is established methodology \cite{BalabanV9}; no result for a
different averaging rule is substituted for the present calculation.

\section{The actual finite constraint and its stationary branch}
\label{sec:v11-branch}

Let $L=2^n\ge2$ and let the fine torus be $(\mathbb Z/L\mathbb Z)^4$, with one
coarse periodic cell after $n$ blocking steps. Coordinates are
$x_\nu\in\{0,\ldots,L-1\}$. Use quaternion coefficient norms on
$\mathfrak{su}(2)$, so
\[
 [a\cdot E,b\cdot E]=2(a\times b)\cdot E,
 \qquad \mathcal S_2(A)=\tfrac12\|dA\|_2^2.
\]
All fine counting norms are unnormalized. Write
$\langle w\rangle_L=L^{-4}\sum_xw(x)$ and work on the logarithmic Coulomb
slice $d^*A=0$. Its detail space is
\[
 V_L=\{w:d^*w=0,\ \langle w\rangle_L=0\},
 \qquad \langle w,\mathcal J_Lw\rangle=\|dw\|_2^2.
\]
Here $\mathcal J_L$ is the positive detail Riesz operator. On $V_L$, the Hodge
identity identifies it with the componentwise periodic negative Laplacian.
Its smallest eigenvalue in the unscaled counting norm is allowed to tend to
zero; no uniform counting-norm inverse bound is asserted.

Let $\mathcal F_L$ be V10's exact balanced nested root logarithm, restricted to
this slice. Then
\begin{equation}\label{eq:v11-zero}
 \mathcal F_L(0)=0,\qquad D\mathcal F_L(0)w=L\langle w\rangle_L,
 \qquad \mathcal F_L(a/L)=\Phi_n(a).
\end{equation}
The four Lie components of $a$ and of the coarse source $v$ form a real
12-dimensional space. Put
\begin{equation}\label{eq:v11-E4}
 E_4(a)=\frac12\sum_{\mu<\nu}|[a_\mu,a_\nu]|^2,
 \qquad
 g_\mu(a)=\partial_{a_\mu}E_4(a)
       =-\sum_\nu[a_\nu,[a_\nu,a_\mu]].
\end{equation}
The subscript four is field degree; it is not a loop index.

\begin{proposition}[Local classical minimum for the literal root orbit]
\label{prop:v11-existence}
At every fixed $L$, there is a real-analytic constrained stationary branch
$A_L^*(v)$ near zero for the Wilson action and $\mathcal F_L(A)=v$ in the
Coulomb slice. It is a strict local minimum on that fibre. Define
\[
 \mathcal G_L(v)=\mathcal S(A_L^*(v)),\qquad
 A_L^{\rm hom}(v)=\Phi_n^{-1}(v)/L,\qquad
 E_L^{\rm hom}(v)=\mathcal S(A_L^{\rm hom}(v)).
\]
These are finite-regulator germs. The scalar minimum is a local root-orbit
quantity, while the representative-valued density still needs its own
Jacobian and coarea treatment.
\end{proposition}
\begin{proof}
Every Coulomb field decomposes uniquely as $A=s/L+w$ with $w\in V_L$.
The derivative in $s$ of $\mathcal F_L(s/L+w)$ is the identity at zero.
The implicit function theorem solves its constraint as
$s=s(v,w)$ near zero. The resulting chart action has vanishing first
variation and detail Hessian $\mathcal J_L\succ0$ at zero. A second implicit
solve gives its stationary point and positivity of the fibre Hessian on a
smaller neighbourhood. The neighbourhood may depend on $L$.

For the orbit interpretation, on every fixed finite torus a logarithmic
Coulomb condition is a local slice modulo constant gauges: its linearization
in a mean-zero infinitesimal gauge is the scalar negative Laplacian, which
is invertible on that mean-zero space. This is another finite implicit solve.
V10's exact orbit identity preserves the coarse root orbit under that change.
At the one-cell coarse level the remaining coarse gauge is simultaneous
conjugation of the four source groups. A constant fine gauge preserves the
Coulomb slice and implements that same conjugation. Hence a nearby fine orbit
whose coarse orbit is that of $v$ can be represented with the exact source
$v$ on this slice. Wilson gauge invariance identifies its scalar energy.
This is not an equality of arbitrary gauge-fixed source densities, nor a
uniform global gauge-fixing theorem.
\end{proof}

\section{The mixed root Hessian is an exact first spatial moment}
\label{sec:v11-mixed}

The following computation uses the entire spatial stencil, not just its
homogeneous restriction. For an even side $M$, let $\mathfrak T$ be the
one-step balanced logarithmic root map. A homogeneous input $a$ in the next
lemma is not divided by $M$.

\begin{lemma}[One-step averaged mixed derivative]
\label{lem:v11-one}
For every fine Lie field $w$ on the side-$M$ torus,
\begin{equation}\label{eq:v11-one}
 \left\langle D^2\mathfrak T(0)[a,w]_\mu\right\rangle_{M/2}
 =\sum_{\nu\ne\mu}
   [a_\nu,\langle (2(x_\nu\bmod2)-1)w_\mu(x)\rangle_M].
\end{equation}
No Coulomb or zero-mean assumption on $w$ is needed for this identity.
\end{lemma}
\begin{proof}
For one cell let $\theta_{r,s}$ be the first-order ordered-tree logarithm,
$\phi_r=16^{-1}\sum_s\theta_{r,s}$, and
$\psi_{r,s}=\theta_{r,s}-\phi_r$. For an edge of orientation $\mu$ let
\[
 b_s=A_\mu(2r+s)+A_\mu(2r+s+e_\mu),\qquad
 b_s^{[2]}=\tfrac12[A_\mu(2r+s),A_\mu(2r+s+e_\mu)].
\]
The logarithm of the distinguished-root average agrees with the average of
its path logarithms through degree two. Expanding the three path factors
and the two coarse balancing factors by the ordered second-degree BCH
identity gives the quadratic homogeneous coefficient
\begin{equation}\label{eq:v11-quadratic-stencil}
 E^{[2]}_{r,\mu}(A)
 =\overline{b_s^{[2]}}
  +\tfrac12\overline{[\psi_{r,s}+\psi_{r+e_\mu,s},b_s]}
  -\tfrac12\overline{[\psi_{r,s},\psi_{r+e_\mu,s}]}.
\end{equation}
The bars in this formula average only $s$. Second-order tree logarithms
cancel with their means in the balancing factors; none has been dropped.

For the constant leg $a$,
$\psi_{r,s}(a)=\sum_\nu(s_\nu-1/2)a_\nu$, independent of $r$, and
$b_s(a)=2a_\mu$. In the mixed coefficient, the average of
$b_s^{[2]}$ is a periodic forward difference and sums to zero. Terms pairing
$\psi(w)$ with $2a_\mu$ vanish by the zero $s$-mean of $\psi(w)$.
The last term of \cref{eq:v11-quadratic-stencil} becomes a coarse difference
of $\psi(w)$ and vanishes after the $r$ average. What remains is
\[
 \left\langle\sum_\nu
   [(x_\nu\bmod2-1/2)a_\nu,
             w_\mu(x)+w_\mu(x+e_\mu)]\right\rangle_M.
\]
If $\nu=\mu$, translation by $e_\mu$ reverses the parity weight, so its
contribution is zero. If $\nu\ne\mu$, that translation leaves the weight
unchanged and the two terms agree. This is \cref{eq:v11-one}. The mixed
coefficient of $E^{[2]}(ta+sw)$ is the true bilinear Hessian, with no extra
factor of two in the displayed identity.
\end{proof}

Define the centered coordinate, periodically continued, by
\begin{equation}\label{eq:v11-f}
 f_L(x)=x-\frac{L-1}{2},\qquad 0\le x<L,
 \qquad M_\nu(w_\mu)=\langle f_L(x_\nu)w_\mu(x)\rangle_L.
\end{equation}
It is a periodic sawtooth, not an unbounded coordinate on the torus.

\begin{theorem}[All-level homogeneous--detail root Hessian]
\label{thm:v11-mixed}
For the exact nested root source,
\begin{equation}\label{eq:v11-mixed}
 \boxed{\displaystyle
 \mathcal L_a w:=D^2\mathcal F_L(0)[a/L,w],\qquad
 (\mathcal L_a w)_\mu=\sum_{\nu\ne\mu}[a_\nu,M_\nu(w_\mu)].}
\end{equation}
This includes every fine detail mode and holds before imposing the Coulomb
constraint. Moreover
\begin{equation}\label{eq:v11-mixedbound}
 |M_\nu(w_\mu)|\le\tfrac12\|w_\mu\|_4,
 \qquad
 |\mathcal L_aw|\le C|a|\|w\|_4.
\end{equation}
On mean-zero Coulomb details, the last bound is at most
$CC_S|a|\|dw\|_2$, uniformly in $L$.
\end{theorem}
\begin{proof}
The second derivative of a composition is a sum with one nonlinear
second-derivative insertion. At stage $s$ its constant input is $2^sa/L$,
and its other input is $B^{\circ s}w$. Subsequent linear blocks multiply the
spatial mean by $2^{n-1-s}$. Therefore \cref{eq:v11-one} gives
\[
 (\mathcal L_aw)_\mu
 =\tfrac12\sum_{s=0}^{n-1}\sum_{\nu\ne\mu}
 [a_\nu,\langle\xi_\nu B^{\circ s}w_\mu\rangle_{L/2^s}],
 \quad \xi_\nu(r)=2(r_\nu\bmod2)-1.
\]
This is the exact second chain rule. No higher root coefficient enters it.

For $p=2^s$, the actual composed linear block is
\[
 (B_p w)_\mu(r)=p^{-4}\sum_{z\in\{0,\ldots,p-1\}^4}
                       \sum_{t=0}^{p-1}w_\mu(pr+z+te_\mu).
\]
When $\nu\ne\mu$, translation in the longitudinal summation leaves the
$\nu$-th block parity unchanged. Thus
\[
 \langle\xi_\nu B_pw_\mu\rangle_{L/p}
 =p\left\langle
    \bigl(2(\lfloor x_\nu/p\rfloor\bmod2)-1\bigr)w_\mu(x)
                                            \right\rangle_L.
\]
The binary digits of $x\in\{0,\ldots,L-1\}$ satisfy
\[
 \tfrac12\sum_{s=0}^{n-1}2^s
       \bigl(2(\lfloor x/2^s\rfloor\bmod2)-1\bigr)
                       =x-(L-1)/2.
\]
Substitution proves \cref{eq:v11-mixed}.

Finally $|f_L|\le L/2$ and Holder in unnormalized counting norm give
\[
 L^{-4}\left(\sum_x|f_L(x_\nu)|^{4/3}\right)^{3/4}\le\tfrac12.
\]
The bracket estimate and V9's lattice Sobolev inequality prove
\cref{eq:v11-mixedbound}. Unlike a Cauchy bound in $L\|w\|_\infty$, this
mixed estimate is valid for arbitrary details of small critical energy.
It does not give all higher root derivatives on that energy domain.
\end{proof}

\section{The leading Euler source and its complete periodic Green solution}
\label{sec:v11-Green}

For the homogeneous coarse Lie tuple $a$ put
\begin{equation}\label{eq:v11-Jtensor}
 J_{\mu\nu}(a)=[g_\mu(a),a_\nu]\quad(\mu\ne\nu),
 \qquad
 \mathscr P_8(a)=\sum_\mu\sum_{\nu\ne\mu}|J_{\mu\nu}(a)|^2.
\end{equation}
The two orientation indices are ordered. Lie contractions use the same
quaternion coefficient normalization as the action. Define a fine source
\begin{equation}\label{eq:v11-force}
 \mathfrak f_{a,\mu}(x)=L^{-4}\sum_{\nu\ne\mu}f_L(x_\nu)J_{\mu\nu}(a).
\end{equation}
Ad-invariance of the inner product gives the exact identity
\begin{equation}\label{eq:v11-force-dual}
 \langle g(a),\mathcal L_aw\rangle
                   =\sum_{x,\mu}\langle\mathfrak f_{a,\mu}(x),w_\mu(x)\rangle.
\end{equation}
Each component of this source is independent of its own link coordinate
$x_\mu$. It is therefore divergence free and has zero mean: the source lies
in the actual detail space, without an uncomputed Coulomb projection.

\begin{lemma}[An exact periodic inverse of the centered coordinate]
\label{lem:v11-Green}
The unique mean-zero solution of
\[
 2q_L(x)-q_L(x+1)-q_L(x-1)=f_L(x)
\]
on the periodic length-$L$ cycle is
\begin{equation}\label{eq:v11-q}
 \boxed{\displaystyle
 q_L(x)=\frac{f_L(x)}{24}\bigl(L^2+3-4f_L(x)^2\bigr).}
\end{equation}
At $L=2$ both neighbour terms are present even though their sites coincide.
Moreover
\begin{equation}\label{eq:v11-capacity}
 \boxed{\displaystyle
 \sum_{x=0}^{L-1}f_L(x)q_L(x)
      =\frac{L(L^2-1)(L^2+11)}{720},\qquad
 c_L:=L^{-5}\sum_xf_L(x)q_L(x)
      =\frac{1+10L^{-2}-11L^{-4}}{720}.}
\end{equation}
For $L\ge2$, $1/720<c_L\le1/256$, and
$0<c_L-1/720\le1/(72L^2)$.
\end{lemma}
\begin{proof}
The second difference of $-r^3/6+\alpha r$ is $r$. Setting
$r=x-(L-1)/2$, the periodic endpoint equation fixes
$\alpha=(L^2+3)/24$. Direct substitution at $x=-1,L$ gives the periodic
values; antisymmetry about the midpoint gives zero mean. Positivity of the
cycle Laplacian on mean-zero vectors proves uniqueness.

The centered sums are
\[
 \sum_x f_L(x)^2=\frac{L(L^2-1)}{12},\qquad
 \sum_x f_L(x)^4=\frac{L(L^2-1)(3L^2-7)}{240}.
\]
Insert these into \cref{eq:v11-q} to obtain \cref{eq:v11-capacity}.
The stated bounds follow by writing $u=L^{-2}\in(0,1/4]$ and inspecting
$(1+10u-11u^2)/720$.
\end{proof}

\begin{theorem}[Full detail profile and energy of the root forcing]
\label{thm:v11-profile}
The solution of $\mathcal J_Lv_4=\mathfrak f_a$ on $V_L$ is
\begin{equation}\label{eq:v11-v4}
 \boxed{\displaystyle
 (v_4(a))_\mu(x)=L^{-4}\sum_{\nu\ne\mu}q_L(x_\nu)J_{\mu\nu}(a).}
\end{equation}
It obeys
\begin{equation}\label{eq:v11-energy}
 \boxed{\displaystyle
 \|dv_4(a)\|_2^2
 =\langle\mathfrak f_a,\mathcal J_L^{-1}\mathfrak f_a\rangle
 =c_L\mathscr P_8(a).}
\end{equation}
In particular, with $J_*(a)=\max_\mu\sum_{\nu\ne\mu}|J_{\mu\nu}(a)|$,
\begin{equation}\label{eq:v11-scaled}
 \boxed{\displaystyle
 L\|v_4(a)\|_\infty\le J_*(a)/24,\qquad
 L^2\|\nabla v_4(a)\|_\infty\le J_*(a)/12.}
\end{equation}
These are estimates for the ordinary degree-four displacement coefficient.
They are not bounds for its unknown finite-amplitude remainder.
\end{theorem}
\begin{proof}
The displayed field is mean zero and divergence free for the same reasons as
its source. On it the four-dimensional Laplacian separates into the one
coordinate in each summand, so \cref{lem:v11-Green} proves the equation.
For two different indices $\nu,\rho$, the spatial sum of
$f_L(x_\nu)q_L(x_\rho)$ vanishes by their zero means. A diagonal term gives
$L^3\sum_x f_L(x)q_L(x)$. The two factors $L^{-4}$ in source and solution
therefore give precisely $L^{-5}\sum_xf_Lq_L=c_L$. This proves the complete
energy, with no cross mode omitted.

For $L\ge2$, \cref{eq:v11-q} implies $\max|q_L|\le L^3/24$.
Its forward differences have maximum absolute value at most $L^2/12$;
one may verify this directly from the quadratic first difference and the
periodic endpoint difference $-(L^2-1)/12$. Multiplying by $L^{-4}$ gives
\cref{eq:v11-scaled}.
\end{proof}

There is an exact continuum profile for this \emph{coefficient}. With midpoint
coordinate $r_x=(x+1/2)/L-1/2$ and
$q_\infty(r)=r(1-4r^2)/24$, periodically joined on $[-1/2,1/2]$, one has
\begin{equation}\label{eq:v11-profile-limit}
 \frac{q_L(x)}{L^3}=q_\infty(r_x)+\frac{r_x}{8L^2},\qquad
 \frac{q_L(x+1)-q_L(x)}{L^2}
       =q_\infty'(r_x+1/(2L))+\frac1{12L^2}.
\end{equation}
The second formula includes the wrapping edge, using the common endpoint
value of the periodic first derivative. Thus the scaled displacement and its
first difference have exact $O(L^{-2})$ errors at their natural staggered
coordinates. This identifies a scale-aware elliptic response rather than
extrapolating its size from an energy bound.

\begin{corollary}[Only commuting homogeneous tuples avoid this leading response]
\label{cor:v11-nonzero}
For real compact Lie data,
\begin{equation}\label{eq:v11-nonzero}
 \mathscr P_8(a)=0
 \quad\Longleftrightarrow\quad
 [a_\mu,a_\nu]=0\ \text{for all }\mu,\nu.
\end{equation}
\end{corollary}
\begin{proof}
Put $A_\mu=-\sum_{\nu\ne\mu}\operatorname{ad}_{a_\nu}^2\succeq0$ on Lie
space. Then $g_\mu=A_\mu a_\mu$ and
$\sum_{\nu\ne\mu}|[g_\mu,a_\nu]|^2=\langle g_\mu,A_\mu g_\mu\rangle$.
If this vanishes, $g_\mu$ belongs both to $\ker A_\mu$ and to
$\Ran A_\mu$, so $g_\mu=0$. Hence all $g_\mu$ vanish and Euler homogeneity
gives $4E_4(a)=\sum_\mu\langle g_\mu,a_\mu\rangle=0$.
Every summand in $E_4$ is nonnegative, proving commutativity. The converse
follows immediately from \cref{eq:v11-E4}. This is a finite classical
source statement, not a mass gap or a global compact-fibre repair estimate.
\end{proof}

\section{The actual minimum differs from the homogeneous trial at degree eight}
\label{sec:v11-energy-correction}

\begin{theorem}[Leading constrained displacement and complete eighth-order decrease]
\label{thm:v11-main}
For every fixed dyadic $L\ge2$ and homogeneous Lie direction $a$, take the
\emph{actual} source curve $v(t)=\Phi_n(ta)$. On the branch of
\cref{prop:v11-existence},
\begin{align}\label{eq:v11-minimum-profile}
 A_L^*(\Phi_n(ta))&=ta/L+t^4v_4(a)+O_L(t^5),\\
 L\langle A_L^*(\Phi_n(ta))\rangle_L
       &=ta+t^5b_5(a)+O_L(t^6),\notag\\
 b_{5,\mu}(a)&=-c_L\sum_{\nu\ne\mu}[a_\nu,J_{\mu\nu}(a)].\notag
\end{align}
Here $v_4$ is the explicit full detail field in \cref{eq:v11-v4}.
The first line states ordinary Taylor coefficients; for example the fourth
true derivative of that curve is $24v_4$.
Its energy satisfies
\begin{equation}\label{eq:v11-energy-curve}
 \mathcal G_L(\Phi_n(ta))
 =\mathcal S(ta/L)-\frac{t^8c_L}{2}\mathscr P_8(a)+O_L(t^9).
\end{equation}
Equivalently, as germs in the \emph{fixed balanced source} $v$,
\begin{equation}\label{eq:v11-main}
 \boxed{\displaystyle
 \mathcal G_L(v)=E_L^{\rm hom}(v)
      -\frac{1+10L^{-2}-11L^{-4}}{1440}\mathscr P_8(v)
      +O_L(|v|^9).}
\end{equation}
For every noncommuting $a$, the homogeneous field $ta/L$ is therefore not
stationary on its actual root fibre for all sufficiently small nonzero $t$.
The homogeneous trial is exact through degree seven, but not degree eight.
\end{theorem}
\begin{proof}
Write the exact constraint chart for the source $\Phi_n(ta)$ as
$A=s(t,w)/L+w$, $w\in V_L$, with $s(t,0)=ta$. Differentiating in $w$ at
zero and using \cref{eq:v11-mixed} gives
\begin{equation}\label{eq:v11-chart-derivative}
 D_ws(t,0)=-D\Phi_n(ta)^{-1}D\mathcal F_L(ta/L)|_{V_L}
                         =-t\mathcal L_a+O_L(t^2).
\end{equation}
The constant part of $D\mathcal F_L|_{V_L}$ vanishes by
\cref{eq:v11-zero}. The leading homogeneous source derivative is the identity.

At every constant fine connection, the Wilson gradient is constant in
position by translation invariance. It annihilates all mean-zero detail
variations exactly, not merely through a specified order. The constant
connection action begins with $E_4(s)$, with no lower terms. Its mean-direction
gradient at $s=ta$ is therefore $t^3g(a)+O_L(t^4)$ (indeed the first bare
correction is of order five). The derivative of the constrained chart action
in a detail direction is consequently
\[
 -t^4\langle g(a),\mathcal L_aw\rangle+O_L(t^5)\|w\|
       =-t^4\langle\mathfrak f_a,w\rangle+O_L(t^5)\|w\|.
\]
Its detail Hessian at zero is $\mathcal J_L$. The finite analytic stationary
solve then yields $w_*(t)=t^4\mathcal J_L^{-1}\mathfrak f_a+O_L(t^5)$,
which is \cref{eq:v11-minimum-profile}'s first assertion. Substituting into
the constraint chart gives $s(t,w_*)=ta-t^5\mathcal L_av_4+O_L(t^6)$.
By the same separated spatial sums used in \cref{thm:v11-profile},
$M_\nu((v_4)_\mu)=c_LJ_{\mu\nu}$, proving the displayed $b_5$.

Expanding the chart action at $w=0$ and inserting $w_*(t)$ gives
\[
 -t^8\langle\mathfrak f_a,\mathcal J_L^{-1}\mathfrak f_a\rangle
 +\frac{t^8}{2}\langle\mathcal J_L^{-1}\mathfrak f_a,
                       \mathcal J_L\mathcal J_L^{-1}\mathfrak f_a\rangle
 +O_L(t^9).
\]
This is \cref{eq:v11-energy-curve}. All details are allowed in this solve;
the explicit source happens to lie in the separable subspace on which the
Green inverse is computed. No other detail can change the minimum of this
quadratic response problem.

Finally $\Phi_n(a)=a+O(|a|^3)$. Substitution of its inverse into a homogeneous
degree-eight polynomial leaves that degree-eight part unchanged. Polarization
of the analytic germs therefore gives \cref{eq:v11-main}. For noncommuting
$a$, \cref{cor:v11-nonzero} and $c_L>0$ make the leading constrained derivative
nonzero; its nonvanishing persists for sufficiently small nonzero $t$ at fixed
$L$. All uses of $O_L$ retain their regulator-dependent scope.
\end{proof}

\begin{corollary}[Cofinal classical homogeneous-source packet through degree eight]
\label{cor:v11-cofinal}
Let $J^8$ denote the Taylor polynomial at the stationary zero source. There is
a limiting polynomial $\mathcal G_*^{[8]}$ such that, in any fixed tensor norm
on this 12-dimensional source space,
\begin{equation}\label{eq:v11-cofinal}
 \boxed{\displaystyle
 \|J^8\mathcal G_L-\mathcal G_*^{[8]}\|\le C L^{-2},\qquad
 \mathcal G_*^{[8]}=
 J^8\!\left(E_4(\Phi^{-1}(v))-\frac{\mathscr P_8(v)}{1440}\right).}
\end{equation}
The errors are summable for dyadic $L$. This statement includes all homogeneous
coarse source directions, while minimizing over all fine inhomogeneous details.
It is not convergence of the complete finite-amplitude functions $\mathcal G_L$
or a bound for spatially varying coarse source tensors.
\end{corollary}
\begin{proof}
V10 proves convergence of $E_L^{\rm hom}$ to $E_4\circ\Phi^{-1}$, with all
fixed derivatives on smaller common domains, at rate $O(L^{-2})$.
The explicit coefficient $c_L$ converges to $1/720$ at that rate by
\cref{lem:v11-Green}. Equation \cref{eq:v11-main} identifies exactly the
Taylor polynomials through degree eight, irrespective of any size estimate
on its higher remainder. Combine the two coefficient bounds. The source
space has fixed dimension, so all fixed tensor norms differ by fixed
constants. This does not assume a common domain for the true minima.
\end{proof}

\section{An explicit complete two-colour source polynomial, not just a trial value}
\label{sec:v11-plane}

On the source plane $v=(xE_1,yE_2,0,0)$ the invariant in
\cref{eq:v11-Jtensor} is particularly simple:
\begin{equation}\label{eq:v11-plane-force}
 g_1=4xy^2E_1,\quad g_2=4x^2yE_2,\quad
 J_{12}=8xy^3E_3,\quad J_{21}=-8x^3yE_3,
 \qquad \mathscr P_8=64(x^6y^2+x^2y^6).
\end{equation}
All other force tensors vanish. The leading fine detail is therefore
\[
 (v_4)_1(x_{\rm site})=\frac{8xy^3}{L^4}q_L(x_{{\rm site},2})E_3,
 \qquad
 (v_4)_2(x_{\rm site})=-\frac{8x^3y}{L^4}q_L(x_{{\rm site},1})E_3.
\]
The use of a third Lie colour is forced by the actual commutator; it is not a
new polarization inserted to lower the action artificially.

\begin{theorem}[Full planar root-constrained classical germ through degree eight]
\label{thm:v11-plane}
Put $u=L^{-2}$. The actual constrained minimum obeys
\begin{equation}\label{eq:v11-plane}
 \boxed{\begin{aligned}
 \mathcal G_L(xE_1,yE_2,0,0)
 ={}&2x^2y^2+\frac23(1-2u)(x^4y^2+x^2y^4)\\
 &+\frac4{45}(10u^2-10u+1)(x^6y^2+x^2y^6)\\
 &+\frac{2u(9u-8)}9x^4y^4+O_L((|x|+|y|)^{10}).
 \end{aligned}}
\end{equation}
For example at $L=2$ and $x=y=t$ this is
$2t^4+\frac23t^6-\frac{19}{40}t^8+O(t^{10})$.
The cofinal polynomial on this plane is
\[
 2x^2y^2+\frac23(x^4y^2+x^2y^4)
                      +\frac4{45}(x^6y^2+x^2y^6).
\]
This last display is the limiting \emph{Taylor polynomial}, not a claimed
formula for the full limiting potential.
\end{theorem}
\begin{proof}
We specify the finite homogeneous calculation needed beyond V10's sextic
term. Write $T=2I+C_3+C_4+C_5+O(a^6)$ for the literal homogeneous balanced
root map, using ordinary homogeneous coefficients. By V10, $C_3=-g/4$.
Its finite iterate has
\[
 \Phi_n(a)=a+\alpha C_3+\beta_4 C_4+\gamma C_5
                         +\delta DC_3[C_3]+O(a^6),
\]
where
\begin{equation}\label{eq:v11-Phi5}
 \alpha=\frac{1-u}{6},\quad
 \beta_4=\frac{1-L^{-3}}{14},\quad
 \gamma=\frac{1-u^2}{30},\quad
 \delta=\frac{1-5u+4u^2}{180}.
\end{equation}
These follow by substituting $\Phi_{n+1}(a)=T(\Phi_n(a/2))$ and matching
degrees. For example $\delta_{n+1}=\delta_n/16+\alpha_n/32$.
Inversion through degree five gives
\[
 \Phi_n^{-1}(v)=v-\alpha C_3-\beta_4C_4-\gamma C_5
                     +(\alpha^2-\delta)DC_3[C_3]+O(v^6).
\]
Terms beyond this degree cannot affect the eighth-degree constant Wilson
action, which begins in degree four.

Here is a finite exact determination of the only new scalar root contraction
needed on the plane:
\begin{equation}\label{eq:v11-T5scalar}
 g\cdot C_5=\frac13(x^6y^2+x^2y^6)+\frac{52}{3}x^4y^4.
\end{equation}
Homogeneity gives total degree eight. Global conjugations fixing $E_1$ or
$E_2$ make this scalar even separately in $x,y$; it vanishes on each commuting
axis. Thus its only possible monomials are $x^6y^2,x^4y^4,x^2y^6$.
Expanding the \emph{literal} ordered quaternion exponentials and distinguished
root logarithms through degree five gives
\[
 \begin{array}{c|cc|c}
 (x,y)&(C_{5,1})_{E_1}&(C_{5,2})_{E_2}&g\cdot C_5\\ \hline
 (1,1)&7/4&11/4&18\\
 (1,2)&8&43/2&300\\
 (2,1)&27/2&12&300
 \end{array}
\]
These three evaluations are an invertible three-monomial interpolation and
therefore prove \cref{eq:v11-T5scalar}; they are not a numerical extrapolation
of scale dependence. The exact word prescription is V10's
\cref{eq:v10-hompaths,eq:v10-K,eq:v10-balance}, with exponential and logarithm
series truncated at degree five. The supplied checker reconstructs these
values and retains two additional regression ratios.

Let $H_4=D^2E_4$ and let $E_6=-2(x^4y^2+x^2y^4)/3$ be the primitive sixth
constant-connection coefficient. Direct differentiation on the plane gives
\[
 g^TH_4g=64(x^6y^2+x^2y^6)+256x^4y^4,
 \quad
 DE_6[C_3]=\frac43(x^6y^2+x^2y^6)+\frac{16}{3}x^4y^4.
\]
The eighth-degree term of the homogeneous trial is
\[
 \frac{3\alpha^2/2-\delta}{16}\,g^TH_4g
       -\gamma g\cdot C_5-u\alpha DE_6[C_3]+u^2E_8,
 \qquad
 E_8=\frac4{45}(x^6y^2+x^2y^6)+\frac29x^4y^4.
\]
Substitution reduces it to
\[
 \frac2{45}(9u^2-10u+3)(x^6y^2+x^2y^6)
                       +\frac{2u(9u-8)}9x^4y^4.
\]
The sixth-degree term is V10's $2(1-2u)(x^4y^2+x^2y^4)/3$.
Now subtract the \emph{actual} relaxation
$\frac12c_L\mathscr P_8=\frac2{45}(1+10u-11u^2)(x^6y^2+x^2y^6)$ from
\cref{thm:v11-main}. This gives \cref{eq:v11-plane}.

The scalar constrained germ is invariant under the same global conjugations,
by uniqueness and Wilson/root equivariance. It is even in $x$ and in $y$.
Consequently its ninth-degree and all odd-total-degree terms vanish on this
plane, justifying the displayed degree-ten remainder. Its constant may still
depend on $L$. Setting $u=0$ is legitimate for the finite coefficient polynomial
by \cref{cor:v11-cofinal}, not by an unproved exchange of finite-amplitude limits.
\end{proof}

\section{Independent constrained-curve checks and the new proof boundary}
\label{sec:v11-verification}

\paragraph{An independent full-stencil regression.}
For $a=(E_1,E_2,0,0)$, form the fine Coulomb curve
\[
 A_z(t)=ta/L+zt^4v_4(a)+zt^5b_5(a)/L,
 \qquad z\in\R.
\]
Here $z$ is a test multiplier, not a lattice coordinate. Literal nested root
multiplication gives
$\mathcal F_L(A_z(t))=\Phi_n(ta)+O(t^6)$.
The exact full Wilson plaquette words give
\begin{equation}\label{eq:v11-testcurve}
 [t^8]\mathcal S(A_z(t))
 =\frac{2}{5L^4}
     +c_L\mathscr P_8(a)\left(\frac{z^2}{2}-z\right).
\end{equation}
The quadratic in $z$ has its minimum at $z=1$. At $L=2$, the bare coefficient
is $1/40$, $c_L\mathscr P_8=1/2$, and the corrected coefficient along the
source curve $\Phi_n(ta)$ is $-9/40$. This differs from the $-19/40$ in
\cref{thm:v11-plane} because that theorem uses the straight source $v=ta$.
The two parameter conventions are not interchangeable.

Only constraint accuracy through degree five is required to check this eighth
energy coefficient: repairing an $O(t^6)$ source defect changes the constant
mode by $O(t^6)$, and the homogeneous action gradient begins at degree three,
so its first possible energy effect is degree nine. This finite check is
independent of the variational proof of the full minimizing displacement.

The new driver \path{verify_ym_root_relaxation_V11.py} reports 21 passing exact
test families. They include the literal full spatial mixed root Hessian on
sides two and four, the exact compensated source through degree five, and the
full Wilson word coefficient on three constrained test curves. General
noncommuting four-component directions test the four-dimensional Coulomb
constraint, Poisson equation, energy, and source duality. Independent rational
cycle matrices and exact Fourier inversions check the Green polynomial, while
symbolic sums establish its all-length capacity and scaled midpoint identities.
Literal homogeneous source inversion checks the complete planar coefficient
at four dyadic levels and two source ratios. All 21 V10 test families were
also freshly rerun and passed. No floating-point fit or interval sampling is
used to prove an all-level statement.

\paragraph{What is now known about the actual constraint.}
The mixed second root derivative with an arbitrary fine detail is controlled
uniformly in critical energy when its other leg is homogeneous. The root-induced
Euler source is computed and inverted on the full detail space. Its first
nonlinear response has precisely the fine amplitude and first-difference
scales asked for at the end of V10. Most importantly, the actual minimized
classical $J^8$ packet for every one-cell homogeneous source now has a cofinal
limit, rather than being replaced by a homogeneous trial action. The
nonzero degree-eight decrease survives refinement.

\paragraph{What this still does not imply.}
The complete displacement is $t^4v_4+O_L(t^5)$ on the chosen source curve.
The scaled bounds on $v_4$ do not bound $L\|O_L(t^5)\|_\infty$ or its first
difference uniformly in $L$. A common nonlinear background neighbourhood
still requires higher mixed root derivatives and elliptic estimates on their
complete forcing, or a direct common-domain argument. The calculation has
one coarse cell; arbitrary spatially varying coarse sources and their
anchored quantum vertices are not covered by a fixed 12-dimensional coefficient
limit. The measured $J^6Q,J^4B$ packet, analytic tails, local-to-compact
comparison, and physical-time contraction remain separate. No quantum
continuum measure or physical mass is inferred from a negative classical
relaxation term or from the positivity of $\mathscr P_8$.

The useful next acquisition is consequently more specific than another generic
uniformity assumption: estimate the remaining mixed derivatives of the actual
root source against arbitrary detail fields in a scale-aware dual norm. The
centered-coordinate identity and its explicit Green return provide both an
analytic starting point and a regression that such a bound must retain. It
would be incorrect to recover a simpler induction by deleting this return,
freezing the homogeneous lift, or identifying the source-coordinate Jacobian
with a quantum determinant.

\paragraph{Append-only record.}
The complete V1--V10 mathematical body, including prior bibliography entries
and end markers, is retained byte-for-byte. Only version displays and the
front-matter date are advanced. The bundle contains the new appendix, executable
checks, exact reports, portable assembly script, and source/output hashes.
Earlier results are inherited with their stated hypotheses; the fresh rerun
validates their finite V10 regressions, not their entire analytic proof chain.
Any later correction is to be appended explicitly. Future V12 work is to
begin after the marker below.

% END OF VERSION 11 MATHEMATICAL BODY.
% Append subsequent requested versions below this marker; retain V1--V11 above.


\clearpage
\section{V12: a contracting remainder norm closes the actual root-map domain}
\label{sec:v12-context}

V11 determines the first inhomogeneous relaxation and the complete classical
homogeneous-source packet through degree eight.  Its final unresolved
classical input is a common domain for the full nonlinear constrained branch,
not another coefficient of that branch.  The route suggested by the preceding
scaled pointwise estimates would ask for more regularity than is needed to
construct that domain.  This version instead uses two different counting norms:
the exact linear block is nonexpansive in $\ell^4$ and strictly contractive in
$\ell^2$, while its local nonlinear remainder maps two $\ell^4$ inputs into
$\ell^2$.  Keeping that remainder in its contracting norm prevents a loss at
every dyadic step.

The resulting theorem concerns the \emph{literal balanced root map} of V10,
not an alternative group mean or a linear replacement.  It is holomorphic on a
common small global $\ell^4$ ball at every refinement depth and surrounding
volume.  All its fixed-order mixed derivatives are uniformly bounded, including
arbitrary energy-controlled fine details.  For one coarse periodic cell this
populates V9's nonlinear-source hypothesis, and gives a common analytic local
minimum for the actual root constraint and the full Wilson action.  In
particular V11's regulator-dependent higher remainders can now be made uniform
in the critical-energy norm.  The leading degree-eight relaxation remains
present with its original coefficient.

A further application constructs a dyadic continuum limit of the root
\emph{observable} on a small continuum $L^4$ ball.  It permits spatially varying
inputs and is not restricted to homogeneous connections.  This is a
classical deterministic observable limit, not a measure, a limiting minimum,
or a quantum continuum theory.

\paragraph{Dependency audit and scope.}
The primitive analytic stencil, exact orbit equivalence and its linear
cancellation are inherited from \cref{prop:v10-orbit,lem:v10-local}.
V9 supplies the lattice Sobolev estimate and uniform Wilson action estimates;
V11 supplies the exact leading relaxation.  None is reproved as a new
coefficient calculation.  Inspection of the current lineage and the relevant
supplied YM and Grand Tensor passages found no preceding two-norm iteration
estimate of the form proved below.  Semantic file retrieval returned no
matching passages, and was followed by exact-text inspection; this is a
nonduplication audit of the accessed material, not a universal novelty claim.
Analytic lattice averaging and constrained background minimization have a
substantial earlier literature, including \cite{BalabanV9}.  No theorem for a
different averaging rule is substituted for the present estimates.

All norms below use unnormalized counting measure on lattice links.  A small
global $\ell^4$ or critical-energy ball is not a thermodynamic small-field
maximum-norm domain.  On more than one coarse cell the balanced source is not
automatically in logarithmic Coulomb gauge.  The general source estimate and
the one-cell variational theorem are therefore stated separately.

\section{Exact block norm scaling and the local quadratic estimate}
\label{sec:v12-local}

Let $M\ge2$ be even.  Write $\mathcal X_M$ for Lie-valued link arrays on
$(\mathbb Z/M\mathbb Z)^4$, with the complexified quaternion coefficient norm
on each link.  Set
\[
 \|A\|_p=\left(\sum_{x,\mu}|A_\mu(x)|^p\right)^{1/p},
 \qquad 1\le p<\infty.
\]
Complex coefficient and matrix operator norms are uniformly equivalent on one
link; this fixed conversion is included in every primitive radius below.
Let $\mathfrak T_M:\mathcal X_M\to\mathcal X_{M/2}$ be the actual balanced
root logarithm \cref{eq:v10-path,eq:v10-K,eq:v10-balance}.  Its derivative is
\begin{equation}\label{eq:v12-B}
 (B_M A)_{r,\mu}=\frac1{16}\sum_{s\in\{0,1\}^4}
       \bigl(A_{2r+s,\mu}+A_{2r+s+e_\mu,\mu}\bigr).
\end{equation}
Repeated links on a short torus retain their multiplicity in this formula.

\begin{lemma}[Exact counting-norm scale of the actual linear block]
\label{lem:v12-B}
For every $1\le p\le\infty$,
\begin{equation}\label{eq:v12-Bnorm}
 \boxed{\ \|B_M\|_{p\to p}=2^{1-4/p}.\ }
\end{equation}
In particular, $\|B_M\|_{4\to4}=1$ and $\|B_M\|_{2\to2}=1/2$.
These constants are independent of $M$ and of the number of Lie components.
\end{lemma}
\begin{proof}
Each output row has nonnegative weights of sum two.  Each fine link appears
twice among all the chronological path occurrences for its orientation, so
the corresponding column sum is $1/8$.  Jensen's inequality gives
\[
 |B_MA(r,\mu)|^p\le 2^{p-1}\frac1{16}
       \sum_s\bigl(|A_{2r+s,\mu}|^p+|A_{2r+s+e_\mu,\mu}|^p\bigr).
\]
Sum over $r,\mu$ to obtain $\|B_MA\|_p^p\le2^{p-4}\|A\|_p^p$.
The endpoint estimates follow directly from the row and column sums.
A constant nonzero field in one orientation attains equality: its amplitude
is doubled and the number of sites is divided by sixteen.  Thus these are
operator norms, not just upper bounds.
\end{proof}

Write $E_M(A)=\mathfrak T_M(A)-B_MA$.  It has no constant or linear term.
V10's primitive holomorphic neighbourhood is a local maximum-norm
neighbourhood, so it contains a fixed small $\ell^4$ ball at every volume.

\begin{lemma}[The local remainder improves the counting exponent]
\label{lem:v12-E}
There exist $r_{\rm p}>0$ and $C\ge1$, independent of $M$, such that whenever
$\|A\|_\infty,\|A'\|_\infty<r_{\rm p}$,
\begin{align}
 \|E_M(A)\|_2&\le C\|A\|_4^2,\label{eq:v12-E2}\\
 \|E_M(A)-E_M(A')\|_2
  &\le C(\|A\|_4+\|A'\|_4)\|A-A'\|_4.\label{eq:v12-Elip}
\end{align}
Also, with a fixed constant $C_8$,
\begin{equation}\label{eq:v12-E4}
                  \|E_M(A)\|_4\le C_8\|A\|_8^2.
\end{equation}
All these statements hold on the complexified spaces, with holomorphic $E_M$.
\end{lemma}
\begin{proof}
An output edge uses the ordered trees in its two adjacent cells and its
sixteen two-link paths.  It depends on at most $54$ fine links, and a fine
link belongs to at most ten such output supports.  These are safe upper
bounds: the tree union has at most $30$ links and the straight-path union at
most $24$; a tree link can be used at at most eight incident outputs and a
straight link at at most two.  Coincidences on a small torus reduce supports,
or are counted with bounded multiplicity.

Let $K$ be the corresponding positive incidence matrix.  Uniform Cauchy
bounds on the fixed primitive polydisc, together with $E_M(0)=DE_M(0)=0$,
give the pointwise bounds on a smaller polydisc
\[
 |E_M(A)_e|\le C_0(K|A|)_e^2,
\]
\[
 |E_M(A)_e-E_M(A')_e|
 \le C_0(K(|A|+|A'|))_e(K|A-A'|)_e.
\]
The integral first-derivative formula along the segment between the two
fields proves the second estimate from the same local second-derivative
bound.  Every segment remains in the chosen polydisc.
For $p\ge1$, Jensen and the support counts give
\[
                  \|K\|_{p\to p}\le54^{1-1/p}10^{1/p}.
\]
Use Holder with two $\ell^4$ factors to obtain
\cref{eq:v12-E2,eq:v12-Elip}; use two $\ell^8$ factors for
\cref{eq:v12-E4}.  All constants concern one finite stencil and its bounded
incidence, not the total number of cells.  The branch and holomorphy are the
ones already proved in \cref{lem:v10-local}.
\end{proof}

\section{A uniform nonlinear root map on a global critical counting ball}
\label{sec:v12-iteration}

Let $L=2^n$ and let the fine torus have side $NL$.  Set
\[
 A_0=A,\qquad A_{s+1}=\mathfrak T_{NL/2^s}(A_s),\qquad
 B^{[s]}=B_{NL/2^{s-1}}\cdots B_{NL},
\]
with $B^{[0]}=I$.  Define the nonlinear remainder \emph{before} estimating it:
\begin{equation}\label{eq:v12-Rs}
 R_s=A_s-B^{[s]}A,
 \qquad
 R_{s+1}=B_{NL/2^s}R_s+E_{NL/2^s}(A_s).
\end{equation}
The norms at stage $s$ are those of its current lattice; no cell-volume
normalization is inserted between these identities.

\begin{theorem}[Uniform holomorphic nested root map by two-norm induction]
\label{thm:v12-root}
Let
\begin{equation}\label{eq:v12-radius}
              r_*:=\min\{r_{\rm p}/2,\,(16C)^{-1}\}.
\end{equation}
For $\rho=\|A\|_4<r_*$, every nested step uses the admitted identity
logarithm branch, and for all $0\le s\le n$,
\begin{equation}\label{eq:v12-rootbounds}
 \boxed{\ \|R_s\|_2\le4C\rho^2,\qquad
               \|A_s\|_4\le\rho+4C\rho^2\le\tfrac54\rho.\ }
\end{equation}
In particular, the actual balanced composite $\mathcal F_{L,N}(A)=A_n$ is
holomorphic on this common ball and
\begin{equation}\label{eq:v12-terminal}
 \boxed{\ \|\mathcal F_{L,N}(A)-B_LA\|_2\le4C\|A\|_4^2,\qquad
 D\mathcal F_{L,N}(0)=B_L.\ }
\end{equation}
The constants and domain are independent of $L,N$.  For real compact inputs
this composite represents the same coarse gauge orbit as the literal nested
root block, with the qualifications on measured coordinates in V10 unchanged.
\end{theorem}
\begin{proof}
The embedding $\ell^2\subset\ell^4$ has norm one for unnormalized counting
measure.  By \cref{lem:v12-B}, $\|B^{[s]}A\|_4\le\rho$.  If the remainder
bound holds at stage $s$, then
\[
 \|A_s\|_\infty\le\|A_s\|_4\le\rho+4C\rho^2
           \le\tfrac54\rho<r_{\rm p}.
\]
The actual primitive step is therefore defined, and its remainder satisfies
\[
 \|R_{s+1}\|_2\le\tfrac12\|R_s\|_2
                     +C(\rho+\|R_s\|_2)^2
 \le2C\rho^2+C\rho^2(1+4C\rho)^2
 <4C\rho^2,
\]
since $4C\rho\le1/4$ and $2+(5/4)^2<4$.  Starting at $R_0=0$ proves
the induction.  It applies without change on complex coefficient spaces,
so each finite composition is holomorphic on the same ball.  Its first
jet is the product of the primitive first jets.  Exact real orbit equivalence
is \cref{prop:v10-orbit} applied at the now justified intermediate fields.
No gauge-fixed density equality or Jacobian cancellation is inferred from
that orbit statement.
\end{proof}

A same-norm estimate $\|A_{s+1}\|_4\le\|A_s\|_4+C\|A_s\|_4^2$ would
obscure this conclusion.  Its linear marginal part is retained as
$B^{[s]}A$; only the nonlinear remainder is estimated in the norm in which
the subsequent block contracts.  The factor $1/2$ is a property of the
actual four-dimensional average, not a hypothesis of nonlinear RG stability.
This argument controls the observable, not the measured effective action.

\begin{corollary}[Uniform Lipschitz and all finite mixed-derivative bounds]
\label{cor:v12-jets}
Put $\mathcal Z_{L,N}=\mathcal F_{L,N}-B_L$.  If
$\|A\|_4,\|A'\|_4\le r<r_*$, then
\begin{equation}\label{eq:v12-Rlip}
 \|\mathcal Z_{L,N}(A)-\mathcal Z_{L,N}(A')\|_2
                      \le10Cr\,\|A-A'\|_4.
\end{equation}
Consequently the whole composite has $\ell^4$ Lipschitz bound $1+10Cr$.
For $\|A\|_4<r_*/2$, every $m\ge2$ obeys, for example,
\begin{equation}\label{eq:v12-jets}
 \boxed{\
 \|D^m\mathcal F_{L,N}(A)[W_1,\ldots,W_m]\|_2
 \le4Cr_*^2\left(\frac{2m}{r_*}\right)^m\prod_i\|W_i\|_4.\ }
\end{equation}
Also $\|D\mathcal Z_{L,N}(A)\|_{4\to2}\le10C\|A\|_4$.
All bounds are independent of the refinement depth and surrounding volume.
\end{corollary}
\begin{proof}
Use the exact remainder recursion on two inputs, and set
$d_s=\|R_s(A)-R_s(A')\|_2$, $d=\|A-A'\|_4$.
The preceding theorem and \cref{eq:v12-Elip} give
\[
 d_{s+1}\le(1/2+(5/2)Cr)d_s+(5/2)Cr\,d.
\]
Its first coefficient is at most $3/4$ because $Cr\le1/16$.
Sum the geometric series, using $d_0=0$, to obtain $d_s\le10Cr\,d$.
The derivative bound follows by taking the limit $A'\to A$ and then decreasing
$r$ to $\|A\|_4$; at zero it follows from the exact first jet.

For higher derivatives, apply multivariate Cauchy to the remainder on
$A+\sum_i z_iW_i$, with normalized directions and
$|z_i|<r_*/(2m)$.  Its value has $\ell^2$ norm at most $4Cr_*^2$ on that
polydisc.  Cauchy's bound gives \cref{eq:v12-jets}; the linear map has no
higher derivatives.  Use slightly smaller discs and then a limit at boundary
radii if needed.  This is a Banach-valued finite-dimensional Cauchy estimate,
not a claim of a numerically optimized analytic radius.
\end{proof}

\section{Critical-energy consequences, and the precise gauge distinction}
\label{sec:v12-energy}

Use the actual fine Coulomb space and norm from V9,
$\|A\|_{\mathcal E}=\|dA\|_2+\|A\|_4$.
On the coarse lattice use the same expression even when the coarse field is
not Coulomb.  The norm of the discrete curl from $\ell^2$ to $\ell^2$ is at
most four, by its Fourier symbol and $\lambda\le16$.

\begin{proposition}[The actual nonlinear source is controlled in critical energy]
\label{prop:v12-Esource}
For fine Coulomb fields in the common small critical ball,
\begin{equation}\label{eq:v12-Esource}
 \boxed{\
 \|\mathcal F_{L,N}(A)-B_LA\|_{\mathcal E_c}
                       \le20C\|A\|_4^2,\qquad
 \|\mathcal F_{L,N}(A)\|_{\mathcal E_c}
                       \le\sqrt5\|A\|_{\mathcal E}+20C\|A\|_4^2.\ }
\end{equation}
The higher mixed derivatives of the remainder have analogous uniform
energy-output bounds.  In particular arbitrary mean-zero fine Coulomb
detail directions may be used in \cref{eq:v12-jets}, replacing each
$\|W_i\|_4$ by $C_S\|dW_i\|_2$.
\end{proposition}
\begin{proof}
For any coarse remainder $R$, $\|dR\|_2+\|R\|_4\le5\|R\|_2$.
Apply \cref{thm:v12-root,cor:v12-jets}.
The linear block has $\|B_LA\|_4\le\|A\|_4$.
Furthermore $dB_LA=d\Pi_cB_LA$, including the constant coarse modes.
V4's relative kernel comparison and the variational identity used in
\cref{prop:v9-splitting} give
$\|dB_LA\|_2\le\sqrt5\|dA\|_2$.
Combine these two linear estimates and use the triangle inequality.
V9's discrete Sobolev estimate supplies the final substitution.
\end{proof}

This proves all fixed-order mixed root estimates in the indicated global
norms, rather than only the homogeneous--detail Hessian of V11.  It does not
assert $L\|W_i\|_\infty$ bounds, which are unnecessary in this argument.
In particular the thin-line family of V9 has
$\|A\|_4=\pi L^{-3/4}\to0$ and eventually belongs to this domain even
though $L\|A\|_\infty=\pi$.  The theorem justifies its nested short logarithms;
it never places the single unaveraged long holonomy $-I$ in an identity chart.
Its exact abelian output remains $\pi E_3/L^3$, as proved in V10.

On multiple coarse cells, $\mathcal F_{L,N}(A)$ need not satisfy a logarithmic
Coulomb condition.  Applying a \emph{linear} transverse projector to its
nonlinear logarithm is not an exact coarse gauge transformation.  Thus
\cref{prop:v12-Esource} does not, by itself, populate a gauge-identified
onto source map for every prescribed multi-cell Coulomb source.  This
qualification disappears for the following one-cell source problem: there
the coarse divergence is zero and the output consists of twelve constant
Lie coordinates.  Uniform gauge fixing on larger coarse volumes is not
silently assumed.

\section{A common classical branch for the actual one-cell root constraint}
\label{sec:v12-minimum}

Set $N=1$, $L=2^n$, and retain the definitions of V11:
\[
 V_L=\{w:d^*w=0,\ \langle w\rangle_L=0\},\qquad
 \|w\|_V=\|dw\|_2,
 \qquad A=s/L+w.
\]
The twelve-dimensional $s$ and source $v$ use a fixed Euclidean norm.
Constants have $\|s/L\|_4\le|s|$, and
$\|w\|_4\le C_S\|w\|_V$.  On this space $B_LA=s$ exactly.  All inner
products in complex formulas have their bilinear holomorphic extensions;
positivity assertions refer to real fields.

\begin{lemma}[Uniform chart for the literal nonlinear constraint]
\label{lem:v12-chart}
There are fixed positive radii, independent of $L$, on which
$\mathcal F_L(s/L+w)=v$ has a unique solution $s=s_L(v,w)$ near zero.
It is holomorphic and
\begin{equation}\label{eq:v12-chart}
 s_L(v,w)=v+O((|v|+\|w\|_V)^2),\qquad
 A_L(v,w)=s_L(v,w)/L+w.
\end{equation}
These bounds, and every fixed derivative of the chart on smaller balls,
have constants independent of $L$.  For real fields this is an actual
constraint chart, not just a formal jet chart.
\end{lemma}
\begin{proof}
The exact equation is
\[
       s=v-\mathcal Z_{L,1}(s/L+w).
\]
By \cref{thm:v12-root,cor:v12-jets}, its correction has size at most
$C_1(|s|+C_S\|w\|_V)^2$ and its $s$ derivative has norm at most
$C_2(|s|+C_S\|w\|_V)$.  On fixed sufficiently small balls it is a contraction
with constant at most $1/2$, mapping the chosen $s$ ball into itself.
This constructs the chart, including its quadratic error, by uniform Picard
iteration on a complex domain.  The limit is holomorphic; Cauchy estimates
give its uniform finite derivative bounds.  Norms of constant inclusions
and of the detail embedding into $\ell^4$ are independent of $L$, so no
small counting-norm Laplacian eigenvalue enters the argument.
\end{proof}

\begin{theorem}[Uniform analytic local minimum for the actual root orbit]
\label{thm:v12-minimum}
There exist $\delta,R,K>0$, independent of dyadic $L$, such that the full
Wilson action on $A_L(v,w)$ has, for $|v|<\delta$, a unique stationary point
$w_L(v)$ in $\|w\|_V<R$.  It is a strict local minimum on this constraint
chart.  The branch and
\[
             A_L^*(v)=A_L(v,w_L(v)),\qquad
             \mathcal G_L(v)=\mathcal S(A_L^*(v))
\]
have holomorphic extensions on common smaller complex balls, and
\begin{equation}\label{eq:v12-minbounds}
 \boxed{\ \|w_L(v)\|_V\le K|v|^2,\qquad
                 \|A_L^*(v)\|_{\mathcal E}\le K|v|.\ }
\end{equation}
On a common smaller real chart, the constrained Hessian satisfies
\begin{equation}\label{eq:v12-Hessian}
 \boxed{\ \tfrac12\|\xi\|_V^2
 \le D_w^2[\mathcal S(A_L(v,w))][\xi,\xi]
 \le\tfrac32\|\xi\|_V^2.\ }
\end{equation}
For every fixed field order $m$, all source derivatives of $A_L^*$ in
$\mathcal E$ and of $\mathcal G_L$ are bounded independently of $L$ on
smaller common source balls.
\end{theorem}
\begin{proof}
The constant part of a one-cell lift has zero curl.  Therefore the exact
chart action, using V9's full Wilson nonlinearity $\mathcal N$, is
\[
           \mathcal S(A_L(v,w))
               =\tfrac12\|dw\|_{\rm bil}^2+
                         \mathcal N(s_L(v,w)/L+w).
\]
This equality does not delete the noncommuting constant-connection energy:
that energy belongs to $\mathcal N$.  V9's analytic bounds and the uniform
chart estimates imply, for $q=|v|+\|w\|_V$ small,
\[
 \|D_w\mathcal N(A_L(v,w))\|_{V^*}\le Cq^2,
 \qquad \|D_w^2\mathcal N(A_L(v,w))\|\le Cq.
\]
The second chain rule includes both $D^2\mathcal N[DA_L,DA_L]$ and
$D\mathcal N[D^2A_L]$; both are controlled, not omitted.  The detail Riesz
map $\mathcal J_L:V_L\to V_L^*$ has inverse norm one.  Hence the exact
stationary equation
\[
 w=-\mathcal J_L^{-1}D_w\mathcal N(A_L(v,w))
\]
is a contraction of one fixed small detail ball, with constant at most
$1/2$, after fixing $R$ and then $\delta$.  It also maps a ball of radius
$K|v|^2$ into itself, proving \cref{eq:v12-minbounds}.  The same estimates
on a complex ball give a holomorphic branch.  The Hessian perturbation has
norm at most $1/2$, which gives \cref{eq:v12-Hessian} and strict convexity
along every line segment in the detail ball.  Cauchy estimates on the
uniform complex domain give the final derivative assertion.

For each finite $L$ this is the germ already constructed in
\cref{prop:v11-existence}, by uniqueness.  Its local gauge-orbit
interpretation is consequently the same one.  This does not prove a
uniform global gauge slice or exclude other compact critical points.
\end{proof}

The hypothesis on the nonlinear observable in V9 is thus genuinely populated
for its one-cell actual root problem; it is no longer left as an unspecified
higher-mixed-derivative stock.  Positivity is in the constrained detail
energy, not the unscaled physical Hamiltonian.  The retained constant sector
still has no quadratic mass term, and the result is entirely classical.

\section{Uniform higher remainders while retaining the V11 relaxation}
\label{sec:v12-remainders}

Let $v=\Phi_n(a)=\mathcal F_L(a/L)$, so the reference constant field has
exactly the desired source.  Write
\[
 A_L^*(\Phi_n(a))=s_L(a)/L+w_L^{\rm hom}(a),\qquad
                       w_L^{\rm hom}(a)\in V_L.
\]
The notation $s_L(a)$ in this section abbreviates the mean coordinate of the
stationary point, rather than the two-argument chart in
\cref{lem:v12-chart}.  Use $v_4(a),b_5(a),c_L,\mathscr P_8(a)$ exactly as
in V11.

\begin{theorem}[The leading response now has a uniform energy-norm remainder]
\label{thm:v12-remainders}
On one common small real or complex $a$ ball,
\begin{align}
 \|w_L^{\rm hom}(a)\|_V&\le K|a|^4,\label{eq:v12-w4}\\
 \|w_L^{\rm hom}(a)-v_4(a)\|_V&\le K|a|^5,\label{eq:v12-wrem}\\
 |s_L(a)-a-b_5(a)|&\le K|a|^6.\label{eq:v12-srem}
\end{align}
Moreover
\begin{equation}\label{eq:v12-actionrem}
 \boxed{\
 \left|\mathcal G_L(\Phi_n(a))-\mathcal S(a/L)
                     +\tfrac12c_L\mathscr P_8(a)\right|\le K|a|^9.\ }
\end{equation}
Equivalently, on a common small fixed-source ball,
\begin{equation}\label{eq:v12-fixedrem}
 \boxed{\
 \left|\mathcal G_L(v)-E_L^{\rm hom}(v)
                     +\tfrac12c_L\mathscr P_8(v)\right|\le K|v|^9.\ }
\end{equation}
The constants are independent of $L$.  For each fixed derivative order
$m\le9$, the last Taylor remainder has derivative bound
$K_m|v|^{9-m}$ on a smaller ball.  On the separately even source plane
$(xE_1,yE_2,0,0)$, its order-nine bound improves to order ten, with the
corresponding fixed derivative bounds.
\end{theorem}
\begin{proof}
At $w=0$ on the source $\Phi_n(a)$, the chart has $s=a$.
V10, or \cref{cor:v12-jets}, gives a uniform inverse of $D\Phi_n(a)$ on a
smaller ball.  Since $B_L$ annihilates every $w\in V_L$,
\[
 D_ws(\Phi_n(a),0)
   =-D\Phi_n(a)^{-1}D\mathcal F_L(a/L)|_{V_L}=O(|a|)
\]
in $V\to\mathbb C^{12}$ norm, by \cref{eq:v12-Rlip} and the Sobolev
embedding.  The Wilson gradient at a constant connection is constant in
position and annihilates all mean-zero directions.  Its mean-coordinate
gradient is $D_a\mathcal S(a/L)=O(|a|^3)$ uniformly in $L$: the constant
Wilson action starts in degree four, and its ordered exponential series
has a common bound $C|a|^4$ on a small complex ball.  The constrained
chart gradient at $w=0$ is consequently $O(|a|^4)$ in $V^*$.
The uniform Hessian inverse, or the stationary contraction, now gives
\cref{eq:v12-w4}.  The constraint chart derivative then gives
$s_L(a)-a=O(|a|^5)$, uniformly.

All these functions are holomorphic on common complex domains by the preceding
theorem and the bounded homogeneous source maps.  Their first nonzero
coefficients are exactly the $v_4,b_5$ already determined in
\cref{thm:v11-main}; uniqueness identifies the germs.  Uniform Cauchy
estimates and summation of the higher Taylor tail therefore give
\cref{eq:v12-wrem,eq:v12-srem}.  This use of V11 retains its full
inhomogeneous Green response rather than recomputing another coefficient.

The uniform $O(|a|^4)$ chart gradient and stationary displacement show that
the action difference is $O(|a|^8)$ also on a smaller complex domain.
V11 identifies its degree-eight term as $-c_L\mathscr P_8(a)/2$.
A uniform Cauchy tail gives \cref{eq:v12-actionrem}.  The homogeneous maps
and inverses have common domains and $\Phi_n(a)=a+O(|a|^3)$; replacing $a$
by $\Phi_n^{-1}(v)$ changes $\mathscr P_8$ only in degree ten and above,
proving \cref{eq:v12-fixedrem}.  A holomorphic function whose Taylor
terms below degree $k$ vanish has, on a smaller common ball, its $m$th
multilinear derivative bounded by $C_m|v|^{k-m}$ for $m\le k$:
expand its normally convergent homogeneous series and differentiate, or
apply Cauchy on a ball of radius proportional to $|v|$ after the Taylor-tail
bound.  This proves the derivative assertion.  The exact global-conjugation
symmetries on the two-colour plane make the scalar action even separately
in $x,y$, as in V11; the ninth-degree part is zero there.
\end{proof}

Thus the planar polynomial in \cref{thm:v11-plane} now has a genuine common
$O((|x|+|y|)^{10})$ remainder, not just $O_L$ at each regulator.  The
cofinal $J^8$ estimate from V11 remains $O(L^{-2})$.  These two conclusions
\emph{do not} prove convergence of the full finite-amplitude minima: higher
Taylor coefficients are uniformly bounded but have not yet all been shown to
converge.  Nor does \cref{eq:v12-wrem} imply
$L\|w_L^{\rm hom}-v_4\|_\infty=O(|a|^5)$.  The norm is critical energy;
the stronger pointwise scaling is deliberately not used to justify the
nonlinear domain.

\section{A dyadic continuum limit of the actual root observable on \texorpdfstring{$L^4$}{L4} fields}
\label{sec:v12-continuum}

The uniform composite estimate also has a useful nonhomogeneous application.
This section constructs only an observable, without asserting convergence of
an action, a minimum, a Gibbs state, or a physical Hamiltonian.
Let $\mathbb T_N^4=(\mathbb R/N\mathbb Z)^4$, with coarse cells
$Q_r=r+[0,1)^4$.  For a continuum Lie-valued one-form $A\in L^4$ define the
fine cell-average link coordinates
\begin{equation}\label{eq:v12-sampling}
 (\mathsf S_L A)_\mu(x)
  =\frac1L\,L^4\int_{x/L+[0,1/L)^4}A_\mu(y)\,dy.
\end{equation}
The sampling convention, root origin, ordered tree, and dyadic path are fixed.
It is not a Wilson line integral.  It specifies the link field whose literal
root average is evaluated.  Jensen gives
\begin{equation}\label{eq:v12-samplebound}
 \|\mathsf S_L A\|_4\le\|A\|_{L^4},\qquad
 \|\mathsf S_L A\|_8\le L^{-1/2}\|A\|_{L^8}.
\end{equation}
Put $\mathcal O_{L,N}(A)=\mathcal F_{L,N}(\mathsf S_L A)$.

\begin{lemma}[One refinement defect for a smooth input]
\label{lem:v12-consistency}
If $A\in W^{1,4}(\mathbb T_N^4)\cap L^8$ and
$\|A\|_{L^4}<r_*/4$, then
\begin{equation}\label{eq:v12-consistency}
 \|\mathfrak T(\mathsf S_{2L}A)-\mathsf S_L A\|_4
 \le\frac{C'}L\left(\|\nabla A\|_{L^4}+\|A\|_{L^8}^2\right).
\end{equation}
The constant is independent of $L,N$.  Both inputs to the remaining $n$
steps belong to the same ball used in \cref{cor:v12-jets}.
\end{lemma}
\begin{proof}
For each orientation, the two chronological terms in the linear block give
\[
 (B\mathsf S_{2L}A)_\mu(r)
  =\frac1{2L}\left(\operatorname{avg}_{Q_{r,L}}A_\mu+
        \operatorname{avg}_{Q_{r,L}+e_\mu/(2L)}A_\mu\right),
 \quad Q_{r,L}=r/L+[0,1/L)^4.
\]
Its difference from $\mathsf S_LA$ is half the sampled translation
difference.  Jensen, translation invariance of $L^4$, and the fundamental
theorem of calculus give an $\ell^4$ bound
$C L^{-1}\|\nabla A\|_{L^4}$.
The nonlinear term satisfies
\[
 \|E(\mathsf S_{2L}A)\|_4
          \le C_8\|\mathsf S_{2L}A\|_8^2
          \le(C_8/(2L))\|A\|_{L^8}^2.
\]
This proves \cref{eq:v12-consistency}.  The sampling bound and
\cref{thm:v12-root} put both $\mathsf S_LA$ and
$\mathfrak T(\mathsf S_{2L}A)$ in the ball of radius $r_*/2$.
No pointwise smallness of the continuum field is assumed.
\end{proof}

\begin{theorem}[An actual continuum root observable on a small global $L^4$ ball]
\label{thm:v12-continuum}
For each fixed coarse torus and every $\|A\|_{L^4}<r_*/4$, the dyadic limit
\begin{equation}\label{eq:v12-obslimit}
             \mathcal O_{\infty,N}(A)
                  :=\lim_{L=2^n\to\infty}\mathcal O_{L,N}(A)
\end{equation}
exists.  It is locally Lipschitz and holomorphic on the complex $L^4$ ball,
with constants independent of $N$.  Convergence is uniform on compact
subsets of that ball.  For $A\in W^{1,4}\cap L^8$,
\begin{equation}\label{eq:v12-obsrate}
 \boxed{\
 \|\mathcal O_{L,N}(A)-\mathcal O_{\infty,N}(A)\|_4
 \le\frac{C''}L
       \left(\|\nabla A\|_{L^4}+\|A\|_{L^8}^2\right).\ }
\end{equation}
The derivative at zero is the explicit linear continuum average
\begin{equation}\label{eq:v12-Binfty}
 (\mathcal B_{\infty,N}A)_{r,\mu}
             =\int_0^1\int_{Q_r}A_\mu(x+t e_\mu)\,dx\,dt.
\end{equation}
Moreover
\begin{equation}\label{eq:v12-inftyR}
 \|\mathcal O_{\infty,N}(A)-\mathcal B_{\infty,N}A\|_2
                           \le4C\|A\|_{L^4}^2.
\end{equation}
For fixed smooth affine families in the stated ball, all their fixed
parameter derivatives converge at rate $O(L^{-1})$ on smaller parameter
polydiscs.
\end{theorem}
\begin{proof}
The exact composition gives
\[
 \mathcal O_{2L,N}(A)
             =\mathcal F_{L,N}(\mathfrak T(\mathsf S_{2L}A)).
\]
Apply the uniform composite Lipschitz bound to this and
$\mathcal F_{L,N}(\mathsf S_LA)$, then use
\cref{lem:v12-consistency}.  The successive differences are bounded by
$C L^{-1}(\|\nabla A\|_4+\|A\|_8^2)$, whose dyadic tail is at most
twice the first term.  This proves existence and \cref{eq:v12-obsrate} on
the dense smooth class.

Every $\mathcal O_{L,N}$ has one common local Lipschitz bound in $L^4$ by
\cref{eq:v12-samplebound,cor:v12-jets}.  Approximate an arbitrary $L^4$ field
by smooth fields within a smaller ball.  The triangle inequality, with the
two approximation errors controlled independently of $L$, proves the Cauchy
property and a unique Lipschitz extension of the limit.  The same argument
using a finite net gives uniform convergence on compact subsets.  There is
no asserted quantitative rate for all fields in an $L^4$ ball.

The finite maps are holomorphic and uniformly bounded.  On any finite
complex affine polydisc in the ball, compact convergence and Cauchy's
integral formula pass holomorphy and all finite derivatives to the limit.
The local uniform bound yields the usual uniformly convergent homogeneous
Cauchy series on a smaller Banach ball (the multilinear bounds follow by
polarization).  Thus the limit is Banach-holomorphic, not merely
one-dimensional analytic.  For smooth affine directions the norms in
\cref{eq:v12-obsrate} are uniformly bounded on each compact complex
polydisc, and Cauchy gives its differentiated rate.

The direct linear composition is
\[
 (B_L\mathsf S_LA)_{r,\mu}
       =\frac1L\sum_{t=0}^{L-1}\int_{Q_r}A_\mu(x+t e_\mu/L)\,dx.
\]
Translation is strongly continuous on $L^4$, so its Riemann sums converge to
\cref{eq:v12-Binfty}.  Alternatively prove this first on smooth fields
and extend by the norm-one linear bounds.  The finite nonlinear remainder
has vanishing first jet and satisfies \cref{eq:v12-terminal}; passing to
the fixed finite-dimensional output limit proves both the derivative and
\cref{eq:v12-inftyR}.  All norm constants used in this argument were
independent of $N$, although the limit is constructed at each fixed coarse
volume.
\end{proof}

\subsection{A nonconstant exact derivative of the limiting source}

For $N=1$, $\mathcal B_{\infty,1}$ is simply the spatial mean.  If $a$ is a
constant continuum tuple and $w\in L^4(\mathbb T^4)$, V11 and the preceding
limit give
\begin{equation}\label{eq:v12-contmoment}
 \boxed{\
 \bigl(D^2\mathcal O_{\infty,1}(0)[a,w]\bigr)_\mu
   =\sum_{\nu\ne\mu}
       \left[a_\nu,\int_{[0,1]^4}(x_\nu-\tfrac12)w_\mu(x)\,dx\right].\ }
\end{equation}
Indeed $\mathsf S_La=a/L$ and the exact finite centered-coordinate
functional in V11 equals integration against the cell-midpoint step
approximation to $x_\nu-1/2$.  That weight differs from its continuum value
by at most $1/(2L)$.  Holder proves convergence for every $L^4$ $w$,
consistent with differentiated convergence on smooth inputs and continuous
extension.  At constant input $a$, the entire limit is precisely V10's
$\Phi(a)$; the theorem extends that source to nonconstant fields, not to a
new homogeneous normalization.

An exact rational regression uses
\[
 p(x)=x^2(1-x)^2(x-\tfrac12),\qquad
       w_\mu(x)=p(x_\nu)b,\quad \mu\ne\nu.
\]
This periodic field is in $W^{1,4}\cap L^8$, has zero mean, and is divergence
free in its single active component.  Direct integration gives
\begin{equation}\label{eq:v12-polymoment}
 \int_0^1(x-\tfrac12)p(x)\,dx=\frac1{840},\qquad
 m_L=\frac1{840}+\frac1{360L^4}-\frac1{252L^6},
\end{equation}
where $m_L$ is the actual finite midpoint-weighted cell-average moment.
Thus its finite mixed root response is $m_L[a_\nu,b]$ and its limit is
$[a_\nu,b]/840$.  The summation formula is checked by integrating the
polynomial in each fine cell and summing powers of the cell index.  Literal
noncommuting root products provide independent small-lattice regressions.
This uses the inherited mixed Hessian as a check on the new continuum map;
it is not counted as a new derivation of that Hessian.

\section{What is resolved, and what still requires a different estimate}
\label{sec:v12-next}

The actual root observable now has one common holomorphic domain in the fine
global $\ell^4$ norm, including all finite mixed derivatives.  In the
one-coarse-cell problem, the uniform small critical-energy minimum and its
higher remainders are proved for the actual nonlinear root constraint.
This removes the unspecified common-domain obligation at the end of V11
without requiring an additional pointwise scale bound on every detail.
The deterministic continuum source limit further gives a well-defined
inhomogeneous root observable in a specified sampling and dyadic convention.

Several limits must remain distinct.  The source limit does not prove that
$\mathcal G_L(v)$ has a unique finite-amplitude limit.  Bounded analytic
classical branches and convergence through degree eight alone leave higher
coefficients undetermined; subsequential compactness in the fixed source
space does not identify all subsequential limits.  Passing a constrained
variational problem to a limit also requires convergence of the action and
of the relevant derivative response on admissible minimizing fields, not
only pointwise convergence of the source on fixed continuum fields.
The present continuum observable is not a construction of the four-dimensional
interacting measure.

The global domain is genuinely restrictive for statistical mechanics.  A
fixed nonzero constant field over more and more coarse cells has a growing
$L^4$ norm.  Typical Gaussian fluctuations at fixed temperature do not lie
in one common global $\ell^4$ ball as volume tends to infinity.  There is
therefore no transfer here to V7's thermodynamic local-box hypotheses or to
its compact complement.  The balance is an exact statement about gauge
orbits; a source-dependent change of representative still needs its actual
coarea/Jacobian when constructing the measured $J^6Q,J^4B$ packet.

The next useful classical target is convergence of the \emph{constrained
response}, not another homogeneous coefficient: compare the discrete detail
stationary equations with the continuum source functional in a common
energy-dual topology, retaining the V11 nonzero return.  The next useful
quantum target is different: localize the two-norm mechanism to bounded
physical regions with explicit boundary/source terms, so that global
$\ell^4$ smallness is not imposed on all fluctuations.  Neither task is
replaced by the contraction of $B$ on deterministic $\ell^2$ remainders.
The V17 transition-cutoff and high-occupancy coverage problem, noncollapse,
and calibrated physical-time contraction remain outside these conclusions.

\section{Exact checks and append-only record}
\label{sec:v12-verification}

The new verifier reports thirteen passing exact test families, and all
21 V11 exact test families were freshly rerun and passed.  The new verifier
reconstructs the actual linear path incidence, including
periodic multiplicity, and checks the $\ell^2$, $\ell^4$ and $\ell^8$ power
inequalities with exact rational vector data.  It enumerates the primitive
root dependency supports and their overlap bounds, and tests the nonlinear
remainder's vanishing linear jet by literal ordered root multiplication.
Independent positive-kernel nonlinear examples check the contracting-remainder
recursion and the need to distinguish the two counting norms; these examples
are labelled as verification models, not Yang--Mills action data.

The all-length polynomial in \cref{eq:v12-polymoment} is checked by symbolic
integration and summation, then against the literal spatial root Hessian.
The exact finite continuum linear averaging identity is checked independently
on polynomial cell averages.  The V11 verifier is rerun to preserve the
nonzero source-forcing, Green-return and eighth-order constrained-energy
regressions.  The reports state which checks were executed and their elapsed
times.  No finite set of examples is used to infer the holomorphic-domain,
all-volume, or continuum-limit theorems: those conclusions rest on the norm
and convergence proofs above.  No Lean formalization, numerical enclosure of
the common nonlinear radius, or quantum mass-gap certificate is asserted.

The entire V1--V11 mathematical body, all prior labels, bibliographies and end
markers are preserved byte-for-byte.  Only the version display and the front-matter table-of-contents label spacing
are changed in the preamble.  The bundle contains the new appendix, cumulative source and
PDF, executable checks and reports, a portable assembler, and source/output
hashes.  Future V13 work is to append after the following marker.  Any
correction to a preceding proof must be identified explicitly rather than
silently rewriting its statement.

% END OF VERSION 12 MATHEMATICAL BODY.
% Append subsequent requested versions below this marker; retain V1--V12 above.


\clearpage
\section{V13: weak source continuity and the full local classical continuum minimum}
\label{sec:v13-context}

V12 constructs a common analytic root-constrained classical branch and a
continuum root observable on fixed $L^4$ inputs.  Its final boundary is a
variational one: the minimizing fine fields themselves vary with the cutoff,
and the critical Sobolev embedding does not make their $L^4$ norms compact.
Neither convergence through degree eight nor pointwise source convergence on
smooth test fields identifies their finite-amplitude limit.  This version
addresses that boundary, rather than computing another homogeneous coefficient.

The additional estimate is that the \emph{remaining root steps contract
$\ell^2$ perturbations}.  Fine nonlinear contributions can consequently be
removed before a fixed intermediate scale at a uniform small cost.  The
remaining source is a continuous function of finitely many linear averages.
This gives weak $L^4$ continuity, a uniform approximation rate for the source,
and source convergence on varying critical-energy fields.  The exact Wilson
plaquette energy then supplies a lower-limit inequality without assuming that
concentrations are absent.  Exact-source recovery fields and the already
proved constrained convexity identify the entire local classical minimum and
give strong energy convergence of its fields.

\paragraph{Scope and dependencies.}
The literal balanced root stencil and its exact gauge-orbit interpretation are
inherited from V10.  Its two-norm bounds, finite mixed derivatives, and continuum
sampling convention are inherited from V12; the actual one-cell analytic
minimum is \cref{thm:v12-minimum}.  V9 supplies the discrete Coulomb identity,
Sobolev estimate, and full Wilson analytic bounds.  These inputs are not
recertified by the finite checks below.  The new conclusions concern the
\emph{actual root constraint and full Wilson action on their common small
classical branch}, with one final coarse periodic cell for the variational
conclusion.  They concern neither the measured loop determinants nor a Gibbs
or Schwinger functional.  The continuum observable retains its specified
cell-average sampling, ordered tree, root origin, and dyadic convention.

The relevant passages of the supplied f12 and Grand Tensor sources and the
present lineage were searched for the new weak-source and variational
comparison; no corresponding result was located in those passages.  This is
not a universal novelty claim.  Constrained lattice background minima are
established methodology, including \cite{BalabanV9}; no estimate for another
averaging rule is used as a substitute for the arguments below.

\section{Contraction of perturbations and a finite-resolution source approximation}
\label{sec:v13-stop}

All discrete norms still use unnormalized counting measure.  Let
$\mathfrak T_M=B_M+E_M$ be the actual balanced logarithmic root map on an even
side $M$.  Constants in this section are independent of $M$ and the surrounding
volume.  Choose any $0<\alpha<1$.  Radii may be decreased depending on $\alpha$,
but not on a refinement depth.

\begin{lemma}[A contracting derivative on the error space]
\label{lem:v13-l2}
There is $D<\infty$ on a common primitive maximum-norm polydisc such that
\begin{equation}\label{eq:v13-local2}
 \|DE_M(A)\|_{2\to2}\le D\|A\|_\infty.
\end{equation}
There is $r_\alpha>0$ such that the composite of $m$ actual steps, denoted
$\mathcal F_{K,N}$ with $K=2^m$, satisfies
\begin{equation}\label{eq:v13-contract}
 \boxed{\
 \|\mathcal F_{K,N}(A)-\mathcal F_{K,N}(A')\|_2
       \le K^{-\alpha}\|A-A'\|_2
 \quad(\|A\|_4,\|A'\|_4<r_\alpha).\ }
\end{equation}
This also holds for complex fields in the admitted identity germs.
\end{lemma}
\begin{proof}
The local Cauchy and incidence argument of \cref{lem:v12-E} gives
$|DE_M(A)W|\le C_0(K|A|)(K|W|)$, where the positive stencil incidence
matrix has uniformly bounded row and column sums.  Its $\ell^\infty$ and
$\ell^2$ operator norms are therefore bounded, proving
\cref{eq:v13-local2}.  By \cref{lem:v12-B},
$\|D\mathfrak T_M(A)\|_{2\to2}\le1/2+D\|A\|_\infty$.
Choose $r_\alpha$ below half of V12's radius, with
\[
                 \frac12+\frac54Dr_\alpha\le2^{-\alpha}<1.
\]
Every intermediate field on a trajectory starting in this $\ell^4$ ball has
maximum norm at most $5r_\alpha/4$, by \cref{thm:v12-root}.  Thus the chain
rule bounds the $\ell^2$ norm of the composite derivative by
$(2^{-\alpha})^m$.  Apply it along the straight segment between the two
initial inputs.  That segment stays in the $\ell^4$ ball, and V12 controls
each of its intermediate trajectories.  This proves the displayed difference
bound; intermediate straight segments need not be invariant under the map.
\end{proof}

The linear marginal sector still has $\ell^4$ operator norm one.  The lemma
does not contradict this fact: input and output $\ell^2$ norms count different
numbers of fine links.  It is not a contraction of a physical time evolution.

\begin{theorem}[Removal of the finest nonlinear steps]
\label{thm:v13-stop}
Let $L=2^n$, $K=2^m\le L$, and $p=L/K$.  For a fine field on side $NL$ with
$\rho=\|A\|_4<r_\alpha/2$,
\begin{equation}\label{eq:v13-stop}
 \boxed{\
 \|\mathcal F_{L,N}(A)-\mathcal F_{K,N}(B_pA)\|_2
                   \le4C\rho^2 K^{-\alpha}.\ }
\end{equation}
Here $B_p$ is the \emph{actual composed linear average} from side $NL$ to
side $NK$, not a decimation or an independently chosen interpolation.
\end{theorem}
\begin{proof}
Perform the first $n-m$ actual steps.  Their result is
$B_pA+R_{n-m}$, with $\|R_{n-m}\|_2\le4C\rho^2$ by V12.  Both this
result and $B_pA$ lie inside the radius-$r_\alpha$ ball.  The last $m$
steps are exactly $\mathcal F_{K,N}$.  Apply \cref{lem:v13-l2} to these two
intermediate inputs.  No nonlinear interaction has been deleted from the
actual source; the displayed expression is a quantified approximation to it.
\end{proof}

\section{A uniform source limit on the whole small \texorpdfstring{$L^4$}{L4} ball}
\label{sec:v13-source}

Let $\mathbb T_N^4=(\mathbb R/N\mathbb Z)^4$.  Write
$Q_{r,K}=r/K+[0,1/K)^4$, and introduce the finite-rank linear map
\begin{equation}\label{eq:v13-AK}
 (\mathsf A_K A)_{r,\mu}
   =\frac1K\int_0^1
       \operatorname{avg}_{Q_{r,K}+t e_\mu/K}A_\mu(y)\,dt.
\end{equation}
Its output is a link array on side $NK$.  This is a box average followed by
one normalized longitudinal average.  In particular
$\|\mathsf A_K\|_{L^4\to\ell^4}\le1$ by Jensen and translation invariance.
Let $\mathsf S_L$ be exactly V12's fine cell-average sampling, and put
$\mathcal O_{L,N}=\mathcal F_{L,N}\mathsf S_L$.

\begin{lemma}[Operator-norm Riemann comparison without a derivative of the field]
\label{lem:v13-quadrature}
For dyadic $K\le L$ and $p=L/K$,
\begin{equation}\label{eq:v13-quadrature}
 \boxed{\
 \|B_p\mathsf S_L-\mathsf A_K\|_{L^4\to\ell^4}\le\frac4p.\ }
\end{equation}
The constant is independent of $K,L,N$.  No $W^{1,4}$ or $L^8$ hypothesis
on the field is required.
\end{lemma}
\begin{proof}
The literal linear path sum gives
\[
 (B_p\mathsf S_L A)_{r,\mu}
   =\frac1K\frac1p\sum_{t=0}^{p-1}
       \operatorname{avg}_{Q_{r,K}+t e_\mu/(Kp)} A_\mu.
\]
The two kernels multiplying $A_\mu(y)$ are $K^3$ times respectively the
Riemann mean and the integral mean of
$1_{Q_{r,K}+t e_\mu/K}(y)$ for $0\le t\le1$.  For almost every fixed $y$
this indicator has at most two jumps as a function of $t$, including periodic
wraps.  Its Riemann error is at most $2/p$.  The error kernel is supported in
the union swept out by this box, of volume at most $2K^{-4}$.  Its row
absolute integral is at most $4/(Kp)$.  At a fixed $y$ at most two such swept
boxes of a given orientation occur, so the column sum of the absolute kernel
is at most $4K^3/p$.  Short-period coincidences reduce these bounds or give
zero error when the shifted interval is the whole coordinate circle.
For a nonnegative majorant kernel $H$, Holder yields
\[
 \sum_r\left|\int H_r(y)f(y)\,dy\right|^4
 \le \left(\sup_r\int H_r\right)^3
       \left(\sup_y\sum_r H_r(y)\right)\|f\|_4^4.
\]
The stated row and column bounds give $4^4p^{-4}$.  Orientations do not mix,
and Jensen covers the fixed Lie component norm.  This proves
\cref{eq:v13-quadrature}, also for merely $L^4$ inputs.
\end{proof}

Define the finite-resolution nonlinear function
\begin{equation}\label{eq:v13-PK}
                    \mathcal P_{K,N}=\mathcal F_{K,N}\circ\mathsf A_K.
\end{equation}
It depends on finitely many bounded linear functionals of the continuum field.
Its rank here refers to the input averages; it is not an approximation by a
linear output map.

\begin{theorem}[Uniform approximation, weak continuity, and a quantitative rate]
\label{thm:v13-weaksource}
On a common small complex $L^4(\mathbb T_N^4)$ ball, the continuum root
observable of \cref{thm:v12-continuum} satisfies
\begin{equation}\label{eq:v13-finite-rank}
 \boxed{\
 \|\mathcal O_{\infty,N}(A)-\mathcal P_{K,N}(A)\|_4
       \le4C\|A\|_4^2K^{-\alpha}.\ }
\end{equation}
For every fixed $N$ it is weakly sequentially continuous on any smaller
closed ball.  Moreover, for each radius $\rho<r_\alpha/2$,
\begin{equation}\label{eq:v13-rate}
 \boxed{\
 \sup_{\|A\|_{L^4}\le\rho}
  \|\mathcal O_{L,N}(A)-\mathcal O_{\infty,N}(A)\|_4
                  \le C_{\rho,\alpha}L^{-\alpha/(1+\alpha)}.\ }
\end{equation}
The constants in these norm estimates are independent of $N$.  On smaller
balls all fixed Frechet derivatives have the same rate in multilinear
$L^4$ operator norm.  For example, a fixed sufficiently small radius permits
$\alpha=1/2$ and the rate $O(L^{-1/3})$.
\end{theorem}
\begin{proof}
Sampling is contractive from $L^4$ to $\ell^4$.  Combine
\cref{thm:v13-stop,lem:v13-quadrature} with the uniform $\ell^4$ Lipschitz
bound of the remaining composite to obtain
\begin{equation}\label{eq:v13-two-errors}
 \|\mathcal O_{L,N}(A)-\mathcal P_{K,N}(A)\|_4
 \le4C\rho^2K^{-\alpha}+C_1\rho K/L.
\end{equation}
The first error is in $\ell^2$ and hence in $\ell^4$; the second uses the
field-norm Lipschitz bound, not a dimension-growing conversion from $\ell^4$
to $\ell^2$.  First let $L\to\infty$ with $K$ fixed, using the actual
continuum observable already constructed in V12.  This gives
\cref{eq:v13-finite-rank}.

For fixed $K,N$, weak $L^4$ convergence gives norm convergence under the
finite-rank $\mathsf A_K$, and therefore under its continuous composition
$\mathcal F_{K,N}$.  Approximate both the varying field and its weak limit
by \cref{eq:v13-finite-rank}, then let $K\to\infty$.  This proves weak
sequential continuity; weak lower semicontinuity of the $L^4$ norm keeps the
limit in the ball.  The argument does not assume compactness of the critical
Sobolev embedding.

Combining \cref{eq:v13-two-errors,eq:v13-finite-rank} gives a uniform bound
$8C\rho^2K^{-\alpha}+C_1\rho K/L$.  Choose a dyadic $K$ within a factor two
of $L^{1/(1+\alpha)}$.  This proves \cref{eq:v13-rate}.  To differentiate,
use this uniform estimate on a slightly larger complex ball and apply
multivariate Cauchy on $A+\sum z_iW_i$, with a fixed total radius strictly
inside that ball.  This gives the same power of $L$ for each fixed derivative,
with its own radius-dependent constant.  No claim at the limiting exponent
$1/2$ is made on one fixed nonzero ball: the proof uses $\alpha<1$.
\end{proof}

\begin{corollary}[Source passage on varying fine fields]
\label{cor:v13-varying}
For a fine array $a_L$ on side $NL$, define its physical reconstruction by
\begin{equation}\label{eq:v13-reconstruct}
 (\mathcal I_La_L)_\mu(y)=L(a_L)_\mu(x),
                       \qquad y\in Q_{x,L}.
\end{equation}
Suppose $\|a_L\|_4\le\rho<r_\alpha/2$ and
$\mathcal I_La_L\rightharpoonup A$ in continuum $L^4$.  Then
\begin{equation}\label{eq:v13-varying}
                 \boxed{\mathcal F_{L,N}(a_L)\longrightarrow
                                      \mathcal O_{\infty,N}(A).}
\end{equation}
\end{corollary}
\begin{proof}
Exactly $\mathsf S_L\mathcal I_La_L=a_L$, and
$\|\mathcal I_La_L\|_{L^4}=\|a_L\|_4$.  Apply the uniform estimate
\cref{eq:v13-rate} to this varying field and use the weak continuity just
proved.  This is stronger than evaluating a source limit on a fixed smooth
field and interchanging it with an unspecified minimizing sequence.
\end{proof}

\section{The continuum constrained problem in the same local branch}
\label{sec:v13-cont-min}

For the rest of the variational argument take $N=1$.  All limiting statements
are in a fixed one-cell coarse source space, not in growing physical volume.
Let
\[
 V=\{w\in H^1(\mathbb T^4;\mathfrak{su}(2)^4):
                  \operatorname{div}w=0,\ \int w=0\},
                 \qquad \|w\|_V=\|dw\|_{L^2}.
\]
The Fourier Hodge identity and Sobolev--Poincare estimate make this a Hilbert
norm equivalent to the mean-zero $H^1$ norm.  Constants $s$ have the same
12-dimensional normalization as V11--V12.  Define
\begin{equation}\label{eq:v13-YM}
 \mathcal S_\infty(A)
    =\frac12\sum_{\mu<\nu}\int_{\mathbb T^4}
       |\partial_\mu A_\nu-\partial_\nu A_\mu+[A_\mu,A_\nu]|^2.
\end{equation}
This expression uses the quaternion Lie normalization $[a,b]=2a\times b$.
For complex fields, the square has its bilinear holomorphic extension.
The continuum constraint is \emph{the actual limiting source functional}
\begin{equation}\label{eq:v13-cont-constraint}
                         \mathcal O_{\infty,1}(s+w)=v.
\end{equation}
It is not replaced by fixing the mean of $A$.

\begin{theorem}[A unique local continuum stationary branch]
\label{thm:v13-cont-min}
On common small source and detail balls, \cref{eq:v13-cont-constraint} has a
holomorphic chart $s=s_\infty(v,w)=v+O((|v|+\|w\|_V)^2)$.
The chart action has a unique stationary point $w_\infty(v)$ in its small
detail ball, and this point is a strict local minimum for real data.  Put
\[
 A_\infty^*(v)=s_\infty(v,w_\infty(v))+w_\infty(v),\qquad
 \mathcal G_\infty(v)=\mathcal S_\infty(A_\infty^*(v)).
\]
The functions are holomorphic on common smaller complex source balls.  The
constrained Hessian is between $\frac12 I$ and $\frac32 I$ in the detail
energy norm, and $\|w_\infty(v)\|_V\le C|v|^2$.
\end{theorem}
\begin{proof}
The linear continuum source on one cell is the mean.  By V12 and
\cref{thm:v13-weaksource}, the nonlinear remainder and its first derivative
are respectively $O(\|A\|_4^2)$ and $O(\|A\|_4)$, uniformly on a complex
ball.  Sobolev gives $\|w\|_4\le C_S\|w\|_V$.  The equation
$s=v-(\mathcal O_{\infty,1}(s+w)-s)$ is therefore a contraction on fixed
small balls and gives the holomorphic chart, just as in V12.

Write $\mathcal S_\infty(A)=\frac12\|dA\|_{\rm bil}^2+\mathcal N_\infty(A)$.
Holder with $2,4,4$ bounds the cubic term, and four $L^4$ factors bound the
quartic term.  Thus $D\mathcal N_\infty=O(\|A\|_{\mathcal E}^2)$ and
$D^2\mathcal N_\infty=O(\|A\|_{\mathcal E})$ in the continuum critical
norm $\mathcal E=\|dA\|_2+\|A\|_4$.  In the chart, $d(s+w)=dw$; the
noncommuting constant part remains in $\mathcal N_\infty$.  The detail
stationary equation is an identity Riesz inverse applied to the small
nonlinear gradient.  The estimates, including the second chart derivative,
give a contraction with constant at most $1/2$ on a fixed small ball and a
solution of size $C|v|^2$.  The Hessian bounds and uniqueness follow on that
ball.  The complex contraction and Cauchy estimates give holomorphy and
uniform fixed source derivatives.  This is a continuum local construction;
no global minimizer or compact gauge-sector classification is inferred.
\end{proof}

\section{Weak compactness and the exact Wilson lower-limit inequality}
\label{sec:v13-liminf}

Write $h_L=L^{-1}$ only for mesh size in this section; it is not the inverse
coupling used in V8.  For a fine array $a_L$ let $A_L=\mathcal I_La_L$ and
let its physical forward difference, reconstructed on cells, be
\[
 D^L_\nu A_{L,\mu}(y)
      =L^2\bigl(a_{L,\mu}(x+e_\nu)-a_{L,\mu}(x)\bigr).
\]
Then exactly
\begin{equation}\label{eq:v13-scale-norms}
 \|A_L\|_{L^4}=\|a_L\|_4,\qquad
 \sum_{\mu,\nu}\|D^L_\nu A_{L,\mu}\|_2^2
    =\sum_{\mu,\nu}\|\nabla_\nu a_{L,\mu}\|_2^2.
\end{equation}

\begin{lemma}[Discrete critical compactness with the correct weak topology]
\label{lem:v13-compact}
Suppose $d^*a_L=0$, the means $L\langle a_L\rangle_L$ are bounded, and
$\|da_L\|_2$ is bounded.  A subsequence has
\[
 A_L\to A\ \hbox{in }L^p\ (1\le p<4),\qquad
 A_L\rightharpoonup A\ \hbox{in }L^4,\qquad
 D^L A_L\rightharpoonup\nabla A\ \hbox{in }L^2.
\]
The limit is an $H^1$ Coulomb field, and
$\|dA\|_2\le\liminf\|da_L\|_2$.  Strong $L^4$ convergence is not asserted.
\end{lemma}
\begin{proof}
The discrete Hodge identity turns the curl bound into the bound for every
forward difference in \cref{eq:v13-scale-norms}.  Periodically interpolate
the physical vertex values $L a_L$ by multilinear functions on mesh cubes.
The unit-cube estimates used in V9, after scaling, give a bounded $H^1$
interpolant and an $L^2$ difference from $A_L$ at most $C L^{-1}\|\nabla a_L\|_2$.
Its mean is the displayed physical mean.  Compactness in $L^2$ can be proved
on the fixed torus by Fourier truncation: the $H^1$ bound makes the squared
$L^2$ tail above frequency $R$ at most $C/R^2$, while the retained Fourier
space is finite dimensional.  Thus a subsequence converges strongly in
$L^2$.  The uniform $L^4$ bound from discrete Sobolev, and interpolation,
give strong $L^p$ for $p<4$ and weak $L^4$.

The reconstructed forward differences are bounded in $L^2$.
Discrete summation by parts against cell averages of a smooth test function
identifies their distributional limits with $\nabla A$.  The same argument
with the exact discrete divergence shows $\operatorname{div}A=0$.
Weak lower semicontinuity of the curl norm proves the last inequality.
No endpoint Sobolev compactness has been used.
\end{proof}

For each oriented plaquette let $U_{L,\mu\nu}(x)$ be the literal ordered
product of $e^{a_L}$ in the Wilson action, and set
\begin{equation}\label{eq:v13-Pcurv}
 \mathcal C_{L,\mu\nu}(y)=L^2(U_{L,\mu\nu}(x)-I),\qquad y\in Q_{x,L}.
\end{equation}
The norm on a quaternion includes its real and three imaginary components.
For a real unit quaternion $U$, $|U-I|^2=2(1-\Re U)$, so
\begin{equation}\label{eq:v13-exact-energy}
                 \boxed{\mathcal S(a_L)=\tfrac12\|\mathcal C_L\|_2^2.}
\end{equation}
This exact square is used instead of declaring a Taylor action uniformly
consistent on arbitrary concentrating sequences.

\begin{theorem}[The full Wilson action has the required weak lower limit]
\label{thm:v13-liminf}
In addition to the compactness hypotheses, suppose $\|a_L\|_4$ stays in the
common small ball.  For every subsequential limit from
\cref{lem:v13-compact},
\begin{equation}\label{eq:v13-liminf}
 \boxed{\mathcal S_\infty(A)\le\liminf_{L\to\infty}\mathcal S(a_L).}
\end{equation}
If the energies are bounded, $\mathcal C_L$ converges weakly in $L^2$ to the
purely imaginary continuum curvature $dA+[A,A]$ along that subsequence.
\end{theorem}
\begin{proof}
For a finite lower limit choose the subsequence with bounded energies.
Equation \cref{eq:v13-exact-energy} bounds $\mathcal C_L$ in $L^2$.
Expand one ordered plaquette through degree two in its four oriented Lie
entries $X_0,\ldots,X_3$.  The linear term is the discrete curl.  Its
quadratic real part is $-\frac12|\sum_i X_i|^2$, and its imaginary part is
$\sum_{i<j}x_i\times x_j$, where $x_i$ are the quaternion Lie vectors.
If both shifted components are replaced by their values at the initial
vertex, this imaginary part is $[a_{L,\mu},a_{L,\nu}]$.
The differences consist of products of one field and one lattice difference.
All cubic and higher terms are bounded pointwise by
$C(\sum_i|X_i|)^3$ on the common small maximum-norm ball, which follows
from $\|a_L\|_\infty\le\|a_L\|_4$.

After physical reconstruction, the cubic error has $L^1$ norm at most
$C L^{-1}\|A_L\|_3^3$; bounded local incidence changes only $C$.
The quadratic shift error has $L^1$ norm at most
$C L^{-1}\|A_L\|_2\|D^L A_L\|_2$.
The real quadratic part has $L^1$ norm at most
$C L^{-2}\|da_L\|_2^2$.  All tend to zero.
Consequently, in distributions,
\[
 \mathcal C_{L,\mu\nu}
 =D^L_\mu A_{L,\nu}-D^L_\nu A_{L,\mu}
                       +[A_{L,\mu},A_{L,\nu}]+o(1).
\]
The derivative terms have their weak $L^2$ limits.  Strong $L^2$ convergence
of $A_L$ gives strong $L^1$ convergence of the bracket; its $L^2$ norm is
bounded by the $L^4$ stock, so the same bracket converges weakly in $L^2$.
The distributional limit of every bounded $L^2$ subsequence of
$\mathcal C_L$ is therefore $dA+[A,A]$, with zero real component.  This
identifies its weak $L^2$ limit.  Applying weak lower semicontinuity to the
exact square proves \cref{eq:v13-liminf}.  Concentration may contribute
additional nonnegative energy; the proof has not assumed it vanishes.
\end{proof}

\section{Recovery with the exact nonlinear source}
\label{sec:v13-recovery}

\begin{lemma}[Strong recovery on the actual discrete root fibres]
\label{lem:v13-recovery}
Let $A=s_\infty(v,w)+w$ lie strictly inside the small continuum constraint
chart, with $w\in V$.  There are real discrete Coulomb fields $\widetilde a_L$
in the common V12 constraint chart such that
\begin{equation}\label{eq:v13-recovery}
 \boxed{\
 \mathcal F_{L,1}(\widetilde a_L)=v,\quad
 \mathcal I_L\widetilde a_L\to A\text{ in }L^4,\quad
 D^L\mathcal I_L\widetilde a_L\to\nabla A\text{ in }L^2,\quad
 \mathcal S(\widetilde a_L)\to\mathcal S_\infty(A).\ }
\end{equation}
Their detail-energy norms remain in a fixed smaller admissible ball whenever
$w$ is in that ball.  No uniform $L^4$ bound for a moving Fourier projector
is needed in this construction.
\end{lemma}
\begin{proof}
First let $w$ be a real trigonometric polynomial with mean zero and continuum
divergence zero.  For each of its finitely many nonzero Fourier modes
$k\in2\pi\mathbb Z^4$, use
\[
 q_L(k)_\mu=e^{ik_\mu/L}-1,\qquad
 \Pi_L(k)=I-q_L(k)q_L(k)^*/|q_L(k)|^2.
\]
For large $L$ these modes are distinct modulo fine aliases.  Define its fine
coefficients as $L^{-1}\Pi_L(k)\widehat w(k)$.  This is exactly discrete
Coulomb and mean zero.  Conjugate modes retain reality, and
$\Pi_L(k)\widehat w(k)\to\widehat w(k)$ because $k\cdot\widehat w(k)=0$.
With constant part $s/L$, its reconstructions and physical differences
converge uniformly to $s+w$ and its derivatives.  Smooth plaquette expansion
gives convergence of the full Wilson energy.  Its source tends to
$\mathcal O_{\infty,1}(s+w)=v$ by \cref{cor:v13-varying}.

Keep this fine detail fixed and solve only for its twelve constant
coordinates using V12's chart.  The constant-coordinate derivative of the
source is uniformly close to the identity.  The required correction to $s$
therefore tends to zero, bounded by a fixed multiple of the source defect.
It preserves discrete Coulomb and the detail energy.  The uniform critical
Wilson first-derivative bound makes its energy change tend to zero.  Thus
all displayed conclusions hold with the \emph{exact} source, not an
approximately constrained recovery.

For a general $w\in V$, approximate it strongly in $H^1$ by real Coulomb
trigonometric polynomials $w^{(j)}$.  Sobolev gives strong $L^4$ convergence.
Set $s^{(j)}=s_\infty(v,w^{(j)})$ so these continuum approximants have the
exact source.  Chart continuity gives $s^{(j)}\to s$ and continuity of
\cref{eq:v13-YM} in $H^1$ gives convergence of their energies.  For each
fixed $j$ use the preceding discrete construction, then choose a diagonal
$j=j(L)\to\infty$ slowly enough to obtain all conclusions simultaneously.
The strict interior margin leaves every chosen detail and its corrected
source in the common local chart.  This argument never applies the Fourier
projection to an arbitrary varying $L^4$ field with an unproved operator
bound.
\end{proof}

\section{The actual finite-amplitude classical minima converge}
\label{sec:v13-main}

Let $A_L^*(v),\mathcal G_L(v)$ be precisely the V12 branch, and let
$A_\infty^*(v),\mathcal G_\infty(v)$ be the branch of
\cref{thm:v13-cont-min}.  Choose one common smaller source ball so all these
fields and the recovery constructions lie strictly inside the domains above.

\begin{theorem}[Full local classical continuum limit with strong constrained response]
\label{thm:v13-main}
For every real source $v$ in that common ball,
\begin{equation}\label{eq:v13-main}
 \boxed{\
 \mathcal G_L(v)\longrightarrow\mathcal G_\infty(v),\qquad
 \mathcal I_L A_L^*(v)\longrightarrow A_\infty^*(v)\text{ in }L^4,\qquad
 D^L\mathcal I_L A_L^*(v)\longrightarrow\nabla A_\infty^*(v)
                                                  \text{ in }L^2.\ }
\end{equation}
The scalar functions converge locally uniformly on a common smaller complex
source ball, and every fixed source derivative converges there.  In particular
\emph{all} their classical Taylor coefficients converge, not just those through
degree eight.  No rate for \cref{eq:v13-main} is asserted by this proof.
\end{theorem}
\begin{proof}
V12 bounds the means and detail energies of its minimizing fields uniformly,
with $\|w_L(v)\|_V\le K|v|^2$ and small global $\ell^4$ norm.
\Cref{lem:v13-compact} supplies a subsequential limit $A=s+w$ with
$w\in V$ inside the continuum small ball.  The source passage
\cref{cor:v13-varying} gives $\mathcal O_{\infty,1}(A)=v$, so uniqueness
of the constant-coordinate chart identifies $s=s_\infty(v,w)$.
The lower-limit theorem yields
\[
 \mathcal S_\infty(A)\le\liminf\mathcal G_L(v).
\]
For the continuum minimum, \cref{lem:v13-recovery} supplies exactly constrained
competitors in the finite detail ball.  The finite chart is strictly convex,
so its stationary point has no larger action than each competitor.  Therefore
\[
 \limsup\mathcal G_L(v)\le\mathcal S_\infty(A_\infty^*(v)).
\]
The continuum constrained action is also strictly convex on its detail ball;
its unique minimum is $A_\infty^*(v)$.  These inequalities force
$A=A_\infty^*(v)$ and convergence of the values.  They apply to every
subsequence, giving convergence of the entire sequence in the weak topologies.

To strengthen it, take a recovery sequence $\widetilde a_L$ for
$A_\infty^*(v)$ and let $\widetilde w_L,w_L$ be its and the minimizer's
fine details.  They have the \emph{same exact source}.  Taylor's integral
formula along their segment in the fixed detail ball and V12's Hessian lower
bound imply
\begin{equation}\label{eq:v13-strong-defect}
 \boxed{\
 \frac14\|d(\widetilde w_L-w_L)\|_2^2
 \le\mathcal S(\widetilde a_L)-\mathcal G_L(v)\longrightarrow0.\ }
\end{equation}
Discrete Sobolev then gives $\|\widetilde w_L-w_L\|_4\to0$.  The uniform
constant-coordinate chart Lipschitz bound gives convergence of their mean
coordinates as well.  Their physical reconstructions therefore differ by
$o(1)$ in $L^4$; the discrete Hodge identity gives $o(1)$ for all physical
first differences.  Combine with the strong recovery to prove
\cref{eq:v13-main}.  This is where possible concentration is excluded for
minimizers.  Source convergence alone did not exclude it.

Finally the scalar $\mathcal G_L$ are holomorphic and uniformly bounded on a
common complex source ball by V12.  On every smaller complex polydisc,
Cauchy bounds and finite-dimensional Arzela--Ascoli give subsequential compact
convergence.  Every such limit agrees with $\mathcal G_\infty$ on the real
open source ball, by the real convergence just proved.  The holomorphic
identity theorem, applied successively in the twelve coordinates, identifies
it with the continuum branch's holomorphic extension.  Thus the entire
sequence converges locally uniformly.  Cauchy's formula gives convergence of
every fixed source derivative.  Qualitative weak compactness and diagonal
recovery give no uniform mesh rate for the full minimum; the source rate alone
must not be advertised as that rate.
\end{proof}

In particular the classical source multiplier, with the sign convention
$D_v\mathcal G=-\eta$, also converges locally, as do the finite-dimensional
source Hessians and all higher source tensors.  The theorem does not assert
norm-resolvent convergence of infinite-dimensional fluctuation operators or
convergence of their determinants.
The exact reconstructed plaquette curvatures also converge strongly in $L^2$:
the lower-limit proof gives weak convergence, and the convergence of the
minimum values gives convergence of their squared norms in
\cref{eq:v13-exact-energy}.  This is a field-strength statement, not a
fluctuation-determinant estimate.

\subsection{The nonzero V11 return remains in the limiting functional}

By the scalar derivative convergence and the inherited V11 coefficient,
\begin{equation}\label{eq:v13-keep-return}
 J^8\mathcal G_\infty(v)
   =J^8\left[E_4(\Phi^{-1}(v))-\frac{\mathscr P_8(v)}{1440}\right].
\end{equation}
This is now the Taylor polynomial of an identified finite-amplitude
continuum constrained minimum, not merely a candidate coefficient limit.
It can also be checked directly at the continuum stationary equation.
V12's continuum mixed root Hessian has weight $x_\nu-1/2$, so its degree-four
force is
\[
 f_{a,\mu}(x)=\sum_{\nu\ne\mu}(x_\nu-1/2)[g_\mu(a),a_\nu].
\]
It is mean zero and divergence free.  The periodic function
$q(r)=r(1-4r^2)/24$, $-1/2\le r\le1/2$, has matching value and first
derivative at the endpoints and satisfies $-q''=r$ weakly.  Hence
\[
 \int_{-1/2}^{1/2}r q(r)\,dr=\frac1{720},\qquad
 \frac12\langle f_a,(-\Delta)^{-1}f_a\rangle
                                  =\frac{\mathscr P_8(a)}{1440}.
\]
The continuum variational expansion thus retains exactly the same decrease.
These values are inherited regressions, not newly counted eighth-order
coefficients.  In particular the full limit is not identified with the
homogeneous trial $E_4\circ\Phi^{-1}$.

\section{A concentrating Coulomb loop separates source convergence from energy convergence}
\label{sec:v13-loop}

Here is an exact diagnostic of the compactness issue addressed in the proof.
For $L\ge4$, fix a sufficiently small real $a$, one Lie generator $E$, and
place a signed loop on the four edges of a single elementary $(1,2)$ square:
\[
 a_{L,1}(0)=aE,\quad a_{L,2}(e_1)=aE,\quad
 a_{L,1}(e_2)=-aE,\quad a_{L,2}(0)=-aE,
\]
with all other links zero.  It is discrete Coulomb and has zero mean in each
orientation.  All its Lie entries commute.  V10's exact abelian reduction
therefore gives, on the justified small branches,
\begin{equation}\label{eq:v13-loop-source}
                         \mathcal F_{L,1}(a_L)=B_La_L=0.
\end{equation}
The scalar curl equals $4a$ on the central plaquette and has magnitude $a$
on twenty other plaquettes.  Consequently
\begin{equation}\label{eq:v13-loop-energy}
 \boxed{\begin{aligned}
 \|a_L\|_4&=\sqrt2|a|,\qquad \|da_L\|_2^2=36a^2,\\
 \mathcal S(a_L)&=1-\cos(4a)+20(1-\cos a)>0
                         \quad(a\ne0\text{ small}).
 \end{aligned}}
\end{equation}
These counts retain the full four-dimensional plaquette incidence; the
sixteen plaquettes in the other transverse planes are part of the twenty.
The physical reconstruction is supported on four cells of total volume
$4L^{-4}$, has constant nonzero $L^4$ norm, and tends weakly to zero in $L^4$.
Indeed pairing with an $L^{4/3}$ test function is bounded by that norm times
the test function's $L^{4/3}$ norm on the shrinking support.  Its $L^2$ norm
is $2|a|/L$.

Thus exact source convergence and weak critical convergence permit a positive
energy defect.  The lower-limit inequality allows it, while
\cref{eq:v13-strong-defect} excludes it for minimizing sequences.  This example
is not a hidden-diffusion gap test, a continuum instanton, or a competitor
contradicting the zero-source local minimum: at that source the zero field
has strictly smaller energy.

\section{Verification, changed boundary, and append-only record}
\label{sec:v13-next}

The new proof replaces the final common-space assumption of V12 by three
explicit steps: finite-resolution approximation of the actual source,
weak lower-limit control from the exact Wilson square, and recovery with the
\emph{same nonlinear source}.  It then uses constrained convexity to upgrade
weak convergence to strong energy convergence.  The one-cell classical
continuum functional is thereby identified at finite amplitude, with all
fixed source derivatives, rather than only through $J^8$.

The next quantum target is not supplied by this result.  A convergent classical
background does not control the renormalized trace of its fluctuation operator,
the measured $J^6Q,J^4B$ packet, or the analytic tails required for a uniform
integral comparison.  The new finite-resolution estimate is a natural starting
point for comparing \emph{source-dependent fluctuation forms} around these
backgrounds, but their loop traces and ultraviolet subtractions have to be
estimated separately.  Localizing a global $L^4$ ball to thermodynamic field
regions, with the boundary and coarea terms retained, also remains necessary.

The coarse volume in the variational theorem is one fixed periodic cell.
There is no claim here of uniform multi-cell nonlinear gauge fixing, arbitrary
rough retained-source control, compact-complement suppression, transition-cutoff
coverage, a quantum state, or a physical Hamiltonian mass gap.  The continuum
source and its local constrained action are in the declared deterministic
dyadic and sampling convention.  A comparison with a different microscopic
sampling or renormalization route is a separate theorem, not part of uniqueness
on this one prescribed route.

\paragraph{Executed checks.}
The new driver \path{verify_ym_classical_continuum_V13.py} reports 22
passing exact test families.  It checks the contraction and truncation arithmetic,
the exact box-translation kernel comparison, the plaquette second-order
expansion and its continuum normalization, the discrete Coulomb Fourier
recovery, the concentrating-loop incidence and literal root regression,
and the continuum Green regression.  Exact rational and symbolic checks are
reported as such.  They validate finite identities and normalization, not
weak compactness or a continuum theorem by sampling.  All 13 V12 exact test families were freshly rerun and passed; that new
report is distinguished from prior archival reports.
The written proofs, with the stated inherited hypotheses, establish the
analytic conclusions.  No Lean formalization or numerical enclosure of the
common nonlinear radius is asserted.

\paragraph{Append-only record.}
The complete V1--V12 mathematical body, including all labels, bibliographies
and end markers, is retained byte-for-byte.  Only preamble version displays
change.  The bundle records the original and preserved-body hashes, new
appendix, code, verification reports, build checks, and final output hashes.
A subsequent V14 is to append below the next marker.  Any correction to an
earlier result must be stated in an appended correction rather than silently
rewriting that result.

% END OF VERSION 13 MATHEMATICAL BODY.
% Append subsequent requested versions below this marker; retain V1--V13 above.


\clearpage
\section{V14: the constrained fluctuation form and a trace-class determinant remainder}
\label{sec:v14-context}

V13 identifies the small one-cell classical continuum minimum.  Its final
paragraph correctly leaves the fluctuation operator and its determinant
uncontrolled.  Convergence of the minimized scalar action, even with all source
derivatives, is not a trace-ideal estimate on the infinitely many detail
fluctuations.  The present continuation treats that next object.  It keeps the
second derivative of the \emph{nonlinear root constraint} and its classical
multiplier, rather than using the unconstrained Wilson Hessian as a replacement.

The main new observation is quantitative.  The finite-resolution approximation
of V13 factors through at most $12K^4$ linear averages on one coarse cell.  After
two field derivatives it is therefore a finite-rank approximation to the
energy-normalized source Hessian.  Its error is $O(K^{-\alpha})$ for any fixed
$0<\alpha<1$ on a sufficiently small common ball.  This gives singular values
$O(n^{-\alpha/4})$.  Taking $\alpha>4/5$ puts the actual root-Hessian term in
$\mathfrak S_5$, including all its fixed source derivatives.  This is not a
consequence merely of the source having twelve output coordinates: a nonlinear
scalar-valued function can have an infinite-rank Hessian.

Combining that estimate with a stated Cwikel multiplier theorem gives a fifth
modified determinant for the \emph{whole continuum constrained fluctuation
form} on the V13 branch.  The first four Born traces and the finite normal
coarea factor are displayed separately.  No claim is made that these Born
traces have already been matched to local Yang--Mills counterterms, or that the
actual lattice Wilson Hessians converge in $\mathfrak S_5$.  The latter comparison
is proved here for the actual \emph{source-Hessian pullbacks}, not for every
term of the lattice fluctuation form.

\paragraph{Sources and nonduplication.}
The common source ball, finite-resolution root approximation, and continuum
background are inherited from \cref{thm:v13-weaksource,thm:v13-main}.
The measured constrained-Hessian sign is already present in V4.  The supplied
SM--Einstein monograph already explains fifth modified determinants and the
need to retain the first four Born contributions \cite[Section ``Relative
determinants and rough-field control'']{SM6}; the Grand Tensor has the
corresponding three-dimensional reservoir example
\cite[\texttt{thm:SMST-reservoir-det4}]{GT76V6}.  Those general devices are not
relabelled as new.  What is acquired is the trace ideal of the \emph{actual
nonlocal root curvature}, its dyadic transfer estimate, and the construction
of the complete constrained continuum determinant remainder on this branch.
An external analytic input is Frank's operator-valued Cwikel theorem
\cite[Theorem 3.1]{FrankV14}.  Its hypotheses and its specialization to the unit
four-torus are stated below; this is the only new imported analytic theorem.

\section{Finite-resolution root estimates imply Schatten bounds}
\label{sec:v14-source}

All continuum spaces in this version are on the fixed unit four-torus.  Put
$X=L^4(\T^4;\mathfrak{su}(2)^4)$ and use its complexification for holomorphic
statements.  Let
\[
 V=\{w\in H^1:\operatorname{div}w=0,\ \textstyle\int w=0\},\qquad
 \|w\|_V=\|dw\|_2,
 \qquad \widehat V=\R^{12}\oplus V.
\]
The fixed Hilbert norm on $\widehat V$ is $(|s|^2+\|w\|_V^2)^{1/2}$.
The inclusion $j:\widehat V\to X$, $j(s,w)=s+w$, is bounded by Sobolev.
Let $\mathcal O=\mathcal O_{\infty,1}$ and
$\mathcal O_L=\mathcal F_{L,1}\mathsf S_L$ be the actual continuum source and
its dyadic sampled versions.  For $\eta\in\C^{12}$ define the bilinear form
\begin{equation}\label{eq:v14-rootform}
 t_{A,\eta}(u,z)=\langle\eta,D^2\mathcal O(A)[ju,jz]\rangle,
 \qquad u,z\in\widehat V,
\end{equation}
and denote its operator in the fixed Hilbert coordinates by $\mathsf T(A,\eta)$.
For real parameters it is self-adjoint.  For complex parameters the transpose,
not a parameter-dependent Hermitian conjugation, defines its holomorphic
extension.  Write $\mathsf T_L(A,\eta)$ for the same expression with
$\mathcal O_L$.  These are source curvatures, not covariances.

\begin{lemma}[Rank versus approximation rate]
\label{lem:v14-rank}
Suppose a bounded Hilbert-space operator $T$ has, for dyadic $K\ge1$,
approximants $T_K$ with
\[
 \operatorname{rank}T_K\le cK^d,\qquad
 \|T-T_K\|\le MK^{-\alpha},\qquad \|T\|\le M.
\]
Then $T$ is compact and $s_n(T)\le C M n^{-\alpha/d}$, with $C$ depending only
on $c,d,\alpha$.  For $p>d/\alpha$,
\begin{equation}\label{eq:v14-ranktail}
 \|T\|_{\mathfrak S_p}\le C_pM,\qquad
 \|T-T_K\|_{\mathfrak S_p}\le C_pM K^{d/p-\alpha}.
\end{equation}
If $T,T'$ have a common singular-value bound $M n^{-\alpha/d}$ and
$\|T-T'\|\le\delta$, then
\begin{equation}\label{eq:v14-interpolation}
 \|T-T'\|_{\mathfrak S_p}
 \le C_p M^{d/(\alpha p)}\delta^{1-d/(\alpha p)}
\end{equation}
for $0\le\delta\le M$, after increasing the common constant if necessary.
\end{lemma}
\begin{proof}
Norm approximation by finite-rank operators proves compactness.  The
min--max characterization of singular values gives
$s_{\lfloor cK^d\rfloor+1}(T)\le MK^{-\alpha}$.  Choosing dyadic $K$ just
below the corresponding power of $n$ proves the singular-value bound, with
the first finitely many indices covered by $\|T\|\le M$.
For the tail estimate telescope $T_K,T_{2K},T_{4K},\ldots$.
Each difference has rank at most $c(1+2^d)K^d$ and norm at most
$M(1+2^{-\alpha})K^{-\alpha}$.  For a rank-$r$ matrix,
$\|D\|_p\le r^{1/p}\|D\|$.  The resulting series is geometric with ratio
$2^{d/p-\alpha}<1$ and converges to $T$ in operator norm, proving
\cref{eq:v14-ranktail}.  The singular-value inequality
$s_{2n-1}(T-T')\le s_n(T)+s_n(T')$, together with the norm bound, gives
$s_n(T-T')\le\min\{\delta,C M n^{-\alpha/d}\}$.
Split its $p$th-power sum at $n=(CM/\delta)^{d/\alpha}$ and compare the
remaining power tail with its integral.  This proves
\cref{eq:v14-interpolation}; $\delta=0$ is immediate.
\end{proof}

\begin{theorem}[Trace ideal of the actual nonlinear source curvature]
\label{thm:v14-rootideal}
Fix $0<\alpha<1$ and a smaller common complex source ball supplied by V13.
On this ball,
\begin{equation}\label{eq:v14-rootsingular}
 s_n(\mathsf T(A,\eta)),\ s_n(\mathsf T_L(A,\eta))
            \le C_\alpha|\eta|n^{-\alpha/4},
\end{equation}
uniformly in dyadic $L$.  For every $p>4/\alpha$ these are
$\mathfrak S_p$ operators.  They are holomorphic in $A,\eta$ in that norm,
and all fixed derivatives in $A$ have the corresponding local bounds.
In particular $\alpha=9/10$ gives $\mathfrak S_5$ on one fixed sufficiently
small ball.

Let $\mathcal P_K=\mathcal F_{K,1}\mathsf A_K$ be V13's finite-resolution
source and let $\mathsf T^{[K]}$ use $D^2\mathcal P_K$ in
\cref{eq:v14-rootform}.  Then
\begin{equation}\label{eq:v14-roottail}
 \operatorname{rank}\mathsf T^{[K]}\le12K^4,\qquad
 \|\mathsf T-\mathsf T^{[K]}\|_{\mathfrak S_p}
                   \le C_{p,\alpha}|\eta|K^{4/p-\alpha}.
\end{equation}
The same estimates hold with any fixed number of background derivatives and
their direction norms included.
\end{theorem}
\begin{proof}
V13 proves $\|\mathcal O-\mathcal P_K\|=O(K^{-\alpha})$ uniformly on a
complex $X$ ball.  Cauchy differentiation on a larger ball than the one used
in the statement gives the same error for $D^{r+2}$ in multilinear $X$
operator norm, for every fixed $r$.  Composing the two free slots with the
bounded $j$ gives operator-norm error $C_r|\eta|K^{-\alpha}$ on
$\widehat V$.

Since $\mathsf A_K$ is linear and has $12K^4$ scalar output coordinates,
\[
 D^2(\mathcal F_{K,1}\mathsf A_K)(A)[ju,jz]
 =D^2\mathcal F_{K,1}(\mathsf A_KA)[\mathsf A_Kju,\mathsf A_Kjz].
\]
The associated operator factors through that finite-dimensional space, so its
rank is at most $12K^4$.  Extra background derivatives do not remove either
free factor $\mathsf A_Kj$, and obey the same rank bound.  Apply
\cref{lem:v14-rank} with $d=4$.

For finite $L$ and $K\le L$ use instead
$\mathcal F_{K,1}B_{L/K}\mathsf S_L$.  The stopping estimate
\cref{thm:v13-stop}, followed by Cauchy, gives exactly the same rank and
error bound.  For $K\ge L$ the operator is already of rank at most $12L^4$
and can be used as its own approximant.  Thus the constants are uniform in
$L$ as claimed.  The dyadic telescoping series in the proof of
\cref{lem:v14-rank} converges locally uniformly in $\mathfrak S_p$ and each
term is holomorphic.  Cauchy's formula in this Banach ideal proves
holomorphy and the differentiated assertions.  No critical Sobolev compactness
has been assumed; the finite-resolution structure of the source provides the
compactness.
\end{proof}

\begin{corollary}[Dyadic source-Hessian transfer, including varying backgrounds]
\label{cor:v14-roottransfer}
On smaller balls and bounded multiplier sets, put
$\gamma=\alpha/(1+\alpha)$.  Then
\begin{equation}\label{eq:v14-rootrate}
 \|\mathsf T_L(A,\eta)-\mathsf T(A,\eta)\|_{\mathfrak S_p}
       \le C L^{-(\alpha-4/p)/(1+\alpha)},\qquad p>4/\alpha.
\end{equation}
At $\alpha=9/10,p=5$ the exponent is $1/19$, while the finite-resolution
$K$-tail in \cref{eq:v14-roottail} is $K^{-1/10}$.
More generally, if $A_L\to A$ strongly in $L^4$ and $\eta_L\to\eta$, then
\begin{equation}\label{eq:v14-movingroot}
 \|\mathsf T_L(A_L,\eta_L)-\mathsf T(A,\eta)\|_{\mathfrak S_p}
 \le C\bigl(L^{-\gamma}+\|A_L-A\|_4+|\eta_L-\eta|\bigr)^{1-4/(\alpha p)}
\end{equation}
when the expression in parentheses is small.  In particular this convergence
holds along V13's reconstructed real minimizing fields and their multipliers.
\end{corollary}
\begin{proof}
The uniform differentiated source rate in \cref{thm:v13-weaksource} gives
operator-norm error $O(L^{-\gamma})$.  The same Cauchy bounds give Lipschitz
control in the base field and multiplier.  Apply
\cref{eq:v14-interpolation} to the uniform singular-value estimates.  V13 gives
strong $L^4$ convergence of the reconstructed fields and convergence of the
classical source multipliers.  Their higher rates are not asserted.
\end{proof}

The finite operator in this corollary is the \emph{cell-average pullback of the
literal lattice source Hessian}: its arguments are $\mathsf S_Lju$ and
$\mathsf S_Ljz$.  It includes the nonlinear root curvature and has a fixed
common carrier.  Cell-average sampling does not by itself identify the full
lattice Coulomb energy basis.  The corollary is therefore not already a
comparison of the entire lattice constrained Wilson precision with its
continuum counterpart.

\section{The complete constrained fluctuation form, with the multiplier retained}
\label{sec:v14-form}

Let $A(v)=A_\infty^*(v)$ be V13's continuum minimum, and abbreviate the source
by $\mathcal O$.  Let $\iota:\R^{12}\to\widehat V$ be constant inclusion and
put
\begin{equation}\label{eq:v14-tangent}
 M_v=D\mathcal O(A(v))\iota,\qquad
 C_vu=-M_v^{-1}D\mathcal O(A(v))[u],\qquad
 T_vu=u+\iota C_vu\quad(u\in V).
\end{equation}
In the first expression $\iota$ is also identified with its constant field in
$X$.  These notational inclusions are fixed, not background-dependent
changes of measure.  The $12$-by-$12$ matrix $M_v$ is close to the identity,
$T_v$ maps into the actual tangent $\ker D\mathcal O(A(v))$, and
$\operatorname{rank}C_v\le12$.  Use the inherited multiplier convention
\[
 D\mathcal S_\infty(A)+D\mathcal O(A)^*\eta(v)=0,
                    \qquad \eta(v)=-D_v\mathcal G_\infty(v).
\]
The equality is on the fixed Coulomb ambient space, including its constants.

For a one-form $b$ and a pair $b,c$ define
\[
 (d_Ab)_{\mu\nu}=\partial_\mu b_\nu-\partial_\nu b_\mu
                       +[A_\mu,b_\nu]+[b_\mu,A_\nu],
 \qquad
 \Gamma(b,c)_{\mu\nu}=[b_\mu,c_\nu]+[c_\mu,b_\nu],
\]
and let $F_A=dA+[A,A]$ with the convention of \cref{eq:v13-YM}.

\begin{proposition}[Actual root-constrained Hessian on a fixed detail space]
\label{prop:v14-Hessian}
The second derivative of the continuum constrained chart action at its
minimum is the bilinear form
\begin{equation}\label{eq:v14-Hessian}
 \boxed{\begin{aligned}
 h_v(u,z)={}&\langle d_AT_vu,d_AT_vz\rangle_{L^2}
       +\langle F_A,\Gamma(T_vu,T_vz)\rangle_{L^2}\\
       &+\langle\eta(v),D^2\mathcal O(A)[T_vu,T_vz]\rangle .
 \end{aligned}}
\end{equation}
On real fields it is between $\frac12\langle du,du\rangle$ and
$\frac32\langle du,du\rangle$ on the common smaller branch.  Its operator in
the detail energy inner product is $I+\mathsf K(v)$, with
$\mathsf K(0)=0$ and $\|\mathsf K(v)\|\le C|v|$ locally.  These operators
have a holomorphic extension in the fixed real Hilbert coordinates.
\end{proposition}
\begin{proof}
The first derivative of the constraint chart in a detail direction is exactly
$T_v$.  Differentiating its constraint a second time gives
\[
 D\mathcal O(A)[D^2A_{\rm chart}(u,z)]
                  =-D^2\mathcal O(A)[T_vu,T_vz].
\]
The second chain rule for the action contains
$D\mathcal S_\infty(A)[D^2A_{\rm chart}(u,z)]$.
By stationarity it is the \emph{plus} multiplier term in
\cref{eq:v14-Hessian}.  The other part follows by twice differentiating
$\frac12\|F_A\|^2$.  Holder with $2,4,4$ makes every term bounded on $V$.
The Hessian bounds and common analytic branch are the already proved
\cref{thm:v13-cont-min}; no new physical spectral gap is claimed.
At zero, $A=0$, $\eta=0$, and the source derivative is the mean, so $T_0=I$
on $V$ and $h_0(u,z)=\langle du,dz\rangle$.  Uniform analytic derivatives
give the displayed norm estimate and the holomorphic extension.
\end{proof}

The term containing $D^2\mathcal O$ is not generally finite rank despite its
finite-dimensional multiplier.  Nor does stationarity permit it to be dropped.
For example, with $S(x,y)=(a x^2+y^2)/2$ and
$F(x,y)=x(1+\kappa y^2/2)$, the stationary point at source $c$ is $(c,0)$,
its multiplier is $-ac$, and its constrained Hessian is $1-a\kappa c^2$.
The unconstrained $y$ Hessian is one.  This is a sign regression inherited
from V4, not a new Yang--Mills model.

\section{Critical Sobolev coefficients and the continuum trace ideal}
\label{sec:v14-localideal}

For a compact operator $D$, write
$\|D\|_{p,\infty}=\sup_{n\ge1} n^{1/p}s_n(D)$; this weak-ideal quantity is
used only for estimates.  Strong Schatten spaces are Banach ideals when
$p\ge1$.

\paragraph{The external theorem used.}
Frank's Theorem 3.1 \cite{FrankV14} states that if a unitary transform
$\mathcal U:L^2(X)\to L^2(Y)$ is bounded from $L^1$ to $L^\infty$, then,
for $p>2$, multiplication by $f\in L^p(X)$ followed by the inverse transform
and a multiplier $g\in L^{p,\infty}(Y)$ belongs to $\mathfrak S_{p,\infty}$,
with norm controlled by those two function norms and the transform bound.
It includes finite-dimensional matrix multipliers.  We use $X=\T^4$,
$Y=\mathbb Z^4$, and the Fourier series transform, with normalized torus measure.
Its $L^1\to\ell^\infty$ norm is one.  Thus this is an explicit application
of an established theorem, not an assumed regularity gain of the root map.

\begin{lemma}[Multiplier bounds at the actual critical regularity]
\label{lem:v14-Cwikel}
Let $J_0=-\Delta$ on mean-zero one-forms, with the Coulomb projection when
restricted to that subspace.  For $f\in L^4$,
\begin{equation}\label{eq:v14-Cwikel}
 \|M_f J_0^{-1/2}\|_{4,\infty}\le C\|f\|_4.
\end{equation}
Consequently the operator of the unconstrained continuum Hessian difference
\[
 D^2\mathcal S_\infty(A)[u,z]-\langle du,dz\rangle
\]
in the $V$ energy metric belongs to $\mathfrak S_p$ for every $p>4$, and
\begin{equation}\label{eq:v14-localideal}
 \|\mathsf K_{\rm W}(A)\|_{\mathfrak S_p}
 \le C_p\{\|A\|_4+\|A\|_4^2+\|F_A\|_2\}.
\end{equation}
It depends continuously, and complex-holomorphically, on $A$ in $H^1$ on the
fixed Coulomb space.  In particular,
\begin{equation}\label{eq:v14-localdiff}
 \|\mathsf K_{\rm W}(A)-\mathsf K_{\rm W}(B)\|_{\mathfrak S_p}
 \le C_p\{(1+\|A\|_4+\|B\|_4)\|A-B\|_4+\|F_A-F_B\|_2\}.
\end{equation}
\end{lemma}
\begin{proof}
The Fourier sequence $g(k)=(2\pi|k|)^{-1}$ for $k\ne0$, and $g(0)=0$,
has weak $\ell^4$ norm bounded by a dimensional constant, since the number
of lattice points of radius $R$ is at most $C R^4$.  Apply the stated theorem
with $p=4$.  The Coulomb projection is a bounded right factor and does not
increase singular values.  Finite orientation and Lie components change only
a fixed constant.

Identify $V$ with its mean-zero Coulomb $L^2$ space by $u=J_0^{-1/2}\xi$.
Then $dJ_0^{-1/2}$ is bounded, and the operator $u\mapsto[A,u]$ composed with
$J_0^{-1/2}$ has weak $\mathfrak S_4$ norm at most $C\|A\|_4$.
The two first-order cross terms in $d_A^*d_A-d^*d$ are consequently weak
$\mathfrak S_4$; the term with two $A$ factors is weak $\mathfrak S_2$.
For the curvature term, each scalar coefficient in $F_A$ is in $L^2$.
Factor its multiplication as two $L^4$ square-root multipliers and a bounded
phase between them.  Each square-root multiplier composed with
$J_0^{-1/2}$ is weak $\mathfrak S_4$ with bound $C\|F_A\|_2^{1/2}$.
The inequality $s_{2n-1}(XY)\le s_n(X)s_n(Y)$ proves the weak
$\mathfrak S_2$ bound for their product.  Finite sums preserve the indicated
weak singular-value exponents, with a fixed change of constant.

Summing $n^{-p/4}$ for $p>4$ gives \cref{eq:v14-localideal}.  Subtract the
bilinear expressions at $A,B$ and apply the same estimates to each difference
to get \cref{eq:v14-localdiff}.  The form itself is polynomial in $A,dA$ on
$H^1$ and its complex bilinear extension.  The estimates show that this
polynomial takes values continuously in $\mathfrak S_p$, which proves its
holomorphy.  Square roots and phases used for the estimates are not used to
define that holomorphic polynomial.
\end{proof}

\begin{theorem}[The full actual continuum constrained precision is $I+\mathfrak S_5$]
\label{thm:v14-fullideal}
Shrink the fixed one-cell source ball so V13's source approximation holds with
$\alpha=9/10$.  The operator in \cref{prop:v14-Hessian} satisfies
\begin{equation}\label{eq:v14-fullideal}
 \boxed{\mathsf K(v)\in\mathfrak S_5(V),\qquad
        \|\mathsf K(v)\|_{\mathfrak S_5}\le C|v|.}
\end{equation}
It is holomorphic in $\mathfrak S_5$ on a common smaller complex source ball,
including every fixed source derivative.  For real sources it is self-adjoint
and $I+\mathsf K(v)\succeq\frac12I$.  After another fixed shrink one may take
$\|\mathsf K(v)\|\le r<1$ throughout the complex ball.
\end{theorem}
\begin{proof}
First restrict the local action Hessian to mean-zero directions without the
constant correction $C_v$.  \Cref{lem:v14-Cwikel} applies to the actual
$A(v)\in H^1$ and $F_{A(v)}\in L^2$.  Replacing each direction $u$ by
$T_vu=u+\iota C_vu$ adds only a finite-rank form of rank at most $24$:
every extra term has a factor through the twelve constants in at least one
slot, and the ranges of the two possible factors have total dimension at most
$24$.  Its coefficients are bounded and holomorphic by the common chart and
background estimates.  At zero this correction is zero.

The last line of \cref{eq:v14-Hessian} is
$T_v^{\mathsf t}\mathsf T(A(v),\eta(v))T_v$, restricted to $V$.  It is
$\mathfrak S_5$ by \cref{thm:v14-rootideal}, and retains the actual multiplier.
Composition of the holomorphic maps and the ideal inequality prove
holomorphy in $\mathfrak S_5$.  All parts vanish at zero.  The local
$\mathfrak S_5$ analytic bound therefore yields $C|v|$ on a smaller ball.
The real positivity and operator-norm smallness follow from
\cref{prop:v14-Hessian}.  The argument uses the continuum background from V13,
not a smoothness assumption on it beyond its proved $H^1$ regularity.
\end{proof}

\section{The determinant remainder exists, but its Born terms remain in the problem}
\label{sec:v14-det}

For $K\in\mathfrak S_5$ with $\|K\|\le r<1$ define
\begin{equation}\label{eq:v14-logdet5}
 \mathcal L_5(K)
 =\Tr\left[\log(I+K)-K+\tfrac12K^2-\tfrac13K^3+\tfrac14K^4\right]
 =\sum_{m=5}^{\infty}\frac{(-1)^{m+1}}m\Tr K^m.
\end{equation}
The bracket is defined by its trace-norm convergent combined series, not by
subtracting individually existing infinite traces.  Its exponential is the
fifth modified determinant on the branch connected to one.

\begin{theorem}[A source-differentiable determinant remainder for the actual form]
\label{thm:v14-det}
On the common smaller source ball the function
\begin{equation}\label{eq:v14-Qtail}
 \boxed{\mathcal Q_{\mathrm{tail},5}(v)
                     =\tfrac12\mathcal L_5(\mathsf K(v))}
\end{equation}
is holomorphic.  It has all fixed source derivatives, including the six
required in V8's one-loop packet, and
\begin{equation}\label{eq:v14-detbound}
 |\mathcal L_5(K)|\le\frac{\|K\|_5^5}{5(1-r)}.
\end{equation}
For a differentiable $\mathfrak S_5$ curve $K(t)$, with $R=(I+K)^{-1}$,
\begin{align}
 \frac{d}{dt}\mathcal L_5(K)&=\Tr(K^4RK'),\label{eq:v14-detD1}\\
 \frac{d^2}{dt^2}\mathcal L_5(K)&=
 \Tr(K^4RK'')+
 \sum_{j=0}^3\Tr(K^jK'K^{3-j}RK')-\Tr(K^4RK'RK').
 \label{eq:v14-detD2}
\end{align}
No commutation of $K$ and $K'$ is required.

For any finite-rank orthogonal projections $P_m\to I$ strongly on $V$, put
$K_m(v)=P_m\mathsf K(v)P_m$.  Then $\mathcal L_5(K_m(v))$ and all fixed source
derivatives converge locally uniformly to \cref{eq:v14-Qtail}'s corresponding
quantity (with the factor $1/2$ when appropriate).
\end{theorem}
\begin{proof}
Schatten Holder gives $\|K^m\|_1\le\|K\|_5^5r^{m-5}$ for $m\ge5$.
Summation proves the bound and locally normal trace-norm convergence.
Differentiate the trace series, retaining cyclicity of the trace.  Its scalar
coefficient series is
$\sum_{m\ge5}(-1)^{m+1}K^{m-1}=K^4(I+K)^{-1}$,
which proves \cref{eq:v14-detD1}.  The ordinary noncommutative product rule
for $K^4$ and $R'=-RK'R$ gives \cref{eq:v14-detD2}.
Every displayed product is trace class: the first derivative contains four
$K$ factors and one $K'$, the first sum in the second derivative has three
$K$ factors and two $K'$, and its last term has six ideal factors.
Holomorphy also follows directly from the normally convergent series and
\cref{thm:v14-fullideal}.

Finite-rank approximation in $\mathfrak S_5$ proves
$\|P_m K P_m-K\|_5\to0$.  A compact set of $\mathfrak S_5$ operators has the
same uniform compression property by a finite-net argument.  The image of a
compact source polydisc under the holomorphic map $\mathsf K$ is such a set.
Continuity of the trace series, or its first derivative bound, gives uniform
convergence of the scalar tails there.  Cauchy differentiation on larger
compact polydiscs proves convergence of every fixed source derivative.
This is exhaustion independence on the fixed continuum energy carrier; it is
not yet a theorem equating a literal lattice regulator to these compressions.
\end{proof}

For later comparisons a useful quantitative form is
\begin{equation}\label{eq:v14-detlip}
 |\mathcal L_5(K)-\mathcal L_5(L)|
 \le\frac{(\|K\|_5+\|L\|_5)^4}{1-r}\,\|K-L\|_5
\end{equation}
when both operator norms are at most $r<1$.  Integrate
\cref{eq:v14-detD1} along the straight segment; Schatten Holder bounds its
fourth power and the inverse.  This bound applies in particular to the
convergent source-Hessian contribution from \cref{cor:v14-roottransfer}.
It does not supply a missing Schatten bound for other regulator terms.

\subsection{Finite-dimensional normal coarea is a separate nonzero contribution}

At $A(v)$, restrict detail energy coordinates to any finite-dimensional
subspace, keep all twelve normal constant coordinates, and write the source
Jacobian in these coordinates as $(L_m,M_v)$.  The tangent chart is
$u\mapsto(u,-M_v^{-1}L_mu)$.  Let $\mathcal A_m$ be the restriction of the
\emph{Lagrangian} Hessian $S''+\eta\cdot\mathcal O''$ on this ambient space.
The elementary bordered determinant identity is
\begin{equation}\label{eq:v14-coarea}
 \boxed{\left|\det\begin{pmatrix}
 \mathcal A_m&(L_m,M_v)^{\mathsf t}\\(L_m,M_v)&0
 \end{pmatrix}\right|
     =(\det M_v)^2\det(I+K_m(v)).}
\end{equation}
Indeed the change to tangent and constant coordinates has determinant one,
the constraint becomes $(0,M_v)$, and expansion in the normal/border rows
leaves exactly the tangent Hessian and two factors $M_v$.  Reference energy
coordinates make the zero-source tangent Hessian the identity.
Equivalently, the map $(w,s)\mapsto(w,v)$ has Jacobian $\det M_v$, giving
$\frac12\log\det(I+K_m)+\log\det M_v$ in the negative logarithm.
There is no second copy of the same coarea factor to add.

On V13's actual branch, V12 gives $w(v)=O(|v|^4)$ and
$s(v)=\Phi^{-1}(v)+O(|v|^5)$.  The uniform source derivatives imply
$M_v=D\Phi(\Phi^{-1}(v))+O(|v|^4)$.  Thus V10's source Jacobian gives the
\emph{actual stationary normal factor}
\begin{equation}\label{eq:v14-normaljet}
 \boxed{\log\det M_v=-\sum_{\mu=1}^4|v_\mu|^2+O(|v|^3).}
\end{equation}
This is a finite normal-coordinate contribution, not a physical mass term.
Other measured and fluctuation contributions are not identified with it.

At finite rank the exact relative precision contribution is consequently
\begin{equation}\label{eq:v14-born}
 \boxed{\begin{aligned}
 Q_m^{\rm precision}(v)
 ={}&\log\det M_v+
       \frac12\sum_{r=1}^4\frac{(-1)^{r+1}}r\Tr K_m(v)^r\\
    &+\frac12\mathcal L_5(K_m(v)).
 \end{aligned}}
\end{equation}
The last line converges by the theorem.  The first four traces have \emph{not}
been declared finite or local in the limit.  Haar, gauge-slice, and other
coordinate densities of the original compact law also remain to be matched.
Since $K(0)=0$ and $\|K(v)\|_5=O(|v|)$,
\begin{equation}\label{eq:v14-tailjet}
                  J^4\mathcal Q_{\mathrm{tail},5}=0.
\end{equation}
In particular this tail alone cannot determine the quadratic running coupling.
In the finite diagnostic $K(c)=-a\kappa c^2$, its first nonzero field term is
of degree ten, while the complete Hessian contribution already has a nonzero
second derivative at zero.  Ultraviolet subtraction is not permission to lose
those low-order terms.

\section{What inverse convergence does and does not imply}
\label{sec:v14-resolvent}

\begin{proposition}[A continuum Galerkin resolvent, not the physical transfer operator]
\label{prop:v14-resolvent}
The closed form $h_v$ defines a positive operator $\mathscr A_v$ in mean-zero
Coulomb $L^2$, with compact inverse.  Under the energy isometry its inverse is
\begin{equation}\label{eq:v14-inverse}
 \mathscr A_v^{-1}
       =J_0^{-1/2}(I+\widetilde K(v))^{-1}J_0^{-1/2},
\end{equation}
where $\widetilde K$ is the unitary image of $\mathsf K$ in $L^2$.
Replacing $K$ by $P_mKP_m$ gives norm-convergent inverses, locally uniformly
in the source, and similarly for every fixed source derivative.  Finite
Fourier Galerkin inverses extended by zero have the same limit.
\end{proposition}
\begin{proof}
The form norm is equivalent to mean-zero $H^1$ by
\cref{prop:v14-Hessian}.  Fourier series show that its injection into $L^2$
is compact.  The representation theorem, or direct Riesz solution of the
form equation, gives \cref{eq:v14-inverse}.  The inverse difference in the
middle is at most $(1-r)^{-2}\|K_m-K\|$ and tends to zero, proving operator
norm convergence.  The source derivatives follow on complex polydiscs by
Cauchy.  For Fourier projections, the free high-frequency inverse omitted by
zero extension has norm $O(R^{-2})$ at frequency cutoff $R$; the projections
commute with $J_0$.  Thus the finite Galerkin inverse has the same limit.
This statement is on a unit spatial torus and on constrained elliptic
fluctuations.  It is neither a fixed-physical-time contraction nor a mass gap
of a reconstructed Hamiltonian.
\end{proof}

\begin{example}[Operator-norm convergence alone does not control the modified determinant]
\label{ex:v14-no-trace}
Let $P_N$ be a rank-$N$ orthogonal projection on one Hilbert space and set
$K_N(t)=tN^{-1/5}P_N$.  For $|t|<1/2$, these are uniformly coercive relative
precisions and $\|K_N(t)\|\to0$, uniformly in $t$.  Nevertheless
\begin{equation}\label{eq:v14-notrace}
 \mathcal L_5(K_N(t))
 =N\left[\log(1+tN^{-1/5})-tN^{-1/5}
   +\frac{t^2N^{-2/5}}2-\frac{t^3N^{-3/5}}3
   +\frac{t^4N^{-4/5}}4\right]
 \longrightarrow\frac{t^5}{5}.
\end{equation}
The logarithmic modified determinant of the norm limit is zero.
Indeed the fifth series coefficient is exactly $t^5/5$ for every $N$;
the absolute remaining series is $O_t(N^{-1/5})$.
Here $\|K_N(t)\|_5=|t|$ does not tend to zero.  This is a spectral diagnostic,
not a counterexample to the actual source transfer theorem.
\end{example}

A second diagonal diagnostic clarifies the four retained Born traces.
Give the modes of maximum Fourier radius $r\ge1$ eigenvalue $\rho/r$.
There are exactly $64r^3+16r$ such modes in four dimensions.  Hence this
operator belongs to $\mathfrak S_p$ precisely when $p>4$; its first four
cutoff traces grow respectively as $R^3,R^2,R,\log R$.  Its fifth-power tail
above $R$ is bounded by
$\rho^5(64/R+16/(3R^3))$.  The modified logarithm therefore has an absolutely
convergent tail although none of its first four positive trace terms has a
finite limit.  This model fixes the ideal bookkeeping; no claim is made that
these four divergence coefficients equal the gauge-invariant Yang--Mills ones.

\section{The acquired quantum interface and its remaining identification}
\label{sec:v14-next}

The actual continuum root-constrained precision is now a defined coercive
operator in a fixed energy space, with its nonlinear multiplier curvature
retained.  Its relative perturbation is in $\mathfrak S_5$, its fifth modified
determinant is source differentiable, and the source-curvature part has an
actual dyadic $\mathfrak S_5$ transfer estimate along the V13 backgrounds.
The normal coarea factor is also explicit and is not zero.  None of these
conclusions is an inference from a finite list of eigenvalue samples.

The next decisive estimate has two distinct parts.  First compare the
\emph{full literal lattice Lagrangian Hessians}, with compatible Coulomb/detail
identifications and measured coordinate conventions, in a common trace ideal.
The cell-average source pullback proved here is only one part of that
comparison.  Strong background convergence plus the continuum estimate
\cref{eq:v14-localdiff} does not automatically handle lattice symbols on
arbitrary ultraviolet detail directions.  Second, form the signed, complete
Born packet in \cref{eq:v14-born}, add the actual compact density and gauge
factors once, and establish its Ward-compatible local subtraction and finite
nonlocal remainder.  Calling the fifth determinant the whole one-loop action
would erase precisely the low-order terms needed for coupling matching.

This version does not yet evaluate $J^6Q$ or the two-loop $J^4B$, prove
ultraviolet matching for the original regulator, construct the interacting
measure, globalize the small-field domain, or supply physical-time
contraction.  All spaces here remain on the fixed one-cell classical branch.
The unrestricted compact well and transition-cutoff problem is not bypassed by
an elliptic determinant construction.

\paragraph{Verification and append-only record.}
The accompanying driver verifies exact finite determinant identities,
noncommuting differentiated trace formulas, the normal coarea sign and source
Jacobian, the critical Fourier counting, and the rank/rate arithmetic.  It also
recomputes the literal root mixed-Hessian regression inherited from V11 in its
finite test carrier.  These checks validate algebra and normalization; they do
not establish the analytic estimates by sampling.  The new trace-ideal theorems
use the written finite-rank argument and the explicitly cited external Cwikel
input.  The report lists only tests actually executed, and distinguishes any
fresh V13 rerun from archival reports.  No Lean formalization or numerical
value for the full continuum determinant is claimed.

V1--V13's complete mathematical body and end markers are retained byte-for-byte.
Only the preamble's version displays change.  The bundle contains the new
appendix, portable assembler, verification code, reports, and input/output
hashes.  A later continuation is to append after the following end marker;
any correction to a prior proof must be stated explicitly.

\begin{thebibliography}{99}
\bibitem[F14]{FrankV14}
R.~L. Frank, \emph{Cwikel's theorem and the CLR inequality}, Journal of
Spectral Theory \textbf{4} (2014), 1--21, DOI: 10.4171/JST/59;
arXiv:1206.3325.  Theorem 3.1, the operator-valued bounded-transform version,
is used explicitly in \cref{lem:v14-Cwikel}.
\url{https://arxiv.org/abs/1206.3325}.
The theorem was checked in the primary full-text version.  No regulator
estimate for the present root rule is imported from that reference.
\end{thebibliography}

% END OF VERSION 14 MATHEMATICAL BODY.
% Append subsequent requested versions below this marker; retain V1--V14 above.


\clearpage
\section{V15: the literal lattice fluctuation operator reaches the determinant ideal}
\label{sec:v15-context}

V14 constructs a fifth modified determinant for the continuum constrained
fluctuation form.  Its last unresolved comparison is quite specific: its
finite source-Hessian pullback is not the complete lattice Wilson Hessian in
lattice Coulomb energy coordinates.  Replacing the latter by continuum
Galerkin compressions would leave that comparison unproved.  This version
uses the exact plaquette derivatives to supply the missing lattice estimate.

On the small one-coarse-cell branch already constructed in V12--V13, the
complete lattice constrained relative precisions, embedded by their actual
Fourier energy coordinates, converge in $\mathfrak S_5$ to V14's continuum
operator.  The associated Born-subtracted determinant and the normal source
Jacobian converge with every fixed source derivative.  No convergence rate
for the full moving-background comparison is asserted.  An exact fourth-field
calculation also shows why this does not complete the Born packet: a particular
literal lattice kinetic-weight term vanishes in $\mathfrak S_5$ while its
first trace has a nonzero limiting fourth-field coefficient.

\paragraph{Inherited inputs and the new comparison.}
We use the actual real backgrounds of \cref{thm:v13-main}, including strong
physical $L^4$ convergence and strong $L^2$ convergence of their reconstructed
plaquette curvatures.  Their common holomorphic critical-energy bounds are
those of V12.  The literal balanced root observable, its finite-resolution
approximation, and its differentiated bounds are inherited from V10--V14.
The only external analytic theorem used is the same bounded-transform Cwikel
theorem as V14, \cite[Theorem 3.1]{FrankV14}; its applicability is checked
below for normalized finite-grid Fourier transforms.  This is not an
independent recertification of those inherited theorems.

The new argument has four components: an exact plaquette-Hessian identity;
a lattice-uniform multiplier estimate; convergence on the actual lattice
Coulomb carrier, rather than only on sampled continuum directions; and the
resulting determinant comparison.  The sources and present lineage were
searched for this full lattice comparison.  Semantic retrieval gave no matches;
the relevant exact source passages were then inspected.  This is an audit of
the accessed material, not a universal novelty assertion about lattice spectral
approximation.

Starred formulas below are adjoints for real backgrounds.  Their holomorphic
continuations are obtained by continuing the fixed real-coordinate matrices
(and their real transposes), not by conjugating the source parameter.
Fourier-coordinate changes are fixed, independent of that parameter.

All continuum spaces in this version are on the fixed unit four-torus.  No
physical-volume limit, compact-complement comparison, or Hamiltonian mass gap
is implied.  Haar and gauge-slice densities not contained in the displayed
normal coarea factor still belong to the measured one-loop problem.

\section{A common carrier for the actual lattice Coulomb energy}
\label{sec:v15-carrier}

Let $L=2^n\ge2$.  Give the grid $X_L=(L^{-1}\mathbb Z/\mathbb Z)^4$ the
probability counting measure $L^{-4}\sum_x$.  Norms in this measure are written
$\|\cdot\|_{p,L}$, to distinguish them from the unnormalized fine-link norms.
For a fine Lie array $a$, its physical field is $A_L(x/L)=La(x)$, extended
constantly on the mesh cubes.  Thus
\[
       \|A_L\|_{4,L}=\|a\|_{\ell^4},\qquad
       \|D_L A_L\|_{2,L}=\|\nabla a\|_{\ell^2},
       \quad D_{L,\nu}=L(\mathsf t_\nu-I).
\]
Here $\mathsf t_\nu$ is the actual one-mesh translation.  In this section
$D_L$ with an orientation index denotes a spatial difference, not a source
or noise derivative.

Choose centered Fourier representatives
$\mathcal K_L=\{-L/2,\ldots,L/2-1\}^4$.  For $k\ne0$ put
\begin{equation}\label{eq:v15-symbols}
 \begin{aligned}
 q_{L,\mu}(k)&=L(e^{2\pi i k_\mu/L}-1),&
 \lambda_L(k)&=\sum_\mu|q_{L,\mu}(k)|^2,\\
 P_L(k)&=I-\frac{q_L(k)q_L(k)^*}{\lambda_L(k)},&
 g_L(k)&=\lambda_L(k)^{-1/2}.
 \end{aligned}
\end{equation}
All these symbols are zero outside $\mathcal K_L\setminus\{0\}$ when used
as truncated maps; $P_L$ always means the truncated projection in this
convention.  Tensor the orientation maps with the three Lie components.
The common carrier is
\[
 \mathcal H=\ell^2(\mathbb Z^4\setminus\{0\};\mathbb C^{12}).
\]
On it let $P(k)=I-kk^{\mathsf t}/|k|^2$, $g(k)=(2\pi|k|)^{-1}$ and
$j=gP$ followed by continuum Fourier inversion.  Let $j_L$ be Fourier
multiplication by $g_LP_L$, inverse discrete Fourier transform, and
piecewise-constant physical reconstruction.  The unscaled fine variation
corresponding to $\xi\in\mathcal H$ is
\begin{equation}\label{eq:v15-variation}
                         \delta a=L^{-1}j_L\xi.
\end{equation}
It is exactly discrete Coulomb and mean zero.  The nonzero subspace
$P_L\mathcal H$ has complex dimension $9(L^4-1)$, the complexification of
the real mean-zero Coulomb detail dimension.  Nyquist representatives are
kept with their actual discrete reality convention; no extra real field is
counted at an aliased boundary frequency.

\begin{lemma}[Uniform energy embeddings and strong low-frequency identification]
\label{lem:v15-embedding}
The maps $j_L:\mathcal H\to L^4$ are uniformly bounded, and
\begin{equation}\label{eq:v15-embedding}
 \|d_Lj_L\xi\|_{2,L}^2=\|P_L\xi\|_2^2,\qquad
 \|j_L-j\|_{\mathcal H\to L^2}\longrightarrow0.
\end{equation}
The normalized curl multipliers $R_L=d_Lg_LP_L$ and their adjoints,
embedded on Fourier carriers for links and plaquettes, converge strongly to
$R=dgP$ and its adjoint; their norms are at most one.  Every fixed finite
mesh translation converges strongly to the identity on these carriers.
\end{lemma}
\begin{proof}
The first identity is the discrete Hodge identity on the Coulomb subspace.
Apply V9's lattice Sobolev estimate to the fine variation
$L^{-1}j_L\xi$ to obtain the uniform $L^4$ bound.  The elementary sine
inequalities give, for centered representatives,
\begin{equation}\label{eq:v15-dispersion}
          16|k|^2\le\lambda_L(k)\le4\pi^2|k|^2.
\end{equation}
For every fixed nonzero $k$, $q_L(k)\to2\pi i k$, $P_L(k)\to P(k)$ and
$g_L(k)\to g(k)$.  On finitely many fixed modes the reconstructed plane
waves converge in $L^2$ to their continuum waves.  On the complementary
modes $|k|>R$, both embedding norms into $L^2$ are at most $C/R$, by
\cref{eq:v15-dispersion} and Parseval.  First choose $R$ large and then
$L$ large; this proves the operator-norm $L^2$ convergence.  The bounded
curl and translation symbols converge on each fixed mode, as do their
adjoints.  Density of finite Fourier sequences proves the strong assertions.
This argument does not claim operator-norm convergence $\mathcal H\to L^4$.
\end{proof}

When a grid multiplication or curl has a Fourier output, it is embedded by
its discrete Fourier coefficients, not by identifying its point samples
with a continuum multiplication operator.  The following simple observation
will be used repeatedly: if piecewise-constant grid fields $f_L\to f$ in
$L^2$, their embedded discrete Fourier coefficients converge in $\ell^2$
to the Fourier coefficients of $f$.  Indeed each fixed coefficient converges
by testing against a smooth plane wave and estimating its within-cell error;
Parseval gives convergence of the norms.  Weak convergence and convergence
of Hilbert norms then give strong convergence.  Finite Fourier approximation
provides an equivalent proof.

\section{A uniform discrete multiplier estimate and its consistency}
\label{sec:v15-multipliers}

We use $\|T\|_{p,\infty}=\sup_{m\ge1}m^{1/p}s_m(T)$ for weak Schatten
estimates.  Constants below are uniform in $L$; fixed orientation, Lie, and
plaquette component counts only change those constants.

\begin{lemma}[Cwikel on the actual finite grid]
\label{lem:v15-Cwikel}
Let $M_f$ be grid multiplication and let $\mathcal D_L$ denote the map
$g_LP_L$ followed by inverse discrete Fourier transform, before physical
reconstruction.  Then
\begin{equation}\label{eq:v15-Cwikel}
 \boxed{\quad \|M_f\mathcal D_L\|_{4,\infty}
                  \le C\|f\|_{4,L}.\quad}
\end{equation}
The estimate remains valid with a fixed mesh translation before
$\mathcal D_L$.  Bilinear forms with an $L^2$ coefficient satisfy
\begin{equation}\label{eq:v15-potential}
 \|\mathcal D_L^*\mathsf t^*M_V\mathsf t'\mathcal D_L\|_{2,\infty}
                          \le C\|V\|_{2,L}.
\end{equation}
In particular these operators are uniformly bounded in every
$\mathfrak S_p$, $p>4$, with the corresponding right sides.
\end{lemma}
\begin{proof}
The discrete Fourier transform from normalized grid $L^2$ to counting
Fourier $\ell^2$ is unitary, and its $L^1\to\ell^\infty$ norm is one.
The theorem of \cite[Theorem 3.1]{FrankV14} therefore applies on these two
finite measure spaces with exactly the same transform bound as for the
unit continuum torus.  Equation \cref{eq:v15-dispersion} and four-dimensional
lattice-point counting give
\[
        \sup_{t>0}t^4\#\{k:g_L(k)>t\}\le C.
\]
Apply that theorem with exponent four and append the bounded projection
$P_L$ to prove \cref{eq:v15-Cwikel}.  A translation commutes with the
scalar multiplier and projection and is unitary.  Alternatively move it
across $M_f$, translating $f$ without changing its norm.

For a scalar coefficient $V$, write $V=|V|^{1/2}\omega|V|^{1/2}$ with
$|\omega|\le1$.  Each square-root multiplier is in $L^4$ and gives a
weak $\mathfrak S_4$ factor.  The singular-value product inequality gives
weak $\mathfrak S_2$ for their product.  Matrix coefficients are finite
sums of scalar entries.  Summing $m^{-p/4}$ or $m^{-p/2}$ proves the last
assertion.  Square roots are only an estimate; the bilinear operator is
linear in $V$ and keeps its original sign.
\end{proof}

\begin{lemma}[Consistency of the local operator class on varying rough coefficients]
\label{lem:v15-consistency}
Suppose the physical reconstructions of $b_L$ converge strongly in $L^4$
to $b$, and those of $V_L$ converge strongly in $L^2$ to $V$.  Embed all
Fourier operators as above.  For every $p>4$, the operators
\[
 R_L^*M_{b_L}\mathsf t\mathcal D_L,\
 \mathcal D_L^*\mathsf t^*M_{V_L}\mathsf t'\mathcal D_L
\]
converge in $\mathfrak S_p$ to their continuum analogues, with mesh
translations replaced by the identity.  Products of two factors of the
first smoothing type have the analogous convergence.  Fixed finite sums
and bounded strongly convergent Fourier multipliers, with their adjoints
strongly convergent when needed, may be inserted.
\end{lemma}
\begin{proof}
First let a coefficient be a fixed trigonometric polynomial sampled on the
grid.  Its multiplication shifts Fourier indices by only finitely many fixed
vectors; a wrap can occur only in the high-frequency tail.  On a fixed finite
set of input modes all entries converge to the continuum ones.  For the
remaining input modes,
\[
 \sum_{\substack{k\in\mathcal K_L\\|k|>R}}g_L(k)^p
                    \le C_pR^{4-p},\qquad p>4.
\]
Thus the smoothing multiplier restricted to that tail has uniformly small
$\mathfrak S_p$ norm.  This proves convergence for bounded polynomial
coefficients, including the zero-order sandwich.  The bounded curl factors
converge strongly together with their adjoints by
\cref{lem:v15-embedding}.  Multiplication of a fixed Schatten operator by
such a family converges in Schatten norm: approximate the fixed operator
by finite rank and then use strong convergence.  This proves the assertions
with the extra bounded factors.

For general $b$, choose a smooth polynomial approximation in $L^4$.
Its cell averages have the same approximation error in grid $L^4$ by Jensen;
point samples of that fixed polynomial differ from its cell averages by a
quantity tending uniformly to zero.  Strong convergence of $b_L$ and
\cref{lem:v15-Cwikel} bound the replacement error uniformly in $L$.
The same approximation for $V$ uses its $L^2$ norm and
\cref{eq:v15-potential}, applied to the difference coefficient.  This avoids
assuming continuity of a chosen square-root phase.  The triangle inequality
finishes the proof.  Products are handled by the ideal inequality, for example
$\|XY-X'Y'\|_p\le\|X-X'\|_p\|Y\|+\|X'\|\|Y-Y'\|_p$.
This is a comparison of the actual grid multiplications, including their
aliases, not an assumption that the lattice matrices are Galerkin compressions.
\end{proof}

\section{An exact plaquette identity separates the lattice kinetic weight}
\label{sec:v15-plaquette}

For one plaquette, write its four oriented logarithms as $X_0,\ldots,X_3$,
including the signs on its inverse links, and put
\[
  P=e^{X_0}e^{X_1}e^{X_2}e^{X_3}=p_0+\boldsymbol p\cdot E.
\]
For fine variations $b,c$ use the corresponding oriented $Y_i,Z_i$.  Define
right-trivialized derivatives
\begin{equation}\label{eq:v15-q-r}
 q_b=(DP[b])P^{-1},\qquad
 r_{b,c}=\operatorname{Im}\bigl((D^2P[b,c])P^{-1}\bigr).
\end{equation}
The dot product on imaginary quaternion vectors is the coefficient inner
product, complex-bilinearly extended when appropriate.

\begin{lemma}[Literal plaquette Hessian, with no background truncation]
\label{lem:v15-plaquette}
On the identity chart,
\begin{equation}\label{eq:v15-plaquette}
 \boxed{\quad
 D^2(1-\operatorname{Re}P)[b,c]
      =p_0\,q_b\cdot q_c+\boldsymbol p\cdot r_{b,c}.
 \quad}
\end{equation}
Here $q_b$ is purely imaginary and
\begin{equation}\label{eq:v15-dexp}
 q_b=\sum_{i=0}^3\operatorname{Ad}_{e^{X_0}\cdots e^{X_{i-1}}}
                \mathscr E(X_i)Y_i,\qquad
 \mathscr E(X)Y=\int_0^1e^{sX}Ye^{-sX}\,ds.
\end{equation}
The maps defining $r_{b,c}$ are analytic local bilinear maps of the four
variation entries, uniformly bounded on a fixed small polydisc.  At zero,
\begin{equation}\label{eq:v15-rzero}
 r_{b,c}(0)=\frac12\sum_{i<j}\bigl([Y_i,Z_j]+[Z_i,Y_j]\bigr).
\end{equation}
\end{lemma}
\begin{proof}
Differentiate the ordered product and multiply on the right by its inverse.
The derivative of one exponential on the right is
$D(e^X)[Y]e^{-X}=\mathscr E(X)Y$, proving
\cref{eq:v15-dexp}.  It is imaginary by Lie algebra closure.
Twice differentiating $P\overline P=1$ shows that the scalar part of
$(D^2P[b,c])P^{-1}$ equals $-q_b\cdot q_c$.  Multiplying this quaternion
by $P$ and taking its negative scalar part proves
\cref{eq:v15-plaquette}.  The ordered product and inverse are analytic
on the fixed finite stencil, giving the bounded bilinear coefficients of $r$.
At zero, the same-factor second exponential derivatives are anticommutators
and have no imaginary part.  The different-factor terms, in their actual
order, have imaginary part \cref{eq:v15-rzero}.  This proves the formula
including its factors and sign.  For complex fields ``Re'' and ``Im'' in
these identities mean scalar and imaginary quaternion components, not
complex conjugation of a parameter.
\end{proof}

Now let $b=L^{-1}u$, $c=L^{-1}z$, with physical variations $u,z$ on $X_L$.
Set
\[
 \mathcal Q_{a,L}u=L^2q_b,
 \qquad \mathcal R_{a,L}(u,z)=L^2r_{b,c},
 \qquad F_L=L^2\boldsymbol p.
\]
Then, exactly,
\begin{equation}\label{eq:v15-physical-Hessian}
 D^2\mathcal S(a)[L^{-1}u,L^{-1}z]
 =\int_{X_L}\left\{
       p_0\,\mathcal Q_{a,L}u\cdot\mathcal Q_{a,L}z
       +F_L\cdot\mathcal R_{a,L}(u,z)\right\}.
\end{equation}
The integral sums plaquette orientations as well as normalized spatial cells.

\begin{proposition}[The exact local coefficients have the critical scaling]
\label{prop:v15-coefficients}
The first map has the exact form
\begin{equation}\label{eq:v15-Qform}
             \mathcal Q_{a,L}u=d_Lu+\sum_{e}M_{b_{L,e}}\mathsf t_eu,
             \qquad \sum_e\|b_{L,e}\|_{4,L}\le C\|A_L\|_4,
\end{equation}
with a fixed number of local shifts and component contractions.  Suppose
$A_L\to A$ strongly in $L^4$, and the actual reconstructed curvatures
$L^2(P-I)\to F_A$ strongly in $L^2$.  Then the coefficients in
\cref{eq:v15-Qform} converge in $L^4$ to their linear-in-$A$ limits; their
combined co-located map is $u\mapsto[A,u]$, as in $d_A=d+[A,\cdot]$.
The coefficients of $F_L\cdot\mathcal R_{a,L}$ converge in $L^2$ to those
of $F_A\cdot\Gamma(u,z)$.

The remaining factor $\omega_L=p_0-1$ satisfies
\begin{equation}\label{eq:v15-scalarweight}
 \boxed{\quad
 \|M_{\omega_L}\|_{\mathfrak S_5}
       \le\|\omega_L\|_\infty^{4/5}
                         \Bigl(\sum_{x,\mu<\nu}|\omega_L(x)|\Bigr)^{1/5}
       \longrightarrow0.
 \quad}
\end{equation}
These are singular values of grid multiplication, so the sum in this display
is \emph{unnormalized}.  Finite Lie multiplicities only change its constant.
\end{proposition}
\begin{proof}
In \cref{eq:v15-dexp} the coefficient of a variation at an oriented link is
its signed identity plus an analytic correction of size $C\max|a_e|$.
After the physical factors are restored, the correction in
\cref{eq:v15-Qform} is $L$ times this local correction.  Bounded stencil
incidence proves its $L^4$ estimate.

Strong $L^4$ convergence of the physical reconstructions implies
\begin{equation}\label{eq:v15-linksmall}
                           \max_{x,\mu}|a_{L,\mu}(x)|\longrightarrow0.
\end{equation}
Indeed $|a_{L,\mu}(x)|^4$ is the integral of the corresponding reconstructed
fourth power on one cell.  Strong $L^4$ convergence implies uniform
integrability of those fourth powers, and the cell volumes tend to zero.
Every fixed fine translation of $A_L$ has the same strong limit, by translation
continuity of $A$ and the triangle inequality.  The quadratic-and-higher part
of each local coefficient, multiplied by $L$, has $L^4$ norm at most
$C\max|a_L|\,\|A_L\|_4$, and tends to zero.  The first-order coefficient
therefore has the asserted limit.  Expanding one plaquette with co-located
values gives its combined first variation
$[A_\mu,u_\nu]+[u_\mu,A_\nu]$; this also follows by differentiating its
quadratic imaginary term.  Equation \cref{eq:v15-rzero} gives the combined
second variation $\Gamma(u,z)$ after co-location.  Its remaining analytic
coefficients differ uniformly from their zero values by $O(\max|a_L|)$.
Multiply by the strongly convergent $F_L$ to obtain the $L^2$ assertion.

For real unit quaternions, $\omega_L=-|P-I|^2/2$, so
$\sum|\omega_L|=\mathcal S(a_L)$ is bounded.  Equation
\cref{eq:v15-linksmall} and the fixed plaquette word imply
$\|\omega_L\|_\infty\to0$.  The singular values of multiplication are
the absolute values of its entries.  Bounding their fifth powers by
$\|\omega_L\|_\infty^4\sum|\omega_L|$ proves
\cref{eq:v15-scalarweight}.  The normalized grid measure does not change
these eigenvalues.
\end{proof}

For uniform complex estimates, one uses $P\overline P=1$ in complex quaternion
algebra: near identity, $p_0-1=-\boldsymbol p\cdot\boldsymbol p/(1+p_0)$.
Also the first-order plaquette expansion with its actual curl gives
\begin{equation}\label{eq:v15-complexcurv}
 \|L^2(P-I)\|_{2,L}
                 \le C(\|da\|_2+\|a\|_4^2).
\end{equation}
Hence $\sum|p_0-1|$ has the same uniform bound by a fixed multiple of the
squared right side.  These inequalities are bounds for the holomorphic
polynomials; no conjugated source parameter is inserted in their definition.

\section{Schatten convergence of the full literal local Wilson Hessian}
\label{sec:v15-Wilson}

Define $K_{{\rm W},L}$ on $\mathcal H$ as the local Wilson Hessian in the
variations \cref{eq:v15-variation}, minus its free value $P_L$.  It vanishes
on $(P_L\mathcal H)^\perp$.  The continuum operator $K_{\rm W}$ is V14's
energy-normalized local Wilson perturbation, extended by zero off the
continuum Coulomb subspace.

\begin{theorem}[The lattice local precision converges at the actual background regularity]
\label{thm:v15-Wilson}
Under the strong background and curvature hypotheses in
\cref{prop:v15-coefficients},
\begin{equation}\label{eq:v15-Wilsonlimit}
                   \boxed{\quad
 K_{{\rm W},L}\longrightarrow K_{\rm W}(A)
                           \quad\hbox{in }\mathfrak S_5(\mathcal H).
 \quad}
\end{equation}
Uniform small critical-energy bounds on complex lattice backgrounds also give
uniform $\mathfrak S_5$ bounds, of order $O(\|a\|_{\mathcal E})$ on a
fixed small ball.  No continuum smoothness beyond $A\in H^1$, $A\in L^4$
and $F_A\in L^2$ is required for the real convergence statement.
\end{theorem}
\begin{proof}
Insert \cref{eq:v15-Qform} into \cref{eq:v15-physical-Hessian} and first
replace $p_0$ by one.  Subtract the free term $R_L^*R_L=P_L$.  The result
is the sum of two first-order cross terms, products of two smoothing
multipliers, and the curvature bilinear terms.  By
\cref{prop:v15-coefficients} their coefficients converge in the $L^4$ and
$L^2$ norms used in \cref{lem:v15-consistency}.  That lemma proves their
$\mathfrak S_5$ convergence, retaining every shifted variation and every
ultraviolet Fourier mode.  The co-located limiting form is exactly
\[
       \langle d_Au,d_Az\rangle-\langle du,dz\rangle
                           +\langle F_A,\Gamma(u,z)\rangle.
\]

The part not included in those terms is the exact kinetic-weight correction
\begin{equation}\label{eq:v15-kineticdefect}
  E_{{\rm kin},L}=
  (\mathcal Q_{a,L}\mathcal D_L)^*M_{p_0-1}
                         (\mathcal Q_{a,L}\mathcal D_L).
\end{equation}
The outside operators are uniformly bounded, since their free parts are
$R_L$ and their remaining parts have the bound
\cref{eq:v15-Cwikel}.  Equation \cref{eq:v15-scalarweight} and the ideal
inequality give $\|E_{{\rm kin},L}\|_5\to0$.  This proves
\cref{eq:v15-Wilsonlimit} without discarding this term from the finite
operator.  The same estimates, using \cref{eq:v15-complexcurv}, give a
uniform complex bound.  The local first-order terms are $O(\|a\|_4)$;
the other coefficients are controlled by the critical-energy norm or its
square.  This gives the asserted small-ball bound.  Holomorphy of the
finite matrices is their literal logarithmic-coordinate holomorphy.
\end{proof}

This theorem is not an application of the continuum Cwikel estimate to an
unidentified lattice symbol.  The grid normalization, its dispersion inequality,
its aliases, the precise scalar kinetic-weight term, and the actual shifted
plaquette derivatives have all entered the comparison.  It is also not
trace-norm convergence of the first four powers; the distinction is tested
explicitly below.

\section{The nonlinear root term on lattice energy coordinates}
\label{sec:v15-root}

The source-Hessian pullback of V14 has continuum arguments $j\xi$.  The
lattice problem instead has $j_L\xi$.  The difference is significant at
critical regularity; \cref{lem:v15-embedding} did not claim $L^4$
operator-norm convergence of those embeddings.

Adjoin the twelve constants and write
$\widehat j_L(s,\xi)=s+j_L\xi$, $\widehat j(s,\xi)=s+j\xi$.
For physical reconstructed backgrounds $A_L$, the exact source Hessian in
these ambient coordinates is
\begin{equation}\label{eq:v15-roothessian}
 \widehat T_L[(s,\xi),(t,\zeta)]
 =\left\langle\eta_L,
 D^2\mathcal O_L(A_L)
          [\widehat j_L(s,\xi),\widehat j_L(t,\zeta)]\right\rangle.
\end{equation}
Since $\mathsf S_Lj_L\xi=L^{-1}j_L\xi$ as a fine array, this is the
literal lattice root Hessian, not a newly chosen source.

\begin{theorem}[Actual lattice-energy source transport]
\label{thm:v15-root}
Suppose $A_L\to A$ in $L^4$ in the common small ball and
$\eta_L\to\eta$.  For $\alpha>4/5$ as admitted by V13,
\begin{equation}\label{eq:v15-rootlimit}
              \boxed{\quad
       \widehat T_L\longrightarrow\widehat T
                                      \quad\hbox{in }\mathfrak S_5.
              \quad}
\end{equation}
Here $\widehat T$ is the continuum source Hessian with $\widehat j$.
Uniform complex-background bounds hold in this ideal.  Moreover the normal
matrices and detail source rows obey
\begin{equation}\label{eq:v15-normalconvergence}
 D\mathcal O_L(A_L)\iota\to D\mathcal O(A)\iota,
 \qquad
 D\mathcal O_L(A_L)j_L\to D\mathcal O(A)j
\end{equation}
in their finite-matrix and Hilbert-to-$\mathbb C^{12}$ operator norms.
\end{theorem}
\begin{proof}
Fix a dyadic intermediate scale $K\le L$ and replace $\mathcal O_L$ by
$\mathcal F_{K,1}B_{L/K}\mathsf S_L$.  V13's uniform stopping estimate,
Cauchy differentiation, and the uniform $L^4$ bounds for $\widehat j_L$
give a rank-at-most-$12K^4$ Hessian approximant with operator error
$CK^{-\alpha}$.  For $K\ge L$ use the finite Hessian itself.  The rank/rate
lemma of V14 therefore gives a uniform $\mathfrak S_5$ tail
$CK^{4/5-\alpha}$ for this approximation on the actual lattice energy
carrier.

For fixed $K$, the linear maps carrying its two free arguments converge in
operator norm:
\begin{equation}\label{eq:v15-fixedK}
 B_{L/K}\mathsf S_L\widehat j_L
                         \longrightarrow\mathsf A_K\widehat j.
\end{equation}
Indeed V13's quadrature error is at most $4K/L$ in $L^4\to\ell^4$ norm,
and the embeddings are uniformly bounded in $L^4$.  The remaining difference
is $\mathsf A_K(\widehat j_L-\widehat j)$, which tends to zero by
\cref{lem:v15-embedding}: for this fixed $K$, $\mathsf A_K$ is bounded
on $L^2$ into a finite-dimensional space.  The finite base point
$B_{L/K}\mathsf S_LA_L$ also converges to $\mathsf A_KA$.  The ordinary
finite-dimensional derivatives of $\mathcal F_{K,1}$ are continuous.
Thus the fixed-$K$ Hessian approximants converge in operator norm and, by their
common rank bound, in $\mathfrak S_5$ norm.  First choose $K$ large to
control the two tails, then send $L$ to infinity.  This proves
\cref{eq:v15-rootlimit}.

For the first derivatives use the same stopping approximation and
\cref{eq:v15-fixedK}; its differentiated remainder is $O(K^{-\alpha})$
in the Hilbert-to-source norm.  Constants are unchanged by the embeddings.
This proves \cref{eq:v15-normalconvergence}.  The uniform complex bounds
use the identical rank approximation and Cauchy estimates before taking any
real limit.  No operator-norm convergence of $j_L$ into $L^4$ has been used.
\end{proof}

\section{The complete constrained precision and determinant now transfer}
\label{sec:v15-main}

Evaluate all quantities at V13's actual real minimizing field $a_L^*(v)$,
with physical reconstruction $A_L(v)$ and multiplier
$\eta_L(v)=-D\mathcal G_L(v)$.  Set
\[
 M_L=D\mathcal O_L(A_L)\iota,\quad
 C_L=-M_L^{-1}D\mathcal O_L(A_L)j_L,\quad
 \mathcal T_L\xi=j_L\xi+\iota C_L\xi.
\]
These are the derivatives of the actual Coulomb/constant constraint chart.
The complete lattice Lagrangian Hessian in energy coordinates is
\begin{equation}\label{eq:v15-hfull}
 h_L(\xi,\zeta)=
 D^2\mathcal S(a_L^*)[L^{-1}\mathcal T_L\xi,L^{-1}\mathcal T_L\zeta]
 +\langle\eta_L,D^2\mathcal O_L(A_L)
                    [\mathcal T_L\xi,\mathcal T_L\zeta]\rangle.
\end{equation}
Define $K_L=h_L-P_L$ and extend it by zero off $P_L\mathcal H$.
Both free argument maps annihilate that orthogonal complement.  The
continuum $K$ is the extension by zero of V14's energy-normalized constrained
operator.  The plus multiplier sign in \cref{eq:v15-hfull} is the same
second-constraint-chain-rule sign as V4 and V14.

\begin{theorem}[Literal lattice-to-continuum constrained Schatten transport]
\label{thm:v15-main}
On a common sufficiently small one-cell source ball,
\begin{equation}\label{eq:v15-main}
                         \boxed{\quad K_L(v)\to K(v)
                              \quad\hbox{in }\mathfrak S_5(\mathcal H).
                         \quad}
\end{equation}
The convergence is locally uniform on a smaller common complex source ball,
and every fixed source derivative converges in the same ideal norm.
The normal matrices $M_L(v)$ and their derivatives converge as well.
For real sources the finite and limiting relative precisions are positive;
after shrinking the complex ball one has uniformly $\|K_L(v)\|,\|K(v)\|<r<1$.
No mesh rate for this full statement is asserted.
\end{theorem}
\begin{proof}
For each real $v$, V13 gives strong $A_L\to A$ in $L^4$, strong exact
curvature convergence in $L^2$, and convergence of $\eta_L$.  The local
Wilson contribution on mean-zero details converges by
\cref{thm:v15-Wilson}.  The source curvature converges on the ambient
constant-plus-detail carrier by \cref{thm:v15-root}.  Its constant and
source-row assertions give $M_L\to M$ and $C_L\to C$ in operator norm;
the normal inverse stays uniformly bounded.

It remains to account for the constant legs in the local Wilson Hessian,
rather than infer them from its mean-zero restriction.  Their number is
finite.  In \cref{eq:v15-physical-Hessian}, a constant has zero free curl,
and its $\mathcal Q$ image is a finite sum of the $b_{L,e}$ coefficients.
These converge in $L^4$ and hence in $L^2$.  Cross terms with $R_L$ therefore
converge as Hilbert functionals, using strong convergence of $R_L^*$ and the
Fourier observation following \cref{lem:v15-embedding}.  Terms with an
additional $b_{L,e}\mathcal D_L$ use
\cref{lem:v15-consistency}.  Curvature terms with one constant are handled
by the same $L^2$ coefficient estimate: approximate that coefficient in
$L^2$ by a smooth one, use the uniform $L^4$ detail embedding for the error
(on a unit volume, $L^4\subset L^2$), and use low-frequency convergence
for the smooth term.  Terms with two constants converge directly in the
finite normal matrix.  The kinetic-weight corrections tend to zero in
operator norm by their uniform maximum weight and the bounded constant-leg
maps.  Thus all additional tangent terms converge in operator norm and are
of uniformly finite rank.  They consequently converge in $\mathfrak S_5$.
This proves \cref{eq:v15-main} for real sources.

For complex sources, V12 bounds the actual analytic lattice backgrounds in
critical energy on a common domain.  The estimate
\cref{eq:v15-complexcurv}, the discrete multiplier bounds, and the root
rank estimates give uniform $\mathfrak S_5$ bounds for $K_L(v)$ there.
The normal inverses and constant corrections are uniformly analytic.
Every $K_L(0)=0$, so Cauchy bounds give $\|K_L(v)\|_5\le C|v|$ after
one shrink, uniformly in $L$.  The same holds for the continuum operator.

For completeness, norm convergence of derivatives does not follow by invoking
scalar compactness for operator-valued maps.  At a real source point,
uniform Banach-valued Cauchy estimates make a fixed real finite-difference
formula for any derivative accurate in $\mathfrak S_5$ with a uniform error
tending to zero as its step tends to zero.  For each fixed step its finitely
many values converge in $\mathfrak S_5$ by the already proved real result.
Thus every Taylor coefficient converges in that norm.  On a smaller complex
polydisc the Taylor-series tail is uniformly geometrically bounded by the
common Cauchy bound.  Convergence of finitely many coefficients and then
control of that tail prove local uniform norm convergence; another Cauchy
formula gives all fixed differentiated statements.  This argument also
applies to $M_L$.  Real positivity is the inherited constrained Hessian
bound.  The uniform $C|v|$ estimate provides the stated complex operator
radius.  V13's moving-background convergence is qualitative, so no rate
is manufactured from a rate for only the source term.
\end{proof}

\begin{corollary}[A regulator limit of the actual Born-subtracted precision contribution]
\label{cor:v15-det}
Let the determinant and traces below be on $P_L\mathcal H$, and normalize
the zero-source detail precision to the identity.  Define
\begin{equation}\label{eq:v15-subtracted}
 \begin{split}
 Q^{\rm prec}_{5,L}(v)={}&\log\det M_L(v)
       +\tfrac12\log\det(I+K_L(v))\\
   &-\tfrac12\sum_{r=1}^4\frac{(-1)^{r+1}}r\Tr K_L(v)^r.
 \end{split}
\end{equation}
Then, locally uniformly with every fixed source derivative,
\begin{equation}\label{eq:v15-detlimit}
 \boxed{\quad
 Q^{\rm prec}_{5,L}(v)\longrightarrow
       \log\det M_v+\tfrac12\mathcal L_5(K(v)).\quad}
\end{equation}
The finite expression uses the literal lattice constrained Hessian, not a
continuum Galerkin matrix.  The normal coarea factor is counted once.
\end{corollary}
\begin{proof}
The exact finite trace identity makes
\cref{eq:v15-subtracted} equal to
$\log\det M_L+\frac12\mathcal L_5(K_L)$.  Extension by zero adds only
identity eigenvalues to $I+K_L$, so the modified determinant is unchanged
on the common carrier.  Apply \cref{thm:v15-main} and the trace-ideal
continuity estimate \cref{eq:v14-detlip}.  The normal determinant has its
branch near one.  Uniform holomorphic convergence gives the differentiated
conclusion.  In the chosen constant-plus-energy coordinates the bordered
identity is exactly $(\det M_L)^2\det(I+K_L)$, by the same block
calculation as \cref{eq:v14-coarea}.  Coordinate scale factors independent
of $v$ belong to the separately retained vacuum normalization.
\end{proof}

The theorem does not identify the four subtracted traces with physically
admissible local counterterms.  Nor does the coordinate formula for $M_L$
include the full compact Haar density or the Jacobian of changing the original
tree slice to the present Coulomb slice.  Those measured terms still have to
be combined in one convention.  What is now closed is the actual lattice
precision tail and normal factor on this fixed small classical branch.

\section{A literal ultraviolet contribution invisible to the fifth determinant limit}
\label{sec:v15-Born}

The small ideal norm of the kinetic-weight term
\cref{eq:v15-kineticdefect} does not make its first Born trace negligible.
This can be seen in the actual Wilson/root problem, not only in a diagonal
model.  Denote by $[4]$ the homogeneous fourth coefficient in the source, not
the fourth true derivative.  Evaluate $E_{{\rm kin},L}$ on the actual
constrained tangent and actual minimizing branch.  Its degree-four part is
unchanged by the higher tangent and background corrections.

\begin{theorem}[A nonzero fourth-field first-trace limit of a vanishing ideal term]
\label{thm:v15-Born}
For a homogeneous source direction $v=(v_1,\ldots,v_4)$, set
\[
                E_4(v)=\tfrac12\sum_{\mu<\nu}|[v_\mu,v_\nu]|^2.
\]
On the literal finite lattice constrained branch,
\begin{equation}\label{eq:v15-Borntrace}
 \boxed{\quad
       \Tr E_{{\rm kin},L}^{[4]}(v)
              =-\frac32(1-L^{-4})E_4(v).
       \quad}
\end{equation}
In contrast,
\begin{equation}\label{eq:v15-Bornideal}
          \|E_{{\rm kin},L}^{[4]}(v)\|_5
                          \le C|v|^4L^{-16/5}\longrightarrow0.
\end{equation}
Its contribution to the fourth-field term of the negative-log first Born
trace therefore tends to $-3E_4(v)/4$, even though it is absent from the
limiting $\mathfrak S_5$ operator.
\end{theorem}
\begin{proof}
The actual background has leading term $v/L$, its mean correction begins
at degree three, and its inhomogeneous detail begins at degree four by V11--V12.
The leading plaquette imaginary part is consequently
$[v_\mu,v_\nu]/L^2$, of field degree two.  The quaternion identity gives
\[
             (p_0-1)^{[4]}_{\mu\nu}
                     =-\frac{|[v_\mu,v_\nu]|^2}{2L^4},
\]
constant over the fine lattice.  The degree-zero tangent is $j_L$ and
$\mathcal Q_{0,L}\mathcal D_L=R_L$.  Since the weight begins at degree four,
all its other tangent and $\mathcal Q$ corrections enter only higher degrees.
Thus its fourth coefficient is $R_L^*M_{(p_0-1)^{[4]}}R_L$.

For a single plaquette orientation $(\mu,\nu)$ the squared norm of its
normalized curl row at $k\ne0$ is
\[
                  \frac{|q_{L,\mu}(k)|^2+|q_{L,\nu}(k)|^2}{\lambda_L(k)}.
\]
That row already annihilates the gauge vector, so $P_L$ does not change it.
Permuting the four lattice axes shows that the sum of each individual
$|q_{L,\mu}|^2/\lambda_L$ over nonzero modes is $(L^4-1)/4$.
The three Lie colours therefore give
\[
                   \Tr(R_{L,\mu\nu}^*R_{L,\mu\nu})
                               =\frac32(L^4-1).
\]
Multiply by the displayed constant weight and sum the six orientations.
This proves \cref{eq:v15-Borntrace}, with every fine mode included.
The operator norm of this fourth coefficient is at most $C|v|^4L^{-4}$,
and its rank is at most $9(L^4-1)$.  The elementary rank bound for its
$\mathfrak S_5$ norm gives \cref{eq:v15-Bornideal}.  Finally the first Born
trace in the Gaussian negative logarithm has prefactor $1/2$.
\end{proof}

For $v=(xE_1,yE_2,0,0)$ one has $E_4=2x^2y^2$, so the isolated trace is
$-3(1-L^{-4})x^2y^2$.  In particular its coefficients at $L=2,4$ are
$-45/16$ and $-765/256$.  This is not the complete quartic one-loop
coefficient: the other Wilson terms, the nonlinear constraint, Haar and
coordinate factors can contribute and cancel.  It is not a two-loop
Laplace coefficient or a mass.  Its role is to prevent the false inference
that a term vanishing in the determinant ideal may be removed from the
separately retained low Born traces.

\section{Verification, proof boundary, and append-only record}
\label{sec:v15-next}

The actual lattice source and local Wilson contributions have now been
transported on the same Fourier energy carrier, including their constant
normal corrections and the nonlinear multiplier.  This closes V14's stated
lattice-precision comparison on the one-cell small classical branch.
The fifth determinant tail, its normal coarea factor and all fixed source
derivatives have a regulator limit.  The proof does not rely on treating the
lattice as a continuum spectral compression or on excluding ultraviolet
fluctuation directions from its matrices.

The next task should be the signed, complete first-four-Born packet with its
compact density and gauge factors, not a repeat of the fifth-determinant
argument.  The explicitly computed surviving term in
\cref{thm:v15-Born} supplies one necessary regression for that calculation.
Any Ward-compatible subtraction must use the complete measured expression
and retain its finite nonlocal part.  This version does not show that the
full packet is local, compute the running coupling, evaluate $J^6Q$ or
$J^4B$, or obtain a quantum Gibbs or Schwinger functional.  Localizing the
global critical-energy branch to thermodynamic field regions, the compact
well/transition-cutoff coverage, noncollapse and physical-time contraction
remain separate.

\paragraph{Verification scope.}
The new exact checker expands literal noncommuting quaternion plaquettes in
independent background and variation parameters, verifies the right-trivialized
Hessian against direct log-coordinate differentiation, checks its co-located
continuum limits, reconstructs full four-dimensional discrete curl and
Coulomb incidence, and evaluates the finite Fourier first-trace coefficient.
It separately checks normalization, grid multiplier counting, normal coarea,
and the distinction between a vanishing Schatten term and its retained trace.
The new driver passed 23 exact test families.  All 22 V14 exact families
were freshly rerun and passed; their report is separate from the older
archival reports.  These finite checks are not a numerical
proof of Cwikel's theorem, Schatten convergence or the earlier classical
background construction.  Those conclusions use the written estimates and
the explicitly inherited analytic inputs.

The V1--V14 mathematical body is retained byte-for-byte; only preamble version
displays change.  The bundle contains this appendix, the cumulative source
and PDF, the exact checker and report, an assembler, and source/output hashes.
No Lean formalization, certified numerical nonlinear radius, complete quantum
one-loop coefficient, or Yang--Mills mass-gap proof is claimed.  A later V16
is to append below this version's end marker.  Any correction must be stated
explicitly rather than silently changing a preceding proof.

% END OF VERSION 15 MATHEMATICAL BODY.
% Append subsequent requested versions below this marker; retain V1--V15 above.


\clearpage
\section{V16: the missing measure and a complete flat-root one-loop limit}
\label{sec:v16-context}

V15 compares the literal lattice constrained precision with its continuum
counterpart. Its last paragraph correctly leaves the compact density and the
first four Born traces unassembled. Another estimate for the fifth modified
determinant would not determine those terms. This version changes coordinates
only after writing the exact local Haar disintegration, distinguishes its gauge
Jacobian from the normal root Jacobian, and evaluates the resulting
\emph{complete measured one-loop coefficient} on a flat one-holonomy branch.
The fine lattice is still four dimensional and every nonconstant momentum and
both charged colours are retained. The root source is the literal balanced
nested source, not a constraint on the mean.

There are two main conclusions. First, in the same local Coulomb chart as V15,
the missing density is an explicitly defined logarithmic Faddeev--Popov
determinant times the full fine-link Haar density. Its relative operator has
an actual lattice-to-continuum fifth-determinant limit. Second, for
$v=(xE_3,0,0,0)$, all gauge, Haar, root-normal, and tangent-return terms can be
combined before a limit is taken. The answer is a finite three-dimensional
sum with an exponentially convergent continuum series and an $O(L^{-2})$
coefficient-function error. This is a coefficient of the local stationary
expansion, not an exact interacting integral or a proof that its remainder is
uniform in a simultaneous coupling/refinement limit.

\paragraph{Audit and conventions.}
The Grand Tensor already has an exact Haar quotient in tree coordinates,
without an added Faddeev--Popov factor
\cite[\texttt{thm:YM-relative-tree-gauge}]{GT76V6}. The determinant below is
\emph{the change to logarithmic Coulomb coordinates}; it is not an additional
factor to multiply into that tree quotient. Logarithmic lattice gauge operators
are standard; \cite[Section IV.B, equations (34)--(37)]{LogGaugeV16} supplies
primary context. Their exact finite formula is derived below in the present
anti-Hermitian quaternion normalization. The earlier flat-background shift
identity in f12 V9 is also inherited, not relabelled as a new principle. What is
new here is the measured Coulomb assembly on the present nested root fibre,
its all-refinement flat evaluation, and its root-specific normal return.

All gauge-fixing identities are local finite-regulator identities on a regular
branch. Constant gauge rotations are retained, not divided out by a singular
orbit volume. The small classical branch, the lattice precision $K_L$, and its
normal root matrix $M_L$ are those of V12--V15. No other compact wells or global
Gribov regions are classified. Throughout this appendix $L=2^n$ is refinement,
$x$ is a retained holonomy angle, and $\beta^{-1}$ is the noise variance.

\section{The local Haar law in the actual logarithmic Coulomb slice}
\label{sec:v16-measure}

Let $a_e\in\mathfrak{su}(2)$ be logarithmic link coordinates, with
$[u,v]=2u\mathbin\times v$. On vertices set
$(d\omega)_{xy}=\omega_y-\omega_x$, $\Delta=d^*d$, and let
$\mathcal C=\ker d^*$ be the fixed linear Coulomb slice. A prime denotes the
mean-zero vertex space, omitting exactly three constant gauge directions.
The full fine Haar density in logarithmic coordinates, up to a constant, is
\begin{equation}\label{eq:v16-Haar}
 J_{\rm H}(a)=\prod_e\left(\frac{\sin|a_e|}{|a_e|}\right)^2,
 \qquad \ell_{\rm H}(a)=\log J_{\rm H}(a).
\end{equation}
The product is over \emph{all} fine links in this coordinate presentation.
It must not be replaced by the product over only the non-tree links while
retaining the Coulomb Jacobian derived in this section.

Use the gauge action $U_{xy}\mapsto g_x^{-1}U_{xy}g_y$, with
$g_x=\exp(t\omega_x)$. Define $R_a\omega$ to be the derivative of its link
logarithm at $t=0$. Introduce the analytic identity germs
\[
 f(z)=\frac{z}{e^z-1},\qquad
 \sigma(z)=\frac z2\coth\frac z2
       =1+\frac{z^2}{12}-\frac{z^4}{720}+\frac{z^6}{30240}+\cdots.
\]
The apparent values at zero are removable.

\begin{lemma}[Exact logarithmic gauge derivative]
\label{lem:v16-gauge}
For each oriented link $e=(x,y)$,
\begin{align}\label{eq:v16-R}
 (R_a\omega)_e
 &=-f(\operatorname{ad}_{a_e})\omega_x
       +f(-\operatorname{ad}_{a_e})\omega_y\notag\\
 &=\sigma(\operatorname{ad}_{a_e})(\omega_y-\omega_x)
       +\tfrac12[a_e,\omega_x+\omega_y].
\end{align}
On $\mathcal C$ the operator
\begin{equation}\label{eq:v16-FP}
                       \mathscr M_a=d^*R_a
\end{equation}
annihilates constant vertex fields and is symmetric on real coordinates.
Its restriction to the primed space is positive near zero.
\end{lemma}
\begin{proof}
The right-trivialized derivative of $e^a$ is
$\mathscr E(a)=(e^{\operatorname{ad}a}-1)/\operatorname{ad}a$.
The gauge derivative of $U$ on the right is
$-\omega_x+\operatorname{Ad}_{e^{a_e}}\omega_y$. Applying
$\mathscr E(a_e)^{-1}$ gives the first line of \cref{eq:v16-R}.
The identities $f(z)=\sigma(z)-z/2$ and
$f(-z)=\sigma(z)+z/2$ give the second line.

For a constant $\omega$, $R_a\omega=[a,\omega]$, so
$\mathscr M_a\omega=[d^*a,\omega]=0$ on the slice.
The even function $\sigma(\operatorname{ad}a)$ is symmetric for real compact
Lie data. The antisymmetric part of the remaining bilinear form is
\[
 \langle\phi,\mathscr M_a\omega\rangle
 -\langle\mathscr M_a\phi,\omega\rangle
 =\sum_x\langle\phi_x,[(d^*a)_x,\omega_x]\rangle=0.
\]
This identity follows by pairing the two endpoints of every link; no boundary
term occurs on the torus. At zero the primed operator is $\Delta\succ0$.
Continuity proves finite-regulator positivity near zero. A uniform energy
version is proved in the next section.
\end{proof}

\begin{proposition}[The local quotient density is fixed, not optional]
\label{prop:v16-density}
On a regular local Coulomb branch, integration of gauge-invariant functions
against fine product Haar measure reduces, up to a field-independent constant,
to the density
\begin{equation}\label{eq:v16-density}
 \boxed{\quad
 \rho_{\mathcal C}(a)
 =J_{\rm H}(a)\,
       \frac{\det{}'\mathscr M_a}{\det{}'\Delta}
 \quad}
\end{equation}
relative to fixed Euclidean volume on $\mathcal C$. The ratio is positive on
its branch through one. Its logarithm is counted once.
\end{proposition}
\begin{proof}
Fix the gauge at one vertex and use the based gauge group. It acts freely on
links of the connected finite graph. The derivative at a slice point of the
local map from slice coordinates and based gauge coordinates to link
logarithms has columns $[I_{\mathcal C},R_a]$. Projecting its second block
onto the fixed normal space gives $\Delta^{-1/2}d^*R_a$, with the appropriate
fixed isomorphism from based to mean-zero vertex coordinates. At zero the same
matrix has $R_0=d$. The ratio of the two Jacobians is consequently
$\det{}'\mathscr M_a/\det{}'\Delta$.
Indeed $d^*R_a$ kills constants on the slice, so replacing a based gauge
parameter by its mean-zero representative changes this determinant only by a
fixed normalization. Multiplication by the full logarithmic Haar density gives
\cref{eq:v16-density}. Gauge invariance and Haar invariance make the remaining
based-group integration a constant, when the local slice is saturated in the
chosen orbit neighbourhood. Equivalently the calculation gives the local
quotient volume form, with compatible localizations used for its germ.
This argument neither inserts a second gauge volume into the tree quotient
nor makes a global uniqueness claim for Coulomb gauge.
\end{proof}

There is a small but important source issue. The raw rooted one-cell output is
invariant under based fine gauges, whereas balancing can conjugate that output
by a field-dependent \emph{constant coarse} group element. The Coulomb quotient
law, on a constant-conjugation-invariant local branch, is invariant under
constant conjugation. Both source maps are equivariant for those rotations and
have the same output orbit pointwise. Their pushforward laws are therefore
equal: average any test function over constant conjugations first, and its
averaged value is identical on the two outputs. This justifies the use of the
balanced source for the invariant local density. It does not justify dropping
a Jacobian from a representative-valued coordinate change without this
measure argument. A non-invariant Cartesian localization has its own boundary
comparison; stationary germs may be represented by a compatible invariant
localization equal to one near the branch.

\section{The actual gauge determinant also has a continuum tail}
\label{sec:v16-ghost}

On the primed scalar-gauge energy space set
\begin{equation}\label{eq:v16-G}
 I+G_L^{\rm gh}(a)=\Delta^{-1/2}\mathscr M_a\Delta^{-1/2}.
\end{equation}
A fixed Fourier carrier is used exactly as in V15, but with scalar vertex
orientations and three Lie components; the free zero-mode sector is omitted.
No physical degree of freedom is removed by adding a small ghost mass.

\begin{theorem}[Uniform gauge coercivity and actual Schatten transport]
\label{thm:v16-ghost}
For one sufficiently small global fine $\ell^4$ ball, independently of $L$,
\begin{equation}\label{eq:v16-ghost-bound}
 \|G_L^{\rm gh}(a)\|\le C(\|a\|_4+\|a\|_4^2)<\tfrac12.
\end{equation}
On the Coulomb slice the real primed gauge form is thus at least
$\frac12\Delta$. Along the actual V13 backgrounds, the Fourier-embedded
operators converge in $\mathfrak S_5$ to
\begin{equation}\label{eq:v16-ghostlimit}
 G^{\rm gh}(A)=J_0^{-1/2}d^*\operatorname{ad}_A J_0^{-1/2},
\end{equation}
where $J_0=-\sum_\mu\partial_\mu^2$ is the positive continuum scalar
Laplacian on mean-zero functions, and $(\operatorname{ad}_A\omega)_\mu=[A_\mu,\omega]$. The maps and all fixed source derivatives converge locally uniformly
on smaller complex source domains. In particular
$\mathcal L_5(G_L^{\rm gh})\to\mathcal L_5(G^{\rm gh})$ with those derivatives.
\end{theorem}
\begin{proof}
The odd term of \cref{eq:v16-R}, paired with $d\omega$, is bounded by
$C\|a\|_4\|d\phi\|_2\|\omega\|_4$ plus its symmetric counterpart.
The lattice Sobolev inequality on mean-zero scalar fields gives the linear
term of \cref{eq:v16-ghost-bound}. The even remainder has
$\|\sigma(\operatorname{ad}a)-I\|_\infty\le C\|a\|_\infty^2$ and therefore
the quadratic bound. This proves uniform form positivity, not a lower bound
for the unscaled smallest Laplacian eigenvalue.

In V15's physical grid convention, the odd relative term is a bounded
normalized gradient on the left, multiplication by $A_L/2$, a fixed finite
sum of translations, and the free inverse-square-root on the right. The
actual finite-grid Cwikel bound and its consistency proof
\cref{lem:v15-Cwikel,lem:v15-consistency} apply unchanged to scalar vertex
gauge fields. Strong physical $L^4$ convergence gives convergence in
$\mathfrak S_p$, $p>4$, to \cref{eq:v16-ghostlimit}.

The remaining exact term is
\[
 E_L^{\rm gh}=(d\Delta^{-1/2})^*
       M_{\sigma(\operatorname{ad}a)-I}(d\Delta^{-1/2}).
\]
It has $\mathfrak S_2$ norm at most $C\|a\|_4^2$, since its multiplier
has squared eigenvalue sum at most $C\sum_e|a_e|^4$. For $p>2$,
\begin{equation}\label{eq:v16-evenideal}
 \|E_L^{\rm gh}\|_p^p
       \le C\|a\|_\infty^{2p-4}\sum_e|a_e|^4\longrightarrow0.
\end{equation}
Here $\max|a_L|\to0$ follows from the strong physical $L^4$ convergence as in
\cref{eq:v15-linksmall}. This proves the real $\mathfrak S_5$ statement.
The same estimates give uniform ideal bounds on the inherited complex
classical branches. Real norm convergence, uniform Banach-valued Cauchy bounds,
and the coefficient argument in \cref{thm:v15-main} give local complex and
source-derivative convergence. V14's trace-series theorem then transfers the
modified determinant. No new external ideal theorem beyond the already
specified Cwikel input is being assumed.
\end{proof}

The even gauge term has a low-trace contribution despite
\cref{eq:v16-evenideal}. At the homogeneous fine connection
$a_1=xE_3/L$, $a_2=a_3=a_4=0$, its quadratic coefficient has
\begin{equation}\label{eq:v16-even-trace}
 \boxed{\quad
 \Tr(E_L^{\rm gh})^{[2]}=-\frac{L^4-1}{6L^2}x^2,
 \qquad \|(E_L^{\rm gh})^{[2]}\|_5=O(x^2L^{-6/5}).\quad}
\end{equation}
Indeed $\Tr_{\mathfrak{su}(2)}(\operatorname{ad}_{xE_3})^2=-8x^2$,
the coefficient of $\sigma-I$ is $1/12$, and
$\sum_{k\ne0}|q_1(k)|^2/\lambda(k)=(L^4-1)/4$.
The rank is at most $3(L^4-1)$ and the operator norm is $O(x^2/L^2)$,
which proves the ideal bound. This is another literal low-Born term which
must not be discarded when its high-order determinant tail disappears.

\section{The complete finite one-loop expression in one measured convention}
\label{sec:v16-assembly}

Write $\ell_c(v)=\sum_{\mu=1}^4 2\log(\sin|v_\mu|/|v_\mu|)$ for the coarse
product-Haar logarithm. All functions in this section are relative to their
zero-source values, so field-independent Gaussian and coordinate-scale factors
are separated. Retain V15's actual constrained precision $K_L$ and normal
source matrix $M_L$.

\begin{theorem}[The missing gauge and Haar factors complete the finite formula]
\label{thm:v16-assembly}
The one-loop coefficient of the local rooted pushforward, expressed relative
to coarse product Haar, is
\begin{equation}\label{eq:v16-assembly}
 \boxed{\begin{aligned}
 \widehat Q_L(v)={}&\log\det M_L(v)
        +\tfrac12\log\det(I+K_L(v))\\
    &-\log\det(I+G_L^{\rm gh}(a_L^*(v)))
        -\ell_{\rm H}(a_L^*(v))+\ell_c(v).
 \end{aligned}}
\end{equation}
This is a finite local coefficient identity, with $\widehat Q_L(0)=0$.
The gauge determinant in the second line and the normal determinant in the
first line have different domains and different roles.
\end{theorem}
\begin{proof}
Use \cref{prop:v16-density} on the based-gauge quotient. In fixed
constant-plus-detail Coulomb coordinates impose the balanced source using
V12's chart. Its normal Jacobian is $M_L$, so its coarea density is
$(\det M_L)^{-1}$. Laplace integration over the detail coordinates gives
$\det(I+K_L)^{-1/2}$ after the source-independent free energy normalization.
Evaluate the amplitude \cref{eq:v16-density} at the classical minimum; its
derivatives first enter the next Laplace coefficient and are not extra
one-loop Hessian terms. Taking the negative logarithm gives the first four
terms of \cref{eq:v16-assembly}. Conversion from coarse Lebesgue to coarse
Haar adds $\ell_c$. The preceding orbit argument identifies the invariant
balanced and raw rooted pushforwards on the same local branch.
\end{proof}

There is a useful algebraic check on this expression. Let $E$ be the full
link right-exponential derivative, and let $\mathbb B_{\log}$ be a bordered
stationary Hessian with all gauge and root constraints in log coordinates.
Under the linear tangent-coordinate change to right Lie variations,
\[
 \mathbb B_{\log}
 =\operatorname{diag}(E^{\mathsf t},I)\,
   \mathbb B_{\rm right}\,
   \operatorname{diag}(E,I),\qquad
 \det\mathbb B_{\log}=(\det E)^2\det\mathbb B_{\rm right}.
\]
At a stationary Lagrangian the ordinary Hessian chain rule has no additional
first-gradient term. Since $J_{\rm H}=\det E$, the corresponding Haar factor
cancels in this \emph{complete} bordered change. All constraint rows also
change. Deleting $\ell_{\rm H}$ while leaving V15's old log-coordinate
Hessian and constraint rows fixed would not implement this identity.

Let $c_r=(-1)^{r+1}/r$. The exact Born decomposition is now
\begin{align}\label{eq:v16-headtail}
 \widehat Q_L={}&\log\det M_L+
     \tfrac12\mathcal L_5(K_L)-\mathcal L_5(G_L^{\rm gh})+\mathcal H_L,\\
 \mathcal H_L={}&
   \sum_{r=1}^4 c_r\bigl\{\tfrac12\Tr K_L^r-\Tr(G_L^{\rm gh})^r\bigr\}
                -\ell_{\rm H}(a_L^*)+\ell_c.\notag
\end{align}
The displayed tail and normal factor converge with all fixed derivatives by
V15 and \cref{thm:v16-ghost}. The head is now a \emph{specified measured
finite expression}, not an unspecified gauge factor. Its general
noncommuting ultraviolet subtraction is still to be proved. The next sections
evaluate the entire expression, without Born truncation, on a nontrivial flat
source. This directly tests the cancellations which separate absolute trace
bounds cannot establish.

\section{Exact charged-mode cancellation at a flat one-holonomy background}
\label{sec:v16-flatbulk}

Take the actual source and flat fine connection
\begin{equation}\label{eq:v16-flat}
 v=(xE_3,0,0,0),\qquad a_1=xE_3/L,\qquad a_2=a_3=a_4=0.
\end{equation}
All plaquettes equal the identity and the literal nested source is exactly
$v$. The full Wilson gradient and its source multiplier vanish. The flat
connection is therefore the local constrained minimum on the inherited branch.
The following formulas are analytic on a fixed neighbourhood of $x=0$; for
the explicit formula one may take a sufficiently small subdisc of $|x|<1/4$.
They define the local one-loop coefficient, not a sum over all flat sectors.

Put $\varepsilon=x/L$. For $k\in(2\pi/L)\{0,\ldots,L-1\}^4$ set
\[
 q_\mu=e^{ik_\mu}-1,\qquad \lambda_0=\sum_\mu|q_\mu|^2,
 \qquad q_{+,\mu}=e^{i(k_\mu+2\varepsilon\delta_{\mu1})}-1,
 \qquad \lambda_+=|q_+|^2.
\]
The other charge has $\varepsilon\mapsto-\varepsilon$. Neutral fields have
no background dependence in these flat mode matrices. Only $k\ne0$ are
integrated; the twelve constant link coordinates are normal source coordinates,
not fluctuation modes to be silently included.

\begin{lemma}[The complete nonconstant-mode density reduces to a shifted scalar symbol]
\label{lem:v16-mode}
With a linear mean constraint in place of the root constraint only for this
intermediate calculation, the relative one-loop coefficient is
\begin{equation}\label{eq:v16-meanQ}
 \boxed{\quad Q_L^{\rm mean}(x)
   =2\sum_{k\ne0}\log\frac{\lambda_+(k)}{\lambda_0(k)}
                 -2\log\frac{\sin\varepsilon}{\varepsilon}.\quad}
\end{equation}
It includes the full fine Haar density and the logarithmic gauge determinant.
The last term is the uncancelled constant-coordinate exponential volume.
\end{lemma}
\begin{proof}
In one charged complex component the logarithmic-to-right tangent matrix is
$E=\operatorname{diag}(e^{i\varepsilon}\sin\varepsilon/\varepsilon,1,1,1)$.
The right-coordinate flat Wilson Hessian is
$H_+=\lambda_+I-q_+q_+^*$; its kernel is the gauge vector $q_+$.
The logarithmic gauge scalar is
\begin{equation}\label{eq:v16-fpsymbol}
 m(k)=q^*E^{-1}q_+
 =|q_1|^2\varepsilon\cot\varepsilon
       +2\varepsilon\sin k_1+\sum_{\nu>1}|q_\nu|^2.
\end{equation}
It is positive at small real $x$ for each nonzero mode and has its analytic
branch through $\lambda_0$.

Choose an orthonormal basis $C_0$ of $q^\perp$. The elementary codimension-one
determinant identity gives
\begin{equation}\label{eq:v16-mode-det}
 \det(C_0^*E^*H_+EC_0)
 =|\det E|^2\lambda_+^2\frac{|m(k)|^2}{\lambda_0}.
\end{equation}
One proof completes $C_0$ by $q/|q|$ and the target transverse basis by
$q_+/|q_+|$. The transverse volume of $EC_0$ is
$|\det E|\,|q^*E^{-1}q_+|/(|q||q_+|)$. The three positive eigenvalues of
$H_+$ are $\lambda_+$, yielding \cref{eq:v16-mode-det}.
This proof retains the Coulomb plane actually used in V15, not a replacement
by the covariant plane without its volume factor.

In the negative logarithm, half of \cref{eq:v16-mode-det}, the fine Haar
volume, and the gauge determinant cancel mode by mode to leave
$\log(\lambda_+/\lambda_0)$ per real charge. The two charged spectra have the
same sum by $k\mapsto-k$, giving the factor two in \cref{eq:v16-meanQ}.
There are $L^4$ fine exponential volume factors but only $L^4-1$ nonconstant
Fourier modes. Their remaining factor is
$(\sin\varepsilon/\varepsilon)^2$, giving the last term. All other
normalizations are independent of $x$ and have been removed relative to zero.
\end{proof}

This cancellation is exact before taking a mesh limit. In particular it
contains the divergent quadratic even-gauge term
\cref{eq:v16-even-trace}, the fine Haar quadratic term
$L^2x^2/3$, and their Wilson partners. It cannot be inferred by taking the
absolute value of each separately.

\section{The actual root constraint has a six-dimensional Gaussian return}
\label{sec:v16-rootreturn}

The mean constraint in \cref{lem:v16-mode} is not the problem's source.
The correction to it can be computed exactly at \cref{eq:v16-flat}.
Define
\begin{equation}\label{eq:v16-rh}
 r_L(x)=\frac{\sin x}{L\sin(x/L)},\qquad
 h_L(x)=1-\frac23(1-L^{-2})\sin^2x.
\end{equation}
Both have value one at zero.

\begin{lemma}[Literal transverse source derivative on the flat branch]
\label{lem:v16-rootrow}
For every fine logarithmic variation $w$ and every $\mu\ne1$,
\begin{equation}\label{eq:v16-rootrow}
 \boxed{\quad
 (D\mathcal F_L(a)w)_\mu
 =L\left\langle
  \operatorname{Ad}_{\exp((x_1-(L-1)/2)\varepsilon E_3)}w_\mu(x_{\rm site})
                                            \right\rangle_L.\quad}
\end{equation}
Here $x_1$ is the integer first coordinate of the site, while the unadorned
$x$ in $\varepsilon=x/L$ is the source angle. The normal root Jacobian has
\begin{equation}\label{eq:v16-normalflat}
                            \det M_L=r_L(x)^6.
\end{equation}
\end{lemma}
\begin{proof}
At one step, the background tree is
$t_s=\exp(s_1\varepsilon E_3)$. For an output orientation $\mu\ne1$ its
background path is the identity. The derivative of the distinguished-root
average at identical identity paths is their ordinary mean, exactly, regardless
of higher derivatives of that average. Differentiating the two tree factors
produces a coarse difference in orientation $\mu$. The balancing factors give
another such difference. The remaining term is the two-link variation averaged
with the conjugation weight
$\operatorname{Ad}_{\exp((s_1-1/2)\varepsilon E_3)}$.

In the final spatial average, all these differences telescope: their later
weights depend on the first coordinate only and $\mu\ne1$. Under another
step the background angle doubles, and the conjugation weights multiply.
Their exponent on a fine site is
\[
 \sum_{j=0}^{n-1}2^j\bigl((\lfloor x_1/2^j\rfloor\bmod2)-1/2\bigr)
                   =x_1-(L-1)/2.
\]
The longitudinal two-link sums in orientation $\mu$ contribute the total factor
$L$. This proves \cref{eq:v16-rootrow}, including arbitrary variations in
other orientations, whose tree-only contributions have telescoped.

For a constant variation $s_\mu/L$, the charged average in
\cref{eq:v16-rootrow} is $r_L(x)$ and the neutral average is one. For changes
in the first constant component alone, the homogeneous source is exactly that
component, so the corresponding $3$-by-$3$ diagonal block is the identity.
Other constant directions can rotate the first output, but the transverse
outputs have no derivative from them except their own diagonal block.
The full normal matrix is block triangular with six charged diagonal entries
$r_L$ and all other diagonal entries one. This gives \cref{eq:v16-normalflat}.
\end{proof}

\begin{proposition}[The complete tangent correction is finite rank and evaluated]
\label{prop:v16-return}
Replacing the mean constraint by the actual root constraint adds to the
negative-log Gaussian determinant
\begin{equation}\label{eq:v16-return}
 \boxed{\quad \frac12\log\frac{\det h_L^{\rm root}}{\det h_L^{\rm mean}}
                  =3\log d_L(x),\qquad d_L(x)=r_L(x)^{-4}h_L(x).\quad}
\end{equation}
This is in addition to the normal coarea contribution $6\log r_L$.
The correction begins at fourth source order, but it is retained exactly.
\end{proposition}
\begin{proof}
At the constant flat connection the full Wilson gradient is zero, and the
multiplier is zero. Fourier modes do not mix in its Hessian. The constant
Hessian vanishes on the first-orientation charged components and on all neutral
constants. On the six charged transverse constants, expressed in physical
mean coordinates $s$, it is
$\tau_0 I$ with $\tau_0=4L^2\sin^2\varepsilon$.
The actual constraint tangent is the mean-zero detail plus its normal
correction $-M_L^{-1}D\mathcal F_L|_{V_L}$. Thus the constrained Hessian is
the mean-fixed Hessian plus a positive term of rank at most six.

The row in \cref{eq:v16-rootrow} depends only on axial Fourier modes.
Those modes have $w_1=0$ in the free Coulomb slice, so their remaining
charged logarithmic covariance is exactly
$[L^4\,4\sin^2(\varepsilon+\pi m/L)]^{-1}$ in the convention
$w(x)=\sum_m\widehat w_m e^{2\pi imx_1/L}$.
The squared normalized Fourier weight of the conjugation average is
$\sin^2x/[L^2\sin^2(\varepsilon+\pi m/L)]$.
Consequently the detail source covariance in one charged transverse component is
\begin{equation}\label{eq:v16-sigma}
 \Sigma_L(x)=\frac{\sin^2x}{4L^4}
           \sum_{m=1}^{L-1}\csc^4(\varepsilon+\pi m/L).
\end{equation}
The omitted $m=0$ is a retained constant, not an integrated detail. All six
components have this same covariance and their cross terms vanish by
orientation and charge symmetry. The determinant lemma therefore gives
$3\log(1+\tau_0r_L^{-2}\Sigma_L)$.

Differentiate the elementary cotangent product identity twice to obtain
\[
 \sum_{m=0}^{L-1}\csc^4(z+\pi m/L)
 =L^4\csc^4(Lz)-\frac23L^2(L^2-1)\csc^2(Lz).
\]
For example it follows from
$\sum_m\csc^2(z+\pi m/L)=L^2\csc^2(Lz)$ and
$(f''+4f)/6=f^2$ for $f=\csc^2z$.
Restoring the missing term in \cref{eq:v16-sigma} gives
\[
 1+\tau_0r_L^{-2}\Sigma_L
 =\sin^4\varepsilon\sum_{m=0}^{L-1}\csc^4(\varepsilon+\pi m/L)
 =r_L^{-4}h_L.
\]
This proves \cref{eq:v16-return}. The expression is understood by its
removable limit at $x=0$. No additional real charge is introduced by a
self-conjugate Nyquist mode; the determinant count is the real dimension six.
\end{proof}

\section{A closed expression and a summable continuum limit for the measured coefficient}
\label{sec:v16-main}

For $m\in\mathcal K_L^3$, the centered three-dimensional momentum cube, put
\begin{equation}\label{eq:v16-D}
 c_L(m)=2\sum_{\nu=1}^3\sin^2(\pi m_\nu/L),\qquad
 t_L(m)=L\operatorname{arcosh}(1+c_L(m)),\qquad
 D_L(m)=\cosh t_L(m)-1.
\end{equation}
This $c_L(m)$ is a transverse dispersion, unrelated to V11's scalar Green
coefficient $c_L$. Only the notation within this section is used for it.
Equivalently $D_L(m)=T_L(1+c_L(m))-1$, where $T_L$ is the Chebyshev polynomial.

\begin{theorem}[Complete measured flat-root one-loop coefficient]
\label{thm:v16-flat}
With all factors in \cref{eq:v16-assembly}, the local one-loop coefficient on
\cref{eq:v16-flat}, relative to its value at zero, is exactly
\begin{equation}\label{eq:v16-flatQ}
 \boxed{\quad
 \widehat Q_L(xE_3,0,0,0)
 =3\log h_L(x)
  +2\sum_{m\in\mathcal K_L^3\setminus\{0\}}
       \log\left(1+\frac{1-\cos(2x)}{D_L(m)}\right).\quad}
\end{equation}
The identities refer to the actual local Gaussian/Laplace coefficient at every
finite $L$, not to an arbitrary continuum regularization. Neutral modes,
nonconstant charged modes, fine and coarse Haar, the logarithmic gauge
determinant, and both root determinant factors have all been accounted for.
\end{theorem}
\begin{proof}
At fixed transverse momentum $m\ne0$, use the polynomial identity
\[
 \prod_{j=0}^{L-1}2\bigl(z-\cos(2\pi j/L+\delta)\bigr)
                     =2\bigl(T_L(z)-\cos(L\delta)\bigr).
\]
It follows by comparing the roots and leading coefficients of the two
polynomials. Set $z=1+c_L(m)$ and $\delta=2x/L$ to evaluate the longitudinal
part of the shifted determinant. Its ratio to the unshifted product is
$1+(1-\cos 2x)/D_L(m)$.

For $m=0$ omit the longitudinal zero mode as required in
\cref{eq:v16-meanQ}. The full shifted product is $4\sin^2x$, the omitted
factor is $4\sin^2(x/L)$, and the unshifted nonzero product is $L^2$.
The ratio is $r_L^2$. The two charges therefore contribute $4\log r_L$.
Equation \cref{eq:v16-meanQ} becomes
\[
 4\log r_L+2\sum_{m\ne0}\log\left(1+\frac{1-\cos2x}{D_L(m)}\right)
                           -2\log\frac{\sin(x/L)}{x/L}.
\]
Add the actual normal root contribution $6\log r_L$, the tangent return
$3\log(r_L^{-4}h_L)$, and the coarse Haar logarithm
$2\log(\sin x/x)$. The logarithms of $r_L$ and of the two sine ratios cancel
exactly, leaving \cref{eq:v16-flatQ}. In particular the single uncancelled
fine constant-mode volume in \cref{eq:v16-meanQ} is essential.
As a check, at $L=1$ the sum is empty and $h_1=1$, giving zero for the identity
block. Without that constant-mode volume the identity-block answer would be
incorrect.
\end{proof}

\begin{theorem}[An identified continuum coefficient with geometric refinement error]
\label{thm:v16-limit}
On every smaller complex disc $|x|\le\rho<1/4$ admitted by the local branch,
\begin{align}\label{eq:v16-Qlimit}
 \widehat Q_\infty(x)
 ={}&3\log\left(1-\tfrac23\sin^2x\right)\notag\\
 &+2\sum_{m\in\mathbb Z^3\setminus\{0\}}
       \log\left(1+\frac{1-\cos(2x)}{\cosh(2\pi|m|)-1}\right)
\end{align}
converges normally, and
\begin{equation}\label{eq:v16-rate}
 \boxed{\quad
 \sup_{|x|\le\rho}|\widehat Q_L(x)-\widehat Q_\infty(x)|\le C_\rho L^{-2}.
 \quad}
\end{equation}
Every fixed $x$ derivative has the same rate on smaller discs. The error is
summable for dyadic refinement. This is convergence of the complete measured
one-loop coefficient on this flat family, without deleting its low Born terms.
\end{theorem}
\begin{proof}
For centered $m$, put $u=(\sum_\nu\sin^2(\pi m_\nu/L))^{1/2}$.
Then $t_L(m)=2L\operatorname{arsinh}u$. Since $u\le\sqrt3$ and
$\operatorname{arsinh}u\ge u/\sqrt{1+u^2}\ge u/2$, the elementary sine
bound gives
\begin{equation}\label{eq:v16-decay}
                         2|m|\le t_L(m)\le2\pi|m|.
\end{equation}
Thus the summands in \cref{eq:v16-flatQ}, and their fixed complex-$x$
derivatives on smaller discs, are bounded by $C_\rho e^{-2|m|}$.
The denominator is uniformly separated from zero for $m\ne0$, and
$|1-\cos 2x|$ is uniformly small on the stated disc, so all logarithms use
their common branch through zero.

For $|m|\le\varepsilon L$ with a fixed small $\varepsilon$, the sine and
inverse-hyperbolic-sine Taylor estimates give
\[
                  |t_L(m)-2\pi|m||\le C|m|^3L^{-2}.
\]
The derivative with respect to $t$ of
$\log(1+(1-\cos2x)/(\cosh t-1))$ is bounded by $C_\rho e^{-t}$ for
$t\ge2$. Equation \cref{eq:v16-decay} therefore makes the summed low-frequency
error at most
$C_\rho L^{-2}\sum_m |m|^3e^{-2|m|}$.
The complementary modes have exponentially small total contribution, both
for the lattice and the continuum series. The missing modes outside the
centered cube obey the same bound. Finally $h_L-h_\infty=O_\rho(L^{-2})$
and their logarithms have a common nonzero neighbourhood. This proves
\cref{eq:v16-rate}. Normal convergence proves holomorphy; Cauchy on a slightly
larger disc gives the differentiated rate. No asymptotic statement for the
full integral with $\beta$ and $L$ moving together has been used.
\end{proof}

\subsection{Exact low-order consequences and a finite complete regression}

Let
\[
 S_{j,L}=\sum_{m\in\mathcal K_L^3\setminus\{0\}}D_L(m)^{-j},\qquad
 S_{j,\infty}=\sum_{m\in\mathbb Z^3\setminus\{0\}}
                       (\cosh(2\pi|m|)-1)^{-j}.
\]
The complete coefficient has the expansion
\begin{equation}\label{eq:v16-Qjets}
 \widehat Q_L(x)=q_{2,L}x^2+q_{4,L}x^4+O(x^6),
 \quad
 \begin{cases}
 q_{2,L}=-2(1-L^{-2})+4S_{1,L},\\
 q_{4,L}=\dfrac23L^{-2}(1-L^{-2})-\dfrac43S_{1,L}-4S_{2,L}.
 \end{cases}
\end{equation}
These are ordinary coefficients: the true second and fourth derivatives are
$2q_{2,L}$ and $24q_{4,L}$. In particular,
\begin{equation}\label{eq:v16-finitejets}
 \boxed{\begin{aligned}
 \widehat Q_2(x)&=-\frac{11}{24}x^2-\frac{211}{768}x^4+O(x^6),\\
 \widehat Q_4(x)&=-\frac{6887}{4704}x^2
             -\frac{3343129451}{33191424000}x^4+O(x^6).
 \end{aligned}}
\end{equation}
For $L=2$ the three nonzero transverse energy classes give
$S_{1,2}=25/96$ and $S_{2,2}=121/9216$. At $L=4$ they give
$S_{1,4}=1933/18816$ and
$S_{2,4}=93253451/132765696000$. These are full finite mode sums, not corner
samples standing in for an integral.

The continuum values are
\begin{equation}\label{eq:v16-limitjets}
 q_{2,\infty}=-2+4S_{1,\infty},\qquad
 q_{4,\infty}=-\frac43S_{1,\infty}-4S_{2,\infty}<0.
\end{equation}
The normal-return term $3\log(1-\frac23\sin^2x)$ has zero fourth
coefficient, explaining the simple last expression. The surviving source
potential is a finite-volume holonomy effect on a flat connection. Since the
curvature is zero, it cannot determine the noncommuting $F^2$ counterterm or
be interpreted as a local gluon mass. Its finite nonlocal part is precisely
why a complete Born subtraction must not discard every low-field term.

For comparison, the independent side-two coefficient assembly in the report is
\[
\begin{array}{c|cc}
 \text{contribution}&[x^2]&[x^4]\\\hline
 \tfrac12\log\det H_{\rm Coulomb,detail}^{\rm rel}&-4/3&-1577/3456\\
 -\log\det{}'\mathscr M^{\rm rel}&5/8&175/6912\\
 -\ell_{\rm H}(a)&4/3&1/90\\
 \log\det M_L&-3/4&-1/32\\
 \text{root tangent return}&0&3/16\\
 \ell_c&-1/3&-1/90\\\hline
 \text{complete measured sum}&-11/24&-211/768
\end{array}
\]
The first line here uses the \emph{mean-fixed} detail block only; the root
tangent return is therefore added separately. Adding that return to an already
root-constrained determinant would double-count it. In the general expression
\cref{eq:v16-assembly}, $K_L$ is already constrained and no separate return is
added. The table fixes this distinction as an explicit executable regression.

\section{What is now assembled and the remaining noncommuting calculation}
\label{sec:v16-next}

The missing local Coulomb density is now derived from the same compact Haar
law, and its gauge determinant has a controlled actual trace-ideal tail. The
finite one-loop formula is fully specified in the same convention as V15's
precision. On the flat one-holonomy branch it has been evaluated completely,
including the low Born contributions. The result has an identified continuum
limit and a summable coefficient error, with the normal source and its
six-dimensional tangent return retained.

This resolves a substantive measured calculation but not the general
noncommuting Born problem. V15's nonzero kinetic-weight quartic contribution
is invisible on the flat line, where $E_4=0$; it remains a required regression
on $v=(xE_1,yE_2,0,0)$. The next acquisition should compute the complete
$x^2y^2$ coefficient on that source plane using \cref{eq:v16-assembly}, including
the actual background/source jets and the new gauge density, before estimating
its logarithmic cutoff growth or assigning a running-coupling normalization.
A scalar flat holonomy potential cannot replace that mixed-curvature calculation.

No conclusion here gives a common nonperturbative $\beta,L$ integral remainder,
a thermodynamic small-field measure, suppression of the other compact wells,
coverage of transition-cutoff configurations, noncollapse, or calibrated
physical-time contraction. The deterministic relative gauge floor is not the
already excluded uniform hidden-diffusion gap. The local stationary coefficient
and its refinement limit are the proved objects.

\paragraph{Executed verification and preservation.}
The new driver \path{verify_ym_measured_Born_V16.py} passes twenty exact test
families. An independent dual-number quaternion implementation differentiates
the \emph{literal} spatial root words at the flat background through fourth
background order on side two and through second order on side four, with
arbitrary inhomogeneous variation legs. It checks the transverse row and the
normal matrix. Direct Bernoulli and exponential derivatives verify the
noncommuting logarithmic gauge formula and the Haar Jacobian. On the complete
side-two and side-four momentum grids, exact $\mathbb Q(i)$ matrix/series
arithmetic reconstructs the mean-fixed charged Hessian, gauge determinant,
fine Haar, root coarea, root return, and coarse Haar separately before summing.
These calculations do not use \cref{eq:v16-flatQ} to produce both sides.
Further checks prove finite Chebyshev products, the returning covariance,
identity-block normalization, and the displayed low coefficients.

The finite checks validate algebra and normalization, not a continuum theorem
by sampling. The analytic limit rests on \cref{eq:v16-decay} and the summed
dispersion estimates; the general ideal statement uses V15's explicitly
inherited multiplier inputs. All twenty-three V15 test families were freshly rerun and passed; that report
is separately identified in the provenance record. Older reports remain archival. No Lean formalization or numerical certification
of a mass gap is claimed.

V1--V15's mathematical body, including prior bibliography entries and version
markers, is preserved byte-for-byte. Only preamble version displays and the
table-of-contents page-number reserve change.
The new appendix, checker, exact report, portable assembly script, output PDF,
and hashes are included in the bundle. Any future correction is to be stated
in an appended section rather than silently editing this lineage.

\begin{thebibliography}{99}
\bibitem[LM11]{LogGaugeV16}
E.-M. Ilgenfritz, C. Menz, M. M\"uller-Preussker, A. Schiller and A. Sternbeck,
\emph{SU(3) Landau gauge gluon and ghost propagators using the logarithmic
lattice gluon field definition}, Physical Review D \textbf{83} (2011), 054506,
DOI: \nolinkurl{10.1103/PhysRevD.83.054506}, arXiv:1010.5120v2.
\url{https://arxiv.org/abs/1010.5120}.
Section IV.B, equations (34)--(37), was checked in the primary full text.
The standard logarithmic gauge operator is contextual prior work; the finite
slice measure and present root-specific determinant identities are proved here.
\end{thebibliography}

% END OF VERSION 16 MATHEMATICAL BODY.
% Append subsequent requested versions below this marker; retain V1--V16 above.

\end{document}
